\chapt{बालकाण्डः}

\sect{प्रथमः सर्गः}

\twolineshloka
{तपःस्वाध्यायनिरतं तपस्वी वाग्विदां वरम्}
{नारदं परिपप्रच्छ वाल्मीकिर्मुनिपुङ्गवम्} %||1-1-1||

\twolineshloka
{को न्वस्मिन्साम्प्रतं लोके गुणवान्कश्च वीर्यवान्}
{धर्मज्ञश्च कृतज्ञश्च सत्यवाक्यो दृढव्रतः} %||1-1-2||

\twolineshloka
{चारित्रेण च को युक्तः सर्वभूतेषु को हितः}
{विद्वान्कः कः समर्थश्च कश्चैकप्रियदर्शनः} %||1-1-3||

\twolineshloka
{आत्मवान्को जितक्रोधो मतिमान्कोऽनसूयकः}
{कस्य बिभ्यति देवाश्च जातरोषस्य संयुगे} %||1-1-4||

\twolineshloka
{एतदिच्छाम्यहं श्रोतुं परं कौतूहलं हि मे}
{महर्षे त्वं समर्थोऽसि ज्ञातुमेवंविधं नरम्} %||1-1-5||

\twolineshloka
{श्रुत्वा चैतत्त्रिलोकज्ञो वाल्मीकेर्नारदो वचः}
{श्रूयतामिति चामन्त्र्य प्रहृष्टो वाक्यमब्रवीत्} %||1-1-6||

\twolineshloka
{बहवो दुर्लभाश्चैव ये त्वया कीर्तिता गुणाः}
{मुने वक्ष्याम्यहं बुद्ध्वा तैर्युक्तः श्रूयतां नरः} %||1-1-7||

\twolineshloka
{इक्ष्वाकुवंशप्रभवो रामो नाम जनैः श्रुतः}
{नियतात्मा महावीर्यो द्युतिमान्धृतिमान्वशी} %||1-1-8||

\twolineshloka
{बुद्धिमान्नीतिमान्वाग्मी श्रीमाञ्शत्रुनिबर्हणः}
{विपुलांसो महाबाहुः कम्बुग्रीवो महाहनुः} %||1-1-9||

\twolineshloka
{महोरस्को महेष्वासो गूढजत्रुररिन्दमः}
{आजानुबाहुः सुशिराः सुललाटः सुविक्रमः} %||1-1-10||

\twolineshloka
{समः समविभक्ताङ्गः स्निग्धवर्णः प्रतापवान्}
{पीनवक्षा विशालाक्षो लक्ष्मीवाञ्शुभलक्षणः} %||1-1-11||

\twolineshloka
{धर्मज्ञः सत्यसन्धश्च प्रजानां च हिते रतः}
{यशस्वी ज्ञानसम्पन्नः शुचिर्वश्यः समाधिमान्} %||1-1-12||

\twolineshloka
{रक्षिता जीवलोकस्य धर्मस्य परिरक्षिता}
{वेदवेदाङ्गतत्त्वज्ञो धनुर्वेदे च निष्ठितः} %||1-1-13||

\twolineshloka
{सर्वशास्त्रार्थतत्त्वज्ञो स्मृतिमान्प्रतिभानवान्}
{सर्वलोकप्रियः साधुरदीनात्मा विचक्षणः} %||1-1-14||

\twolineshloka
{सर्वदाभिगतः सद्भिः समुद्र इव सिन्धुभिः}
{आर्यः सर्वसमश्चैव सदैकप्रियदर्शनः} %||1-1-15||

\twolineshloka
{स च सर्वगुणोपेतः कौसल्यानन्दवर्धनः}
{समुद्र इव गाम्भीर्ये धैर्येण हिमवानिव} %||1-1-16||

\twolineshloka
{विष्णुना सदृशो वीर्ये सोमवत्प्रियदर्शनः}
{कालाग्निसदृशः क्रोधे क्षमया पृथिवीसमः} %||1-1-17||

\twolineshloka
{धनदेन समस्त्यागे सत्ये धर्म इवापरः}
{तमेवङ्गुणसम्पन्नं रामं सत्यपराक्रमम्} %||1-1-18||

\twolineshloka
{ज्येष्ठं श्रेष्ठगुणैर्युक्तं प्रियं दशरथः सुतम्}
{यौवराज्येन संयोक्तुमैच्छत्प्रीत्या महीपतिः} %||1-1-19||

\threelineshloka
{तस्याभिषेकसम्भारान्दृष्ट्वा भार्याथ कैकयी}
{पूर्वं दत्तवरा देवी वरमेनमयाचत}
{विवासनं च रामस्य भरतस्याभिषेचनम्} %||1-1-20||

\twolineshloka
{स सत्यवचनाद्राजा धर्मपाशेन संयतः}
{विवासयामास सुतं रामं दशरथः प्रियम्} %||1-1-21||

\twolineshloka
{स जगाम वनं वीरः प्रतिज्ञामनुपालयन्}
{पितुर्वचननिर्देशात्कैकेय्याः प्रियकारणात्} %||1-1-22||

\twolineshloka
{तं व्रजन्तं प्रियो भ्राता लक्ष्मणोऽनुजगाम ह}
{स्नेहाद्विनयसम्पन्नः सुमित्रानन्दवर्धनः} %||1-1-23||

\twolineshloka
{सर्वलक्षणसम्पन्ना नारीणामुत्तमा वधूः}
{सीताप्यनुगता रामं शशिनं रोहिणी यथा} %||1-1-24||

\twolineshloka
{पौरैरनुगतो दूरं पित्रा दशरथेन च}
{शृङ्गवेरपुरे सूतं गङ्गाकूले व्यसर्जयत्} %||1-1-25||

\twolineshloka
{ते वनेन वनं गत्वा नदीस्तीर्त्वा बहूदकाः}
{चित्रकूटमनुप्राप्य भरद्वाजस्य शासनात्} %||1-1-26||

\twolineshloka
{रम्यमावसथं कृत्वा रममाणा वने त्रयः}
{देवगन्धर्वसङ्काशास्तत्र ते न्यवसन्सुखम्} %||1-1-27||

\twolineshloka
{चित्रकूटं गते रामे पुत्रशोकातुरस्तदा}
{राजा दशरथः स्वर्गं जगाम विलपन्सुतम्} %||1-1-28||

\threelineshloka
{मृते तु तस्मिन्भरतो वसिष्ठप्रमुखैर्द्विजैः}
{नियुज्यमानो राज्याय नैच्छद्राज्यं महाबलः}
{स जगाम वनं वीरो रामपादप्रसादकः} %||1-1-29||

\twolineshloka
{पादुके चास्य राज्याय न्यासं दत्त्वा पुनः पुनः}
{निवर्तयामास ततो भरतं भरताग्रजः} %||1-1-30||

\twolineshloka
{स काममनवाप्यैव रामपादावुपस्पृशन्}
{नन्दिग्रामेऽकरोद्राज्यं रामागमनकाङ्क्षया} %||1-1-31||

\twolineshloka
{रामस्तु पुनरालक्ष्य नागरस्य जनस्य च}
{तत्रागमनमेकाग्रे दण्डकान्प्रविवेश ह} %||1-1-32||

\twolineshloka
{विराधं राक्षसं हत्वा शरभङ्गं ददर्श ह}
{सुतीक्ष्णं चाप्यगस्त्यं च अगस्त्य भ्रातरं तथा} %||1-1-33||

\twolineshloka
{अगस्त्यवचनाच्चैव जग्राहैन्द्रं शरासनम्}
{खड्गं च परमप्रीतस्तूणी चाक्षयसायकौ} %||1-1-34||

\twolineshloka
{वसतस्तस्य रामस्य वने वनचरैः सह}
{ऋषयोऽभ्यागमन्सर्वे वधायासुररक्षसाम्} %||1-1-35||

\twolineshloka
{तेन तत्रैव वसता जनस्थाननिवासिनी}
{विरूपिता शूर्पणखा राक्षसी कामरूपिणी} %||1-1-36||

\twolineshloka
{ततः शूर्पणखावाक्यादुद्युक्तान्सर्वराक्षसान्}
{खरं त्रिशिरसं चैव दूषणं चैव राक्षसम्} %||1-1-37||

\twolineshloka
{निजघान रणे रामस्तेषां चैव पदानुगान्}
{रक्षसां निहतान्यासन्सहस्राणि चतुर्दश} %||1-1-38||

\twolineshloka
{ततो ज्ञातिवधं श्रुत्वा रावणः क्रोधमूर्छितः}
{सहायं वरयामास मारीचं नाम राक्षसम्} %||1-1-39||

\twolineshloka
{वार्यमाणः सुबहुशो मारीचेन स रावणः}
{न विरोधो बलवता क्षमो रावण तेन ते} %||1-1-40||

\twolineshloka
{अनादृत्य तु तद्वाक्यं रावणः कालचोदितः}
{जगाम सहमारीचस्तस्याश्रमपदं तदा} %||1-1-41||

\twolineshloka
{तेन मायाविना दूरमपवाह्य नृपात्मजौ}
{जहार भार्यां रामस्य गृध्रं हत्वा जटायुषम्} %||1-1-42||

\twolineshloka
{गृध्रं च निहतं दृष्ट्वा हृतां श्रुत्वा च मैथिलीम्}
{राघवः शोकसन्तप्तो विललापाकुलेन्द्रियः} %||1-1-43||

\twolineshloka
{ततस्तेनैव शोकेन गृध्रं दग्ध्वा जटायुषम्}
{मार्गमाणो वने सीतां राक्षसं सन्ददर्श ह} %||1-1-44||

\twolineshloka
{कबन्धं नाम रूपेण विकृतं घोरदर्शनम्}
{तं निहत्य महाबाहुर्ददाह स्वर्गतश्च सः} %||1-1-45||

\threelineshloka
{स चास्य कथयामास शबरीं धर्मचारिणीम्}
{श्रमणीं धर्मनिपुणामभिगच्छेति राघव}
{सोऽभ्यगच्छन्महातेजाः शबरीं शत्रुसूदनः} %||1-1-46||

\twolineshloka
{शबर्या पूजितः सम्यग्रामो दशरथात्मजः}
{पम्पातीरे हनुमता सङ्गतो वानरेण ह} %||1-1-47||

\twolineshloka
{हनुमद्वचनाच्चैव सुग्रीवेण समागतः}
{सुग्रीवाय च तत्सर्वं शंसद्रामो महाबलः} %||1-1-48||

\threelineshloka
{ततो वानरराजेन वैरानुकथनं प्रति}
{रामायावेदितं सर्वं प्रणयाद्दुःखितेन च}
{वालिनश्च बलं तत्र कथयामास वानरः} %||1-1-49||

\twolineshloka
{प्रतिज्ञातं च रामेण तदा वालिवधं प्रति}
{सुग्रीवः शङ्कितश्चासीन्नित्यं वीर्येण राघवे} %||1-1-50||

\twolineshloka
{राघवः प्रत्ययार्थं तु दुन्दुभेः कायमुत्तमम्}
{पादाङ्गुष्ठेन चिक्षेप सम्पूर्णं दशयोजनम्} %||1-1-51||

\twolineshloka
{बिभेद च पुनः सालान्सप्तैकेन महेषुणा}
{गिरिं रसातलं चैव जनयन्प्रत्ययं तदा} %||1-1-52||

\twolineshloka
{ततः प्रीतमनास्तेन विश्वस्तः स महाकपिः}
{किष्किन्धां रामसहितो जगाम च गुहां तदा} %||1-1-53||

\twolineshloka
{ततोऽगर्जद्धरिवरः सुग्रीवो हेमपिङ्गलः}
{तेन नादेन महता निर्जगाम हरीश्वरः} %||1-1-54||

\twolineshloka
{ततः सुग्रीववचनाद्धत्वा वालिनमाहवे}
{सुग्रीवमेव तद्राज्ये राघवः प्रत्यपादयत्} %||1-1-55||

\twolineshloka
{स च सर्वान्समानीय वानरान्वानरर्षभः}
{दिशः प्रस्थापयामास दिदृक्षुर्जनकात्मजाम्} %||1-1-56||

\twolineshloka
{ततो गृध्रस्य वचनात्सम्पातेर्हनुमान्बली}
{शतयोजनविस्तीर्णं पुप्लुवे लवणार्णवम्} %||1-1-57||

\twolineshloka
{तत्र लङ्कां समासाद्य पुरीं रावणपालिताम्}
{ददर्श सीतां ध्यायन्तीमशोकवनिकां गताम्} %||1-1-58||

\twolineshloka
{निवेदयित्वाभिज्ञानं प्रवृत्तिं च निवेद्य च}
{समाश्वास्य च वैदेहीं मर्दयामास तोरणम्} %||1-1-59||

\twolineshloka
{पञ्च सेनाग्रगान्हत्वा सप्त मन्त्रिसुतानपि}
{शूरमक्षं च निष्पिष्य ग्रहणं समुपागमत्} %||1-1-60||

\twolineshloka
{अस्त्रेणोन्मुहमात्मानं ज्ञात्वा पैतामहाद्वरात्}
{मर्षयन्राक्षसान्वीरो यन्त्रिणस्तान्यदृच्छया} %||1-1-61||

\twolineshloka
{ततो दग्ध्वा पुरीं लङ्कामृते सीतां च मैथिलीम्}
{रामाय प्रियमाख्यातुं पुनरायान्महाकपिः} %||1-1-62||

\twolineshloka
{सोऽभिगम्य महात्मानं कृत्वा रामं प्रदक्षिणम्}
{न्यवेदयदमेयात्मा दृष्टा सीतेति तत्त्वतः} %||1-1-63||

\twolineshloka
{ततः सुग्रीवसहितो गत्वा तीरं महोदधेः}
{समुद्रं क्षोभयामास शरैरादित्यसंनिभैः} %||1-1-64||

\twolineshloka
{दर्शयामास चात्मानं समुद्रः सरितां पतिः}
{समुद्रवचनाच्चैव नलं सेतुमकारयत्} %||1-1-65||

\twolineshloka
{तेन गत्वा पुरीं लङ्कां हत्वा रावणमाहवे}
{अभ्यषिञ्चत्स लङ्कायां राक्षसेन्द्रं विभीषणम्} %||1-1-66||

\twolineshloka
{कर्मणा तेन महता त्रैलोक्यं सचराचरम्}
{सदेवर्षिगणं तुष्टं राघवस्य महात्मनः} %||1-1-67||

\twolineshloka
{तथा परमसन्तुष्टैः पूजितः सर्वदैवतैः}
{कृतकृत्यस्तदा रामो विज्वरः प्रमुमोद ह} %||1-1-68||

\twolineshloka
{देवताभ्यो वरान्प्राप्य समुत्थाप्य च वानरान्}
{पुष्पकं तत्समारुह्य नन्दिग्रामं ययौ तदा} %||1-1-69||

\twolineshloka
{नन्दिग्रामे जटां हित्वा भ्रातृभिः सहितोऽनघः}
{रामः सीतामनुप्राप्य राज्यं पुनरवाप्तवान्} %||1-1-70||

\twolineshloka
{प्रहृष्टमुदितो लोकस्तुष्टः पुष्टः सुधार्मिकः}
{निरायमो अरोगश्च दुर्भिक्षभयवर्जितः} %||1-1-71||

\twolineshloka
{न पुत्रमरणं केचिद्द्रक्ष्यन्ति पुरुषाः क्वचित्}
{नार्यश्चाविधवा नित्यं भविष्यन्ति पतिव्रताः} %||1-1-72||

\twolineshloka
{न वातजं भयं किञ्चिन्नाप्सु मज्जन्ति जन्तवः}
{न चाग्रिजं भयं किञ्चिद्यथा कृतयुगे तथा} %||1-1-73||

\twolineshloka
{अश्वमेधशतैरिष्ट्वा तथा बहुसुवर्णकैः}
{गवां कोट्ययुतं दत्त्वा विद्वद्भ्यो विधिपूर्वकम्} %||1-1-74||

\twolineshloka
{राजवंशाञ्शतगुणान्स्थापयिष्यति राघवः}
{चातुर्वर्ण्यं च लोकेऽस्मिन्स्वे स्वे धर्मे नियोक्ष्यति} %||1-1-75||

\twolineshloka
{दशवर्षसहस्राणि दशवर्षशतानि च}
{रामो राज्यमुपासित्वा ब्रह्मलोकं गमिष्यति} %||1-1-76||

\twolineshloka
{इदं पवित्रं पापघ्नं पुण्यं वेदैश्च सम्मितम्}
{यः पठेद्रामचरितं सर्वपापैः प्रमुच्यते} %||1-1-77||

\twolineshloka
{एतदाख्यानमायुष्यं पठन्रामायणं नरः}
{सपुत्रपौत्रः सगणः प्रेत्य स्वर्गे महीयते} %||1-1-78||

\fourlineindentedshloka
{पठन्द्विजो वागृषभत्वमीयात्}
{ स्यात्क्षत्रियो भूमिपतित्वमीयात्}
{वणिग्जनः पण्यफलत्वमीयात्}
{ जनश्च शूद्रोऽपि महत्त्वमीयात्} %||1-2-79||


{॥इत्यार्षे श्रीमद्रामायणे वाल्मीकीये आदिकाव्ये बालकाण्डे प्रथमः सर्गः॥}

\sect{द्वितीयः सर्गः}

\twolineshloka
{नारदस्य तु तद्वाक्यं श्रुत्वा वाक्यविशारदः}
{पूजयामास धर्मात्मा सहशिष्यो महामुनिः} %||1-2-1||

\twolineshloka
{यथावत्पूजितस्तेन देवर्षिर्नारदस्तदा}
{आपृष्ट्वैवाभ्यनुज्ञातः स जगाम विहायसम्} %||1-2-2||

\twolineshloka
{स मुहूर्तं गते तस्मिन्देवलोकं मुनिस्तदा}
{जगाम तमसातीरं जाह्नव्यास्त्वविदूरतः} %||1-2-3||

\twolineshloka
{स तु तीरं समासाद्य तमसाया महामुनिः}
{शिष्यमाह स्थितं पार्श्वे दृष्ट्वा तीर्थमकर्दमम्} %||1-2-4||

\twolineshloka
{अकर्दममिदं तीर्थं भरद्वाज निशामय}
{रमणीयं प्रसन्नाम्बु सन्मनुष्यमनो यथा} %||1-2-5||

\twolineshloka
{न्यस्यतां कलशस्तात दीयतां वल्कलं मम}
{इदमेवावगाहिष्ये तमसातीर्थमुत्तमम्} %||1-2-6||

\twolineshloka
{एवमुक्तो भरद्वाजो वाल्मीकेन महात्मना}
{प्रायच्छत मुनेस्तस्य वल्कलं नियतो गुरोः} %||1-2-7||

\twolineshloka
{स शिष्यहस्तादादाय वल्कलं नियतेन्द्रियः}
{विचचार ह पश्यंस्तत्सर्वतो विपुलं वनम्} %||1-2-8||

\twolineshloka
{तस्याभ्याशे तु मिथुनं चरन्तमनपायिनम्}
{ददर्श भगवांस्तत्र क्रौञ्चयोश्चारुनिःस्वनम्} %||1-2-9||

\twolineshloka
{तस्मात्तु मिथुनादेकं पुमांसं पापनिश्चयः}
{जघान वैरनिलयो निषादस्तस्य पश्यतः} %||1-2-10||

\twolineshloka
{तं शोणितपरीताङ्गं वेष्टमानं महीतले}
{भार्या तु निहतं दृष्ट्वा रुराव करुणां गिरम्} %||1-2-11||

\twolineshloka
{तथा तु तं द्विजं दृष्ट्वा निषादेन निपातितम्}
{ऋषेर्धर्मात्मनस्तस्य कारुण्यं समपद्यत} %||1-2-12||

\twolineshloka
{ततः करुणवेदित्वादधर्मोऽयमिति द्विजः}
{निशाम्य रुदतीं क्रौञ्चीमिदं वचनमब्रवीत्} %||1-2-13||

\twolineshloka
{मा निषाद प्रतिष्ठां त्वमगमः शाश्वतीः समाः}
{यत्क्रौञ्चमिथुनादेकमवधीः काममोहितम्} %||1-2-14||

\twolineshloka
{तस्यैवं ब्रुवतश्चिन्ता बभूव हृदि वीक्षतः}
{शोकार्तेनास्य शकुनेः किमिदं व्याहृतं मया} %||1-2-15||

\twolineshloka
{चिन्तयन्स महाप्राज्ञश्चकार मतिमान्मतिम्}
{शिष्यं चैवाब्रवीद्वाक्यमिदं स मुनिपुङ्गवः} %||1-2-16||

\twolineshloka
{पादबद्धोऽक्षरसमस्तन्त्रीलयसमन्वितः}
{शोकार्तस्य प्रवृत्तो मे श्लोको भवतु नान्यथा} %||1-2-17||

\twolineshloka
{शिष्यस्तु तस्य ब्रुवतो मुनेर्वाक्यमनुत्तमम्}
{प्रतिजग्राह संहृष्टस्तस्य तुष्टोऽभवद्गुरुः} %||1-2-18||

\twolineshloka
{सोऽभिषेकं ततः कृत्वा तीर्थे तस्मिन्यथाविधि}
{तमेव चिन्तयन्नर्थमुपावर्तत वै मुनिः} %||1-2-19||

\twolineshloka
{भरद्वाजस्ततः शिष्यो विनीतः श्रुतवान्गुरोः}
{कलशं पूर्णमादाय पृष्ठतोऽनुजगाम ह} %||1-2-20||

\twolineshloka
{स प्रविश्याश्रमपदं शिष्येण सह धर्मवित्}
{उपविष्टः कथाश्चान्याश्चकार ध्यानमास्थितः} %||1-2-21||

\twolineshloka
{आजगाम ततो ब्रह्मा लोककर्ता स्वयं प्रभुः}
{चतुर्मुखो महातेजा द्रष्टुं तं मुनिपुङ्गवम्} %||1-2-22||

\twolineshloka
{वाल्मीकिरथ तं दृष्ट्वा सहसोत्थाय वाग्यतः}
{प्राञ्जलिः प्रयतो भूत्वा तस्थौ परमविस्मितः} %||1-2-23||

\twolineshloka
{पूजयामास तं देवं पाद्यार्घ्यासनवन्दनैः}
{प्रणम्य विधिवच्चैनं पृष्ट्वानामयमव्ययम्} %||1-2-24||

\twolineshloka
{अथोपविश्य भगवानासने परमार्चिते}
{वाल्मीकये महर्षये सन्दिदेशासनं ततः} %||1-2-25||

\twolineshloka
{उपविष्टे तदा तस्मिन्साक्षाल्लोकपितामहे}
{तद्गतेनैव मनसा वाल्मीकिर्ध्यानमास्थितः} %||1-2-26||

\twolineshloka
{पापात्मना कृतं कष्टं वैरग्रहणबुद्धिना}
{यस्तादृशं चारुरवं क्रौञ्चं हन्यादकारणात्} %||1-2-27||

\twolineshloka
{शोचन्नेव मुहुः क्रौञ्चीमुपश्लोकमिमं पुनः}
{जगावन्तर्गतमना भूत्वा शोकपरायणः} %||1-2-28||

\twolineshloka
{तमुवाच ततो ब्रह्मा प्रहसन्मुनिपुङ्गवम्}
{श्लोक एव त्वया बद्धो नात्र कार्या विचारणा} %||1-2-29||

\twolineshloka
{मच्छन्दादेव ते ब्रह्मन्प्रवृत्तेयं सरस्वती}
{रामस्य चरितं कृत्स्नं कुरु त्वमृषिसत्तम} %||1-2-30||

\twolineshloka
{धर्मात्मनो गुणवतो लोके रामस्य धीमतः}
{वृत्तं कथय धीरस्य यथा ते नारदाच्छ्रुतम्} %||1-2-31||

\twolineshloka
{रहस्यं च प्रकाशं च यद्वृत्तं तस्य धीमतः}
{रामस्य सह सौमित्रे राक्षसानां च सर्वशः} %||1-2-32||

\twolineshloka
{वैदेह्याश्चैव यद्वृत्तं प्रकाशं यदि वा रहः}
{तच्चाप्यविदितं सर्वं विदितं ते भविष्यति} %||1-2-33||

\twolineshloka
{न ते वागनृता काव्ये काचिदत्र भविष्यति}
{कुरु रामकथां पुण्यां श्लोकबद्धां मनोरमाम्} %||1-2-34||

\twolineshloka
{यावत्स्थास्यन्ति गिरयः सरितश्च महीतले}
{तावद्रामायणकथा लोकेषु प्रचरिष्यति} %||1-2-35||

\twolineshloka
{यावद्रामस्य च कथा त्वत्कृता प्रचरिष्यति}
{तावदूर्ध्वमधश्च त्वं मल्लोकेषु निवत्स्यसि} %||1-2-36||

\twolineshloka
{इत्युक्त्वा भगवान्ब्रह्मा तत्रैवान्तरधीयत}
{ततः सशिष्यो वाल्मीकिर्मुनिर्विस्मयमाययौ} %||1-2-37||

\twolineshloka
{तस्य शिष्यास्ततः सर्वे जगुः श्लोकमिमं पुनः}
{मुहुर्मुहुः प्रीयमाणाः प्राहुश्च भृशविस्मिताः} %||1-2-38||

\twolineshloka
{समाक्षरैश्चतुर्भिर्यः पादैर्गीतो महर्षिणा}
{सोऽनुव्याहरणाद्भूयः शोकः श्लोकत्वमागतः} %||1-2-39||

\twolineshloka
{तस्य बुद्धिरियं जाता वाल्मीकेर्भावितात्मनः}
{कृत्स्नं रामायणं काव्यमीदृशैः करवाण्यहम्} %||1-2-40||

\fourlineindentedshloka
{उदारवृत्तार्थपदैर्मनोरमै}
{स्तदास्य रामस्य चकार कीर्तिमान्}
{समाक्षरैः श्लोकशतैर्यशस्विनो}
{ यशस्करं काव्यमुदारधीर्मुनिः} %||1-3-41||


{॥इत्यार्षे श्रीमद्रामायणे वाल्मीकीये आदिकाव्ये बालकाण्डे द्वितीयः सर्गः॥}

\sect{तृतीयः सर्गः}

\twolineshloka
{श्रुत्वा वस्तु समग्रं तद्धर्मात्मा धर्मसंहितम्}
{व्यक्तमन्वेषते भूयो यद्वृत्तं तस्य धीमतः} %||1-3-1||

\twolineshloka
{उपस्पृश्योदकं संयन्मुनिः स्थित्वा कृताञ्जलिः}
{प्राचीनाग्रेषु दर्भेषु धर्मेणान्वेषते गतिम्} %||1-3-2||

\twolineshloka
{जन्म रामस्य सुमहद्वीर्यं सर्वानुकूलताम्}
{लोकस्य प्रियतां क्षान्तिं सौम्यतां सत्यशीलताम्} %||1-3-3||

\twolineshloka
{नानाचित्राः कथाश्चान्या विश्वामित्रसहायने}
{जानक्याश्च विवाहं च धनुषश्च विभेदनम्} %||1-3-4||

\twolineshloka
{रामरामविवादं च गुणान्दाशरथेस्तथा}
{तथाभिषेकं रामस्य कैकेय्या दुष्टभावताम्} %||1-3-5||

\twolineshloka
{व्याघातं चाभिषेकस्य रामस्य च विवासनम्}
{राज्ञः शोकं विलापं च परलोकस्य चाश्रयम्} %||1-3-6||

\twolineshloka
{प्रकृतीनां विषादं च प्रकृतीनां विसर्जनम्}
{निषादाधिपसंवादं सूतोपावर्तनं तथा} %||1-3-7||

\twolineshloka
{गङ्गायाश्चाभिसन्तारं भरद्वाजस्य दर्शनम्}
{भरद्वाजाभ्यनुज्ञानाच्चित्रकूटस्य दर्शनम्} %||1-3-8||

\twolineshloka
{वास्तुकर्मनिवेशं च भरतागमनं तथा}
{प्रसादनं च रामस्य पितुश्च सलिलक्रियाम्} %||1-3-9||

\twolineshloka
{पादुकाग्र्याभिषेकं च नन्दिग्राम निवासनम्}
{दण्डकारण्यगमनं सुतीक्ष्णेन समागमम्} %||1-3-10||

\twolineshloka
{अनसूयासमस्यां च अङ्गरागस्य चार्पणम्}
{शूर्पणख्याश्च संवादं विरूपकरणं तथा} %||1-3-11||

\twolineshloka
{वधं खरत्रिशिरसोरुत्थानं रावणस्य च}
{मारीचस्य वधं चैव वैदेह्या हरणं तथा} %||1-3-12||

\twolineshloka
{राघवस्य विलापं च गृध्रराजनिबर्हणम्}
{कबन्धदर्शनं चैव पम्पायाश्चापि दर्शनम्} %||1-3-13||

\twolineshloka
{शर्बर्या दर्शनं चैव हनूमद्दर्शनं तथा}
{विलापं चैव पम्पायां राघवस्य महात्मनः} %||1-3-14||

\twolineshloka
{ऋष्यमूकस्य गमनं सुग्रीवेण समागमम्}
{प्रत्ययोत्पादनं सख्यं वालिसुग्रीवविग्रहम्} %||1-3-15||

\twolineshloka
{वालिप्रमथनं चैव सुग्रीवप्रतिपादनम्}
{ताराविलापसमयं वर्षरात्रिनिवासनम्} %||1-3-16||

\twolineshloka
{कोपं राघवसिंहस्य बलानामुपसङ्ग्रहम्}
{दिशः प्रस्थापनं चैव पृथिव्याश्च निवेदनम्} %||1-3-17||

\twolineshloka
{अङ्गुलीयकदानं च ऋक्षस्य बिलदर्शनम्}
{प्रायोपवेशनं चैव सम्पातेश्चापि दर्शनम्} %||1-3-18||

\twolineshloka
{पर्वतारोहणं चैव सागरस्य च लङ्घनम्}
{रात्रौ लङ्काप्रवेशं च एकस्यापि विचिन्तनम्} %||1-3-19||

\twolineshloka
{आपानभूमिगमनमवरोधस्य दर्शनम्}
{अशोकवनिकायानं सीतायाश्चापि दर्शनम्} %||1-3-20||

\twolineshloka
{अभिज्ञानप्रदानं च सीतायाश्चापि भाषणम्}
{राक्षसीतर्जनं चैव त्रिजटास्वप्नदर्शनम्} %||1-3-21||

\twolineshloka
{मणिप्रदानं सीताया वृक्षभङ्गं तथैव च}
{राक्षसीविद्रवं चैव किङ्कराणां निबर्हणम्} %||1-3-22||

\twolineshloka
{ग्रहणं वायुसूनोश्च लङ्कादाहाभिगर्जनम्}
{प्रतिप्लवनमेवाथ मधूनां हरणं तथा} %||1-3-23||

\twolineshloka
{राघवाश्वासनं चैव मणिनिर्यातनं तथा}
{सङ्गमं च समुद्रस्य नलसेतोश्च बन्धनम्} %||1-3-24||

\twolineshloka
{प्रतारं च समुद्रस्य रात्रौ लङ्कावरोधनम्}
{विभीषणेन संसर्गं वधोपायनिवेदनम्} %||1-3-25||

\twolineshloka
{कुम्भकर्णस्य निधनं मेघनादनिबर्हणम्}
{रावणस्य विनाशं च सीतावाप्तिमरेः पुरे} %||1-3-26||

\twolineshloka
{बिभीषणाभिषेकं च पुष्पकस्य च दर्शनम्}
{अयोध्यायाश्च गमनं भरतेन समागमम्} %||1-3-27||

\twolineshloka
{रामाभिषेकाभ्युदयं सर्वसैन्यविसर्जनम्}
{स्वराष्ट्ररञ्जनं चैव वैदेह्याश्च विसर्जनम्} %||1-3-28||

\twolineshloka
{अनागतं च यत्किञ्चिद्रामस्य वसुधातले}
{तच्चकारोत्तरे काव्ये वाल्मीकिर्भगवानृषिः} %||1-4-29||


{॥इत्यार्षे श्रीमद्रामायणे वाल्मीकीये आदिकाव्ये बालकाण्डे तृतीयः सर्गः॥}

\sect{चतुर्थः सर्गः}

\twolineshloka
{प्राप्तराज्यस्य रामस्य वाल्मीकिर्भगवानृषिः}
{चकार चरितं कृत्स्नं विचित्रपदमात्मवान्} %||1-4-1||

\twolineshloka
{कृत्वा तु तन्महाप्राज्ञः सभविष्यं सहोत्तरम्}
{चिन्तयामास को न्वेतत्प्रयुञ्जीयादिति प्रभुः} %||1-4-2||

\twolineshloka
{तस्य चिन्तयमानस्य महर्षेर्भावितात्मनः}
{अगृह्णीतां ततः पादौ मुनिवेषौ कुशीलवौ} %||1-4-3||

\twolineshloka
{कुशीलवौ तु धर्मज्ञौ राजपुत्रौ यशस्विनौ}
{भ्रातरौ स्वरसम्पन्नौ ददर्शाश्रमवासिनौ} %||1-4-4||

\twolineshloka
{स तु मेधाविनौ दृष्ट्वा वेदेषु परिनिष्ठितौ}
{वेदोपबृह्मणार्थाय तावग्राहयत प्रभुः} %||1-4-5||

\twolineshloka
{काव्यं रामायणं कृत्स्नं सीतायाश्चरितं महत्}
{पौलस्त्य वधमित्येव चकार चरितव्रतः} %||1-4-6||

\twolineshloka
{पाठ्ये गेये च मधुरं प्रमाणैस्त्रिभिरन्वितम्}
{जातिभिः सप्तभिर्युक्तं तन्त्रीलयसमन्वितम्} %||1-4-7||

\twolineshloka
{हास्यशृङ्गारकारुण्यरौद्रवीरभयानकैः}
{बीभत्सादिरसैर्युक्तं काव्यमेतदगायताम्} %||1-4-8||

\twolineshloka
{तौ तु गान्धर्वतत्त्वज्ञौ स्थानमूर्च्छनकोविदौ}
{भ्रातरौ स्वरसम्पन्नौ गन्धर्वाविव रूपिणौ} %||1-4-9||

\twolineshloka
{रूपलक्षणसम्पन्नौ मधुरस्वरभाषिणौ}
{बिम्बादिवोद्धृतौ बिम्बौ रामदेहात्तथापरौ} %||1-4-10||

\twolineshloka
{तौ राजपुत्रौ कार्त्स्न्येन धर्म्यमाख्यानमुत्तमम्}
{वाचो विधेयं तत्सर्वं कृत्वा काव्यमनिन्दितौ} %||1-4-11||

\threelineshloka
{ऋषीणां च द्विजातीनां साधूनां च समागमे}
{यथोपदेशं तत्त्वज्ञौ जगतुस्तौ समाहितौ}
{महात्मानौ महाभागौ सर्वलक्षणलक्षितौ} %||1-4-12||

\twolineshloka
{तौ कदाचित्समेतानामृषीणां भावितात्मनाम्}
{आसीनानां समीपस्थाविदं काव्यमगायताम्} %||1-4-13||

\twolineshloka
{तच्छ्रुत्वा मुनयः सर्वे बाष्पपर्याकुलेक्षणाः}
{साधु साध्वित्य्तावूचतुः परं विस्मयमागताः} %||1-4-14||

\twolineshloka
{ते प्रीतमनसः सर्वे मुनयो धर्मवत्सलाः}
{प्रशशंसुः प्रशस्तव्यौ गायमानौ कुशीलवौ} %||1-4-15||

\twolineshloka
{अहो गीतस्य माधुर्यं श्लोकानां च विशेषतः}
{चिरनिर्वृत्तमप्येतत्प्रत्यक्षमिव दर्शितम्} %||1-4-16||

\twolineshloka
{प्रविश्य तावुभौ सुष्ठु तदा भावमगायताम्}
{सहितौ मधुरं रक्तं सम्पन्नं स्वरसम्पदा} %||1-4-17||

\twolineshloka
{एवं प्रशस्यमानौ तौ तपःश्लाघ्यैर्महर्षिभिः}
{संरक्ततरमत्यर्थं मधुरं तावगायताम्} %||1-4-18||

\twolineshloka
{प्रीतः कश्चिन्मुनिस्ताभ्यां संस्थितः कलशं ददौ}
{प्रसन्नो वल्कलं कश्चिद्ददौ ताभ्यां महायशाः} %||1-4-19||

\twolineshloka
{आश्चर्यमिदमाख्यानं मुनिना सम्प्रकीर्तितम्}
{परं कवीनामाधारं समाप्तं च यथाक्रमम्} %||1-4-20||

\twolineshloka
{प्रशस्यमानौ सर्वत्र कदाचित्तत्र गायकौ}
{रथ्यासु राजमार्गेषु ददर्श भरताग्रजः} %||1-4-21||

\twolineshloka
{स्ववेश्म चानीय ततो भ्रातरौ सकुशीलवौ}
{पूजयामास पूजार्हौ रामः शत्रुनिबर्हणः} %||1-4-22||

\twolineshloka
{आसीनः काञ्चने दिव्ये स च सिंहासने प्रभुः}
{उपोपविष्टैः सचिवैर्भ्रातृभिश्च परन्तपः} %||1-4-23||

\twolineshloka
{दृष्ट्वा तु रूपसम्पन्नौ तावुभौ वीणिनौ ततः}
{उवाच लक्ष्मणं रामः शत्रुघ्नं भरतं तथा} %||1-4-24||

\twolineshloka
{श्रूयतामिदमाख्यानमनयोर्देववर्चसोः}
{विचित्रार्थपदं सम्यग्गायतोर्मधुरस्वरम्} %||1-4-25||

\fourlineindentedshloka
{इमौ मुनी पार्थिवलक्ष्मणान्वितौ}
{ कुशीलवौ चैव महातपस्विनौ}
{ममापि तद्भूतिकरं प्रचक्षते}
{ महानुभावं चरितं निबोधत} %||1-4-26||

\fourlineindentedshloka
{ततस्तु तौ रामवचः प्रचोदिता}
{वगायतां मार्गविधानसम्पदा}
{स चापि रामः परिषद्गतः शनैर्-}
{ बुभूषयासक्तमना बभूव} %||1-5-27||


{॥इत्यार्षे श्रीमद्रामायणे वाल्मीकीये आदिकाव्ये बालकाण्डे चतुर्थः सर्गः॥}

\sect{पञ्चमः सर्गः}

\twolineshloka
{सर्वापूर्वमियं येषामासीत्कृत्स्ना वसुन्धरा}
{प्रजापतिमुपादाय नृपाणां जयशालिनाम्} %||1-5-1||

\twolineshloka
{येषां स सगरो नाम सागरो येन खानितः}
{षष्टिः पुत्रसहस्राणि यं यान्तं पर्यवारयन्} %||1-5-2||

\twolineshloka
{इक्ष्वाकूणामिदं तेषां राज्ञां वंशे महात्मनाम्}
{महदुत्पन्नमाख्यानं रामायणमिति श्रुतम्} %||1-5-3||

\twolineshloka
{तदिदं वर्तयिष्यामि सर्वं निखिलमादितः}
{धर्मकामार्थसहितं श्रोतव्यमनसूयया} %||1-5-4||

\twolineshloka
{कोसलो नाम मुदितः स्फीतो जनपदो महान्}
{निविष्टः सरयूतीरे प्रभूतधनधान्यवान्} %||1-5-5||

\twolineshloka
{अयोध्या नाम नगरी तत्रासील्लोकविश्रुता}
{मनुना मानवेन्द्रेण या पुरी निर्मिता स्वयम्} %||1-5-6||

\twolineshloka
{आयता दश च द्वे च योजनानि महापुरी}
{श्रीमती त्रीणि विस्तीर्णा सुविभक्तमहापथा} %||1-5-7||

\twolineshloka
{राजमार्गेण महता सुविभक्तेन शोभिता}
{मुक्तपुष्पावकीर्णेन जलसिक्तेन नित्यशः} %||1-5-8||

\twolineshloka
{तां तु राजा दशरथो महाराष्ट्रविवर्धनः}
{पुरीमावासयामास दिवि देवपतिर्यथा} %||1-5-9||

\twolineshloka
{कपाटतोरणवतीं सुविभक्तान्तरापणाम्}
{सर्वयन्त्रायुधवतीमुपेतां सर्वशिल्पिभिः} %||1-5-10||

\twolineshloka
{सूतमागधसम्बाधां श्रीमतीमतुलप्रभाम्}
{उच्चाट्टालध्वजवतीं शतघ्नीशतसङ्कुलाम्} %||1-5-11||

\twolineshloka
{वधूनाटकसङ्घैश्च संयुक्तां सर्वतः पुरीम्}
{उद्यानाम्रवणोपेतां महतीं सालमेखलाम्} %||1-5-12||

\twolineshloka
{दुर्गगम्भीरपरिघां दुर्गामन्यैर्दुरासदाम्}
{वाजिवारणसम्पूर्णां गोभिरुष्ट्रैः खरैस्तथा} %||1-5-13||

\twolineshloka
{सामन्तराजसङ्घैश्च बलिकर्मभिरावृताम्}
{नानादेशनिवासैश्च वणिग्भिरुपशोभिताम्} %||1-5-14||

\twolineshloka
{प्रसादै रत्नविकृतैः पर्वतैरुपशोभिताम्}
{कूटागारैश्च सम्पूर्णामिन्द्रस्येवामरावतीम्} %||1-5-15||

\twolineshloka
{चित्रामष्टापदाकारां वरनारीगणैर्युताम्}
{सर्वरत्नसमाकीर्णां विमानगृहशोभिताम्} %||1-5-16||

\twolineshloka
{गृहगाढामविच्छिद्रां समभूमौ निवेशिताम्}
{शालितण्डुलसम्पूर्णामिक्षुकाण्डरसोदकाम्} %||1-5-17||

\twolineshloka
{दुन्दुभीभिर्मृदङ्गैश्च वीणाभिः पणवैस्तथा}
{नादितां भृशमत्यर्थं पृथिव्यां तामनुत्तमाम्} %||1-5-18||

\twolineshloka
{विमानमिव सिद्धानां तपसाधिगतं दिवि}
{सुनिवेशितवेश्मान्तां नरोत्तमसमावृताम्} %||1-5-19||

\twolineshloka
{ये च बाणैर्न विध्यन्ति विविक्तमपरापरम्}
{शब्दवेध्यं च विततं लघुहस्ता विशारदाः} %||1-5-20||

\twolineshloka
{सिंहव्याघ्रवराहाणां मत्तानां नदतां वने}
{हन्तारो निशितैः शस्त्रैर्बलाद्बाहुबलैरपि} %||1-5-21||

\twolineshloka
{तादृशानां सहस्रैस्तामभिपूर्णां महारथैः}
{पुरीमावासयामास राजा दशरथस्तदा} %||1-5-22||

\fourlineindentedshloka
{तामग्निमद्भिर्गुणवद्भिरावृताम्}
{ द्विजोत्तमैर्वेदषडङ्गपारगैः}
{सहस्रदैः सत्यरतैर्महात्मभिर्-}
{ महर्षिकल्पैरृषिभिश्च केवलैः} %||1-6-23||


{॥इत्यार्षे श्रीमद्रामायणे वाल्मीकीये आदिकाव्ये बालकाण्डे पञ्चमः सर्गः॥}

\sect{षष्ठः सर्गः}

\twolineshloka
{पुर्यां तस्यामयोध्यायां वेदवित्सर्वसङ्ग्रहः}
{दीर्घदर्शी महातेजाः पौरजानपदप्रियः} %||1-6-1||

\twolineshloka
{इक्ष्वाकूणामतिरथो यज्वा धर्मरतो वशी}
{महर्षिकल्पो राजर्षिस्त्रिषु लोकृषु विश्रुतः} %||1-6-2||

\twolineshloka
{बलवान्निहतामित्रो मित्रवान्विजितेन्द्रियः}
{धनैश्च सञ्चयैश्चान्यैः शक्रवैश्रवणोपमः} %||1-6-3||

\twolineshloka
{यथा मनुर्महातेजा लोकस्य परिरक्षिता}
{तथा दशरथो राजा वसञ्जगदपालयत्} %||1-6-4||

\twolineshloka
{तेन सत्याभिसन्धेन त्रिवर्गमनुतिष्ठता}
{पालिता सा पुरी श्रेष्ठेन्द्रेण इवामरावती} %||1-6-5||

\twolineshloka
{तस्मिन्पुरवरे हृष्टा धर्मात्मना बहुश्रुताः}
{नरास्तुष्टाधनैः स्वैः स्वैरलुब्धाः सत्यवादिनः} %||1-6-6||

\twolineshloka
{नाल्पसंनिचयः कश्चिदासीत्तस्मिन्पुरोत्तमे}
{कुटुम्बी यो ह्यसिद्धार्थोऽगवाश्वधनधान्यवान्} %||1-6-7||

\twolineshloka
{कामी वा न कदर्यो वा नृशंसः पुरुषः क्वचित्}
{द्रष्टुं शक्यमयोध्यायां नाविद्वान्न च नास्तिकः} %||1-6-8||

\twolineshloka
{सर्वे नराश्च नार्यश्च धर्मशीलाः सुसंयताः}
{मुदिताः शीलवृत्ताभ्यां महर्षय इवामलाः} %||1-6-9||

\twolineshloka
{नाकुण्डली नामुकुटी नास्रग्वी नाल्पभोगवान्}
{नामृष्टो नानुलिप्ताङ्गो नासुगन्धश्च विद्यते} %||1-6-10||

\twolineshloka
{नामृष्टभोजी नादाता नाप्यनङ्गदनिष्कधृक्}
{नाहस्ताभरणो वापि दृश्यते नाप्यनात्मवान्} %||1-6-11||

\twolineshloka
{नानाहिताग्निर्नायज्वा विप्रो नाप्यसहस्रदः}
{कश्चिदासीदयोध्यायां न च निर्वृत्तसङ्करः} %||1-6-12||

\twolineshloka
{स्वकर्मनिरता नित्यं ब्राह्मणा विजितेन्द्रियाः}
{दानाध्ययनशीलाश्च संयताश्च प्रतिग्रहे} %||1-6-13||

\twolineshloka
{न नास्तिको नानृतको न कश्चिदबहुश्रुतः}
{नासूयको न चाशक्तो नाविद्वान्विद्यते तदा} %||1-6-14||

\threelineshloka
{न दीनः क्षिप्तचित्तो वा व्यथितो वापि कश्चन}
{कश्चिन्नरो वा नारी वा नाश्रीमान्नाप्यरूपवान्}
{द्रष्टुं शक्यमयोध्यायां नापि राजन्यभक्तिमान्} %||1-6-15||

\twolineshloka
{वर्णेष्वग्र्यचतुर्थेषु देवतातिथिपूजकाः}
{दीर्घायुषो नराः सर्वे धर्मं सत्यं च संश्रिताः} %||1-6-16||

\twolineshloka
{क्षत्रं ब्रह्ममुखं चासीद्वैश्याः क्षत्रमनुव्रताः}
{शूद्राः स्वधर्मनिरतास्त्रीन्वर्णानुपचारिणः} %||1-6-17||

\twolineshloka
{सा तेनेक्ष्वाकुनाथेन पुरी सुपरिरक्षिता}
{यथा पुरस्तान्मनुना मानवेन्द्रेण धीमता} %||1-6-18||

\twolineshloka
{योधानामग्निकल्पानां पेशलानाममर्षिणाम्}
{सम्पूर्णाकृतविद्यानां गुहाकेसरिणामिव} %||1-6-19||

\twolineshloka
{काम्बोजविषये जातैर्बाह्लीकैश्च हयोत्तमैः}
{वनायुजैर्नदीजैश्च पूर्णाहरिहयोपमैः} %||1-6-20||

\twolineshloka
{विन्ध्यपर्वतजैर्मत्तैः पूर्णा हैमवतैरपि}
{मदान्वितैरतिबलैर्मातङ्गैः पर्वतोपमैः} %||1-6-21||

\twolineshloka
{अञ्जनादपि निष्क्रान्तैर्वामनादपि च द्विपैः}
{भद्रमन्द्रैर्भद्रमृगैर्मृगमन्द्रैश्च सा पुरी} %||1-6-22||

\twolineshloka
{नित्यमत्तैः सदा पूर्णा नागैरचलसंनिभैः}
{सा योजने च द्वे भूयः सत्यनामा प्रकाशते} %||1-6-23||

\fourlineindentedshloka
{तां सत्यनामां दृढतोरणार्गला}
{म्गृहैर्विचित्रैरुपशोभितां शिवाम्}
{पुरीमयोध्यां नृसहस्रसङ्कुलाम्}
{ शशास वै शक्रसमो महीपतिः} %||1-7-24||


{॥इत्यार्षे श्रीमद्रामायणे वाल्मीकीये आदिकाव्ये बालकाण्डे षष्ठः सर्गः॥}

\sect{सप्तमः सर्गः}

\twolineshloka
{अष्टौ बभूवुर्वीरस्य तस्यामात्या यशस्विनः}
{शुचयश्चानुरक्ताश्च राजकृत्येषु नित्यशः} %||1-7-1||

\twolineshloka
{धृष्टिर्जयन्तो विजयः सिद्धार्थो अर्थसाधकः}
{अशोको मन्त्रपालश्च सुमन्त्रश्चाष्टमोऽभवत्} %||1-7-2||

\twolineshloka
{ऋत्विजौ द्वावभिमतौ तस्यास्तामृषिसत्तमौ}
{वसिष्ठो वामदेवश्च मन्त्रिणश्च तथापरे} %||1-7-3||

\twolineshloka
{श्रीमन्तश्च महात्मानः शास्त्रज्ञा दृढविक्रमाः}
{कीर्तिमन्तः प्रणिहिता यथावचनकारिणः} %||1-7-4||

\twolineshloka
{तेजःक्षमायशःप्राप्ताः स्मितपूर्वाभिभाषिणः}
{क्रोधात्कामार्थहेतोर्वा न ब्रूयुरनृतं वचः} %||1-7-5||

\twolineshloka
{तेषामविदितं किञ्चित्स्वेषु नास्ति परेषु वा}
{क्रियमाणं कृतं वापि चारेणापि चिकीर्षितम्} %||1-7-6||

\twolineshloka
{कुशला व्यवहारेषु सौहृदेषु परीक्षिताः}
{प्राप्तकालं यथा दण्डं धारयेयुः सुतेष्वपि} %||1-7-7||

\twolineshloka
{कोशसङ्ग्रहणे युक्ता बलस्य च परिग्रहे}
{अहितं चापि पुरुषं न विहिंस्युरदूषकम्} %||1-7-8||

\twolineshloka
{वीराश्च नियतोत्साहा राजशास्त्रमनुष्ठिताः}
{शुचीनां रक्षितारश्च नित्यं विषयवासिनाम्} %||1-7-9||

\twolineshloka
{ब्रह्मक्षत्रमहिंसन्तस्ते कोशं समपूरयन्}
{सुतीक्ष्णदण्डाः सम्प्रेक्ष्य पुरुषस्य बलाबलम्} %||1-7-10||

\twolineshloka
{शुचीनामेकबुद्धीनां सर्वेषां सम्प्रजानताम्}
{नासीत्पुरे वा राष्ट्रे वा मृषावादी नरः क्वचित्} %||1-7-11||

\twolineshloka
{कश्चिन्न दुष्टस्तत्रासीत्परदाररतिर्नरः}
{प्रशान्तं सर्वमेवासीद्राष्ट्रं पुरवरं च तत्} %||1-7-12||

\twolineshloka
{सुवाससः सुवेशाश्च ते च सर्वे सुशीलिनः}
{हितार्थं च नरेन्द्रस्य जाग्रतो नयचक्षुषा} %||1-7-13||

\twolineshloka
{गुरौ गुणगृहीताश्च प्रख्याताश्च पराक्रमैः}
{विदेशेष्वपि विज्ञाताः सर्वतो बुद्धिनिश्चयात्} %||1-7-14||

\twolineshloka
{ईदृशैस्तैरमात्यैस्तु राजा दशरथोऽनघः}
{उपपन्नो गुणोपेतैरन्वशासद्वसुन्धराम्} %||1-7-15||

\twolineshloka
{अवेक्षमाणश्चारेण प्रजा धर्मेण रञ्जयन्}
{नाध्यगच्छद्विशिष्टं वा तुल्यं वा शत्रुमात्मनः} %||1-7-16||

\fourlineindentedshloka
{तैर्मन्त्रिभिर्मन्त्रहितैर्निविष्टैर्-}
{ वृतोऽनुरक्तैः कुशलैः समर्थैः}
{स पार्थिवो दीप्तिमवाप युक्त}
{स्तेजोमयैर्गोभिरिवोदितोऽर्कः} %||1-8-17||


{॥इत्यार्षे श्रीमद्रामायणे वाल्मीकीये आदिकाव्ये बालकाण्डे सप्तमः सर्गः॥}

\sect{अष्टमः सर्गः}

\twolineshloka
{तस्य त्वेवं प्रभावस्य धर्मज्ञस्य महात्मनः}
{सुतार्थं तप्यमानस्य नासीद्वंशकरः सुतः} %||1-8-1||

\twolineshloka
{चिन्तयानस्य तस्यैवं बुद्धिरासीन्महात्मनः}
{सुतार्थं वाजिमेधेन किमर्थं न यजाम्यहम्} %||1-8-2||

\twolineshloka
{स निश्चितां मतिं कृत्वा यष्टव्यमिति बुद्धिमान्}
{मन्त्रिभिः सह धर्मात्मा सर्वैरेव कृतात्मभिः} %||1-8-3||

\twolineshloka
{ततोऽब्रवीदिदं राजा सुमन्त्रं मन्त्रिसत्तमम्}
{शीघ्रमानय मे सर्वान्गुरूंस्तान्सपुरोहितान्} %||1-8-4||

\twolineshloka
{एतच्छ्रुत्वा रहः सूतो राजानमिदमब्रवीत्}
{ऋत्विग्भिरुपदिष्टोऽयं पुरावृत्तो मया श्रुतः} %||1-8-5||

\twolineshloka
{सनत्कुमारो भगवान्पूर्वं कथितवान्कथाम्}
{ऋषीणां संनिधौ राजंस्तव पुत्रागमं प्रति} %||1-8-6||

\twolineshloka
{काश्यपस्य तु पुत्रोऽस्ति विभाण्डक इति श्रुतः}
{ऋष्यशृङ्ग इति ख्यातस्तस्य पुत्रो भविष्यति} %||1-8-7||

\twolineshloka
{स वने नित्यसंवृद्धो मुनिर्वनचरः सदा}
{नान्यं जानाति विप्रेन्द्रो नित्यं पित्रनुवर्तनात्} %||1-8-8||

\twolineshloka
{द्वैविध्यं ब्रह्मचर्यस्य भविष्यति महात्मनः}
{लोकेषु प्रथितं राजन्विप्रैश्च कथितं सदा} %||1-8-9||

\twolineshloka
{तस्यैवं वर्तमानस्य कालः समभिवर्तत}
{अग्निं शुश्रूषमाणस्य पितरं च यशस्विनम्} %||1-8-10||

\twolineshloka
{एतस्मिन्नेव काले तु रोमपादः प्रतापवान्}
{अङ्गेषु प्रथितो राजा भविष्यति महाबलः} %||1-8-11||

\twolineshloka
{तस्य व्यतिक्रमाद्राज्ञो भविष्यति सुदारुणा}
{अनावृष्टिः सुघोरा वै सर्वभूतभयावहा} %||1-8-12||

\twolineshloka
{अनावृष्ट्यां तु वृत्तायां राजा दुःखसमन्वितः}
{ब्राह्मणाञ्श्रुतवृद्धांश्च समानीय प्रवक्ष्यति} %||1-8-13||

\twolineshloka
{भवन्तः श्रुतधर्माणो लोके चारित्रवेदिनः}
{समादिशन्तु नियमं प्रायश्चित्तं यथा भवेत्} %||1-8-14||

\twolineshloka
{वक्ष्यन्ति ते महीपालं ब्राह्मणा वेदपारगाः}
{विभाण्डकसुतं राजन्सर्वोपायैरिहानय} %||1-8-15||

\twolineshloka
{आनाय्य च महीपाल ऋष्यशृङ्गं सुसत्कृतम्}
{प्रयच्छ कन्यां शान्तां वै विधिना सुसमाहितः} %||1-8-16||

\twolineshloka
{तेषां तु वचनं श्रुत्वा राजा चिन्तां प्रपत्स्यते}
{केनोपायेन वै शक्यमिहानेतुं स वीर्यवान्} %||1-8-17||

\twolineshloka
{ततो राजा विनिश्चित्य सह मन्त्रिभिरात्मवान्}
{पुरोहितममात्यांश्च प्रेषयिष्यति सत्कृतान्} %||1-8-18||

\twolineshloka
{ते तु राज्ञो वचः श्रुत्वा व्यथिता वनताननाः}
{न गच्छेम ऋषेर्भीता अनुनेष्यन्ति तं नृपम्} %||1-8-19||

\twolineshloka
{वक्ष्यन्ति चिन्तयित्वा ते तस्योपायांश्च तान्क्षमान्}
{आनेष्यामो वयं विप्रं न च दोषो भविष्यति} %||1-8-20||

\twolineshloka
{एवमङ्गाधिपेनैव गणिकाभिरृषेः सुतः}
{आनीतोऽवर्षयद्देवः शान्ता चास्मै प्रदीयते} %||1-8-21||

\twolineshloka
{ऋष्यशृङ्गस्तु जामाता पुत्रांस्तव विधास्यति}
{सनत्कुमारकथितमेतावद्व्याहृतं मया} %||1-8-22||

\twolineshloka
{अथ हृष्टो दशरथः सुमन्त्रं प्रत्यभाषत}
{यथर्ष्यशृङ्गस्त्वानीतो विस्तरेण त्वयोच्यताम्} %||1-9-23||


{॥इत्यार्षे श्रीमद्रामायणे वाल्मीकीये आदिकाव्ये बालकाण्डे अष्टमः सर्गः॥}

\sect{नवमः सर्गः}

\twolineshloka
{सुमन्त्रश्चोदितो राज्ञा प्रोवाचेदं वचस्तदा}
{यथर्ष्यशृङ्गस्त्वानीतः शृणु मे मन्त्रिभिः सह} %||1-9-1||

\twolineshloka
{रोमपादमुवाचेदं सहामात्यः पुरोहितः}
{उपायो निरपायोऽयमस्माभिरभिचिन्तितः} %||1-9-2||

\twolineshloka
{ऋष्यशृङ्गो वनचरस्तपःस्वाध्यायने रतः}
{अनभिज्ञः स नारीणां विषयाणां सुखस्य च} %||1-9-3||

\twolineshloka
{इन्द्रियार्थैरभिमतैर्नरचित्तप्रमाथिभिः}
{पुरमानाययिष्यामः क्षिप्रं चाध्यवसीयताम्} %||1-9-4||

\twolineshloka
{गणिकास्तत्र गच्छन्तु रूपवत्यः स्वलङ्कृताः}
{प्रलोभ्य विविधोपायैरानेष्यन्तीह सत्कृताः} %||1-9-5||

\twolineshloka
{श्रुत्वा तथेति राजा च प्रत्युवाच पुरोहितम्}
{पुरोहितो मन्त्रिणश्च तथा चक्रुश्च ते तदा} %||1-9-6||

\twolineshloka
{वारमुख्यास्तु तच्छ्रुत्वा वनं प्रविविशुर्महत्}
{आश्रमस्याविदूरेऽस्मिन्यत्नं कुर्वन्ति दर्शने} %||1-9-7||

\twolineshloka
{ऋषिपुत्रस्य घोरस्य नित्यमाश्रमवासिनः}
{पितुः स नित्यसन्तुष्टो नातिचक्राम चाश्रमात्} %||1-9-8||

\twolineshloka
{न तेन जन्मप्रभृति दृष्टपूर्वं तपस्विना}
{स्त्री वा पुमान्वा यच्चान्यत्सत्त्वं नगरराष्ट्रजम्} %||1-9-9||

\twolineshloka
{ततः कदाचित्तं देशमाजगाम यदृच्छया}
{विभाण्डकसुतस्तत्र ताश्चापश्यद्वराङ्गनाः} %||1-9-10||

\twolineshloka
{ताश्चित्रवेषाः प्रमदा गायन्त्यो मधुरस्वरैः}
{ऋषिपुत्रमुपागम्य सर्वा वचनमब्रुवन्} %||1-9-11||

\twolineshloka
{कस्त्वं किं वर्तसे ब्रह्मञ्ज्ञातुमिच्छामहे वयम्}
{एकस्त्वं विजने घोरे वने चरसि शंस नः} %||1-9-12||

\twolineshloka
{अदृष्टरूपास्तास्तेन काम्यरूपा वने स्त्रियः}
{हार्दात्तस्य मतिर्जाता आख्यातुं पितरं स्वकम्} %||1-9-13||

\twolineshloka
{पिता विभाण्डकोऽस्माकं तस्याहं सुत औरसः}
{ऋष्यशृङ्ग इति ख्यातं नाम कर्म च मे भुवि} %||1-9-14||

\twolineshloka
{इहाश्रमपदोऽस्माकं समीपे शुभदर्शनाः}
{करिष्ये वोऽत्र पूजां वै सर्वेषां विधिपूर्वकम्} %||1-9-15||

\twolineshloka
{ऋषिपुत्रवचः श्रुत्वा सर्वासां मतिरास वै}
{तदाश्रमपदं द्रष्टुं जग्मुः सर्वाश्च तेन ह} %||1-9-16||

\twolineshloka
{गतानां तु ततः पूजामृषिपुत्रश्चकार ह}
{इदमर्घ्यमिदं पाद्यमिदं मूलं फलं च नः} %||1-9-17||

\twolineshloka
{प्रतिगृह्य तु तां पूजां सर्वा एव समुत्सुकाः}
{ऋषेर्भीताश्च शीघ्रं तु गमनाय मतिं दधुः} %||1-9-18||

\twolineshloka
{अस्माकमपि मुख्यानि फलानीमानि वै द्विज}
{गृहाण प्रति भद्रं ते भक्षयस्व च मा चिरम्} %||1-9-19||

\twolineshloka
{ततस्तास्तं समालिङ्ग्य सर्वा हर्षसमन्विताः}
{मोदकान्प्रददुस्तस्मै भक्ष्यांश्च विविधाञ्शुभान्} %||1-9-20||

\twolineshloka
{तानि चास्वाद्य तेजस्वी फलानीति स्म मन्यते}
{अनास्वादितपूर्वाणि वने नित्यनिवासिनाम्} %||1-9-21||

\twolineshloka
{आपृच्छ्य च तदा विप्रं व्रतचर्यां निवेद्य च}
{गच्छन्ति स्मापदेशात्ता भीतास्तस्य पितुः स्त्रियः} %||1-9-22||

\twolineshloka
{गतासु तासु सर्वासु काश्यपस्यात्मजो द्विजः}
{अस्वस्थहृदयश्चासीद्दुःखं स्म परिवर्तते} %||1-9-23||

\twolineshloka
{ततोऽपरेद्युस्तं देशमाजगाम स वीर्यवान्}
{मनोज्ञा यत्र ता दृष्टा वारमुख्याः स्वलङ्कृताः} %||1-9-24||

\twolineshloka
{दृष्ट्वैव च तदा विप्रमायान्तं हृष्टमानसाः}
{उपसृत्य ततः सर्वास्तास्तमूचुरिदं वचः} %||1-9-25||

\twolineshloka
{एह्याश्रमपदं सौम्य अस्माकमिति चाब्रुवन्}
{तत्राप्येष विधिः श्रीमान्विशेषेण भविष्यति} %||1-9-26||

\twolineshloka
{श्रुत्वा तु वचनं तासां सर्वासां हृदयङ्गमम्}
{गमनाय मतिं चक्रे तं च निन्युस्तदा स्त्रियः} %||1-9-27||

\twolineshloka
{तत्र चानीयमाने तु विप्रे तस्मिन्महात्मनि}
{ववर्ष सहसा देवो जगत्प्रह्लादयंस्तदा} %||1-9-28||

\twolineshloka
{वर्षेणैवागतं विप्रं विषयं स्वं नराधिपः}
{प्रत्युद्गम्य मुनिं प्रह्वः शिरसा च महीं गतः} %||1-9-29||

\twolineshloka
{अर्घ्यं च प्रददौ तस्मै न्यायतः सुसमाहितः}
{वव्रे प्रसादं विप्रेन्द्रान्मा विप्रं मन्युराविशेत्} %||1-9-30||

\twolineshloka
{अन्तःपुरं प्रविश्यास्मै कन्यां दत्त्वा यथाविधि}
{शान्तां शान्तेन मनसा राजा हर्षमवाप सः} %||1-9-31||

\twolineshloka
{एवं स न्यवसत्तत्र सर्वकामैः सुपूजितः}
{ऋष्यशृङ्गो महातेजाः शान्तया सह भार्यया} %||1-10-32||


{॥इत्यार्षे श्रीमद्रामायणे वाल्मीकीये आदिकाव्ये बालकाण्डे नवमः सर्गः॥}

\sect{दशमः सर्गः}

\twolineshloka
{भूय एव च राजेन्द्र शृणु मे वचनं हितम्}
{यथा स देवप्रवरः कथयामास बुद्धिमान्} %||1-10-1||

\twolineshloka
{इक्ष्वाकूणां कुले जातो भविष्यति सुधार्मिकः}
{राजा दशरथो नाम्ना श्रीमान्सत्यप्रतिश्रवः} %||1-10-2||

\twolineshloka
{अङ्गराजेन सख्यं च तस्य राज्ञो भविष्यति}
{कन्या चास्य महाभागा शान्ता नाम भविष्यति} %||1-10-3||

\twolineshloka
{पुत्रस्त्वङ्गस्य राज्ञस्तु रोमपाद इति श्रुतः}
{तं स राजा दशरथो गमिष्यति महायशाः} %||1-10-4||

\twolineshloka
{अनपत्योऽस्मि धर्मात्मञ्शान्ताभर्ता मम क्रतुम्}
{आहरेत त्वयाज्ञप्तः सन्तानार्थं कुलस्य च} %||1-10-5||

\twolineshloka
{श्रुत्वा राज्ञोऽथ तद्वाक्यं मनसा स विचिन्त्य च}
{प्रदास्यते पुत्रवन्तं शान्ता भर्तारमात्मवान्} %||1-10-6||

\twolineshloka
{प्रतिगृह्य च तं विप्रं स राजा विगतज्वरः}
{आहरिष्यति तं यज्ञं प्रहृष्टेनान्तरात्मना} %||1-10-7||

\twolineshloka
{तं च राजा दशरथो यष्टुकामः कृताञ्जलिः}
{ऋष्यशृङ्गं द्विजश्रेष्ठं वरयिष्यति धर्मवित्} %||1-10-8||

\twolineshloka
{यज्ञार्थं प्रसवार्थं च स्वर्गार्थं च नरेश्वरः}
{लभते च स तं कामं द्विजमुख्याद्विशां पतिः} %||1-10-9||

\twolineshloka
{पुत्राश्चास्य भविष्यन्ति चत्वारोऽमितविक्रमाः}
{वंशप्रतिष्ठानकराः सर्वलोकेषु विश्रुताः} %||1-10-10||

\twolineshloka
{एवं स देवप्रवरः पूर्वं कथितवान्कथाम्}
{सनत्कुमारो भगवान्पुरा देवयुगे प्रभुः} %||1-10-11||

\twolineshloka
{स त्वं पुरुषशार्दूल तमानय सुसत्कृतम्}
{स्वयमेव महाराज गत्वा सबलवाहनः} %||1-10-12||

\twolineshloka
{अनुमान्य वसिष्ठं च सूतवाक्यं निशम्य च}
{सान्तःपुरः सहामात्यः प्रययौ यत्र स द्विजः} %||1-10-13||

\twolineshloka
{वनानि सरितश्चैव व्यतिक्रम्य शनैः शनैः}
{अभिचक्राम तं देशं यत्र वै मुनिपुङ्गवः} %||1-10-14||

\twolineshloka
{आसाद्य तं द्विजश्रेष्ठं रोमपादसमीपगम्}
{ऋषिपुत्रं ददर्शादौ दीप्यमानमिवानलम्} %||1-10-15||

\twolineshloka
{ततो राजा यथान्यायं पूजां चक्रे विशेषतः}
{सखित्वात्तस्य वै राज्ञः प्रहृष्टेनान्तरात्मना} %||1-10-16||

\twolineshloka
{रोमपादेन चाख्यातमृषिपुत्राय धीमते}
{सख्यं सम्बन्धकं चैव तदा तं प्रत्यपूजयत्} %||1-10-17||

\twolineshloka
{एवं सुसत्कृतस्तेन सहोषित्वा नरर्षभः}
{सप्ताष्टदिवसान्राजा राजानमिदमब्रवीत्} %||1-10-18||

\twolineshloka
{शान्ता तव सुता राजन्सह भर्त्रा विशाम्पते}
{मदीयं नगरं यातु कार्यं हि महदुद्यतम्} %||1-10-19||

\twolineshloka
{तथेति राजा संश्रुत्य गमनं तस्य धीमतः}
{उवाच वचनं विप्रं गच्छ त्वं सह भार्यया} %||1-10-20||

\twolineshloka
{ऋषिपुत्रः प्रतिश्रुत्य तथेत्याह नृपं तदा}
{स नृपेणाभ्यनुज्ञातः प्रययौ सह भार्यया} %||1-10-21||

\twolineshloka
{तावन्योन्याञ्जलिं कृत्वा स्नेहात्संश्लिष्य चोरसा}
{ननन्दतुर्दशरथो रोमपादश्च वीर्यवान्} %||1-10-22||

\threelineshloka
{ततः सुहृदमापृच्छ्य प्रस्थितो रघुनन्दनः}
{पौरेभ्यः प्रेषयामास दूतान्वै शीघ्रगामिनः}
{क्रियतां नगरं सर्वं क्षिप्रमेव स्वलङ्कृतम्} %||1-10-23||

\twolineshloka
{ततः प्रहृष्टाः पौरास्ते श्रुत्वा राजानमागतम्}
{तथा प्रचक्रुस्तत्सर्वं राज्ञा यत्प्रेषितं तदा} %||1-10-24||

\twolineshloka
{ततः स्वलङ्कृतं राजा नगरं प्रविवेश ह}
{शङ्खदुन्दुभिनिर्घोषैः पुरस्कृत्य द्विजर्षभम्} %||1-10-25||

\twolineshloka
{ततः प्रमुदिताः सर्वे दृष्ट्वा वै नागरा द्विजम्}
{प्रवेश्यमानं सत्कृत्य नरेन्द्रेणेन्द्रकर्मणा} %||1-10-26||

\twolineshloka
{अन्तःपुरं प्रवेश्यैनं पूजां कृत्वा तु शास्त्रतः}
{कृतकृत्यं तदात्मानं मेने तस्योपवाहनात्} %||1-10-27||

\twolineshloka
{अन्तःपुराणि सर्वाणि शान्तां दृष्ट्वा तथागताम्}
{सह भर्त्रा विशालाक्षीं प्रीत्यानन्दमुपागमन्} %||1-10-28||

\twolineshloka
{पूज्यमाना च ताभिः सा राज्ञा चैव विशेषतः}
{उवास तत्र सुखिता कञ्चित्कालं सह द्विजा} %||1-11-29||


{॥इत्यार्षे श्रीमद्रामायणे वाल्मीकीये आदिकाव्ये बालकाण्डे दशमः सर्गः॥}

\sect{एकादशः सर्गः}

\twolineshloka
{ततः काले बहुतिथे कस्मिंश्चित्सुमनोहरे}
{वसन्ते समनुप्राप्ते राज्ञो यष्टुं मनोऽभवत्} %||1-11-1||

\twolineshloka
{ततः प्रसाद्य शिरसा तं विप्रं देववर्णिनम्}
{यज्ञाय वरयामास सन्तानार्थं कुलस्य च} %||1-11-2||

\twolineshloka
{तथेति च स राजानमुवाच च सुसत्कृतः}
{सम्भाराः सम्भ्रियन्तां ते तुरगश्च विमुच्यताम्} %||1-11-3||

\twolineshloka
{ततो राजाब्रवीद्वाक्यं सुमन्त्रं मन्त्रिसत्तमम्}
{सुमन्त्रावाहय क्षिप्रमृत्विजो ब्रह्मवादिनः} %||1-11-4||

\twolineshloka
{ततः सुमन्त्रस्त्वरितं गत्वा त्वरितविक्रमः}
{समानयत्स तान्विप्रान्समस्तान्वेदपारगान्} %||1-11-5||

\twolineshloka
{सुयज्ञं वामदेवं च जाबालिमथ काश्यपम्}
{पुरोहितं वसिष्ठं च ये चान्ये द्विजसत्तमाः} %||1-11-6||

\twolineshloka
{तान्पूजयित्वा धर्मात्मा राजा दशरथस्तदा}
{इदं धर्मार्थसहितं श्लक्ष्णं वचनमब्रवीत्} %||1-11-7||

\twolineshloka
{मम लालप्यमानस्य पुत्रार्थं नास्ति वै सुखम्}
{तदर्थं हयमेधेन यक्ष्यामीति मतिर्मम} %||1-11-8||

\twolineshloka
{तदहं यष्टुमिच्छामि शास्त्रदृष्टेन कर्मणा}
{ऋषिपुत्रप्रभावेन कामान्प्राप्स्यामि चाप्यहम्} %||1-11-9||

\twolineshloka
{ततः साध्विति तद्वाक्यं ब्राह्मणाः प्रत्यपूजयन्}
{वसिष्ठप्रमुखाः सर्वे पार्थिवस्य मुखाच्च्युतम्} %||1-11-10||

\twolineshloka
{ऋष्यशृङ्गपुरोगाश्च प्रत्यूचुर्नृपतिं तदा}
{सम्भाराः सम्भ्रियन्तां ते तुरगश्च विमुच्यताम्} %||1-11-11||

\twolineshloka
{सर्वथा प्राप्यसे पुत्रांश्चतुरोऽमितविक्रमान्}
{यस्य ते धार्मिकी बुद्धिरियं पुत्रार्थमागता} %||1-11-12||

\twolineshloka
{ततः प्रीतोऽभवद्राजा श्रुत्वा तद्द्विजभाषितम्}
{अमात्यांश्चाब्रवीद्राजा हर्षेणेदं शुभाक्षरम्} %||1-11-13||

\twolineshloka
{गुरूणां वचनाच्छीघ्रं सम्भाराः सम्भ्रियन्तु मे}
{समर्थाधिष्ठितश्चाश्वः सोपाध्यायो विमुच्यताम्} %||1-11-14||

\twolineshloka
{सरय्वाश्चोत्तरे तीरे यज्ञभूमिर्विधीयताम्}
{शान्तयश्चाभिवर्धन्तां यथाकल्पं यथाविधि} %||1-11-15||

\twolineshloka
{शक्यः कर्तुमयं यज्ञः सर्वेणापि महीक्षिता}
{नापराधो भवेत्कष्टो यद्यस्मिन्क्रतुसत्तमे} %||1-11-16||

\twolineshloka
{छिद्रं हि मृगयन्तेऽत्र विद्वांसो ब्रह्मराक्षसाः}
{विधिहीनस्य यज्ञस्य सद्यः कर्ता विनश्यति} %||1-11-17||

\twolineshloka
{तद्यथा विधिपूर्वं मे क्रतुरेष समाप्यते}
{तथाविधानं क्रियतां समर्थाः करणेष्विह} %||1-11-18||

\twolineshloka
{तथेति च ततः सर्वे मन्त्रिणः प्रत्यपूजयन्}
{पार्थिवेन्द्रस्य तद्वाक्यं यथाज्ञप्तमकुर्वत} %||1-11-19||

\twolineshloka
{ततो द्विजास्ते धर्मज्ञमस्तुवन्पार्थिवर्षभम्}
{अनुज्ञातास्ततः सर्वे पुनर्जग्मुर्यथागतम्} %||1-11-20||

\twolineshloka
{गतानां तु द्विजातीनां मन्त्रिणस्तान्नराधिपः}
{विसर्जयित्वा स्वं वेश्म प्रविवेश महाद्युतिः} %||1-12-21||


{॥इत्यार्षे श्रीमद्रामायणे वाल्मीकीये आदिकाव्ये बालकाण्डे एकादशः सर्गः॥}

\sect{द्वादशः सर्गः}

\twolineshloka
{पुनः प्राप्ते वसन्ते तु पूर्णः संवत्सरोऽभवत्}
{अभिवाद्य वसिष्ठं च न्यायतः प्रतिपूज्य च} %||1-12-1||

\twolineshloka
{अब्रवीत्प्रश्रितं वाक्यं प्रसवार्थं द्विजोत्तमम्}
{यज्ञो मे क्रियतां विप्र यथोक्तं मुनिपुङ्गव} %||1-12-2||

\twolineshloka
{यथा न विघ्नः क्रियते यज्ञाङ्गेषु विधीयताम्}
{भवान्स्निग्धः सुहृन्मह्यं गुरुश्च परमो भवान्} %||1-12-3||

\twolineshloka
{वोढव्यो भवता चैव भारो यज्ञस्य चोद्यतः}
{तथेति च स राजानमब्रवीद्द्विजसत्तमः} %||1-12-4||

\twolineshloka
{करिष्ये सर्वमेवैतद्भवता यत्समर्थितम्}
{ततोऽब्रवीद्द्विजान्वृद्धान्यज्ञकर्मसु निष्ठितान्} %||1-12-5||

\twolineshloka
{स्थापत्ये निष्ठितांश्चैव वृद्धान्परमधार्मिकान्}
{कर्मान्तिकाञ्शिल्पकारान्वर्धकीन्खनकानपि} %||1-12-6||

\twolineshloka
{गणकाञ्शिल्पिनश्चैव तथैव नटनर्तकान्}
{तथा शुचीञ्शास्त्रविदः पुरुषान्सुबहुश्रुतान्} %||1-12-7||

\twolineshloka
{यज्ञकर्म समीहन्तां भवन्तो राजशासनात्}
{इष्टका बहुसाहस्री शीघ्रमानीयतामिति} %||1-12-8||

\twolineshloka
{औपकार्याः क्रियन्तां च राज्ञां बहुगुणान्विताः}
{ब्राह्मणावसथाश्चैव कर्तव्याः शतशः शुभाः} %||1-12-9||

\twolineshloka
{भक्ष्यान्नपानैर्बहुभिः समुपेताः सुनिष्ठिताः}
{तथा पौरजनस्यापि कर्तव्या बहुविस्तराः} %||1-12-10||

\twolineshloka
{आवासा बहुभक्ष्या वै सर्वकामैरुपस्थिताः}
{तथा जानपदस्यापि जनस्य बहुशोभनम्} %||1-12-11||

\twolineshloka
{दातव्यमन्नं विधिवत्सत्कृत्य न तु लीलया}
{सर्ववर्णा यथा पूजां प्राप्नुवन्ति सुसत्कृताः} %||1-12-12||

\twolineshloka
{न चावज्ञा प्रयोक्तव्या कामक्रोधवशादपि}
{यज्ञकर्मसु येऽव्यग्राः पुरुषाः शिल्पिनस्तथा} %||1-12-13||

\twolineshloka
{तेषामपि विशेषेण पूजा कार्या यथाक्रमम्}
{यथा सर्वं सुविहितं न किञ्चित्परिहीयते} %||1-12-14||

\twolineshloka
{तथा भवन्तः कुर्वन्तु प्रीतिस्निग्धेन चेतसा}
{ततः सर्वे समागम्य वसिष्ठमिदमब्रुवन्} %||1-12-15||

\twolineshloka
{यथोक्तं तत्करिष्यामो न किञ्चित्परिहास्यते}
{ततः सुमन्त्रमाहूय वसिष्ठो वाक्यमब्रवीत्} %||1-12-16||

\twolineshloka
{निमन्त्रयस्य नृपतीन्पृथिव्यां ये च धार्मिकाः}
{ब्राह्मणान्क्षत्रियान्वैश्याञ्शूद्रांश्चैव सहस्रशः} %||1-12-17||

\twolineshloka
{समानयस्व सत्कृत्य सर्वदेशेषु मानवान्}
{मिथिलाधिपतिं शूरं जनकं सत्यविक्रमम्} %||1-12-18||

\threelineshloka
{निष्ठितं सर्वशास्त्रेषु तथा वेदेषु निष्ठितम्}
{तमानय महाभागं स्वयमेव सुसत्कृतम्}
{पूर्वसम्बन्धिनं ज्ञात्वा ततः पूर्वं ब्रवीमि ते} %||1-12-19||

\twolineshloka
{तथा काशिपतिं स्निग्धं सततं प्रियवादिनम्}
{सद्वृत्तं देवसङ्काशं स्वयमेवानयस्व ह} %||1-12-20||

\twolineshloka
{तथा केकयराजानं वृद्धं परमधार्मिकम्}
{श्वशुरं राजसिंहस्य सपुत्रं तमिहानय} %||1-12-21||

\twolineshloka
{अङ्गेश्वरं महाभागं रोमपादं सुसत्कृतम्}
{वयस्यं राजसिंहस्य तमानय यशस्विनम्} %||1-12-22||

\twolineshloka
{प्राचीनान्सिन्धुसौवीरान्सौराष्ठ्रेयांश्च पार्थिवान्}
{दाक्षिणात्यान्नरेन्द्रांश्च समस्तानानयस्व ह} %||1-12-23||

\twolineshloka
{सन्ति स्निग्धाश्च ये चान्ये राजानः पृथिवीतले}
{तानानय यथाक्षिप्रं सानुगान्सहबान्धवान्} %||1-12-24||

\twolineshloka
{वसिष्ठवाक्यं तच्छ्रुत्वा सुमन्त्रस्त्वरितस्तदा}
{व्यादिशत्पुरुषांस्तत्र राज्ञामानयने शुभान्} %||1-12-25||

\twolineshloka
{स्वयमेव हि धर्मात्मा प्रययौ मुनिशासनात्}
{सुमन्त्रस्त्वरितो भूत्वा समानेतुं महीक्षितः} %||1-12-26||

\twolineshloka
{ते च कर्मान्तिकाः सर्वे वसिष्ठाय च धीमते}
{सर्वं निवेदयन्ति स्म यज्ञे यदुपकल्पितम्} %||1-12-27||

\threelineshloka
{ततः प्रीतो द्विजश्रेष्ठस्तान्सर्वान्पुनरब्रवीत्}
{अवज्ञया न दातव्यं कस्यचिल्लीलयापि वा}
{अवज्ञया कृतं हन्याद्दातारं नात्र संशयः} %||1-12-28||

\twolineshloka
{ततः कैश्चिदहोरात्रैरुपयाता महीक्षितः}
{बहूनि रत्नान्यादाय राज्ञो दशरथस्य ह} %||1-12-29||

\twolineshloka
{ततो वसिष्ठः सुप्रीतो राजानमिदमब्रवीत्}
{उपयाता नरव्याघ्र राजानस्तव शासनात्} %||1-12-30||

\twolineshloka
{मयापि सत्कृताः सर्वे यथार्हं राजसत्तमाः}
{यज्ञियं च कृतं राजन्पुरुषैः सुसमाहितैः} %||1-12-31||

\twolineshloka
{निर्यातु च भवान्यष्टुं यज्ञायतनमन्तिकात्}
{सर्वकामैरुपहृतैरुपेतं वै समन्ततः} %||1-12-32||

\twolineshloka
{तथा वसिष्ठवचनादृष्यशृङ्गस्य चोभयोः}
{शुभे दिवस नक्षत्रे निर्यातो जगतीपतिः} %||1-12-33||

\twolineshloka
{ततो वसिष्ठप्रमुखाः सर्व एव द्विजोत्तमाः}
{ऋष्यशृङ्गं पुरस्कृत्य यज्ञकर्मारभंस्तदा} %||1-13-34||


{॥इत्यार्षे श्रीमद्रामायणे वाल्मीकीये आदिकाव्ये बालकाण्डे द्वादशः सर्गः॥}

\sect{त्रयोदशः सर्गः}

\twolineshloka
{अथ संवत्सरे पूर्णे तस्मिन्प्राप्ते तुरङ्गमे}
{सरय्वाश्चोत्तरे तीरे राज्ञो यज्ञोऽभ्यवर्तत} %||1-13-1||

\twolineshloka
{ऋष्यशृङ्गं पुरस्कृत्य कर्म चक्रुर्द्विजर्षभाः}
{अश्वमेधे महायज्ञे राज्ञोऽस्य सुमहात्मनः} %||1-13-2||

\twolineshloka
{कर्म कुर्वन्ति विधिवद्याजका वेदपारगाः}
{यथाविधि यथान्यायं परिक्रामन्ति शास्त्रतः} %||1-13-3||

\twolineshloka
{प्रवर्ग्यं शास्त्रतः कृत्वा तथैवोपसदं द्विजाः}
{चक्रुश्च विधिवत्सर्वमधिकं कर्म शास्त्रतः} %||1-13-4||

\twolineshloka
{अभिपूज्य ततो हृष्टाः सर्वे चक्रुर्यथाविधि}
{प्रातःसवनपूर्वाणि कर्माणि मुनिपुङ्गवाः} %||1-13-5||

\twolineshloka
{न चाहुतमभूत्तत्र स्खलितं वापि किञ्चन}
{दृश्यते ब्रह्मवत्सर्वं क्षेमयुक्तं हि चक्रिरे} %||1-13-6||

\twolineshloka
{न तेष्वहःसु श्रान्तो वा क्षुधितो वापि दृश्यते}
{नाविद्वान्ब्राह्मणस्तत्र नाशतानुचरस्तथा} %||1-13-7||

\twolineshloka
{ब्राह्मणा भुञ्जते नित्यं नाथवन्तश्च भुञ्जते}
{तापसा भुञ्जते चापि श्रमणा भुञ्जते तथा} %||1-13-8||

\twolineshloka
{वृद्धाश्च व्याधिताश्चैव स्त्रियो बालास्तथैव च}
{अनिशं भुञ्जमानानां न तृप्तिरुपलभ्यते} %||1-13-9||

\twolineshloka
{दीयतां दीयतामन्नं वासांसि विविधानि च}
{इति सञ्चोदितास्तत्र तथा चक्रुरनेकशः} %||1-13-10||

\twolineshloka
{अन्नकूटाश्च बहवो दृश्यन्ते पर्वतोपमाः}
{दिवसे दिवसे तत्र सिद्धस्य विधिवत्तदा} %||1-13-11||

\twolineshloka
{अन्नं हि विधिवत्स्वादु प्रशंसन्ति द्विजर्षभाः}
{अहो तृप्ताः स्म भद्रं ते इति शुश्राव राघवः} %||1-13-12||

\twolineshloka
{स्वलङ्कृताश्च पुरुषा ब्राह्मणान्पर्यवेषयन्}
{उपासते च तानन्ये सुमृष्टमणिकुण्डलाः} %||1-13-13||

\twolineshloka
{कर्मान्तरे तदा विप्रा हेतुवादान्बहूनपि}
{प्राहुः सुवाग्मिनो धीराः परस्परजिगीषया} %||1-13-14||

\twolineshloka
{दिवसे दिवसे तत्र संस्तरे कुशला द्विजाः}
{सर्वकर्माणि चक्रुस्ते यथाशास्त्रं प्रचोदिताः} %||1-13-15||

\twolineshloka
{नाषडङ्गविदत्रासीन्नाव्रतो नाबहुश्रुतः}
{सदस्यस्तस्य वै राज्ञो नावादकुशलो द्विजः} %||1-13-16||

\twolineshloka
{प्राप्ते यूपोच्छ्रये तस्मिन्षड्बैल्वाः खादिरास्तथा}
{तावन्तो बिल्वसहिताः पर्णिनश्च तथापरे} %||1-13-17||

\twolineshloka
{श्लेष्मातकमयो दिष्टो देवदारुमयस्तथा}
{द्वावेव तत्र विहितौ बाहुव्यस्तपरिग्रहौ} %||1-13-18||

\twolineshloka
{कारिताः सर्व एवैते शास्त्रज्ञैर्यज्ञकोविदैः}
{शोभार्थं तस्य यज्ञस्य काञ्चनालङ्कृता भवन्} %||1-13-19||

\twolineshloka
{विन्यस्ता विधिवत्सर्वे शिल्पिभिः सुकृता दृढाः}
{अष्टाश्रयः सर्व एव श्लक्ष्णरूपसमन्विताः} %||1-13-20||

\twolineshloka
{आच्छादितास्ते वासोभिः पुष्पैर्गन्धैश्च भूषिताः}
{सप्तर्षयो दीप्तिमन्तो विराजन्ते यथा दिवि} %||1-13-21||

\threelineshloka
{इष्टकाश्च यथान्यायं कारिताश्च प्रमाणतः}
{चितोऽग्निर्ब्राह्मणैस्तत्र कुशलैः शुल्बकर्मणि}
{स चित्यो राजसिंहस्य सञ्चितः कुशलैर्द्विजैः} %||1-13-22||

\twolineshloka
{गरुडो रुक्मपक्षो वै त्रिगुणोऽष्टादशात्मकः}
{नियुक्तास्तत्र पशवस्तत्तदुद्दिश्य दैवतम्} %||1-13-23||

\twolineshloka
{उरगाः पक्षिणश्चैव यथाशास्त्रं प्रचोदिताः}
{शामित्रे तु हयस्तत्र तथा जल चराश्च ये} %||1-13-24||

\threelineshloka
{ऋत्विग्भिः सर्वमेवैतन्नियुक्तं शास्त्रतस्तदा}
{पशूनां त्रिशतं तत्र यूपेषु नियतं तदा}
{अश्वरत्नोत्तमं तस्य राज्ञो दशरथस्य ह} %||1-13-25||

\twolineshloka
{कौसल्या तं हयं तत्र परिचर्य समन्ततः}
{कृपाणैर्विशशासैनं त्रिभिः परमया मुदा} %||1-13-26||

\twolineshloka
{पतत्रिणा तदा सार्धं सुस्थितेन च चेतसा}
{अवसद्रजनीमेकां कौसल्या धर्मकाम्यया} %||1-13-27||

\twolineshloka
{होताध्वर्युस्तथोद्गाता हयेन समयोजयन्}
{महिष्या परिवृत्थ्याथ वावातामपरां तथा} %||1-13-28||

\twolineshloka
{पतत्रिणस्तस्य वपामुद्धृत्य नियतेन्द्रियः}
{ऋत्विक्परम सम्पन्नः श्रपयामास शास्त्रतः} %||1-13-29||

\twolineshloka
{धूमगन्धं वपायास्तु जिघ्रति स्म नराधिपः}
{यथाकालं यथान्यायं निर्णुदन्पापमात्मनः} %||1-13-30||

\twolineshloka
{हयस्य यानि चाङ्गानि तानि सर्वाणि ब्राह्मणाः}
{अग्नौ प्रास्यन्ति विधिवत्समस्ताः षोडशर्त्विजः} %||1-13-31||

\twolineshloka
{प्लक्षशाखासु यज्ञानामन्येषां क्रियते हविः}
{अश्वमेधस्य चैकस्य वैतसो भाग इष्यते} %||1-13-32||

\twolineshloka
{त्र्यहोऽश्वमेधः सङ्ख्यातः कल्पसूत्रेण ब्राह्मणैः}
{चतुष्टोममहस्तस्य प्रथमं परिकल्पितम्} %||1-13-33||

\twolineshloka
{उक्थ्यं द्वितीयं सङ्ख्यातमतिरात्रं तथोत्तरम्}
{कारितास्तत्र बहवो विहिताः शास्त्रदर्शनात्} %||1-13-34||

\twolineshloka
{ज्योतिष्टोमायुषी चैव अतिरात्रौ च निर्मितौ}
{अभिजिद्विश्वजिच्चैव अप्तोर्यामो महाक्रतुः} %||1-13-35||

\twolineshloka
{प्राचीं होत्रे ददौ राजा दिशं स्वकुलवर्धनः}
{अध्वर्यवे प्रतीचीं तु ब्रह्मणे दक्षिणां दिशम्} %||1-13-36||

\twolineshloka
{उद्गात्रे तु तथोदीचीं दक्षिणैषा विनिर्मिता}
{अश्वमेधे महायज्ञे स्वयम्भुविहिते पुरा} %||1-13-37||

\twolineshloka
{क्रतुं समाप्य तु तदा न्यायतः पुरुषर्षभः}
{ऋत्विग्भ्यो हि ददौ राजा धरां तां क्रतुवर्धनः} %||1-13-38||

\twolineshloka
{ऋत्विजस्त्वब्रुवन्सर्वे राजानं गतकल्मषम्}
{भवानेव महीं कृत्स्नामेको रक्षितुमर्हति} %||1-13-39||

\threelineshloka
{न भूम्या कार्यमस्माकं न हि शक्ताः स्म पालने}
{रताः स्वाध्यायकरणे वयं नित्यं हि भूमिप}
{निष्क्रयं किञ्चिदेवेह प्रयच्छतु भवानिति} %||1-13-40||

\twolineshloka
{गवां शतसहस्राणि दश तेभ्यो ददौ नृपः}
{दशकोटिं सुवर्णस्य रजतस्य चतुर्गुणम्} %||1-13-41||

\twolineshloka
{ऋत्विजस्तु ततः सर्वे प्रददुः सहिता वसु}
{ऋष्यशृङ्गाय मुनये वसिष्ठाय च धीमते} %||1-13-42||

\twolineshloka
{ततस्ते न्यायतः कृत्वा प्रविभागं द्विजोत्तमाः}
{सुप्रीतमनसः सर्वे प्रत्यूचुर्मुदिता भृशम्} %||1-13-43||

\twolineshloka
{ततः प्रीतमना राजा प्राप्य यज्ञमनुत्तमम्}
{पापापहं स्वर्नयनं दुस्तरं पार्थिवर्षभैः} %||1-13-44||

\twolineshloka
{ततोऽब्रवीदृष्यशृङ्गं राजा दशरथस्तदा}
{कुलस्य वर्धनं तत्तु कर्तुमर्हसि सुव्रत} %||1-13-45||

\twolineshloka
{तथेति च स राजानमुवाच द्विजसत्तमः}
{भविष्यन्ति सुता राजंश्चत्वारस्ते कुलोद्वहाः} %||1-14-46||


{॥इत्यार्षे श्रीमद्रामायणे वाल्मीकीये आदिकाव्ये बालकाण्डे त्रयोदशः सर्गः॥}

\sect{चतुर्दशः सर्गः}

\twolineshloka
{मेधावी तु ततो ध्यात्वा स किञ्चिदिदमुत्तमम्}
{लब्धसंज्ञस्ततस्तं तु वेदज्ञो नृपमब्रवीत्} %||1-14-1||

\twolineshloka
{इष्टिं तेऽहं करिष्यामि पुत्रीयां पुत्रकारणात्}
{अथर्वशिरसि प्रोक्तैर्मन्त्रैः सिद्धां विधानतः} %||1-14-2||

\twolineshloka
{ततः प्राक्रमदिष्टिं तां पुत्रीयां पुत्र कारणात्}
{जुहाव चाग्नौ तेजस्वी मन्त्रदृष्टेन कर्मणा} %||1-14-3||

\twolineshloka
{ततो देवाः सगन्धर्वाः सिद्धाश्च परमर्षयः}
{भागप्रतिग्रहार्थं वै समवेता यथाविधि} %||1-14-4||

\twolineshloka
{ताः समेत्य यथान्यायं तस्मिन्सदसि देवताः}
{अब्रुवँल्लोककर्तारं ब्रह्माणं वचनं महत्} %||1-14-5||

\twolineshloka
{भगवंस्त्वत्प्रसादेन रावणो नाम राक्षसः}
{सर्वान्नो बाधते वीर्याच्छासितुं तं न शक्नुमः} %||1-14-6||

\twolineshloka
{त्वया तस्मै वरो दत्तः प्रीतेन भगवन्पुरा}
{मानयन्तश्च तं नित्यं सर्वं तस्य क्षमामहे} %||1-14-7||

\twolineshloka
{उद्वेजयति लोकांस्त्रीनुच्छ्रितान्द्वेष्टि दुर्मतिः}
{शक्रं त्रिदशराजानं प्रधर्षयितुमिच्छति} %||1-14-8||

\twolineshloka
{ऋषीन्यक्षान्सगन्धर्वानसुरान्ब्राह्मणांस्तथा}
{अतिक्रामति दुर्धर्षो वरदानेन मोहितः} %||1-14-9||

\twolineshloka
{नैनं सूर्यः प्रतपति पार्श्वे वाति न मारुतः}
{चलोर्मिमाली तं दृष्ट्वा समुद्रोऽपि न कम्पते} %||1-14-10||

\twolineshloka
{तन्मनन्नो भयं तस्माद्राक्षसाद्घोरदर्शनात्}
{वधार्थं तस्य भगवन्नुपायं कर्तुमर्हसि} %||1-14-11||

\twolineshloka
{एवमुक्तः सुरैः सर्वैश्चिन्तयित्वा ततोऽब्रवीत्}
{हन्तायं विहितस्तस्य वधोपायो दुरात्मनः} %||1-14-12||

\twolineshloka
{तेन गन्धर्वयक्षाणां देवदानवरक्षसाम्}
{अवध्योऽस्मीति वागुक्ता तथेत्युक्तं च तन्मया} %||1-14-13||

\twolineshloka
{नाकीर्तयदवज्ञानात्तद्रक्षो मानुषांस्तदा}
{तस्मात्स मानुषाद्वध्यो मृतुर्नान्योऽस्य विद्यते} %||1-14-14||

\twolineshloka
{एतच्छ्रुत्वा प्रियं वाक्यं ब्रह्मणा समुदाहृतम्}
{देवा महर्षयः सर्वे प्रहृष्टास्तेऽभवंस्तदा} %||1-14-15||

\twolineshloka
{एतस्मिन्नन्तरे विष्णुरुपयातो महाद्युतिः}
{ब्रह्मणा च समागम्य तत्र तस्थौ समाहितः} %||1-14-16||

\twolineshloka
{तमब्रुवन्सुराः सर्वे समभिष्टूय संनताः}
{त्वां नियोक्ष्यामहे विष्णो लोकानां हितकाम्यया} %||1-14-17||

\fourlineshloka
{राज्ञो दशरथस्य त्वमयोध्याधिपतेर्विभो}
{धर्मज्ञस्य वदान्यस्य महर्षिसमतेजसः}
{तस्य भार्यासु तिसृषु ह्रीश्रीकीर्त्युपमासु च}
{विष्णो पुत्रत्वमागच्छ कृत्वात्मानं चतुर्विधम्} %||1-14-18||

\twolineshloka
{तत्र त्वं मानुषो भूत्वा प्रवृद्धं लोककण्टकम्}
{अवध्यं दैवतैर्विष्णो समरे जहि रावणम्} %||1-14-19||

\twolineshloka
{स हि देवान्सगन्धर्वान्सिद्धांश्च ऋषिसत्तमान्}
{राक्षसो रावणो मूर्खो वीर्योत्सेकेन बाधते} %||1-14-20||

\fourlineindentedshloka
{तदुद्धतं रावणमृद्धतेजसम्}
{ प्रवृद्धदर्पं त्रिदशेश्वरद्विषम्}
{विरावणं साधु तपस्विकण्टकम्}
{ तपस्विनामुद्धर तं भयावहम्} %||1-15-21||


{॥इत्यार्षे श्रीमद्रामायणे वाल्मीकीये आदिकाव्ये बालकाण्डे चतुर्दशः सर्गः॥}

\sect{पञ्चदशः सर्गः}

\twolineshloka
{ततो नारायणो विष्णुर्नियुक्तः सुरसत्तमैः}
{जानन्नपि सुरानेवं श्लक्ष्णं वचनमब्रवीत्} %||1-15-1||

\twolineshloka
{उपायः को वधे तस्य राक्षसाधिपतेः सुराः}
{यमहं तं समास्थाय निहन्यामृषिकण्टकम्} %||1-15-2||

\twolineshloka
{एवमुक्ताः सुराः सर्वे प्रत्यूचुर्विष्णुमव्ययम्}
{मानुषीं तनुमास्थाय रावणं जहि संयुगे} %||1-15-3||

\twolineshloka
{स हि तेपे तपस्तीव्रं दीर्घकालमरिन्दम}
{येन तुष्टोऽभवद्ब्रह्मा लोककृल्लोकपूजितः} %||1-15-4||

\twolineshloka
{सन्तुष्टः प्रददौ तस्मै राक्षसाय वरं प्रभुः}
{नानाविधेभ्यो भूतेभ्यो भयं नान्यत्र मानुषात्} %||1-15-5||

\twolineshloka
{अवज्ञाताः पुरा तेन वरदानेन मानवाः}
{तस्मात्तस्य वधो दृष्टो मानुषेभ्यः परन्तप} %||1-15-6||

\twolineshloka
{इत्येतद्वचनं श्रुत्वा सुराणां विष्णुरात्मवान्}
{पितरं रोचयामास तदा दशरथं नृपम्} %||1-15-7||

\twolineshloka
{स चाप्यपुत्रो नृपतिस्तस्मिन्काले महाद्युतिः}
{अयजत्पुत्रियामिष्टिं पुत्रेप्सुररिसूदनः} %||1-15-8||

\twolineshloka
{ततो वै यजमानस्य पावकादतुलप्रभम्}
{प्रादुर्भूतं महद्भूतं महावीर्यं महाबलम्} %||1-15-9||

\twolineshloka
{कृष्णं रक्ताम्बरधरं रक्तास्यं दुन्दुभिस्वनम्}
{स्निग्धहर्यक्षतनुजश्मश्रुप्रवरमूर्धजम्} %||1-15-10||

\twolineshloka
{शुभलक्षणसम्पन्नं दिव्याभरणभूषितम्}
{शैलशृङ्गसमुत्सेधं दृप्तशार्दूलविक्रमम्} %||1-15-11||

\twolineshloka
{दिवाकरसमाकारं दीप्तानलशिखोपमम्}
{तप्तजाम्बूनदमयीं राजतान्तपरिच्छदाम्} %||1-15-12||

\twolineshloka
{दिव्यपायससम्पूर्णां पात्रीं पत्नीमिव प्रियाम्}
{प्रगृह्य विपुलां दोर्भ्यां स्वयं मायामयीमिव} %||1-15-13||

\twolineshloka
{समवेक्ष्याब्रवीद्वाक्यमिदं दशरथं नृपम्}
{प्राजापत्यं नरं विद्धि मामिहाभ्यागतं नृप} %||1-15-14||

\twolineshloka
{ततः परं तदा राजा प्रत्युवाच कृताञ्जलिः}
{भगवन्स्वागतं तेऽस्तु किमहं करवाणि ते} %||1-15-15||

\twolineshloka
{अथो पुनरिदं वाक्यं प्राजापत्यो नरोऽब्रवीत्}
{राजन्नर्चयता देवानद्य प्राप्तमिदं त्वया} %||1-15-16||

\twolineshloka
{इदं तु नरशार्दूल पायसं देवनिर्मितम्}
{प्रजाकरं गृहाण त्वं धन्यमारोग्यवर्धनम्} %||1-15-17||

\twolineshloka
{भार्याणामनुरूपाणामश्नीतेति प्रयच्छ वै}
{तासु त्वं लप्स्यसे पुत्रान्यदर्थं यजसे नृप} %||1-15-18||

\twolineshloka
{तथेति नृपतिः प्रीतः शिरसा प्रतिगृह्यताम्}
{पात्रीं देवान्नसम्पूर्णां देवदत्तां हिरण्मयीम्} %||1-15-19||

\twolineshloka
{अभिवाद्य च तद्भूतमद्भुतं प्रियदर्शनम्}
{मुदा परमया युक्तश्चकाराभिप्रदक्षिणम्} %||1-15-20||

\twolineshloka
{ततो दशरथः प्राप्य पायसं देवनिर्मितम्}
{बभूव परमप्रीतः प्राप्य वित्तमिवाधनः} %||1-15-21||

\twolineshloka
{ततस्तदद्भुतप्रख्यं भूतं परमभास्वरम्}
{संवर्तयित्वा तत्कर्म तत्रैवान्तरधीयत} %||1-15-22||

\twolineshloka
{हर्षरश्मिभिरुद्योतं तस्यान्तःपुरमाबभौ}
{शारदस्याभिरामस्य चन्द्रस्येव नभोंऽशुभिः} %||1-15-23||

\twolineshloka
{सोऽन्तःपुरं प्रविश्यैव कौसल्यामिदमब्रवीत्}
{पायसं प्रतिगृह्णीष्व पुत्रीयं त्विदमात्मनः} %||1-15-24||

\twolineshloka
{कौसल्यायै नरपतिः पायसार्धं ददौ तदा}
{अर्धादर्धं ददौ चापि सुमित्रायै नराधिपः} %||1-15-25||

\twolineshloka
{कैकेय्यै चावशिष्टार्धं ददौ पुत्रार्थकारणात्}
{प्रददौ चावशिष्टार्धं पायसस्यामृतोपमम्} %||1-15-26||

\twolineshloka
{अनुचिन्त्य सुमित्रायै पुनरेव महीपतिः}
{एवं तासां ददौ राजा भार्याणां पायसं पृथक्} %||1-15-27||

\twolineshloka
{तास्त्वेतत्पायसं प्राप्य नरेन्द्रस्योत्तमाः स्त्रियः}
{सम्मानं मेनिरे सर्वाः प्रहर्षोदितचेतसः} %||1-16-28||


{॥इत्यार्षे श्रीमद्रामायणे वाल्मीकीये आदिकाव्ये बालकाण्डे पञ्चदशः सर्गः॥}

\sect{षोडशः सर्गः}

\twolineshloka
{पुत्रत्वं तु गते विष्णौ राज्ञस्तस्य महात्मनः}
{उवाच देवताः सर्वाः स्वयम्भूर्भगवानिदम्} %||1-16-1||

\twolineshloka
{सत्यसन्धस्य वीरस्य सर्वेषां नो हितैषिणः}
{विष्णोः सहायान्बलिनः सृजध्वं कामरूपिणः} %||1-16-2||

\twolineshloka
{मायाविदश्च शूरांश्च वायुवेगसमाञ्जवे}
{नयज्ञान्बुद्धिसम्पन्नान्विष्णुतुल्यपराक्रमान्} %||1-16-3||

\twolineshloka
{असंहार्यानुपायज्ञान्दिव्यसंहननान्वितान्}
{सर्वास्त्रगुणसम्पन्नानमृतप्राशनानिव} %||1-16-4||

\twolineshloka
{अप्सरःसु च मुख्यासु गन्धर्वीणां तनूषु च}
{यक्षपन्नगकन्यासु ऋष्कविद्याधरीषु च} %||1-16-5||

\twolineshloka
{किंनरीणां च गात्रेषु वानरीणां तनूषु च}
{सृजध्वं हरिरूपेण पुत्रांस्तुल्यपराक्रमान्} %||1-16-6||

\twolineshloka
{ते तथोक्ता भगवता तत्प्रतिश्रुत्य शासनम्}
{जनयामासुरेवं ते पुत्रान्वानररूपिणः} %||1-16-7||

\twolineshloka
{ऋषयश्च महात्मानः सिद्धविद्याधरोरगाः}
{चारणाश्च सुतान्वीरान्ससृजुर्वनचारिणः} %||1-16-8||

\twolineshloka
{ते सृष्टा बहुसाहस्रा दशग्रीववधोद्यताः}
{अप्रमेयबला वीरा विक्रान्ताः कामरूपिणः} %||1-16-9||

\twolineshloka
{ते गजाचलसङ्काशा वपुष्मन्तो महाबलाः}
{ऋक्षवानरगोपुच्छाः क्षिप्रमेवाभिजज्ञिरे} %||1-16-10||

\twolineshloka
{यस्य देवस्य यद्रूपं वेषो यश्च पराक्रमः}
{अजायत समस्तेन तस्य तस्य सुतः पृथक्} %||1-16-11||

\twolineshloka
{गोलाङ्गूलीषु चोत्पन्नाः केचित्सम्मतविक्रमाः}
{ऋक्षीषु च तथा जाता वानराः किंनरीषु च} %||1-16-12||

\twolineshloka
{शिलाप्रहरणाः सर्वे सर्वे पादपयोधिनः}
{नखदंष्ट्रायुधाः सर्वे सर्वे सर्वास्त्रकोविदाः} %||1-16-13||

\twolineshloka
{विचालयेयुः शैलेन्द्रान्भेदयेयुः स्थिरान्द्रुमान्}
{क्षोभयेयुश्च वेगेन समुद्रं सरितां पतिम्} %||1-16-14||

\twolineshloka
{दारयेयुः क्षितिं पद्भ्यामाप्लवेयुर्महार्णवम्}
{नभस्तलं विशेयुश्च गृह्णीयुरपि तोयदान्} %||1-16-15||

\twolineshloka
{गृह्णीयुरपि मातङ्गान्मत्तान्प्रव्रजतो वने}
{नर्दमानांश्च नादेन पातयेयुर्विहङ्गमान्} %||1-16-16||

\threelineshloka
{ईदृशानां प्रसूतानि हरीणां कामरूपिणाम्}
{शतं शतसहस्राणि यूथपानां महात्मनाम्}
{बभूवुर्यूथपश्रेष्ठा वीरांश्चाजनयन्हरीन्} %||1-16-17||

\twolineshloka
{अन्ये ऋक्षवतः प्रस्थानुपतस्थुः सहस्रशः}
{अन्ये नानाविधाञ्शैलान्काननानि च भेजिरे} %||1-16-18||

\twolineshloka
{सूर्यपुत्रं च सुग्रीवं शक्रपुत्रं च वालिनम्}
{भ्रातरावुपतस्थुस्ते सर्व एव हरीश्वराः} %||1-16-19||

\fourlineindentedshloka
{तैर्मेघवृन्दाचलतुल्यकायैर्-}
{ महाबलैर्वानरयूथपालैः}
{बभूव भूर्भीमशरीररूपैः}
{ समावृता रामसहायहेतोः} %||1-17-20||


{॥इत्यार्षे श्रीमद्रामायणे वाल्मीकीये आदिकाव्ये बालकाण्डे षोडशः सर्गः॥}

\sect{सप्तदशः सर्गः}

\twolineshloka
{निर्वृत्ते तु क्रतौ तस्मिन्हयमेधे महात्मनः}
{प्रतिगृह्य सुरा भागान्प्रतिजग्मुर्यथागतम्} %||1-17-1||

\twolineshloka
{समाप्तदीक्षानियमः पत्नीगणसमन्वितः}
{प्रविवेश पुरीं राजा सभृत्यबलवाहनः} %||1-17-2||

\twolineshloka
{यथार्हं पूजितास्तेन राज्ञा वै पृथिवीश्वराः}
{मुदिताः प्रययुर्देशान्प्रणम्य मुनिपुङ्गवम्} %||1-17-3||

\twolineshloka
{गतेषु पृथिवीशेषु राजा दशरथः पुनः}
{प्रविवेश पुरीं श्रीमान्पुरस्कृत्य द्विजोत्तमान्} %||1-17-4||

\twolineshloka
{शान्तया प्रययौ सार्धमृष्यशृङ्गः सुपूजितः}
{अन्वीयमानो राज्ञाथ सानुयात्रेण धीमता} %||1-17-5||

\twolineshloka
{कौसल्याजनयद्रामं दिव्यलक्षणसंयुतम्}
{विष्णोरर्धं महाभागं पुत्रमिक्ष्वाकुनन्दनम्} %||1-17-6||

\twolineshloka
{कौसल्या शुशुभे तेन पुत्रेणामिततेजसा}
{यथा वरेण देवानामदितिर्वज्रपाणिना} %||1-17-7||

\twolineshloka
{भरतो नाम कैकेय्यां जज्ञे सत्यपराक्रमः}
{साक्षाद्विष्णोश्चतुर्भागः सर्वैः समुदितो गुणैः} %||1-17-8||

\twolineshloka
{अथ लक्ष्मणशत्रुघ्नौ सुमित्राजनयत्सुतौ}
{वीरौ सर्वास्त्रकुशलौ विष्णोरर्धसमन्वितौ} %||1-17-9||

\twolineshloka
{राज्ञः पुत्रा महात्मानश्चत्वारो जज्ञिरे पृथक्}
{गुणवन्तोऽनुरूपाश्च रुच्या प्रोष्ठपदोपमाः} %||1-17-10||

\twolineshloka
{अतीत्यैकादशाहं तु नाम कर्म तथाकरोत्}
{ज्येष्ठं रामं महात्मानं भरतं कैकयीसुतम्} %||1-17-11||

\threelineshloka
{सौमित्रिं लक्ष्मणमिति शत्रुघ्नमपरं तथा}
{वसिष्ठः परमप्रीतो नामानि कृतवांस्तदा}
{तेषां जन्मक्रियादीनि सर्वकर्माण्यकारयत्} %||1-17-12||

\twolineshloka
{तेषां केतुरिव ज्येष्ठो रामो रतिकरः पितुः}
{बभूव भूयो भूतानां स्वयम्भूरिव सम्मतः} %||1-17-13||

\twolineshloka
{सर्वे वेदविदः शूराः सर्वे लोकहिते रताः}
{सर्वे ज्ञानोपसम्पन्नाः सर्वे समुदिता गुणैः} %||1-17-14||

\twolineshloka
{तेषामपि महातेजा रामः सत्यपराक्रमः}
{बाल्यात्प्रभृति सुस्निग्धो लक्ष्मणो लक्ष्मिवर्धनः} %||1-17-15||

\twolineshloka
{रामस्य लोकरामस्य भ्रातुर्ज्येष्ठस्य नित्यशः}
{सर्वप्रियकरस्तस्य रामस्यापि शरीरतः} %||1-17-16||

\threelineshloka
{लक्ष्मणो लक्ष्मिसम्पन्नो बहिःप्राण इवापरः}
{न च तेन विना निद्रां लभते पुरुषोत्तमः}
{मृष्टमन्नमुपानीतमश्नाति न हि तं विना} %||1-17-17||

\twolineshloka
{यदा हि हयमारूढो मृगयां याति राघवः}
{तदैनं पृष्ठतोऽभ्येति सधनुः परिपालयन्} %||1-17-18||

\twolineshloka
{भरतस्यापि शत्रुघ्नो लक्ष्मणावरजो हि सः}
{प्राणैः प्रियतरो नित्यं तस्य चासीत्तथा प्रियः} %||1-17-19||

\twolineshloka
{स चतुर्भिर्महाभागैः पुत्रैर्दशरथः प्रियैः}
{बभूव परमप्रीतो देवैरिव पितामहः} %||1-17-20||

\twolineshloka
{ते यदा ज्ञानसम्पन्नाः सर्वे समुदिता गुणैः}
{ह्रीमन्तः कीर्तिमन्तश्च सर्वज्ञा दीर्घदर्शिनः} %||1-17-21||

\twolineshloka
{अथ राजा दशरथस्तेषां दारक्रियां प्रति}
{चिन्तयामास धर्मात्मा सोपाध्यायः सबान्धवः} %||1-17-22||

\twolineshloka
{तस्य चिन्तयमानस्य मन्त्रिमध्ये महात्मनः}
{अभ्यागच्छन्महातेजो विश्वामित्रो महामुनिः} %||1-17-23||

\twolineshloka
{स राज्ञो दर्शनाकाङ्क्षी द्वाराध्यक्षानुवाच ह}
{शीघ्रमाख्यात मां प्राप्तं कौशिकं गाधिनः सुतम्} %||1-17-24||

\twolineshloka
{तच्छ्रुत्वा वचनं तस्य राजवेश्म प्रदुद्रुवुः}
{सम्भ्रान्तमनसः सर्वे तेन वाक्येन चोदिताः} %||1-17-25||

\twolineshloka
{ते गत्वा राजभवनं विश्वामित्रमृषिं तदा}
{प्राप्तमावेदयामासुर्नृपायेक्ष्वाकवे तदा} %||1-17-26||

\twolineshloka
{तेषां तद्वचनं श्रुत्वा सपुरोधाः समाहितः}
{प्रत्युज्जगाम संहृष्टो ब्रह्माणमिव वासवः} %||1-17-27||

\twolineshloka
{स दृष्ट्वा ज्वलितं दीप्त्या तापसं संशितव्रतम्}
{प्रहृष्टवदनो राजा ततोऽर्घ्यमुपहारयत्} %||1-17-28||

\twolineshloka
{स राज्ञः प्रतिगृह्यार्घ्यं शास्त्रदृष्ट्तेन कर्मणा}
{कुशलं चाव्ययं चैव पर्यपृच्छन्नराधिपम्} %||1-17-29||

\twolineshloka
{वसिष्ठं च समागम्य कुशलं मुनिपुङ्गवः}
{ऋषींश्च तान्यथा न्यायं महाभागानुवाच ह} %||1-17-30||

\twolineshloka
{ते सर्वे हृष्टमनसस्तस्य राज्ञो निवेशनम्}
{विविशुः पूजितास्तत्र निषेदुश्च यथार्थतः} %||1-17-31||

\twolineshloka
{अथ हृष्टमना राजा विश्वामित्रं महामुनिम्}
{उवाच परमोदारो हृष्टस्तमभिपूजयन्} %||1-17-32||

\fourlineshloka
{यथामृतस्य सम्प्राप्तिर्यथा वर्षमनूदके}
{यथा सदृशदारेषु पुत्रजन्माप्रजस्य च}
{प्रनष्टस्य यथा लाभो यथा हर्षो महोदये}
{तथैवागमनं मन्ये स्वागतं ते महामुने} %||1-17-33||

\threelineshloka
{कं च ते परमं कामं करोमि किमु हर्षितः}
{पात्रभूतोऽसि मे विप्र दिष्ट्या प्राप्तोऽसि धार्मिक}
{अद्य मे सफलं जन्म जीवितं च सुजीवितम्} %||1-17-34||

\twolineshloka
{पूर्वं राजर्षिशब्देन तपसा द्योतितप्रभः}
{ब्रह्मर्षित्वमनुप्राप्तः पूज्योऽसि बहुधा मया} %||1-17-35||

\twolineshloka
{तदद्भुतमिदं विप्र पवित्रं परमं मम}
{शुभक्षेत्रगतश्चाहं तव सन्दर्शनात्प्रभो} %||1-17-36||

\twolineshloka
{ब्रूहि यत्प्रार्थितं तुभ्यं कार्यमागमनं प्रति}
{इच्छाम्यनुगृहीतोऽहं त्वदर्थपरिवृद्धये} %||1-17-37||

\twolineshloka
{कार्यस्य न विमर्शं च गन्तुमर्हसि कौशिक}
{कर्ता चाहमशेषेण दैवतं हि भवान्मम} %||1-17-38||

\fourlineindentedshloka
{इति हृदयसुखं निशम्य वाक्यम्}
{ श्रुतिसुखमात्मवता विनीतमुक्तम्}
{प्रथितगुणयशा गुणैर्विशिष्टः}
{ परम ऋषिः परमं जगाम हर्षम्} %||1-18-39||


{॥इत्यार्षे श्रीमद्रामायणे वाल्मीकीये आदिकाव्ये बालकाण्डे सप्तदशः सर्गः॥}

\sect{अष्टादशः सर्गः}

\twolineshloka
{तच्छ्रुत्वा राजसिंहस्य वाक्यमद्भुतविस्तरम्}
{हृष्टरोमा महातेजा विश्वामित्रोऽभ्यभाषत} %||1-18-1||

\twolineshloka
{सदृशं राजशार्दूल तवैतद्भुवि नान्यतः}
{महावंशप्रसूतस्य वसिष्ठव्यपदेशिनः} %||1-18-2||

\twolineshloka
{यत्तु मे हृद्गतं वाक्यं तस्य कार्यस्य निश्चयम्}
{कुरुष्व राजशार्दूल भव सत्यप्रतिश्रवः} %||1-18-3||

\twolineshloka
{अहं नियममातिष्ठ सिद्ध्यर्थं पुरुषर्षभ}
{तस्य विघ्नकरौ द्वौ तु राक्षसौ कामरूपिणौ} %||1-18-4||

\threelineshloka
{व्रते मे बहुशश्चीर्णे समाप्त्यां राक्षसाविमौ}
{मारीचश्च सुबाहुश्च वीर्यवन्तौ सुशिक्षितौ}
{तौ मांसरुधिरौघेण वेदिं तामभ्यवर्षताम्} %||1-18-5||

\twolineshloka
{अवधूते तथा भूते तस्मिन्नियमनिश्चये}
{कृतश्रमो निरुत्साहस्तस्माद्देशादपाक्रमे} %||1-18-6||

\twolineshloka
{न च मे क्रोधमुत्स्रष्टुं बुद्धिर्भवति पार्थिव}
{तथाभूता हि सा चर्या न शापस्तत्र मुच्यते} %||1-18-7||

\twolineshloka
{स्वपुत्रं राजशार्दूल रामं सत्यपराक्रमम्}
{काकपक्षधरं शूरं ज्येष्ठं मे दातुमर्हसि} %||1-18-8||

\twolineshloka
{शक्तो ह्येष मया गुप्तो दिव्येन स्वेन तेजसा}
{राक्षसा ये विकर्तारस्तेषामपि विनाशने} %||1-18-9||

\twolineshloka
{श्रेयश्चास्मै प्रदास्यामि बहुरूपं न संशयः}
{त्रयाणामपि लोकानां येन ख्यातिं गमिष्यति} %||1-18-10||

\twolineshloka
{न च तौ राममासाद्य शक्तौ स्थातुं कथञ्चन}
{न च तौ राघवादन्यो हन्तुमुत्सहते पुमान्} %||1-18-11||

\twolineshloka
{वीर्योत्सिक्तौ हि तौ पापौ कालपाशवशं गतौ}
{रामस्य राजशार्दूल न पर्याप्तौ महात्मनः} %||1-18-12||

\twolineshloka
{न च पुत्रकृतं स्नेहं कर्तुमर्हसि पार्थिव}
{अहं ते प्रतिजानामि हतौ तौ विद्धि राक्षसौ} %||1-18-13||

\twolineshloka
{अहं वेद्मि महात्मानं रामं सत्यपराक्रमम्}
{वसिष्ठोऽपि महातेजा ये चेमे तपसि स्थिताः} %||1-18-14||

\twolineshloka
{यदि ते धर्मलाभं च यशश्च परमं भुवि}
{स्थिरमिच्छसि राजेन्द्र रामं मे दातुमर्हसि} %||1-18-15||

\twolineshloka
{यद्यभ्यनुज्ञां काकुत्स्थ ददते तव मन्त्रिणः}
{वसिष्ठ प्रमुखाः सर्वे ततो रामं विसर्जय} %||1-18-16||

\twolineshloka
{अभिप्रेतमसंसक्तमात्मजं दातुमर्हसि}
{दशरात्रं हि यज्ञस्य रामं राजीवलोचनम्} %||1-18-17||

\twolineshloka
{नात्येति कालो यज्ञस्य यथायं मम राघव}
{तथा कुरुष्व भद्रं ते मा च शोके मनः कृथाः} %||1-18-18||

\twolineshloka
{इत्येवमुक्त्वा धर्मात्मा धर्मार्थसहितं वचः}
{विरराम महातेजा विश्वामित्रो महामुनिः} %||1-18-19||

\fourlineindentedshloka
{इति हृदयमनोविदारणम्}
{ मुनिवचनं तदतीव शुश्रुवान्}
{नरपतिरगमद्भयं मह}
{द्व्यथितमनाः प्रचचाल चासनात्} %||1-19-20||


{॥इत्यार्षे श्रीमद्रामायणे वाल्मीकीये आदिकाव्ये बालकाण्डे अष्टादशः सर्गः॥}

\sect{एकोनविंशः सर्गः}

\twolineshloka
{तच्छ्रुत्वा राजशार्दूल विश्वामित्रस्य भाषितम्}
{मुहूर्तमिव निःसंज्ञः संज्ञावानिदमब्रवीत्} %||1-19-1||

\twolineshloka
{ऊनषोडशवर्षो मे रामो राजीवलोचनः}
{न युद्धयोग्यतामस्य पश्यामि सह राक्षसैः} %||1-19-2||

\twolineshloka
{इयमक्षौहिणी पूर्णा यस्याहं पतिरीश्वरः}
{अनया संवृतो गत्वा योधाहं तैर्निशाचरैः} %||1-19-3||

\twolineshloka
{इमे शूराश्च विक्रान्ता भृत्या मेऽस्त्रविशारदाः}
{योग्या रक्षोगणैर्योद्धुं न रामं नेतुमर्हसि} %||1-19-4||

\twolineshloka
{अहमेव धनुष्पाणिर्गोप्ता समरमूर्धनि}
{यावत्प्राणान्धरिष्यामि तावद्योत्स्ये निशाचरैः} %||1-19-5||

\twolineshloka
{निर्विघ्ना व्रतचर्या सा भविष्यति सुरक्षिता}
{अहं तत्र गमिष्यामि न राम नेतुमर्हसि} %||1-19-6||

\threelineshloka
{बालो ह्यकृतविद्यश्च न च वेत्ति बलाबलम्}
{न चास्त्रबलसंयुक्तो न च युद्धविशारदः}
{न चासौ रक्षसां योग्यः कूटयुद्धा हि ते ध्रुवम्} %||1-19-7||

\twolineshloka
{विप्रयुक्तो हि रामेण मुहूर्तमपि नोत्सहे}
{जीवितुं मुनिशार्दूल न रामं नेतुमर्हसि} %||1-19-8||

\twolineshloka
{यदि वा राघवं ब्रह्मन्नेतुमिच्छसि सुव्रत}
{चतुरङ्गसमायुक्तं मया सह च तं नय} %||1-19-9||

\twolineshloka
{षष्टिर्वर्षसहस्राणि जातस्य मम कौशिक}
{दुःखेनोत्पादितश्चायं न रामं नेतुमर्हसि} %||1-19-10||

\twolineshloka
{चतुर्णामात्मजानां हि प्रीतिः परमिका मम}
{ज्येष्ठं धर्मप्रधानं च न रामं नेतुमर्हसि} %||1-19-11||

\twolineshloka
{किं वीर्या राक्षसास्ते च कस्य पुत्राश्च के च ते}
{कथं प्रमाणाः के चैतान्रक्षन्ति मुनिपुङ्गव} %||1-19-12||

\twolineshloka
{कथं च प्रतिकर्तव्यं तेषां रामेण रक्षसाम्}
{मामकैर्वा बलैर्ब्रह्मन्मया वा कूटयोधिनाम्} %||1-19-13||

\twolineshloka
{सर्वं मे शंस भगवन्कथं तेषां मया रणे}
{स्थातव्यं दुष्टभावानां वीर्योत्सिक्ता हि राक्षसाः} %||1-19-14||

\twolineshloka
{तस्य तद्वचनं श्रुत्वा विश्वामित्रोऽभ्यभाषत}
{पौलस्त्यवंशप्रभवो रावणो नाम राक्षसः} %||1-19-15||

\twolineshloka
{स ब्रह्मणा दत्तवरस्त्रैलोक्यं बाधते भृशम्}
{महाबलो महावीर्यो राक्षसैर्बहुभिर्वृतः} %||1-19-16||

\twolineshloka
{श्रूयते हि महावीर्यो रावणो राक्षसाधिपः}
{साक्षाद्वैश्रवणभ्राता पुत्रो विश्रवसो मुनेः} %||1-19-17||

\threelineshloka
{यदा स्वयं न यज्ञस्य विघ्नकर्ता महाबलः}
{तेन सञ्चोदितौ तौ तु राक्षसौ सुमहा बलौ}
{मारीचश्च सुबाहुश्च यज्ञविघ्नं करिष्यतः} %||1-19-18||

\twolineshloka
{इत्युक्तो मुनिना तेन राजोवाच मुनिं तदा}
{न हि शक्तोऽस्मि सङ्ग्रामे स्थातुं तस्य दुरात्मनः} %||1-19-19||

\twolineshloka
{स त्वं प्रसादं धर्मज्ञ कुरुष्व मम पुत्रके}
{देवदानवगन्धर्वा यक्षाः पतग पन्नगाः} %||1-19-20||

\twolineshloka
{न शक्ता रावणं सोढुं किं पुनर्मानवा युधि}
{स हि वीर्यवतां वीर्यमादत्ते युधि राक्षसः} %||1-19-21||

\twolineshloka
{तेन चाहं न शक्तोऽस्मि संयोद्धुं तस्य वा बलैः}
{सबलो वा मुनिश्रेष्ठ सहितो वा ममात्मजैः} %||1-19-22||

\twolineshloka
{कथमप्यमरप्रख्यं सङ्ग्रामाणामकोविदम्}
{बालं मे तनयं ब्रह्मन्नैव दास्यामि पुत्रकम्} %||1-19-23||

\twolineshloka
{अथ कालोपमौ युद्धे सुतौ सुन्दोपसुन्दयोः}
{यज्ञविघ्नकरौ तौ ते नैव दास्यामि पुत्रकम्} %||1-19-24||

\twolineshloka
{मारीचश्च सुबाहुश्च वीर्यवन्तौ सुशिक्षितौ}
{तयोरन्यतरेणाहं योद्धा स्यां ससुहृद्गणः} %||1-20-25||


{॥इत्यार्षे श्रीमद्रामायणे वाल्मीकीये आदिकाव्ये बालकाण्डे एकोनविंशः सर्गः॥}

\sect{विंशः सर्गः}

\twolineshloka
{तच्छ्रुत्वा वचनं तस्य स्नेहपर्याकुलाक्षरम्}
{समन्युः कौशिको वाक्यं प्रत्युवच महीपतिम्} %||1-20-1||

\twolineshloka
{पूर्वमर्थं प्रतिश्रुत्य प्रतिज्ञां हातुमिच्छसि}
{रागवाणामयुक्तोऽयं कुलस्यास्य विपर्ययः} %||1-20-2||

\twolineshloka
{यदिदं ते क्षमं राजन्गमिष्यामि यथागतम्}
{मिथ्याप्रतिज्ञः काकुत्स्थ सुखी भव सबान्धवः} %||1-20-3||

\twolineshloka
{तस्य रोषपरीतस्य विश्वामित्रस्य धीमतः}
{चचाल वसुधा कृत्स्ना विवेश च भयं सुरान्} %||1-20-4||

\twolineshloka
{त्रस्तरूपं तु विज्ञाय जगत्सर्वं महानृषिः}
{नृपतिं सुव्रतो धीरो वसिष्ठो वाक्यमब्रवीत्} %||1-20-5||

\twolineshloka
{इक्ष्वाकूणां कुले जातः साक्षाद्धर्म इवापरः}
{धृतिमान्सुव्रतः श्रीमान्न धर्मं हातुमर्हसि} %||1-20-6||

\twolineshloka
{त्रिषु लोकेषु विख्यातो धर्मात्मा इति राघवः}
{स्वधर्मं प्रतिपद्यस्व नाधर्मं वोढुमर्हसि} %||1-20-7||

\twolineshloka
{संश्रुत्यैवं करिष्यामीत्यकुर्वाणस्य राघव}
{इष्टापूर्तवधो भूयात्तस्माद्रामं विसर्जय} %||1-20-8||

\twolineshloka
{कृतास्त्रमकृतास्त्रं वा नैनं शक्ष्यन्ति राक्षसाः}
{गुप्तं कुशिकपुत्रेण ज्वलनेनामृतं यथा} %||1-20-9||

\twolineshloka
{एष विग्रहवान्धर्म एष वीर्यवतां वरः}
{एष बुद्ध्याधिको लोके तपसश्च परायणम्} %||1-20-10||

\twolineshloka
{एषोऽस्त्रान्विविधान्वेत्ति त्रैलोक्ये सचराचरे}
{नैनमन्यः पुमान्वेत्ति न च वेत्स्यन्ति केचन} %||1-20-11||

\twolineshloka
{न देवा नर्षयः केचिन्नासुरा न च राक्षसाः}
{गन्धर्वयक्षप्रवराः सकिंनरमहोरगाः} %||1-20-12||

\twolineshloka
{सर्वास्त्राणि कृशाश्वस्य पुत्राः परमधार्मिकाः}
{कौशिकाय पुरा दत्ता यदा राज्यं प्रशासति} %||1-20-13||

\twolineshloka
{तेऽपि पुत्राः कृशाश्वस्य प्रजापतिसुतासुताः}
{नकरूपा महावीर्या दीप्तिमन्तो जयावहाः} %||1-20-14||

\twolineshloka
{जया च सुप्रभा चैव दक्षकन्ये सुमध्यमे}
{ते सुवातेऽस्त्रशस्त्राणि शतं परम भास्वरम्} %||1-20-15||

\twolineshloka
{पञ्चाशतं सुताँल्लेभे जया नाम वरान्पुरा}
{वधायासुरसैन्यानाममेयान्कामरूपिणः} %||1-20-16||

\twolineshloka
{सुप्रभाजनयच्चापि पुत्रान्पञ्चाशतं पुनः}
{संहारान्नाम दुर्धर्षान्दुराक्रामान्बलीयसः} %||1-20-17||

\twolineshloka
{तानि चास्त्राणि वेत्त्येष यथावत्कुशिकात्मजः}
{अपूर्वाणां च जनने शक्तो भूयश्च धर्मवित्} %||1-20-18||

\twolineshloka
{एवं वीर्यो महातेजा विश्वामित्रो महातपाः}
{न रामगमने राजन्संशयं गन्तुमर्हसि} %||1-21-19||


{॥इत्यार्षे श्रीमद्रामायणे वाल्मीकीये आदिकाव्ये बालकाण्डे विंशः सर्गः॥}

\sect{एकविंशः सर्गः}

\twolineshloka
{तथा वसिष्ठे ब्रुवति राजा दशरथः सुतम्}
{प्रहृष्टवदनो राममाजुहाव सलक्ष्मणम्} %||1-21-1||

\twolineshloka
{कृतस्वस्त्ययनं मात्रा पित्रा दशरथेन च}
{पुरोधसा वसिष्ठेन मङ्गलैरभिमन्त्रितम्} %||1-21-2||

\twolineshloka
{स पुत्रं मूर्ध्न्युपाघ्राय राजा दशरथः प्रियम्}
{ददौ कुशिकपुत्राय सुप्रीतेनान्तरात्मना} %||1-21-3||

\twolineshloka
{ततो वायुः सुखस्पर्शो विरजस्को ववौ तदा}
{विश्वामित्रगतं रामं दृष्ट्वा राजीवलोचनम्} %||1-21-4||

\twolineshloka
{पुष्पवृष्टिर्महत्यासीद्देवदुन्दुभिनिस्वनः}
{शङ्खदुन्दुभिनिर्घोषः प्रयाते तु महात्मनि} %||1-21-5||

\twolineshloka
{विश्वामित्रो ययावग्रे ततो रामो महायशाः}
{काकपक्षधरो धन्वी तं च सौमित्रिरन्वगात्} %||1-21-6||

\threelineshloka
{कलापिनौ धनुष्पाणी शोभयानौ दिशो दश}
{विश्वामित्रं महात्मानं त्रिशीर्षाविव पन्नगौ}
{अनुजग्मतुरक्षुद्रौ पितामहमिवाश्विनौ} %||1-21-7||

\twolineshloka
{बद्धगोधाङ्गुलित्राणौ खड्गवन्तौ महाद्युती}
{स्थाणुं देवमिवाचिन्त्यं कुमाराविव पावकी} %||1-21-8||

\twolineshloka
{अध्यर्धयोजनं गत्वा सरय्वा दक्षिणे तटे}
{रामेति मधुरा वाणीं विश्वामित्रोऽभ्यभाषत} %||1-21-9||

\twolineshloka
{गृहाण वत्स सलिलं मा भूत्कालस्य पर्ययः}
{मन्त्रग्रामं गृहाण त्वं बलामतिबलां तथा} %||1-21-10||

\twolineshloka
{न श्रमो न ज्वरो वा ते न रूपस्य विपर्ययः}
{न च सुप्तं प्रमत्तं वा धर्षयिष्यन्ति नैरृताः} %||1-21-11||

\twolineshloka
{न बाह्वोः सदृशो वीर्ये पृथिव्यामस्ति कश्चन}
{त्रिषु लोकेषु वा राम न भवेत्सदृशस्तव} %||1-21-12||

\twolineshloka
{न सौभाग्ये न दाक्षिण्ये न ज्ञाने बुद्धिनिश्चये}
{नोत्तरे प्रतिपत्तव्यो समो लोके तवानघ} %||1-21-13||

\twolineshloka
{एतद्विद्याद्वये लब्धे भविता नास्ति ते समः}
{बला चातिबला चैव सर्वज्ञानस्य मातरौ} %||1-21-14||

\threelineshloka
{क्षुत्पिपासे न ते राम भविष्येते नरोत्तम}
{बलामतिबलां चैव पठतः पथि राघव}
{विद्याद्वयमधीयाने यशश्चाप्यतुलं भुवि} %||1-21-15||

\twolineshloka
{पितामहसुते ह्येते विद्ये तेजःसमन्विते}
{प्रदातुं तव काकुत्स्थ सदृशस्त्वं हि धार्मिक} %||1-21-16||

\twolineshloka
{कामं बहुगुणाः सर्वे त्वय्येते नात्र संशयः}
{तपसा सम्भृते चैते बहुरूपे भविष्यतः} %||1-21-17||

\threelineshloka
{ततो रामो जलं स्पृष्ट्वा प्रहृष्टवदनः शुचिः}
{प्रतिजग्राह ते विद्ये महर्षेर्भावितात्मनः}
{विद्यासमुदितो रामः शुशुभे भूरिविक्रमः} %||1-21-18||

\twolineshloka
{गुरुकार्याणि सर्वाणि नियुज्य कुशिकात्मजे}
{ऊषुस्तां रजनीं तत्र सरय्वां सुसुखं त्रयः} %||1-22-19||


{॥इत्यार्षे श्रीमद्रामायणे वाल्मीकीये आदिकाव्ये बालकाण्डे एकविंशः सर्गः॥}

\sect{द्वाविंशः सर्गः}

\twolineshloka
{प्रभातायां तु शर्वर्यां विश्वामित्रो महामुनिः}
{अभ्यभाषत काकुत्स्थं शयानं पर्णसंस्तरे} %||1-22-1||

\twolineshloka
{कौसल्या सुप्रजा राम पूर्वा सन्ध्या प्रवर्तते}
{उत्तिष्ठ नरशार्दूल कर्तव्यं दैवमाह्निकम्} %||1-22-2||

\twolineshloka
{तस्यर्षेः परमोदारं वचः श्रुत्वा नृपात्मजौ}
{स्नात्वा कृतोदकौ वीरौ जेपतुः परमं जपम्} %||1-22-3||

\twolineshloka
{कृताह्निकौ महावीर्यौ विश्वामित्रं तपोधनम्}
{अभिवाद्याभिसंहृष्टौ गमनायोपतस्थतुः} %||1-22-4||

\twolineshloka
{तौ प्रयाते महावीर्यौ दिव्यं त्रिपथगां नदीम्}
{ददृशाते ततस्तत्र सरय्वाः सङ्गमे शुभे} %||1-22-5||

\twolineshloka
{तत्राश्रमपदं पुण्यमृषीणामुग्रतेजसाम्}
{बहुवर्षसहस्राणि तप्यतां परमं तपः} %||1-22-6||

\twolineshloka
{तं दृष्ट्वा परमप्रीतौ राघवौ पुण्यमाश्रमम्}
{ऊचतुस्तं महात्मानं विश्वामित्रमिदं वचः} %||1-22-7||

\twolineshloka
{कस्यायमाश्रमः पुण्यः को न्वस्मिन्वसते पुमान्}
{भगवञ्श्रोतुमिच्छावः परं कौतूहलं हि नौ} %||1-22-8||

\twolineshloka
{तयोस्तद्वचनं श्रुत्वा प्रहस्य मुनिपुङ्गवः}
{अब्रवीच्छ्रूयतां राम यस्यायं पूर्व आश्रमः} %||1-22-9||

{कन्दर्पो मूर्तिमानासीत्काम इत्युच्यते बुधैः॥\devanumber{10}॥} %||1-22-10|| (Check)

\addtocounter{shlokacount}{1}

\threelineshloka
{तपस्यन्तमिह स्थाणुं नियमेन समाहितम्}
{कृतोद्वाहं तु देवेशं गच्छन्तं समरुद्गणम्}
{धर्षयामास दुर्मेधा हुङ्कृतश्च महात्मना} %||1-22-11||

\twolineshloka
{दग्धस्य तस्य रौद्रेण चक्षुषा रघुनन्दन}
{व्यशीर्यन्त शरीरात्स्वात्सर्वगात्राणि दुर्मतेः} %||1-22-12||

\twolineshloka
{तस्य गात्रं हतं तत्र निर्दग्धस्य महात्मना}
{अशरीरः कृतः कामः क्रोधाद्देवेश्वरेण ह} %||1-22-13||

\twolineshloka
{अनङ्ग इति विख्यातस्तदा प्रभृति राघव}
{स चाङ्गविषयः श्रीमान्यत्राङ्गं स मुमोच ह} %||1-22-14||

\twolineshloka
{तस्यायमाश्रमः पुण्यस्तस्येमे मुनयः पुरा}
{शिष्या धर्मपरा वीर तेषां पापं न विद्यते} %||1-22-15||

\twolineshloka
{इहाद्य रजनीं राम वसेम शुभदर्शन}
{पुण्ययोः सरितोर्मध्ये श्वस्तरिष्यामहे वयम्} %||1-22-16||

\twolineshloka
{तेषां संवदतां तत्र तपो दीर्घेण चक्षुषा}
{विज्ञाय परमप्रीता मुनयो हर्षमागमन्} %||1-22-17||

\twolineshloka
{अर्घ्यं पाद्यं तथातिथ्यं निवेद्यकुशिकात्मजे}
{रामलक्ष्मणयोः पश्चादकुर्वन्नतिथिक्रियाम्} %||1-22-18||

\twolineshloka
{सत्कारं समनुप्राप्य कथाभिरभिरञ्जयन्}
{न्यवसन्सुसुखं तत्र कामाश्रमपदे तदा} %||1-23-19||


{॥इत्यार्षे श्रीमद्रामायणे वाल्मीकीये आदिकाव्ये बालकाण्डे द्वाविंशः सर्गः॥}

\sect{त्रयोविंशः सर्गः}

\twolineshloka
{ततः प्रभाते विमले कृताह्निकमरिन्दमौ}
{विश्वामित्रं पुरस्कृत्य नद्यास्तीरमुपागतौ} %||1-23-1||

\twolineshloka
{ते च सर्वे महात्मानो मुनयः संशितव्रताः}
{उपस्थाप्य शुभां नावं विश्वामित्रमथाब्रुवन्} %||1-23-2||

\twolineshloka
{आरोहतु भवान्नावं राजपुत्रपुरस्कृतः}
{अरिष्टं गच्छ पन्थानं मा भूत्कालस्य पर्ययः} %||1-23-3||

\twolineshloka
{विश्वामित्रस्तथेत्युक्त्वा तानृषीनभिपूज्य च}
{ततार सहितस्ताभ्यां सरितं सागरं गमाम्} %||1-23-4||

\twolineshloka
{अथ रामः सरिन्मध्ये पप्रच्छ मुनिपुङ्गवम्}
{वारिणो भिद्यमानस्य किमयं तुमुलो ध्वनिः} %||1-23-5||

\twolineshloka
{राघवस्य वचः श्रुत्वा कौतूहल समन्वितम्}
{कथयामास धर्मात्मा तस्य शब्दस्य निश्चयम्} %||1-23-6||

\twolineshloka
{कैलासपर्वते राम मनसा निर्मितं सरः}
{ब्रह्मणा नरशार्दूल तेनेदं मानसं सरः} %||1-23-7||

\twolineshloka
{तस्मात्सुस्राव सरसः सायोध्यामुपगूहते}
{सरःप्रवृत्ता सरयूः पुण्या ब्रह्मसरश्च्युता} %||1-23-8||

\twolineshloka
{तस्यायमतुलः शब्दो जाह्नवीमभिवर्तते}
{वारिसङ्क्षोभजो राम प्रणामं नियतः कुरु} %||1-23-9||

\twolineshloka
{ताभ्यां तु तावुभौ कृत्वा प्रणाममतिधार्मिकौ}
{तीरं दक्षिणमासाद्य जग्मतुर्लघुविक्रमौ} %||1-23-10||

\twolineshloka
{स वनं घोरसङ्काशं दृष्ट्वा नृपवरात्मजः}
{अविप्रहतमैक्ष्वाकः पप्रच्छ मुनिपुङ्गवम्} %||1-23-11||

\twolineshloka
{अहो वनमिदं दुर्गं झिल्लिकागणनादितम्}
{भैरवैः श्वापदैः कीर्णं शकुन्तैर्दारुणारवैः} %||1-23-12||

\twolineshloka
{नानाप्रकारैः शकुनैर्वाश्यद्भिर्भैरवस्वनैः}
{सिंहव्याघ्रवराहैश्च वारणैश्चापि शोभितम्} %||1-23-13||

\twolineshloka
{धवाश्वकर्णककुभैर्बिल्वतिन्दुकपाटलैः}
{सङ्कीर्णं बदरीभिश्च किं न्विदं दारुणं वनम्} %||1-23-14||

\twolineshloka
{तमुवाच महातेजा विश्वामित्रो महामुनिः}
{श्रूयतां वत्स काकुत्स्थ यस्यैतद्दारुणं वनम्} %||1-23-15||

\twolineshloka
{एतौ जनपदौ स्फीतौ पूर्वमास्तां नरोत्तम}
{मलदाश्च करूषाश्च देवनिर्माण निर्मितौ} %||1-23-16||

\twolineshloka
{पुरा वृत्रवधे राम मलेन समभिप्लुतम्}
{क्षुधा चैव सहस्राक्षं ब्रह्महत्या यदाविशत्} %||1-23-17||

\twolineshloka
{तमिन्द्रं स्नापयन्देवा ऋषयश्च तपोधनाः}
{कलशैः स्नापयामासुर्मलं चास्य प्रमोचयन्} %||1-23-18||

\twolineshloka
{इह भूम्यां मलं दत्त्वा दत्त्वा कारुषमेव च}
{शरीरजं महेन्द्रस्य ततो हर्षं प्रपेदिरे} %||1-23-19||

\twolineshloka
{निर्मलो निष्करूषश्च शुचिरिन्द्रो यदाभवत्}
{ददौ देशस्य सुप्रीतो वरं प्रभुरनुत्तमम्} %||1-23-20||

\twolineshloka
{इमौ जनपदौ स्थीतौ ख्यातिं लोके गमिष्यतः}
{मलदाश्च करूषाश्च ममाङ्गमलधारिणौ} %||1-23-21||

\twolineshloka
{साधु साध्विति तं देवाः पाकशासनमब्रुवन्}
{देशस्य पूजां तां दृष्ट्वा कृतां शक्रेण धीमता} %||1-23-22||

\twolineshloka
{एतौ जनपदौ स्थीतौ दीर्घकालमरिन्दम}
{मलदाश्च करूषाश्च मुदितौ धनधान्यतः} %||1-23-23||

\twolineshloka
{कस्यचित्त्वथ कालस्य यक्षी वै कामरूपिणी}
{बलं नागसहस्रस्य धारयन्ती तदा ह्यभूत्} %||1-23-24||

\twolineshloka
{ताटका नाम भद्रं ते भार्या सुन्दस्य धीमतः}
{मारीचो राक्षसः पुत्रो यस्याः शक्रपराक्रमः} %||1-23-25||

\twolineshloka
{इमौ जनपदौ नित्यं विनाशयति राघव}
{मलदांश्च करूषांश्च ताटका दुष्टचारिणी} %||1-23-26||

\twolineshloka
{सेयं पन्थानमावार्य वसत्यत्यर्धयोजने}
{अत एव च गन्तव्यं ताटकाया वनं यतः} %||1-23-27||

\twolineshloka
{स्वबाहुबलमाश्रित्य जहीमां दुष्टचारिणीम्}
{मन्नियोगादिमं देशं कुरु निष्कण्टकं पुनः} %||1-23-28||

\twolineshloka
{न हि कश्चिदिमं देशं शक्रोत्यागन्तुमीदृशम्}
{यक्षिण्या घोरया राम उत्सादितमसह्यया} %||1-23-29||

\twolineshloka
{एतत्ते सर्वमाख्यातं यथैतद्दरुणं वनम्}
{यक्ष्या चोत्सादितं सर्वमद्यापि न निवर्तते} %||1-24-30||


{॥इत्यार्षे श्रीमद्रामायणे वाल्मीकीये आदिकाव्ये बालकाण्डे त्रयोविंशः सर्गः॥}

\sect{चतुर्विंशः सर्गः}

\twolineshloka
{अथ तस्याप्रमेयस्य मुनेर्वचनमुत्तमम्}
{श्रुत्वा पुरुषशार्दूलः प्रत्युवाच शुभां गिरम्} %||1-24-1||

\twolineshloka
{अल्पवीर्या यदा यक्षाः श्रूयन्ते मुनिपुङ्गव}
{कथं नागसहस्रस्य धारयत्यबला बलम्} %||1-24-2||

\twolineshloka
{विश्वामित्रोऽब्रवीद्वाक्यं शृणु येन बलोत्तरा}
{वरदानकृतं वीर्यं धारयत्यबला बलम्} %||1-24-3||

\twolineshloka
{पूर्वमासीन्महायक्षः सुकेतुर्नाम वीर्यवान्}
{अनपत्यः शुभाचारः स च तेपे महत्तपः} %||1-24-4||

\twolineshloka
{पितामहस्तु सुप्रीतस्तस्य यक्षपतेस्तदा}
{कन्यारत्नं ददौ राम ताटकां नाम नामतः} %||1-24-5||

\twolineshloka
{ददौ नागसहस्रस्य बलं चास्याः पितामहः}
{न त्वेव पुत्रं यक्षाय ददौ ब्रह्मा महायशाः} %||1-24-6||

\twolineshloka
{तां तु जातां विवर्धन्तीं रूपयौवनशालिनीम्}
{जम्भपुत्राय सुन्दाय ददौ भार्यां यशस्विनीम्} %||1-24-7||

\twolineshloka
{कस्यचित्त्वथ कालल्स्य यक्षी पुत्रं व्यजायत}
{मारीचं नाम दुर्धर्षं यः शापाद्राक्षसोऽभवत्} %||1-24-8||

\twolineshloka
{सुन्दे तु निहते राम अगस्त्यमृषिसत्तमम्}
{ताटका सह पुत्रेण प्रधर्षयितुमिच्छति} %||1-24-9||

\twolineshloka
{राक्षसत्वं भजस्वेति मारीचं व्याजहार सः}
{अगस्त्यः परमक्रुद्धस्ताटकामपि शप्तवान्} %||1-24-10||

\twolineshloka
{पुरुषादी महायक्षी विरूपा विकृतानना}
{इदं रूपमपहाय दारुणं रूपमस्तु ते} %||1-24-11||

\twolineshloka
{सैषा शापकृतामर्षा ताटका क्रोधमूर्छिता}
{देशमुत्सादयत्येनमगस्त्यचरितं शुभम्} %||1-24-12||

\twolineshloka
{एनां राघव दुर्वृत्तां यक्षीं परमदारुणाम्}
{गोब्राह्मणहितार्थाय जहि दुष्टपराक्रमाम्} %||1-24-13||

\twolineshloka
{न ह्येनां शापसंसृष्टां कश्चिदुत्सहते पुमान्}
{निहन्तुं त्रिषु लोकेषु त्वामृते रघुनन्दन} %||1-24-14||

\twolineshloka
{न हि ते स्त्रीवधकृते घृणा कार्या नरोत्तम}
{चातुर्वर्ण्यहितार्थाय कर्तव्यं राजसूनुना} %||1-24-15||

\twolineshloka
{राज्यभारनियुक्तानामेष धर्मः सनातनः}
{अधर्म्यां जहि काकुत्स्ह धर्मो ह्यस्या न विद्यते} %||1-24-16||

\twolineshloka
{श्रूयते हि पुरा शक्रो विरोचनसुतां नृप}
{पृथिवीं हन्तुमिच्छन्तीं मन्थरामभ्यसूदयत्} %||1-24-17||

\twolineshloka
{विष्णुना च पुरा राम भृगुपत्नी दृढव्रता}
{अनिन्द्रं लोकमिच्छन्ती काव्यमाता निषूदिता} %||1-24-18||

\twolineshloka
{एतैश्चान्यैश्च बहुभी राजपुत्रमहात्मभिः}
{अधर्मनिरता नार्यो हताः पुरुषसत्तमैः} %||1-25-19||


{॥इत्यार्षे श्रीमद्रामायणे वाल्मीकीये आदिकाव्ये बालकाण्डे चतुर्विंशः सर्गः॥}

\sect{पञ्चविंशः सर्गः}

\twolineshloka
{मुनेर्वचनमक्लीबं श्रुत्वा नरवरात्मजः}
{राघवः प्राञ्जलिर्भूत्वा प्रत्युवाच दृढव्रतः} %||1-25-1||

\twolineshloka
{पितुर्वचननिर्देशात्पितुर्वचनगौरवात्}
{वचनं कौशिकस्येति कर्तव्यमविशङ्कया} %||1-25-2||

\twolineshloka
{अनुशिष्टोऽस्म्ययोध्यायां गुरुमध्ये महात्मना}
{पित्रा दशरथेनाहं नावज्ञेयं च तद्वचः} %||1-25-3||

\twolineshloka
{सोऽहं पितुर्वचः श्रुत्वा शासनाद्ब्रह्म वादिनः}
{करिष्यामि न सन्देहस्ताटकावधमुत्तमम्} %||1-25-4||

\twolineshloka
{गोब्राह्मणहितार्थाय देशस्यास्य सुखाय च}
{तव चैवाप्रमेयस्य वचनं कर्तुमुद्यतः} %||1-25-5||

\twolineshloka
{एवमुक्त्वा धनुर्मध्ये बद्ध्वा मुष्टिमरिन्दमः}
{ज्याशब्दमकरोत्तीव्रं दिशः शब्देन पूरयन्} %||1-25-6||

\twolineshloka
{तेन शब्देन वित्रस्तास्ताटका वनवासिनः}
{ताटका च सुसङ्क्रुद्धा तेन शब्देन मोहिता} %||1-25-7||

\twolineshloka
{तं शब्दमभिनिध्याय राक्षसी क्रोधमूर्छिता}
{श्रुत्वा चाभ्यद्रवद्वेगाद्यतः शब्दो विनिःसृतः} %||1-25-8||

\twolineshloka
{तां दृष्ट्वा राघवः क्रुद्धां विकृतां विकृताननाम्}
{प्रमाणेनातिवृद्धां च लक्ष्मणं सोऽभ्यभाषत} %||1-25-9||

\twolineshloka
{पश्य लक्ष्मण यक्षिण्या भैरवं दारुणं वपुः}
{भिद्येरन्दर्शनादस्या भीरूणां हृदयानि च} %||1-25-10||

\twolineshloka
{एनां पश्य दुराधर्षां माया बलसमन्विताम्}
{विनिवृत्तां करोम्यद्य हृतकर्णाग्रनासिकाम्} %||1-25-11||

\twolineshloka
{न ह्येनामुत्सहे हन्तुं स्त्रीस्वभावेन रक्षिताम्}
{वीर्यं चास्या गतिं चापि हनिष्यामीति मे मतिः} %||1-25-12||

\twolineshloka
{एवं ब्रुवाणे रामे तु ताटका क्रोधमूर्छिता}
{उद्यम्य बाहू गर्जन्ती राममेवाभ्यधावत} %||1-25-13||

\twolineshloka
{तामापतन्तीं वेगेन विक्रान्तामशनीमिव}
{शरेणोरसि विव्याध सा पपात ममार च} %||1-25-14||

\twolineshloka
{तां हतां भीमसङ्काशां दृष्ट्वा सुरपतिस्तदा}
{साधु साध्विति काकुत्स्थं सुराश्च समपूजयन्} %||1-25-15||

\twolineshloka
{उवाच परमप्रीतः सहस्राक्षः पुरन्दरः}
{सुराश्च सर्वे संहृष्टा विश्वामित्रमथाब्रुवन्} %||1-25-16||

\twolineshloka
{मुने कौशिके भद्रं ते सेन्द्राः सर्वे मरुद्गणाः}
{तोषिताः कर्मणानेन स्नेहं दर्शय राघवे} %||1-25-17||

\twolineshloka
{प्रजापतेर्भृशाश्वस्य पुत्रान्सत्यपराक्रमान्}
{तपोबलभृतान्ब्रह्मन्राघवाय निवेदय} %||1-25-18||

\twolineshloka
{पात्रभूतश्च ते ब्रह्मंस्तवानुगमने धृतः}
{कर्तव्यं च महत्कर्म सुराणां राजसूनुना} %||1-25-19||

\twolineshloka
{एवमुक्त्वा सुराः सर्वे हृष्टा जग्मुर्यथागतम्}
{विश्वामित्रं पूजयित्वा ततः सन्ध्या प्रवर्तते} %||1-25-20||

\twolineshloka
{ततो मुनिवरः प्रीतिस्ताटका वधतोषितः}
{मूर्ध्नि राममुपाघ्राय इदं वचनमब्रवीत्} %||1-25-21||

\twolineshloka
{इहाद्य रजनीं राम वसेम शुभदर्शन}
{श्वः प्रभाते गमिष्यामस्तदाश्रमपदं मम} %||1-26-22||


{॥इत्यार्षे श्रीमद्रामायणे वाल्मीकीये आदिकाव्ये बालकाण्डे पञ्चविंशः सर्गः॥}

\sect{षड्विंशः सर्गः}

\twolineshloka
{अथ तां रजनीमुष्य विश्वामिरो महायशाः}
{प्रहस्य राघवं वाक्यमुवाच मधुराक्षरम्} %||1-26-1||

\twolineshloka
{पतितुष्टोऽस्मि भद्रं ते राजपुत्र महायशः}
{प्रीत्या परमया युक्तो ददाम्यस्त्राणि सर्वशः} %||1-26-2||

\twolineshloka
{देवासुरगणान्वापि सगन्धर्वोरगानपि}
{यैरमित्रान्प्रसह्याजौ वशीकृत्य जयिष्यसि} %||1-26-3||

\twolineshloka
{तानि दिव्यानि भद्रं ते ददाम्यस्त्राणि सर्वशः}
{दण्डचक्रं महद्दिव्यं तव दास्यामि राघव} %||1-26-4||

\twolineshloka
{धर्मचक्रं ततो वीर कालचक्रं तथैव च}
{विष्णुचक्रं तथात्युग्रमैन्द्रं चक्रं तथैव च} %||1-26-5||

\twolineshloka
{वज्रमस्त्रं नरश्रेष्ठ शैवं शूलवरं तथा}
{अस्त्रं ब्रह्मशिरश्चैव ऐषीकमपि राघव} %||1-26-6||

\twolineshloka
{ददामि ते महाबाहो ब्राह्ममस्त्रमनुत्तमम्}
{गदे द्वे चैव काकुत्स्थ मोदकी शिखरी उभे} %||1-26-7||

\twolineshloka
{प्रदीप्ते नरशार्दूल प्रयच्छामि नृपात्मज}
{धर्मपाशमहं राम कालपाशं तथैव च} %||1-26-8||

\twolineshloka
{वारुणं पाशमस्त्रं च ददान्यहमनुत्तमम्}
{अशनी द्वे प्रयच्छामि शुष्कार्द्रे रघुनन्दन} %||1-26-9||

\twolineshloka
{ददामि चास्त्रं पैनाकमस्त्रं नारायणं तथा}
{आग्नेयमस्त्र दयितं शिखरं नाम नामतः} %||1-26-10||

\twolineshloka
{वायव्यं प्रथमं नाम ददामि तव राघव}
{अस्त्रं हयशिरो नाम क्रौञ्चमस्त्रं तथैव च} %||1-26-11||

\twolineshloka
{शक्ति द्वयं च काकुत्स्थ ददामि तव चानघ}
{कङ्कालं मुसलं घोरं कापालमथ कङ्कणम्} %||1-26-12||

\twolineshloka
{धारयन्त्यसुरा यानि ददाम्येतानि सर्वशः}
{वैद्याधरं महास्त्रं च नन्दनं नाम नामतः} %||1-26-13||

\twolineshloka
{असिरत्नं महाबाहो ददामि नृवरात्मज}
{गान्धर्वमस्त्रं दयितं मानवं नाम नामतः} %||1-26-14||

\twolineshloka
{प्रस्वापनप्रशमने दद्मि सौरं च राघव}
{दर्पणं शोषणं चैव सन्तापनविलापने} %||1-26-15||

\threelineshloka
{मदनं चैव दुर्धर्षं कन्दर्पदयितं तथा}
{पैशाचमस्त्रं दयितं मोहनं नाम नामतः}
{प्रतीच्छ नरशार्दूल राजपुत्र महायशः} %||1-26-16||

\twolineshloka
{तामसं नरशार्दूल सौमनं च महाबलम्}
{संवर्तं चैव दुर्धर्षं मौसलं च नृपात्मज} %||1-26-17||

\twolineshloka
{सत्यमस्त्रं महाबाहो तथा मायाधरं परम्}
{घोरं तेजःप्रभं नाम परतेजोऽपकर्षणम्} %||1-26-18||

\twolineshloka
{सोमास्त्रं शिशिरं नाम त्वाष्ट्रमस्त्रं सुदामनम्}
{दारुणं च भगस्यापि शीतेषुमथ मानवम्} %||1-26-19||

\twolineshloka
{एतान्नाम महाबाहो कामरूपान्महाबलान्}
{गृहाण परमोदारान्क्षिप्रमेव नृपात्मज} %||1-26-20||

\twolineshloka
{स्थितस्तु प्राङ्मुखो भूत्वा शुचिर्निवरतस्तदा}
{ददौ रामाय सुप्रीतो मन्त्रग्राममनुत्तमम्} %||1-26-21||

\twolineshloka
{जपतस्तु मुनेस्तस्य विश्वामित्रस्य धीमतः}
{उपतस्थुर्महार्हाणि सर्वाण्यस्त्राणि राघवम्} %||1-26-22||

\twolineshloka
{ऊचुश्च मुदिता रामं सर्वे प्राञ्जलयस्तदा}
{इमे स्म परमोदार किङ्करास्तव राघव} %||1-26-23||

\twolineshloka
{प्रतिगृह्य च काकुत्स्थः समालभ्य च पाणिना}
{मनसा मे भविष्यध्वमिति तान्यभ्यचोदयत्} %||1-26-24||

\twolineshloka
{ततः प्रीतमना रामो विश्वामित्रं महामुनिम्}
{अभिवाद्य महातेजा गमनायोपचक्रमे} %||1-27-25||


{॥इत्यार्षे श्रीमद्रामायणे वाल्मीकीये आदिकाव्ये बालकाण्डे षड्विंशः सर्गः॥}

\sect{सप्तविंशः सर्गः}

\twolineshloka
{प्रतिगृह्य ततोऽस्त्राणि प्रहृष्टवदनः शुचिः}
{गच्छन्नेव च काकुत्स्थो विश्वामित्रमथाब्रवीत्} %||1-27-1||

\twolineshloka
{गृहीतास्त्रोऽस्मि भगवन्दुराधर्षः सुरैरपि}
{अस्त्राणां त्वहमिच्छामि संहारं मुनिपुङ्गव} %||1-27-2||

\twolineshloka
{एवं ब्रुवति काकुत्स्थे विश्वामित्रो महामुनिः}
{संहारं व्याजहाराथ धृतिमान्सुव्रतः शुचिः} %||1-27-3||

\twolineshloka
{सत्यवन्तं सत्यकीर्तिं धृष्टं रभसमेव च}
{प्रतिहारतरं नाम पराङ्मुखमवाङ्मुखम्} %||1-27-4||

\twolineshloka
{लक्षाक्षविषमौ चैव दृढनाभसुनाभकौ}
{दशाक्षशतवक्त्रौ च दशशीर्षशतोदरौ} %||1-27-5||

\twolineshloka
{पद्मनाभमहानाभौ दुन्दुनाभसुनाभकौ}
{ज्योतिषं कृशनं चैव नैराश्य विमलावुभौ} %||1-27-6||

\twolineshloka
{यौगन्धरहरिद्रौ च दैत्यप्रमथनौ तथा}
{पित्र्यं सौमनसं चैव विधूतमकरावुभौ} %||1-27-7||

\twolineshloka
{करवीरकरं चैव धनधान्यौ च राघव}
{कामरूपं कामरुचिं मोहमावरणं तथा} %||1-27-8||

\twolineshloka
{जृम्भकं सर्वनाभं च सन्तानवरणौ तथा}
{भृशाश्वतनयान्राम भास्वरान्कामरूपिणः} %||1-27-9||

\twolineshloka
{प्रतीच्छ मम भद्रं ते पात्रभूतोऽसि राघव}
{दिव्यभास्वरदेहाश्च मूर्तिमन्तः सुखप्रदाः} %||1-27-10||

\twolineshloka
{रामं प्राञ्जलयो भूत्वाब्रुवन्मधुरभाषिणः}
{इमे स्म नरशार्दूल शाधि किं करवाम ते} %||1-27-11||

\twolineshloka
{गम्यतामिति तानाह यथेष्टं रघुनन्दनः}
{मानसाः कार्यकालेषु साहाय्यं मे करिष्यथ} %||1-27-12||

\twolineshloka
{अथ ते राममामन्त्र्य कृत्वा चापि प्रदक्षिणम्}
{एवमस्त्विति काकुत्स्थमुक्त्वा जग्मुर्यथागतम्} %||1-27-13||

\twolineshloka
{स च तान्राघवो ज्ञात्वा विश्वामित्रं महामुनिम्}
{गच्छन्नेवाथ मधुरं श्लक्ष्णं वचनमब्रवीत्} %||1-27-14||

\twolineshloka
{किं न्वेतन्मेघसङ्काशं पर्वतस्याविदूरतः}
{वृक्षषण्डमितो भाति परं कौतूहलं हि मे} %||1-27-15||

\twolineshloka
{दर्शनीयं मृगाकीर्णं मनोहरमतीव च}
{नानाप्रकारैः शकुनैर्वल्गुभाषैरलङ्कृतम्} %||1-27-16||

\twolineshloka
{निःसृताः स्म मुनिश्रेष्ठ कान्ताराद्रोमहर्षणात्}
{अनया त्ववगच्छामि देशस्य सुखवत्तया} %||1-27-17||

\twolineshloka
{सर्वं मे शंस भगवन्कस्याश्रमपदं त्विदम्}
{सम्प्राप्ता यत्र ते पापा ब्रह्मघ्ना दुष्टचारिणः} %||1-28-18||


{॥इत्यार्षे श्रीमद्रामायणे वाल्मीकीये आदिकाव्ये बालकाण्डे सप्तविंशः सर्गः॥}

\sect{अष्टाविंशः सर्गः}

\twolineshloka
{अथ तस्याप्रमेयस्य तद्वनं परिपृच्छतः}
{विश्वामित्रो महातेजा व्याख्यातुमुपचक्रमे} %||1-28-1||

\twolineshloka
{एष पूर्वाश्रमो राम वामनस्य महात्मनः}
{सिद्धाश्रम इति ख्यातः सिद्धो ह्यत्र महातपाः} %||1-28-2||

\threelineshloka
{एतस्मिन्नेव काले तु राजा वैरोचनिर्बलिः}
{निर्जित्य दैवतगणान्सेन्द्रांश्च समरुद्गणान्}
{कारयामास तद्राज्यं त्रिषु लोकेषु विश्रुतः} %||1-28-3||

\twolineshloka
{बलेस्तु यजमानस्य देवाः साग्निपुरोगमाः}
{समागम्य स्वयं चैव विष्णुमूचुरिहाश्रमे} %||1-28-4||

\twolineshloka
{बलिर्वैरोचनिर्विष्णो यजते यज्ञमुत्तमम्}
{असमाप्ते क्रतौ तस्मिन्स्वकार्यमभिपद्यताम्} %||1-28-5||

\twolineshloka
{ये चैनमभिवर्तन्ते याचितार इतस्ततः}
{यच्च यत्र यथावच्च सर्वं तेभ्यः प्रयच्छति} %||1-28-6||

\twolineshloka
{स त्वं सुरहितार्थाय मायायोगमुपाश्रितः}
{वामनत्वं गतो विष्णो कुरु कल्याणमुत्तमम्} %||1-28-7||

\twolineshloka
{अयं सिद्धाश्रमो नाम प्रसादात्ते भविष्यति}
{सिद्धे कर्मणि देवेश उत्तिष्ठ भगवन्नितः} %||1-28-8||

\twolineshloka
{अथ विष्णुर्महातेजा अदित्यां समजायत}
{वामनं रूपमास्थाय वैरोचनिमुपागमत्} %||1-28-9||

\twolineshloka
{त्रीन्क्रमानथ भिक्षित्वा प्रतिगृह्य च मानतः}
{आक्रम्य लोकाँल्लोकात्मा सर्वभूतहिते रतः} %||1-28-10||

\twolineshloka
{महेन्द्राय पुनः प्रादान्नियम्य बलिमोजसा}
{त्रैलोक्यं स महातेजाश्चक्रे शक्रवशं पुनः} %||1-28-11||

\twolineshloka
{तेनैष पूर्वमाक्रान्त आश्रमः श्रमनाशनः}
{मयापि भक्त्या तस्यैष वामनस्योपभुज्यते} %||1-28-12||

\twolineshloka
{एतमाश्रममायान्ति राक्षसा विघ्नकारिणः}
{अत्र ते पुरुषव्याघ्र हन्तव्या दुष्टचारिणः} %||1-28-13||

\twolineshloka
{अद्य गच्छामहे राम सिद्धाश्रममनुत्तमम्}
{तदाश्रमपदं तात तवाप्येतद्यथा मम} %||1-28-14||

\twolineshloka
{तं दृष्ट्वा मुनयः सर्वे सिद्धाश्रमनिवासिनः}
{उत्पत्योत्पत्य सहसा विश्वामित्रमपूजयन्} %||1-28-15||

\twolineshloka
{यथार्हं चक्रिरे पूजां विश्वामित्राय धीमते}
{तथैव राजपुत्राभ्यामकुर्वन्नतिथिक्रियाम्} %||1-28-16||

\twolineshloka
{मुहूर्तमथ विश्रान्तौ राजपुत्रावरिन्दमौ}
{प्राञ्जली मुनिशार्दूलमूचतू रघुनन्दनौ} %||1-28-17||

\twolineshloka
{अद्यैव दीक्षां प्रविश भद्रं ते मुनिपुङ्गव}
{सिद्धाश्रमोऽयं सिद्धः स्यात्सत्यमस्तु वचस्तव} %||1-28-18||

\twolineshloka
{एवमुक्तो महातेजा विश्वामित्रो महामुनिः}
{प्रविवेश तदा दीक्षां नियतो नियतेन्द्रियः} %||1-28-19||

\twolineshloka
{कुमारावपि तां रात्रिमुषित्वा सुसमाहितौ}
{प्रभातकाले चोत्थाय विश्वामित्रमवन्दताम्} %||1-29-20||


{॥इत्यार्षे श्रीमद्रामायणे वाल्मीकीये आदिकाव्ये बालकाण्डे अष्टाविंशः सर्गः॥}

\sect{एकोनत्रिंशः सर्गः}

\twolineshloka
{अथ तौ देशकालज्ञौ राजपुत्रावरिन्दमौ}
{देशे काले च वाक्यज्ञावब्रूतां कौशिकं वचः} %||1-29-1||

\twolineshloka
{भगवञ्श्रोतुमिच्छावो यस्मिन्काले निशाचरौ}
{संरक्षणीयौ तौ ब्रह्मन्नातिवर्तेत तत्क्षणम्} %||1-29-2||

\twolineshloka
{एवं ब्रुवाणौ काकुत्स्थौ त्वरमाणौ युयुत्सया}
{सर्वे ते मुनयः प्रीताः प्रशशंसुर्नृपात्मजौ} %||1-29-3||

\twolineshloka
{अद्य प्रभृति षड्रात्रं रक्षतं राघवौ युवाम्}
{दीक्षां गतो ह्येष मुनिर्मौनित्वं च गमिष्यति} %||1-29-4||

\twolineshloka
{तौ तु तद्वचनं श्रुत्वा राजपुत्रौ यशस्विनौ}
{अनिद्रौ षडहोरात्रं तपोवनमरक्षताम्} %||1-29-5||

\twolineshloka
{उपासां चक्रतुर्वीरौ यत्तौ परमधन्विनौ}
{ररक्षतुर्मुनिवरं विश्वामित्रमरिन्दमौ} %||1-29-6||

\twolineshloka
{अथ काले गते तस्मिन्षष्ठेऽहनि समागते}
{सौमित्रमब्रवीद्रामो यत्तो भव समाहितः} %||1-29-7||

\twolineshloka
{रामस्यैवं ब्रुवाणस्य त्वरितस्य युयुत्सया}
{प्रजज्वाल ततो वेदिः सोपाध्यायपुरोहिता} %||1-29-8||

\twolineshloka
{मन्त्रवच्च यथान्यायं यज्ञोऽसौ सम्प्रवर्तते}
{आकाशे च महाञ्शब्दः प्रादुरासीद्भयानकः} %||1-29-9||

\twolineshloka
{आवार्य गगनं मेघो यथा प्रावृषि निर्गतः}
{तथा मायां विकुर्वाणौ राक्षसावभ्यधावताम्} %||1-29-10||

\twolineshloka
{मारीचश्च सुबाहुश्च तयोरनुचरास्तथा}
{आगम्य भीमसङ्काशा रुधिरौघानवासृजन्} %||1-29-11||

\twolineshloka
{तावापतन्तौ सहसा दृष्ट्वा राजीवलोचनः}
{लक्ष्मणं त्वभिसम्प्रेक्ष्य रामो वचनमब्रवीत्} %||1-29-12||

\twolineshloka
{पश्य लक्ष्मण दुर्वृत्तान्राक्षसान्पिशिताशनान्}
{मानवास्त्रसमाधूताननिलेन यथाघनान्} %||1-29-13||

\twolineshloka
{मानवं परमोदारमस्त्रं परमभास्वरम्}
{चिक्षेप परमक्रुद्धो मारीचोरसि राघवः} %||1-29-14||

\twolineshloka
{स तेन परमास्त्रेण मानवेन समाहितः}
{सम्पूर्णं योजनशतं क्षिप्तः सागरसम्प्लवे} %||1-29-15||

\twolineshloka
{विचेतनं विघूर्णन्तं शीतेषुबलपीडितम्}
{निरस्तं दृश्य मारीचं रामो लक्ष्मणमब्रवीत्} %||1-29-16||

\twolineshloka
{पश्य लक्ष्मण शीतेषुं मानवं धर्मसंहितम्}
{मोहयित्वा नयत्येनं न च प्राणैर्वियुज्यते} %||1-29-17||

\twolineshloka
{इमानपि वधिष्यामि निर्घृणान्दुष्टचारिणः}
{राक्षसान्पापकर्मस्थान्यज्ञघ्नान्रुधिराशनान्} %||1-29-18||

\twolineshloka
{विगृह्य सुमहच्चास्त्रमाग्नेयं रघुनन्दनः}
{सुबाहुरसि चिक्षेप स विद्धः प्रापतद्भुवि} %||1-29-19||

\twolineshloka
{शेषान्वायव्यमादाय निजघान महायशाः}
{राघवः परमोदारो मुनीनां मुदमावहन्} %||1-29-20||

\twolineshloka
{स हत्वा राक्षसान्सर्वान्यज्ञघ्नान्रघुनन्दनः}
{ऋषिभिः पूजितस्तत्र यथेन्द्रो विजये पुरा} %||1-29-21||

\twolineshloka
{अथ यज्ञे समाप्ते तु विश्वामित्रो महामुनिः}
{निरीतिका दिशो दृष्ट्वा काकुत्स्थमिदमब्रवीत्} %||1-29-22||

\twolineshloka
{कृतार्थोऽस्मि महाबाहो कृतं गुरुवचस्त्वया}
{सिद्धाश्रममिदं सत्यं कृतं राम महायशः} %||1-30-23||


{॥इत्यार्षे श्रीमद्रामायणे वाल्मीकीये आदिकाव्ये बालकाण्डे एकोनत्रिंशः सर्गः॥}

\sect{त्रिंशः सर्गः}

\twolineshloka
{अथ तां रजनीं तत्र कृतार्थौ रामलक्षणौ}
{ऊषतुर्मुदितौ वीरौ प्रहृष्टेनान्तरात्मना} %||1-30-1||

\twolineshloka
{प्रभातायां तु शर्वर्यां कृतपौर्वाह्णिकक्रियौ}
{विश्वामित्रमृषींश्चान्यान्सहितावभिजग्मतुः} %||1-30-2||

\twolineshloka
{अभिवाद्य मुनिश्रेष्ठं ज्वलन्तमिव पावकम्}
{ऊचतुर्मधुरोदारं वाक्यं मधुरभाषिणौ} %||1-30-3||

\twolineshloka
{इमौ स्वो मुनिशार्दूल किङ्करौ समुपस्थितौ}
{आज्ञापय यथेष्टं वै शासनं करवाव किम्} %||1-30-4||

\twolineshloka
{एवमुक्ते ततस्ताभ्यां सर्व एव महर्षयः}
{विश्वामित्रं पुरस्कृत्य रामं वचनमब्रुवन्} %||1-30-5||

\twolineshloka
{मैथिलस्य नरश्रेष्ठ जनकस्य भविष्यति}
{यज्ञः परमधर्मिष्ठस्तत्र यास्यामहे वयम्} %||1-30-6||

\twolineshloka
{त्वं चैव नरशार्दूल सहास्माभिर्गमिष्यसि}
{अद्भुतं च धनूरत्नं तत्र त्वं द्रष्टुमर्हसि} %||1-30-7||

\twolineshloka
{तद्धि पूर्वं नरश्रेष्ठ दत्तं सदसि दैवतैः}
{अप्रमेयबलं घोरं मखे परमभास्वरम्} %||1-30-8||

\twolineshloka
{नास्य देवा न गन्धर्वा नासुरा न च राक्षसाः}
{कर्तुमारोपणं शक्ता न कथञ्चन मानुषाः} %||1-30-9||

\twolineshloka
{धनुषस्तस्य वीर्यं हि जिज्ञासन्तो महीक्षितः}
{न शेकुरारोपयितुं राजपुत्रा महाबलाः} %||1-30-10||

\twolineshloka
{तद्धनुर्नरशार्दूल मैथिलस्य महात्मनः}
{तत्र द्रक्ष्यसि काकुत्स्थ यज्ञं चाद्भुतदर्शनम्} %||1-30-11||

\twolineshloka
{तद्धि यज्ञफलं तेन मैथिलेनोत्तमं धनुः}
{याचितं नरशार्दूल सुनाभं सर्वदैवतैः} %||1-30-12||

\twolineshloka
{एवमुक्त्वा मुनिवरः प्रस्थानमकरोत्तदा}
{सर्षिसङ्घः सकाकुत्स्थ आमन्त्र्य वनदेवताः} %||1-30-13||

\twolineshloka
{स्वस्ति वोऽस्तु गमिष्यामि सिद्धः सिद्धाश्रमादहम्}
{उत्तरे जाह्नवीतीरे हिमवन्तं शिलोच्चयम्} %||1-30-14||

\twolineshloka
{प्रदक्षिणं ततः कृत्वा सिद्धाश्रममनुत्तमम्}
{उत्तरां दिशमुद्दिश्य प्रस्थातुमुपचक्रमे} %||1-30-15||

\twolineshloka
{तं व्रजन्तं मुनिवरमन्वगादनुसारिणाम्}
{शकटी शतमात्रं तु प्रयाणे ब्रह्मवादिनाम्} %||1-30-16||

\twolineshloka
{मृगपक्षिगणाश्चैव सिद्धाश्रमनिवासिनः}
{अनुजग्मुर्महात्मानं विश्वामित्रं महामुनिम्} %||1-30-17||

\twolineshloka
{ते गत्वा दूरमध्वानं लम्बमाने दिवाकरे}
{वासं चक्रुर्मुनिगणाः शोणाकूले समाहिताः} %||1-30-18||

\twolineshloka
{तेऽस्तं गते दिनकरे स्नात्वा हुतहुताशनाः}
{विश्वामित्रं पुरस्कृत्य निषेदुरमितौजसः} %||1-30-19||

\twolineshloka
{रामोऽपि सहसौमित्रिर्मुनींस्तानभिपूज्य च}
{अग्रतो निषसादाथ विश्वामित्रस्य धीमतः} %||1-30-20||

\twolineshloka
{अथ रामो महातेजा विश्वामित्रं महामुनिम्}
{पप्रच्छ मुनिशार्दूलं कौतूहलसमन्वितः} %||1-30-21||

\twolineshloka
{भगवन्को न्वयं देशः समृद्धवनशोभितः}
{श्रोतुमिच्छामि भद्रं ते वक्तुमर्हसि तत्त्वतः} %||1-30-22||

\twolineshloka
{चोदितो रामवाक्येन कथयामास सुव्रतः}
{तस्य देशस्य निखिलमृषिमध्ये महातपाः} %||1-31-23||


{॥इत्यार्षे श्रीमद्रामायणे वाल्मीकीये आदिकाव्ये बालकाण्डे त्रिंशः सर्गः॥}

\sect{एकत्रिंशः सर्गः}

\twolineshloka
{ब्रह्मयोनिर्महानासीत्कुशो नाम महातपाः}
{वैदर्भ्यां जनयामास चतुरः सदृशान्सुतान्} %||1-31-1||

\threelineshloka
{कुशाम्बं कुशनाभं च आधूर्त रजसं वसुम्}
{दीप्तियुक्तान्महोत्साहान्क्षत्रधर्मचिकीर्षया}
{तानुवाच कुशः पुत्रान्धर्मिष्ठान्सत्यवादिनः} %||1-31-2||

\twolineshloka
{कुशस्य वचनं श्रुत्वा चत्वारो लोकसम्मताः}
{निवेशं चक्रिरे सर्वे पुराणां नृवरास्तदा} %||1-31-3||

\twolineshloka
{कुशाम्बस्तु महातेजाः कौशाम्बीमकरोत्पुरीम्}
{कुशनाभस्तु धर्मात्मा परं चक्रे महोदयम्} %||1-31-4||

\twolineshloka
{आधूर्तरजसो राम धर्मारण्यं महीपतिः}
{चक्रे पुरवरं राजा वसुश्चक्रे गिरिव्रजम्} %||1-31-5||

\twolineshloka
{एषा वसुमती राम वसोस्तस्य महात्मनः}
{एते शैलवराः पञ्च प्रकाशन्ते समन्ततः} %||1-31-6||

\twolineshloka
{सुमागधी नदी रम्या मागधान्विश्रुताययौ}
{पञ्चानां शैलमुख्यानां मध्ये मालेव शोभते} %||1-31-7||

\twolineshloka
{सैषा हि मागधी राम वसोस्तस्य महात्मनः}
{पूर्वाभिचरिता राम सुक्षेत्रा सस्यमालिनी} %||1-31-8||

\twolineshloka
{कुशनाभस्तु राजर्षिः कन्याशतमनुत्तमम्}
{जनयामास धर्मात्मा घृताच्यां रघुनन्दन} %||1-31-9||

\twolineshloka
{तास्तु यौवनशालिन्यो रूपवत्यः स्वलङ्कृताः}
{उद्यानभूमिमागम्य प्रावृषीव शतह्रदाः} %||1-31-10||

\twolineshloka
{गायन्त्यो नृत्यमानाश्च वादयन्त्यश्च राघव}
{आमोदं परमं जग्मुर्वराभरणभूषिताः} %||1-31-11||

\twolineshloka
{अथ ताश्चारुसर्वाङ्ग्यो रूपेणाप्रतिमा भुवि}
{उद्यानभूमिमागम्य तारा इव घनान्तरे} %||1-31-12||

\twolineshloka
{ताः सर्वगुणसम्पन्ना रूपयौवनसंयुताः}
{दृष्ट्वा सर्वात्मको वायुरिदं वचनमब्रवीत्} %||1-31-13||

\twolineshloka
{अहं वः कामये सर्वा भार्या मम भविष्यथ}
{मानुषस्त्यज्यतां भावो दीर्घमायुरवाप्स्यथ} %||1-31-14||

\twolineshloka
{तस्य तद्वचनं श्रुत्वा वायोरक्लिष्टकर्मणः}
{अपहास्य ततो वाक्यं कन्याशतमथाब्रवीत्} %||1-31-15||

\twolineshloka
{अन्तश्चरसि भूतानां सर्वेषां त्वं सुरोत्तम}
{प्रभावज्ञाश्च ते सर्वाः किमस्मानवमन्यसे} %||1-31-16||

\twolineshloka
{कुशनाभसुताः सर्वाः समर्थास्त्वां सुरोत्तम}
{स्थानाच्च्यावयितुं देवं रक्षामस्तु तपो वयम्} %||1-31-17||

\twolineshloka
{मा भूत्स कालो दुर्मेधः पितरं सत्यवादिनम्}
{नावमन्यस्व धर्मेण स्वयंवरमुपास्महे} %||1-31-18||

\twolineshloka
{पिता हि प्रभुरस्माकं दैवतं परमं हि सः}
{यस्य नो दास्यति पिता स नो भर्ता भविष्यति} %||1-31-19||

\twolineshloka
{तासां तद्वचनं श्रुत्वा वायुः परमकोपनः}
{प्रविश्य सर्वगात्राणि बभञ्ज भगवान्प्रभुः} %||1-31-20||

\twolineshloka
{ताः कन्या वायुना भग्ना विविशुर्नृपतेर्गृहम्}
{दृष्ट्वा भग्नास्तदा राजा सम्भ्रान्त इदमब्रवीत्} %||1-31-21||

\twolineshloka
{किमिदं कथ्यतां पुत्र्यः को धर्ममवमन्यते}
{कुब्जाः केन कृताः सर्वा वेष्टन्त्यो नाभिभाषथ} %||1-32-22||


{॥इत्यार्षे श्रीमद्रामायणे वाल्मीकीये आदिकाव्ये बालकाण्डे एकत्रिंशः सर्गः॥}

\sect{द्वात्रिंशः सर्गः}

\twolineshloka
{तस्य तद्वचनं श्रुत्वा कुशनाभस्य धीमतः}
{शिरोभिश्चरणौ स्पृष्ट्वा कन्याशतमभाषत} %||1-32-1||

\twolineshloka
{वायुः सर्वात्मको राजन्प्रधर्षयितुमिच्छति}
{अशुभं मार्गमास्थाय न धर्मं प्रत्यवेक्षते} %||1-32-2||

\twolineshloka
{पितृमत्यः स्म भद्रं ते स्वच्छन्दे न वयं स्थिताः}
{पितरं नो वृणीष्व त्वं यदि नो दास्यते तव} %||1-32-3||

\twolineshloka
{तेन पापानुबन्धेन वचनं न प्रतीच्छता}
{एवं ब्रुवन्त्यः सर्वाः स्म वायुना निहता भृषम्} %||1-32-4||

\twolineshloka
{तासां तद्वचनं श्रुत्वा राजा परमधार्मिकः}
{प्रत्युवाच महातेजाः कन्याशतमनुत्तमम्} %||1-32-5||

\twolineshloka
{क्षान्तं क्षमावतां पुत्र्यः कर्तव्यं सुमहत्कृतम्}
{ऐकमत्यमुपागम्य कुलं चावेक्षितं मम} %||1-32-6||

\twolineshloka
{अलङ्कारो हि नारीणां क्षमा तु पुरुषस्य वा}
{दुष्करं तच्च वः क्षान्तं त्रिदशेषु विशेषतः} %||1-32-7||

\twolineshloka
{यादृशीर्वः क्षमा पुत्र्यः सर्वासामविशेषतः}
{क्षमा दानं क्षमा यज्ञः क्षमा सत्यं च पुत्रिकाः} %||1-32-8||

\twolineshloka
{क्षमा यशः क्षमा धर्मः क्षमायां विष्ठितं जगत्}
{विसृज्य कन्याः काकुत्स्थ राजा त्रिदशविक्रमः} %||1-32-9||

\twolineshloka
{मन्त्रज्ञो मन्त्रयामास प्रदानं सह मन्त्रिभिः}
{देशे काले प्रदानस्य सदृशे प्रतिपादनम्} %||1-32-10||

\twolineshloka
{एतस्मिन्नेव काले तु चूली नाम महामुनिः}
{ऊर्ध्वरेताः शुभाचारो ब्राह्मं तप उपागमत्} %||1-32-11||

\twolineshloka
{तप्यन्तं तमृषिं तत्र गन्धर्वी पर्युपासते}
{सोमदा नाम भद्रं ते ऊर्मिला तनया तदा} %||1-32-12||

\twolineshloka
{सा च तं प्रणता भूत्वा शुश्रूषणपरायणा}
{उवास काले धर्मिष्ठा तस्यास्तुष्टोऽभवद्गुरुः} %||1-32-13||

\twolineshloka
{स च तां कालयोगेन प्रोवाच रघुनन्दन}
{परितुष्टोऽस्मि भद्रं ते किं करोमि तव प्रियम्} %||1-32-14||

\twolineshloka
{परितुष्टं मुनिं ज्ञात्वा गन्धर्वी मधुरस्वरम्}
{उवाच परमप्रीता वाक्यज्ञा वाक्यकोविदम्} %||1-32-15||

\twolineshloka
{लक्ष्म्या समुदितो ब्राह्म्या ब्रह्मभूतो महातपाः}
{ब्राह्मेण तपसा युक्तं पुत्रमिच्छामि धार्मिकम्} %||1-32-16||

\twolineshloka
{अपतिश्चास्मि भद्रं ते भार्या चास्मि न कस्यचित्}
{ब्राह्मेणोपगतायाश्च दातुमर्हसि मे सुतम्} %||1-32-17||

\twolineshloka
{तस्याः प्रसन्नो ब्रह्मर्षिर्ददौ पुत्रमनुत्तमम्}
{ब्रह्मदत्त इति ख्यातं मानसं चूलिनः सुतम्} %||1-32-18||

\twolineshloka
{स राजा ब्रह्मदत्तस्तु पुरीमध्यवसत्तदा}
{काम्पिल्यां परया लक्ष्म्या देवराजो यथा दिवम्} %||1-32-19||

\twolineshloka
{स बुद्धिं कृतवान्राजा कुशनाभः सुधार्मिकः}
{ब्रह्मदत्ताय काकुत्स्थ दातुं कन्याशतं तदा} %||1-32-20||

\twolineshloka
{तमाहूय महातेजा ब्रह्मदत्तं महीपतिः}
{ददौ कन्याशतं राजा सुप्रीतेनान्तरात्मना} %||1-32-21||

\twolineshloka
{यथाक्रमं ततः पाणिं जग्राह रघुनन्दन}
{ब्रह्मदत्तो मही पालस्तासां देवपतिर्यथा} %||1-32-22||

\twolineshloka
{स्पृष्टमात्रे ततः पाणौ विकुब्जा विगतज्वराः}
{युक्ताः परमया लक्ष्म्या बभुः कन्याशतं तदा} %||1-32-23||

\twolineshloka
{स दृष्ट्वा वायुना मुक्ताः कुशनाभो महीपतिः}
{बभूव परमप्रीतो हर्षं लेभे पुनः पुनः} %||1-32-24||

\twolineshloka
{कृतोद्वाहं तु राजानं ब्रह्मदत्तं महीपतिः}
{सदारं प्रेषयामास सोपाध्याय गणं तदा} %||1-32-25||

\twolineshloka
{सोमदापि सुसंहृष्टा पुत्रस्य सदृशीं क्रियाम्}
{यथान्यायं च गन्धर्वी स्नुषास्ताः प्रत्यनन्दत} %||1-33-26||


{॥इत्यार्षे श्रीमद्रामायणे वाल्मीकीये आदिकाव्ये बालकाण्डे द्वात्रिंशः सर्गः॥}

\sect{त्रयस्त्रिंशः सर्गः}

\twolineshloka
{कृतोद्वाहे गते तस्मिन्ब्रह्मदत्ते च राघव}
{अपुत्रः पुत्रलाभाय पौत्रीमिष्टिमकल्पयत्} %||1-33-1||

\twolineshloka
{इष्ट्यां तु वर्तमानायां कुशनाभं महीपतिम्}
{उवाच परमप्रीतः कुशो ब्रह्मसुतस्तदा} %||1-33-2||

\twolineshloka
{पुत्रस्ते सदृशः पुत्र भविष्यति सुधार्मिकः}
{गाधिं प्राप्स्यसि तेन त्वं कीर्तिं लोके च शाश्वतीम्} %||1-33-3||

\twolineshloka
{एवमुक्त्वा कुशो राम कुशनाभं महीपतिम्}
{जगामाकाशमाविश्य ब्रह्मलोकं सनातनम्} %||1-33-4||

\twolineshloka
{कस्यचित्त्वथ कालस्य कुशनाभस्य धीमतः}
{जज्ञे परमधर्मिष्ठो गाधिरित्येव नामतः} %||1-33-5||

\twolineshloka
{स पिता मम काकुत्स्थ गाधिः परमधार्मिकः}
{कुशवंशप्रसूतोऽस्मि कौशिको रघुनन्दन} %||1-33-6||

\twolineshloka
{पूर्वजा भगिनी चापि मम राघव सुव्रता}
{नाम्ना सत्यवती नाम ऋचीके प्रतिपादिता} %||1-33-7||

\twolineshloka
{सशरीरा गता स्वर्गं भर्तारमनुवर्तिनी}
{कौशिकी परमोदारा सा प्रवृत्ता महानदी} %||1-33-8||

\twolineshloka
{दिव्या पुण्योदका रम्या हिमवन्तमुपाश्रिता}
{लोकस्य हितकामार्थं प्रवृत्ता भगिनी मम} %||1-33-9||

\twolineshloka
{ततोऽहं हिमवत्पार्श्वे वसामि नियतः सुखम्}
{भगिन्याः स्नेहसंयुक्तः कौशिक्या रघुनन्दन} %||1-33-10||

\twolineshloka
{सा तु सत्यवती पुण्या सत्ये धर्मे प्रतिष्ठिता}
{पतिव्रता महाभागा कौशिकी सरितां वरा} %||1-33-11||

\twolineshloka
{अहं हि नियमाद्राम हित्वा तां समुपागतः}
{सिद्धाश्रममनुप्राप्य सिद्धोऽस्मि तव तेजसा} %||1-33-12||

\twolineshloka
{एषा राम ममोत्पत्तिः स्वस्य वंशस्य कीर्तिता}
{देशस्य च महाबाहो यन्मां त्वं परिपृच्छसि} %||1-33-13||

\twolineshloka
{गतोऽर्धरात्रः काकुत्स्थ कथाः कथयतो मम}
{निद्रामभ्येहि भद्रं ते मा भूद्विघ्नोऽध्वनीह नः} %||1-33-14||

\twolineshloka
{निष्पन्दास्तरवः सर्वे निलीना मृगपक्षिणः}
{नैशेन तमसा व्याप्ता दिशश्च रघुनन्दन} %||1-33-15||

\twolineshloka
{शनैर्वियुज्यते सन्ध्या नभो नेत्रैरिवावृतम्}
{नक्षत्रतारागहनं ज्योतिर्भिरवभासते} %||1-33-16||

\twolineshloka
{उत्तिष्ठति च शीतांशुः शशी लोकतमोनुदः}
{ह्लादयन्प्राणिनां लोके मनांसि प्रभया विभो} %||1-33-17||

\twolineshloka
{नैशानि सर्वभूतानि प्रचरन्ति ततस्ततः}
{यक्षराक्षससङ्घाश्च रौद्राश्च पिशिताशनाः} %||1-33-18||

\twolineshloka
{एवमुक्त्वा महातेजा विरराम महामुनिः}
{साधु साध्विति तं सर्वे मुनयो ह्यभ्यपूजयन्} %||1-33-19||

\twolineshloka
{रामोऽपि सह सौमित्रिः किञ्चिदागतविस्मयः}
{प्रशस्य मुनिशार्दूलं निद्रां समुपसेवते} %||1-34-20||


{॥इत्यार्षे श्रीमद्रामायणे वाल्मीकीये आदिकाव्ये बालकाण्डे त्रयस्त्रिंशः सर्गः॥}

\sect{चतुस्त्रिंशः सर्गः}

\twolineshloka
{उपास्य रात्रिशेषं तु शोणाकूले महर्षिभिः}
{निशायां सुप्रभातायां विश्वामित्रोऽभ्यभाषत} %||1-34-1||

\twolineshloka
{सुप्रभाता निशा राम पूर्वा सन्ध्या प्रवर्तते}
{उत्तिष्ठोत्तिष्ठ भद्रं ते गमनायाभिरोचय} %||1-34-2||

\twolineshloka
{तच्छ्रुत्वा वचनं तस्य कृत्वा पौर्वाह्णिकीं क्रियाम्}
{गमनं रोचयामास वाक्यं चेदमुवाच ह} %||1-34-3||

\twolineshloka
{अयं शोणः शुभजलो गाधः पुलिनमण्डितः}
{कतरेण पथा ब्रह्मन्सन्तरिष्यामहे वयम्} %||1-34-4||

\twolineshloka
{एवमुक्तस्तु रामेण विश्वामित्रोऽब्रवीदिदम्}
{एष पन्था मयोद्दिष्टो येन यान्ति महर्षयः} %||1-34-5||

\twolineshloka
{ते गत्वा दूरमध्वानं गतेऽर्धदिवसे तदा}
{जाह्नवीं सरितां श्रेष्ठां ददृशुर्मुनिसेविताम्} %||1-34-6||

\threelineshloka
{तां दृष्ट्वा पुण्यसलिलां हंससारससेविताम्}
{बभूवुर्मुदिताः सर्वे मुनयः सहराघवाः}
{तस्यास्तीरे ततश्चक्रुस्ते आवासपरिग्रहम्} %||1-34-7||

\twolineshloka
{ततः स्नात्वा यथान्यायं सन्तर्प्य पितृदेवताः}
{हुत्वा चैवाग्निहोत्राणि प्राश्य चामृतवद्धविः} %||1-34-8||

\twolineshloka
{विविशुर्जाह्नवीतीरे शुचौ मुदितमानसाः}
{विश्वामित्रं महात्मानं परिवार्य समन्ततः} %||1-34-9||

\threelineshloka
{सम्प्रहृष्टमना रामो विश्वामित्रमथाब्रवीत्}
{भगवञ्श्रोतुमिच्छामि गङ्गां त्रिपथगां नदीम्}
{त्रैलोक्यं कथमाक्रम्य गता नदनदीपतिम्} %||1-34-10||

\twolineshloka
{चोदितो राम वाक्येन विश्वामित्रो महामुनिः}
{वृद्धिं जन्म च गङ्गाया वक्तुमेवोपचक्रमे} %||1-34-11||

\twolineshloka
{शैलेन्द्रो हिमवान्नाम धातूनामाकरो महान्}
{तस्य कन्या द्वयं राम रूपेणाप्रतिमं भुवि} %||1-34-12||

\twolineshloka
{या मेरुदुहिता राम तयोर्माता सुमध्यमा}
{नाम्ना मेना मनोज्ञा वै पत्नी हिमवतः प्रिया} %||1-34-13||

\twolineshloka
{तस्यां गङ्गेयमभवज्ज्येष्ठा हिमवतः सुता}
{उमा नाम द्वितीयाभूत्कन्या तस्यैव राघव} %||1-34-14||

\twolineshloka
{अथ ज्येष्ठां सुराः सर्वे देवतार्थचिकीर्षया}
{शैलेन्द्रं वरयामासुर्गङ्गां त्रिपथगां नदीम्} %||1-34-15||

\twolineshloka
{ददौ धर्मेण हिमवांस्तनयां लोकपावनीम्}
{स्वच्छन्दपथगां गङ्गां त्रैलोक्यहितकाम्यया} %||1-34-16||

\twolineshloka
{प्रतिगृह्य त्रिलोकार्थं त्रिलोकहितकारिणः}
{गङ्गामादाय तेऽगच्छन्कृतार्थेनान्तरात्मना} %||1-34-17||

\twolineshloka
{या चान्या शैलदुहिता कन्यासीद्रघुनन्दन}
{उग्रं सा व्रतमास्थाय तपस्तेपे तपोधना} %||1-34-18||

\twolineshloka
{उग्रेण तपसा युक्तां ददौ शैलवरः सुताम्}
{रुद्रायाप्रतिरूपाय उमां लोकनमस्कृताम्} %||1-34-19||

\twolineshloka
{एते ते शैल राजस्य सुते लोकनमस्कृते}
{गङ्गा च सरितां श्रेष्ठा उमा देवी च राघव} %||1-34-20||

\twolineshloka
{एतत्ते धर्ममाख्यातं यथा त्रिपथगा नदी}
{खं गता प्रथमं तात गतिं गतिमतां वर} %||1-35-21||


{॥इत्यार्षे श्रीमद्रामायणे वाल्मीकीये आदिकाव्ये बालकाण्डे चतुस्त्रिंशः सर्गः॥}

\sect{पञ्चत्रिंशः सर्गः}

\twolineshloka
{उक्त वाक्ये मुनौ तस्मिन्नुभौ राघवलक्ष्मणौ}
{प्रतिनन्द्य कथां वीरावूचतुर्मुनिपुङ्गवम्} %||1-35-1||

\twolineshloka
{धर्मयुक्तमिदं ब्रह्मन्कथितं परमं त्वया}
{दुहितुः शैलराजस्य ज्येष्ठाय वक्तुमर्हसि} %||1-35-2||

\twolineshloka
{विस्तरं विस्तरज्ञोऽसि दिव्यमानुषसम्भवम्}
{त्रीन्पथो हेतुना केन पावयेल्लोकपावनी} %||1-35-3||

\twolineshloka
{कथं गङ्गां त्रिपथगा विश्रुता सरिदुत्तमा}
{त्रिषु लोकेषु धर्मज्ञ कर्मभिः कैः समन्विता} %||1-35-4||

\twolineshloka
{तथा ब्रुवति काकुत्स्थे विश्वामित्रस्तपोधनः}
{निखिलेन कथां सर्वामृषिमध्ये न्यवेदयत्} %||1-35-5||

\twolineshloka
{पुरा राम कृतोद्वाहः शितिकण्ठो महातपाः}
{दृष्ट्वा च स्पृहया देवीं मैथुनायोपचक्रमे} %||1-35-6||

\twolineshloka
{शितिकण्ठस्य देवस्य दिव्यं वर्षशतं गतम्}
{न चापि तनयो राम तस्यामासीत्परन्तप} %||1-35-7||

\twolineshloka
{ततो देवाः समुद्विग्नाः पितामहपुरोगमाः}
{यदिहोत्पद्यते भूतं कस्तत्प्रतिसहिष्यते} %||1-35-8||

\threelineshloka
{अभिगम्य सुराः सर्वे प्रणिपत्येदमब्रुवन्}
{देवदेव महादेव लोकस्यास्य हिते रत}
{सुराणां प्रणिपातेन प्रसादं कर्तुमर्हसि} %||1-35-9||

\twolineshloka
{न लोका धारयिष्यन्ति तव तेजः सुरोत्तम}
{ब्राह्मेण तपसा युक्तो देव्या सह तपश्चर} %||1-35-10||

\twolineshloka
{त्रैलोक्यहितकामार्थं तेजस्तेजसि धारय}
{रक्ष सर्वानिमाँल्लोकान्नालोकं कर्तुमर्हसि} %||1-35-11||

\twolineshloka
{देवतानां वचः श्रुत्वा सर्वलोकमहेश्वरः}
{बाढमित्यब्रवीत्सर्वान्पुनश्चेदमुवाच ह} %||1-35-12||

\twolineshloka
{धारयिष्याम्यहं तेजस्तेजस्येव सहोमया}
{त्रिदशाः पृथिवी चैव निर्वाणमधिगच्छतु} %||1-35-13||

\twolineshloka
{यदिदं क्षुभितं स्थानान्मम तेजो ह्यनुत्तमम्}
{धारयिष्यति कस्तन्मे ब्रुवन्तु सुरसत्तमाः} %||1-35-14||

\twolineshloka
{एवमुक्तास्ततो देवाः प्रत्यूचुर्वृषभध्वजम्}
{यत्तेजः क्षुभितं ह्येतत्तद्धरा धारयिष्यति} %||1-35-15||

\twolineshloka
{एवमुक्तः सुरपतिः प्रमुमोच महीतले}
{तेजसा पृथिवी येन व्याप्ता सगिरिकानना} %||1-35-16||

\twolineshloka
{ततो देवाः पुनरिदमूचुश्चाथ हुताशनम्}
{प्रविश त्वं महातेजो रौद्रं वायुसमन्वितः} %||1-35-17||

\threelineshloka
{तदग्निना पुनर्व्याप्तं संजातः श्वेतपर्वतः}
{दिव्यं शरवणं चैव पावकादित्यसंनिभम्}
{यत्र जातो महातेजाः कार्तिकेयोऽग्निसम्भवः} %||1-35-18||

\twolineshloka
{अथोमां च शिवं चैव देवाः सर्षि गणास्तदा}
{पूजयामासुरत्यर्थं सुप्रीतमनसस्ततः} %||1-35-19||

\twolineshloka
{अथ शैल सुता राम त्रिदशानिदमब्रवीत्}
{समन्युरशपत्सर्वान्क्रोधसंरक्तलोचना} %||1-35-20||

\threelineshloka
{यस्मान्निवारिता चैव सङ्गता पुत्रकाम्यया}
{अपत्यं स्वेषु दारेषु नोत्पादयितुमर्हथ}
{अद्य प्रभृति युष्माकमप्रजाः सन्तु पत्नयः} %||1-35-21||

\twolineshloka
{एवमुक्त्वा सुरान्सर्वाञ्शशाप पृथिवीमपि}
{अवने नैकरूपा त्वं बहुभार्या भविष्यसि} %||1-35-22||

\twolineshloka
{न च पुत्रकृतां प्रीतिं मत्क्रोधकलुषी कृता}
{प्राप्स्यसि त्वं सुदुर्मेधे मम पुत्रमनिच्छती} %||1-35-23||

\twolineshloka
{तान्सर्वान्व्रीडितान्दृष्ट्वा सुरान्सुरपतिस्तदा}
{गमनायोपचक्राम दिशं वरुणपालिताम्} %||1-35-24||

\twolineshloka
{स गत्वा तप आतिष्ठत्पार्श्वे तस्योत्तरे गिरेः}
{हिमवत्प्रभवे शृङ्गे सह देव्या महेश्वरः} %||1-35-25||

\twolineshloka
{एष ते विस्तरो राम शैलपुत्र्या निवेदितः}
{गङ्गायाः प्रभवं चैव शृणु मे सहलक्ष्मणः} %||1-36-26||


{॥इत्यार्षे श्रीमद्रामायणे वाल्मीकीये आदिकाव्ये बालकाण्डे पञ्चत्रिंशः सर्गः॥}

\sect{षट्त्रिंशः सर्गः}

\twolineshloka
{तप्यमाने तपो देवे देवाः सर्षिगणाः पुरा}
{सेनापतिमभीप्सन्तः पितामहमुपागमन्} %||1-36-1||

\twolineshloka
{ततोऽब्रुवन्सुराः सर्वे भगवन्तं पितामहम्}
{प्रणिपत्य शुभं वाक्यं सेन्द्राः साग्निपुरोगमाः} %||1-36-2||

\twolineshloka
{यो नः सेनापतिर्देव दत्तो भगवता पुरा}
{स तपः परमास्थाय तप्यते स्म सहोमया} %||1-36-3||

\twolineshloka
{यदत्रानन्तरं कार्यं लोकानां हितकाम्यया}
{संविधत्स्व विधानज्ञ त्वं हि नः परमा गतिः} %||1-36-4||

\twolineshloka
{देवतानां वचः श्रुत्वा सर्वलोकपितामहः}
{सान्त्वयन्मधुरैर्वाक्यैस्त्रिदशानिदमब्रवीत्} %||1-36-5||

\twolineshloka
{शैलपुत्र्या यदुक्तं तन्न प्रजास्यथ पत्निषु}
{तस्या वचनमक्लिष्टं सत्यमेव न संशयः} %||1-36-6||

\twolineshloka
{इयमाकाशगा गङ्गा यस्यां पुत्रं हुताशनः}
{जनयिष्यति देवानां सेनापतिमरिन्दमम्} %||1-36-7||

\twolineshloka
{ज्येष्ठा शैलेन्द्रदुहिता मानयिष्यति तं सुतम्}
{उमायास्तद्बहुमतं भविष्यति न संशयः} %||1-36-8||

\twolineshloka
{तच्छ्रुत्वा वचनं तस्य कृतार्था रघुनन्दन}
{प्रणिपत्य सुराः सर्वे पितामहमपूजयन्} %||1-36-9||

\twolineshloka
{ते गत्वा पर्वतं राम कैलासं धातुमण्डितम्}
{अग्निं नियोजयामासुः पुत्रार्थं सर्वदेवताः} %||1-36-10||

\twolineshloka
{देवकार्यमिदं देव समाधत्स्व हुताशन}
{शैलपुत्र्यां महातेजो गङ्गायां तेज उत्सृज} %||1-36-11||

\twolineshloka
{देवतानां प्रतिज्ञाय गङ्गामभ्येत्य पावकः}
{गर्भं धारय वै देवि देवतानामिदं प्रियम्} %||1-36-12||

\twolineshloka
{इत्येतद्वचनं श्रुत्वा दिव्यं रूपमधारयत्}
{स तस्या महिमां दृष्ट्वा समन्तादवकीर्यत} %||1-36-13||

\twolineshloka
{समन्ततस्तदा देवीमभ्यषिञ्चत पावकः}
{सर्वस्रोतांसि पूर्णानि गङ्गाया रघुनन्दन} %||1-36-14||

\threelineshloka
{तमुवाच ततो गङ्गा सर्वदेवपुरोहितम्}
{अशक्ता धारणे देव तव तेजः समुद्धतम्}
{दह्यमानाग्निना तेन सम्प्रव्यथितचेतना} %||1-36-15||

\twolineshloka
{अथाब्रवीदिदं गङ्गां सर्वदेवहुताशनः}
{इह हैमवते पादे गर्भोऽयं संनिवेश्यताम्} %||1-36-16||

\twolineshloka
{श्रुत्वा त्वग्निवचो गङ्गा तं गर्भमतिभास्वरम्}
{उत्ससर्ज महातेजाः स्रोतोभ्यो हि तदानघ} %||1-36-17||

\twolineshloka
{यदस्या निर्गतं तस्मात्तप्तजाम्बूनदप्रभम्}
{काञ्चनं धरणीं प्राप्तं हिरण्यममलं शुभम्} %||1-36-18||

\twolineshloka
{ताम्रं कार्ष्णायसं चैव तैक्ष्ण्यादेवाभिजायत}
{मलं तस्याभवत्तत्र त्रपुसीसकमेव च} %||1-36-19||

{तदेतद्धरणीं प्राप्य नानाधातुरवर्धत॥\devanumber{20}॥} %||1-36-20|| (Check)

\addtocounter{shlokacount}{1}

\twolineshloka
{निक्षिप्तमात्रे गर्भे तु तेजोभिरभिरञ्जितम्}
{सर्वं पर्वतसंनद्धं सौवर्णमभवद्वनम्} %||1-36-21||

\twolineshloka
{जातरूपमिति ख्यातं तदा प्रभृति राघव}
{सुवर्णं पुरुषव्याघ्र हुताशनसमप्रभम्} %||1-36-22||

\twolineshloka
{तं कुमारं ततो जातं सेन्द्राः सहमरुद्गणाः}
{क्षीरसम्भावनार्थाय कृत्तिकाः समयोजयन्} %||1-36-23||

\twolineshloka
{ताः क्षीरं जातमात्रस्य कृत्वा समयमुत्तमम्}
{ददुः पुत्रोऽयमस्माकं सर्वासामिति निश्चिताः} %||1-36-24||

\twolineshloka
{ततस्तु देवताः सर्वाः कार्तिकेय इति ब्रुवन्}
{पुत्रस्त्रैलोक्य विख्यातो भविष्यति न संशयः} %||1-36-25||

\twolineshloka
{तेषां तद्वचनं श्रुत्वा स्कन्नं गर्भपरिस्रवे}
{स्नापयन्परया लक्ष्म्या दीप्यमानमिवानलम्} %||1-36-26||

\twolineshloka
{स्कन्द इत्यब्रुवन्देवाः स्कन्नं गर्भपरिस्रवात्}
{कार्तिकेयं महाभागं काकुत्स्थज्वलनोपमम्} %||1-36-27||

\twolineshloka
{प्रादुर्भूतं ततः क्षीरं कृत्तिकानामनुत्तमम्}
{षण्णां षडाननो भूत्वा जग्राह स्तनजं पयः} %||1-36-28||

\twolineshloka
{गृहीत्वा क्षीरमेकाह्ना सुकुमार वपुस्तदा}
{अजयत्स्वेन वीर्येण दैत्यसैन्यगणान्विभुः} %||1-36-29||

\twolineshloka
{सुरसेनागणपतिं ततस्तममलद्युतिम्}
{अभ्यषिञ्चन्सुरगणाः समेत्याग्निपुरोगमाः} %||1-36-30||

\twolineshloka
{एष ते राम गङ्गाया विस्तरोऽभिहितो मया}
{कुमारसम्भवश्चैव धन्यः पुण्यस्तथैव च} %||1-37-31||


{॥इत्यार्षे श्रीमद्रामायणे वाल्मीकीये आदिकाव्ये बालकाण्डे षट्त्रिंशः सर्गः॥}

\sect{सप्तत्रिंशः सर्गः}

\twolineshloka
{तां कथां कौशिको रामे निवेद्य मधुराक्षरम्}
{पुनरेवापरं वाक्यं काकुत्स्थमिदमब्रवीत्} %||1-37-1||

\twolineshloka
{अयोध्याधिपतिः शूरः पूर्वमासीन्नराधिपः}
{सगरो नाम धर्मात्मा प्रजाकामः स चाप्रजः} %||1-37-2||

\twolineshloka
{वैदर्भदुहिता राम केशिनी नाम नामतः}
{ज्येष्ठा सगरपत्नी सा धर्मिष्ठा सत्यवादिनी} %||1-37-3||

\twolineshloka
{अरिष्टनेमिदुहिता रूपेणाप्रतिमा भुवि}
{द्वितीया सगरस्यासीत्पत्नी सुमतिसंज्ञिता} %||1-37-4||

\twolineshloka
{ताभ्यां सह तदा राजा पत्नीभ्यां तप्तवांस्तपः}
{हिमवन्तं समासाद्य भृगुप्रस्रवणे गिरौ} %||1-37-5||

\twolineshloka
{अथ वर्ष शते पूर्णे तपसाराधितो मुनिः}
{सगराय वरं प्रादाद्भृगुः सत्यवतां वरः} %||1-37-6||

\twolineshloka
{अपत्यलाभः सुमहान्भविष्यति तवानघ}
{कीर्तिं चाप्रतिमां लोके प्राप्स्यसे पुरुषर्षभ} %||1-37-7||

\twolineshloka
{एका जनयिता तात पुत्रं वंशकरं तव}
{षष्टिं पुत्रसहस्राणि अपरा जनयिष्यति} %||1-37-8||

\twolineshloka
{भाषमाणं नरव्याघ्रं राजपत्न्यौ प्रसाद्य तम्}
{ऊचतुः परमप्रीते कृताञ्जलिपुटे तदा} %||1-37-9||

\twolineshloka
{एकः कस्याः सुतो ब्रह्मन्का बहूञ्जनयिष्यति}
{श्रोतुमिच्छावहे ब्रह्मन्सत्यमस्तु वचस्तव} %||1-37-10||

\twolineshloka
{तयोस्तद्वचनं श्रुत्वा भृगुः परम धार्मिकः}
{उवाच परमां वाणीं स्वच्छन्दोऽत्र विधीयताम्} %||1-37-11||

\twolineshloka
{एको वंशकरो वास्तु बहवो वा महाबलाः}
{कीर्तिमन्तो महोत्साहाः का वा कं वरमिच्छति} %||1-37-12||

\twolineshloka
{मुनेस्तु वचनं श्रुत्वा केशिनी रघुनन्दन}
{पुत्रं वंशकरं राम जग्राह नृपसंनिधौ} %||1-37-13||

\twolineshloka
{षष्टिं पुत्रसहस्राणि सुपर्णभगिनी तदा}
{महोत्साहान्कीर्तिमतो जग्राह सुमतिः सुतान्} %||1-37-14||

\twolineshloka
{प्रदक्षिणमृषिं कृत्वा शिरसाभिप्रणम्य च}
{जगाम स्वपुरं राजा सभार्या रघुनन्दन} %||1-37-15||

\twolineshloka
{अथ काले गते तस्मिञ्ज्येष्ठा पुत्रं व्यजायत}
{असमञ्ज इति ख्यातं केशिनी सगरात्मजम्} %||1-37-16||

\twolineshloka
{सुमतिस्तु नरव्याघ्र गर्भतुम्बं व्यजायत}
{षष्टिः पुत्रसहस्राणि तुम्बभेदाद्विनिःसृताः} %||1-37-17||

\twolineshloka
{घृतपूर्णेषु कुम्भेषु धात्र्यस्तान्समवर्धयन्}
{कालेन महता सर्वे यौवनं प्रतिपेदिरे} %||1-37-18||

\twolineshloka
{अथ दीर्घेण कालेन रूपयौवनशालिनः}
{षष्टिः पुत्रसहस्राणि सगरस्याभवंस्तदा} %||1-37-19||

\threelineshloka
{स च ज्येष्ठो नरश्रेष्ठ सगरस्यात्मसम्भवः}
{बालान्गृहीत्वा तु जले सरय्वा रघुनन्दन}
{प्रक्षिप्य प्रहसन्नित्यं मज्जतस्तान्निरीक्ष्य वै} %||1-37-20||

{पौराणामहिते युक्तः पित्रा निर्वासितः पुरात्॥\devanumber{21}॥} %||1-37-21|| (Check)

\addtocounter{shlokacount}{1}

\twolineshloka
{तस्य पुत्रोंऽशुमान्नाम असमञ्जस्य वीर्यवान्}
{सम्मतः सर्वलोकस्य सर्वस्यापि प्रियंवदः} %||1-37-22||

\twolineshloka
{ततः कालेन महता मतिः समभिजायत}
{सगरस्य नरश्रेष्ठ यजेयमिति निश्चिता} %||1-37-23||

\twolineshloka
{स कृत्वा निश्चयं राजा सोपाध्यायगणस्तदा}
{यज्ञकर्मणि वेदज्ञो यष्टुं समुपचक्रमे} %||1-38-24||


{॥इत्यार्षे श्रीमद्रामायणे वाल्मीकीये आदिकाव्ये बालकाण्डे सप्तत्रिंशः सर्गः॥}

\sect{अष्टात्रिंशः सर्गः}

\twolineshloka
{विश्वामित्रवचः श्रुत्वा कथान्ते रघुनन्दन}
{उवाच परमप्रीतो मुनिं दीप्तमिवानलम्} %||1-38-1||

\twolineshloka
{श्रोतुमिछामि भद्रं ते विस्तरेण कथामिमाम्}
{पूर्वको मे कथं ब्रह्मन्यज्ञं वै समुपाहरत्} %||1-38-2||

\twolineshloka
{विश्वामित्रस्तु काकुत्स्थमुवाच प्रहसन्निव}
{श्रूयतां विस्तरो राम सगरस्य महात्मनः} %||1-38-3||

\twolineshloka
{शङ्करश्वशुरो नाम हिमवानचलोत्तमः}
{विन्ध्यपर्वतमासाद्य निरीक्षेते परस्परम्} %||1-38-4||

\twolineshloka
{तयोर्मध्ये प्रवृत्तोऽभूद्यज्ञः स पुरुषोत्तम}
{स हि देशो नरव्याघ्र प्रशस्तो यज्ञकर्मणि} %||1-38-5||

\twolineshloka
{तस्याश्वचर्यां काकुत्स्थ दृढधन्वा महारथः}
{अंशुमानकरोत्तात सगरस्य मते स्थितः} %||1-38-6||

\twolineshloka
{तस्य पर्वणि तं यज्ञं यजमानस्य वासवः}
{राक्षसीं तनुमास्थाय यज्ञियाश्वमपाहरत्} %||1-38-7||

\twolineshloka
{ह्रियमाणे तु काकुत्स्थ तस्मिन्नश्वे महात्मनः}
{उपाध्याय गणाः सर्वे यजमानमथाब्रुवन्} %||1-38-8||

\twolineshloka
{अयं पर्वणि वेगेन यज्ञियाश्वोऽपनीयते}
{हर्तारं जहि काकुत्स्थ हयश्चैवोपनीयताम्} %||1-38-9||

\twolineshloka
{यज्ञच्छिद्रं भवत्येतत्सर्वेषामशिवाय नः}
{तत्तथा क्रियतां राजन्यथाछिद्रः क्रतुर्भवेत्} %||1-38-10||

\twolineshloka
{उपाध्याय वचः श्रुत्वा तस्मिन्सदसि पार्थिवः}
{षष्टिं पुत्रसहस्राणि वाक्यमेतदुवाच ह} %||1-38-11||

\twolineshloka
{गतिं पुत्रा न पश्यामि रक्षसां पुरुषर्षभाः}
{मन्त्रपूतैर्महाभागैरास्थितो हि महाक्रतुः} %||1-38-12||

\twolineshloka
{तद्गच्छत विचिन्वध्वं पुत्रका भद्रमस्तु वः}
{समुद्रमालिनीं सर्वां पृथिवीमनुगच्छत} %||1-38-13||

{एकैकं योजनं पुत्रा विस्तारमभिगच्छत॥\devanumber{14}॥} %||1-38-14|| (Check)

\addtocounter{shlokacount}{1}

\twolineshloka
{यावत्तुरगसन्दर्शस्तावत्खनत मेदिनीम्}
{तमेव हयहर्तारं मार्गमाणा ममाज्ञया} %||1-38-15||

\twolineshloka
{दीक्षितः पौत्रसहितः सोपाध्यायगणो ह्यहम्}
{इह स्थास्यामि भद्रं वो यावत्तुरगदर्शनम्} %||1-38-16||

\twolineshloka
{इत्युक्त्वा हृष्टमनसो राजपुत्रा महाबलाः}
{जग्मुर्महीतलं राम पितुर्वचनयन्त्रिताः} %||1-38-17||

\twolineshloka
{योजनायामविस्तारमेकैको धरणीतलम्}
{बिभिदुः पुरुषव्याघ्र वज्रस्पर्शसमैर्भुजैः} %||1-38-18||

\twolineshloka
{शूलैरशनिकल्पैश्च हलैश्चापि सुदारुणैः}
{भिद्यमाना वसुमती ननाद रघुनन्दन} %||1-38-19||

\twolineshloka
{नागानां वध्यमानानामसुराणां च राघव}
{राक्षसानां च दुर्धर्षः सत्त्वानां निनदोऽभवत्} %||1-38-20||

\twolineshloka
{योजनानां सहस्राणि षष्टिं तु रघुनन्दन}
{बिभिदुर्धरणीं वीरा रसातलमनुत्तमम्} %||1-38-21||

\twolineshloka
{एवं पर्वतसम्बाधं जम्बूद्वीपं नृपात्मजाः}
{खनन्तो नृपशार्दूल सर्वतः परिचक्रमुः} %||1-38-22||

\twolineshloka
{ततो देवाः सगन्धर्वाः सासुराः सहपन्नगाः}
{सम्भ्रान्तमनसः सर्वे पितामहमुपागमन्} %||1-38-23||

\twolineshloka
{ते प्रसाद्य महात्मानं विषण्णवदनास्तदा}
{ऊचुः परमसन्त्रस्ताः पितामहमिदं वचः} %||1-38-24||

\twolineshloka
{भगवन्पृथिवी सर्वा खन्यते सगरात्मजैः}
{बहवश्च महात्मानो वध्यन्ते जलचारिणः} %||1-38-25||

\twolineshloka
{अयं यज्ञहनोऽस्माकमनेनाश्वोऽपनीयते}
{इति ते सर्वभूतानि निघ्नन्ति सगरात्मजः} %||1-39-26||


{॥इत्यार्षे श्रीमद्रामायणे वाल्मीकीये आदिकाव्ये बालकाण्डे अष्टात्रिंशः सर्गः॥}

\sect{एकोनचत्वारिंशः सर्गः}

\twolineshloka
{देवतानां वचः श्रुत्वा भगवान्वै पितामहः}
{प्रत्युवाच सुसन्त्रस्तान्कृतान्तबलमोहितान्} %||1-39-1||

\twolineshloka
{यस्येयं वसुधा कृत्स्ना वासुदेवस्य धीमतः}
{कापिलं रूपमास्थाय धारयत्यनिशं धराम्} %||1-39-2||

\twolineshloka
{पृथिव्याश्चापि निर्भेदो दृष्ट एव सनातनः}
{सगरस्य च पुत्राणां विनाशोऽदीर्घजीविनाम्} %||1-39-3||

\twolineshloka
{पितामहवचः श्रुत्वा त्रयस्त्रिंशदरिन्दमः}
{देवाः परमसंहृष्टाः पुनर्जग्मुर्यथागतम्} %||1-39-4||

\twolineshloka
{सगरस्य च पुत्राणां प्रादुरासीन्महात्मनाम्}
{पृथिव्यां भिद्यमानायां निर्घात सम निःस्वनः} %||1-39-5||

\twolineshloka
{ततो भित्त्वा महीं सर्वां कृत्वा चापि प्रदक्षिणम्}
{सहिताः सगराः सर्वे पितरं वाक्यमब्रुवन्} %||1-39-6||

\twolineshloka
{परिक्रान्ता मही सर्वा सत्त्ववन्तश्च सूदिताः}
{देवदानवरक्षांसि पिशाचोरगकिंनराः} %||1-39-7||

\twolineshloka
{न च पश्यामहेऽश्वं तमश्वहर्तारमेव च}
{किं करिष्याम भद्रं ते बुद्धिरत्र विचार्यताम्} %||1-39-8||

\twolineshloka
{तेषां तद्वचनं श्रुत्वा पुत्राणां राजसत्तमः}
{समन्युरब्रवीद्वाक्यं सगरो रघुनन्दन} %||1-39-9||

\twolineshloka
{भूयः खनत भद्रं वो निर्भिद्य वसुधातलम्}
{अश्वहर्तारमासाद्य कृतार्थाश्च निवर्तथ} %||1-39-10||

\twolineshloka
{पितुर्वचनमास्थाय सगरस्य महात्मनः}
{षष्टिः पुत्रसहस्राणि रसातलमभिद्रवन्} %||1-39-11||

\twolineshloka
{खन्यमाने ततस्तस्मिन्ददृशुः पर्वतोपमम्}
{दिशागजं विरूपाक्षं धारयन्तं महीतलम्} %||1-39-12||

\twolineshloka
{सपर्वतवनां कृत्स्नां पृथिवीं रघुनन्दन}
{शिरसा धारयामास विरूपाक्षो महागजः} %||1-39-13||

\twolineshloka
{यदा पर्वणि काकुत्स्थ विश्रमार्थं महागजः}
{खेदाच्चालयते शीर्षं भूमिकम्पस्तधा भवेत्} %||1-39-14||

\twolineshloka
{तं ते प्रदक्षिणं कृत्वा दिशापालं महागजम्}
{मानयन्तो हि ते राम जग्मुर्भित्त्वा रसातलम्} %||1-39-15||

\twolineshloka
{ततः पूर्वां दिशं भित्त्वा दक्षिणां बिभिदुः पुनः}
{दक्षिणस्यामपि दिशि ददृशुस्ते महागजम्} %||1-39-16||

\twolineshloka
{महापद्मं महात्मानं सुमहापर्वतोपमम्}
{शिरसा धारयन्तं ते विस्मयं जग्मुरुत्तमम्} %||1-39-17||

\twolineshloka
{ततः प्रदक्षिणं कृत्वा सगरस्य महात्मनः}
{षष्टिः पुत्रसहस्राणि पश्चिमां बिभिदुर्दिशम्} %||1-39-18||

\twolineshloka
{पश्चिमायामपि दिशि महान्तमचलोपमम्}
{दिशागजं सौमनसं ददृशुस्ते महाबलाः} %||1-39-19||

\twolineshloka
{तं ते प्रदक्षिणं कृत्वा पृष्ट्वा चापि निरामयम्}
{खनन्तः समुपक्रान्ता दिशं सोमवतीं तदा} %||1-39-20||

\twolineshloka
{उत्तरस्यां रघुश्रेष्ठ ददृशुर्हिमपाण्डुरम्}
{भद्रं भद्रेण वपुषा धारयन्तं महीमिमाम्} %||1-39-21||

\twolineshloka
{समालभ्य ततः सर्वे कृत्वा चैनं प्रदक्षिणम्}
{षष्टिः पुत्रसहस्राणि बिभिदुर्वसुधातलम्} %||1-39-22||

\twolineshloka
{ततः प्रागुत्तरां गत्वा सागराः प्रथितां दिशम्}
{रोषादभ्यखनन्सर्वे पृथिवीं सगरात्मजाः} %||1-39-23||

\twolineshloka
{ददृशुः कपिलं तत्र वासुदेवं सनातनम्}
{हयं च तस्य देवस्य चरन्तमविदूरतः} %||1-39-24||

\twolineshloka
{ते तं यज्ञहनं ज्ञात्वा क्रोधपर्याकुलेक्षणाः}
{अभ्यधावन्त सङ्क्रुद्धास्तिष्ठ तिष्ठेति चाब्रुवन्} %||1-39-25||

\twolineshloka
{अस्माकं त्वं हि तुरगं यज्ञियं हृतवानसि}
{दुर्मेधस्त्वं हि सम्प्राप्तान्विद्धि नः सगरात्मजान्} %||1-39-26||

\twolineshloka
{श्रुत्वा तद्वचनं तेषां कपिलो रघुनन्दन}
{रोषेण महताविष्टो हुङ्कारमकरोत्तदा} %||1-39-27||

\twolineshloka
{ततस्तेनाप्रमेयेन कपिलेन महात्मना}
{भस्मराशीकृताः सर्वे काकुत्स्थ सगरात्मजाः} %||1-40-28||


{॥इत्यार्षे श्रीमद्रामायणे वाल्मीकीये आदिकाव्ये बालकाण्डे एकोनचत्वारिंशः सर्गः॥}

\sect{चत्वारिंशः सर्गः}

\twolineshloka
{पुत्रांश्चिरगताञ्ज्ञात्वा सगरो रघुनन्दन}
{नप्तारमब्रवीद्राजा दीप्यमानं स्वतेजसा} %||1-40-1||

\twolineshloka
{शूरश्च कृतविद्यश्च पूर्वैस्तुल्योऽसि तेजसा}
{पितॄणां गतिमन्विच्छ येन चाश्वोऽपहारितः} %||1-40-2||

\twolineshloka
{अन्तर्भौमानि सत्त्वानि वीर्यवन्ति महान्ति च}
{तेषां त्वं प्रतिघातार्थं सासिं गृह्णीष्व कार्मुकम्} %||1-40-3||

\twolineshloka
{अभिवाद्याभिवाद्यांस्त्वं हत्वा विघ्नकरानपि}
{सिद्धार्थः संनिवर्तस्व मम यज्ञस्य पारगः} %||1-40-4||

\twolineshloka
{एवमुक्तोंऽशुमान्सम्यक्सगरेण महात्मना}
{धनुरादाय खड्गं च जगाम लघुविक्रमः} %||1-40-5||

\twolineshloka
{स खातं पितृभिर्मार्गमन्तर्भौमं महात्मभिः}
{प्रापद्यत नरश्रेष्ठ तेन राज्ञाभिचोदितः} %||1-40-6||

\twolineshloka
{दैत्यदानवरक्षोभिः पिशाचपतगोरगैः}
{पूज्यमानं महातेजा दिशागजमपश्यत} %||1-40-7||

\twolineshloka
{स तं प्रदक्षिणं कृत्वा पृष्ट्वा चैव निरामयम्}
{पितॄन्स परिपप्रच्छ वाजिहर्तारमेव च} %||1-40-8||

\twolineshloka
{दिशागजस्तु तच्छ्रुत्वा प्रीत्याहांशुमतो वचः}
{आसमञ्जकृतार्थस्त्वं सहाश्वः शीघ्रमेष्यसि} %||1-40-9||

\twolineshloka
{तस्य तद्वचनं श्रुत्वा सर्वानेव दिशागजान्}
{यथाक्रमं यथान्यायं प्रष्टुं समुपचक्रमे} %||1-40-10||

\twolineshloka
{तैश्च सर्वैर्दिशापालैर्वाक्यज्ञैर्वाक्यकोविदैः}
{पूजितः सहयश्चैव गन्तासीत्यभिचोदितः} %||1-40-11||

\twolineshloka
{तेषां तद्वचनं श्रुत्वा जगाम लघुविक्रमः}
{भस्मराशीकृता यत्र पितरस्तस्य सागराः} %||1-40-12||

\twolineshloka
{स दुःखवशमापन्नस्त्वसमञ्जसुतस्तदा}
{चुक्रोश परमार्तस्तु वधात्तेषां सुदुःखितः} %||1-40-13||

\twolineshloka
{यज्ञियं च हयं तत्र चरन्तमविदूरतः}
{ददर्श पुरुषव्याघ्रो दुःखशोकसमन्वितः} %||1-40-14||

\twolineshloka
{ददर्श पुरुषव्याघ्रो कर्तुकामो जलक्रियाम्}
{सलिलार्थी महातेजा न चापश्यज्जलाशयम्} %||1-40-15||

\twolineshloka
{विसार्य निपुणां दृष्टिं ततोऽपश्यत्खगाधिपम्}
{पितॄणां मातुलं राम सुपर्णमनिलोपमम्} %||1-40-16||

\twolineshloka
{स चैनमब्रवीद्वाक्यं वैनतेयो महाबलः}
{मा शुचः पुरुषव्याघ्र वधोऽयं लोकसम्मतः} %||1-40-17||

\twolineshloka
{कपिलेनाप्रमेयेन दग्धा हीमे महाबलाः}
{सलिलं नार्हसि प्राज्ञ दातुमेषां हि लौकिकम्} %||1-40-18||

\twolineshloka
{गङ्गा हिमवतो ज्येष्ठा दुहिता पुरुषर्षभ}
{भस्मराशीकृतानेतान्पावयेल्लोकपावनी} %||1-40-19||

\twolineshloka
{तया क्लिन्नमिदं भस्म गङ्गया लोककान्तया}
{षष्टिं पुत्रसहस्राणि स्वर्गलोकं नयिष्यति} %||1-40-20||

\twolineshloka
{गच्छ चाश्वं महाभाग सङ्गृह्य पुरुषर्षभ}
{यज्ञं पैतामहं वीर निर्वर्तयितुमर्हसि} %||1-40-21||

\twolineshloka
{सुपर्णवचनं श्रुत्वा सोंऽशुमानतिवीर्यवान्}
{त्वरितं हयमादाय पुनरायान्महायशाः} %||1-40-22||

\twolineshloka
{ततो राजानमासाद्य दीक्षितं रघुनन्दन}
{न्यवेदयद्यथावृत्तं सुपर्णवचनं तथा} %||1-40-23||

\twolineshloka
{तच्छ्रुत्वा घोरसङ्काशं वाक्यमंशुमतो नृपः}
{यज्ञं निर्वर्तयामास यथाकल्पं यथाविधि} %||1-40-24||

\twolineshloka
{स्वपुरं चागमच्छ्रीमानिष्टयज्ञो महीपतिः}
{गङ्गायाश्चागमे राजा निश्चयं नाध्यगच्छत} %||1-40-25||

\twolineshloka
{अगत्वा निश्चयं राजा कालेन महता महान्}
{त्रिंशद्वर्षसहस्राणि राज्यं कृत्वा दिवं गतः} %||1-41-26||


{॥इत्यार्षे श्रीमद्रामायणे वाल्मीकीये आदिकाव्ये बालकाण्डे चत्वारिंशः सर्गः॥}

\sect{एकचत्वारिंशः सर्गः}

\twolineshloka
{कालधर्मं गते राम सगरे प्रकृतीजनाः}
{राजानं रोचयामासुरंशुमन्तं सुधार्मिकम्} %||1-41-1||

\twolineshloka
{स राजा सुमहानासीदंशुमान्रघुनन्दन}
{तस्य पुत्रो महानासीद्दिलीप इति विश्रुतः} %||1-41-2||

\twolineshloka
{तस्मिन्राज्यं समावेश्य दिलीपे रघुनन्दन}
{हिमवच्छिखरे रम्ये तपस्तेपे सुदारुणम्} %||1-41-3||

\twolineshloka
{द्वाद्त्रिंशच्च सहस्राणि वर्षाणि सुमहायशाः}
{तपोवनगतो राजा स्वर्गं लेभे तपोधनः} %||1-41-4||

\twolineshloka
{दिलीपस्तु महातेजाः श्रुत्वा पैतामहं वधम्}
{दुःखोपहतया बुद्ध्या निश्चयं नाध्यगच्छत} %||1-41-5||

\twolineshloka
{कथं गङ्गावतरणं कथं तेषां जलक्रिया}
{तारयेयं कथं चैतानिति चिन्ता परोऽभवत्} %||1-41-6||

\twolineshloka
{तस्य चिन्तयतो नित्यं धर्मेण विदितात्मनः}
{पुत्रो भगीरथो नाम जज्ञे परमधार्मिकः} %||1-41-7||

\twolineshloka
{दिलीपस्तु महातेजा यज्ञैर्बहुभिरिष्टवान्}
{त्रिंशद्वर्षसहस्राणि राजा राज्यमकारयत्} %||1-41-8||

\twolineshloka
{अगत्वा निश्चयं राजा तेषामुद्धरणं प्रति}
{व्याधिना नरशार्दूल कालधर्ममुपेयिवान्} %||1-41-9||

\twolineshloka
{इन्द्रलोकं गतो राजा स्वार्जितेनैव कर्मणा}
{रम्ये भगीरथं पुत्रमभिषिच्य नरर्षभः} %||1-41-10||

\twolineshloka
{भगीरथस्तु राजर्षिर्धार्मिको रघुनन्दन}
{अनपत्यो महातेजाः प्रजाकामः स चाप्रजः} %||1-41-11||

\twolineshloka
{स तपो दीर्घमातिष्ठद्गोकर्णे रघुनन्दन}
{ऊर्ध्वबाहुः पञ्चतपा मासाहारो जितेन्द्रियः} %||1-41-12||

\twolineshloka
{तस्य वर्षसहस्राणि घोरे तपसि तिष्ठतः}
{सुप्रीतो भगवान्ब्रह्मा प्रजानां पतिरीश्वरः} %||1-41-13||

\twolineshloka
{ततः सुरगणैः सार्धमुपागम्य पितामहः}
{भगीरथं महात्मानं तप्यमानमथाब्रवीत्} %||1-41-14||

\twolineshloka
{भगीरथ महाभाग प्रीतस्तेऽहं जनेश्वर}
{तपसा च सुतप्तेन वरं वरय सुव्रत} %||1-41-15||

\twolineshloka
{तमुवाच महातेजाः सर्वलोकपितामहम्}
{भगीरथो महाभागः कृताञ्जलिरवस्थितः} %||1-41-16||

\twolineshloka
{यदि मे भगवान्प्रीतो यद्यस्ति तपसः फलम्}
{सगरस्यात्मजाः सर्वे मत्तः सलिलमाप्नुयुः} %||1-41-17||

\twolineshloka
{गङ्गायाः सलिलक्लिन्ने भस्मन्येषां महात्मनाम्}
{स्वर्गं गच्छेयुरत्यन्तं सर्वे मे प्रपितामहाः} %||1-41-18||

\twolineshloka
{देया च सन्ततोर्देव नावसीदेत्कुलं च नः}
{इक्ष्वाकूणां कुले देव एष मेऽस्तु वरः परः} %||1-41-19||

\twolineshloka
{उक्तवाक्यं तु राजानं सर्वलोकपितामहः}
{प्रत्युवाच शुभां वाणीं मधुरां मधुराक्षराम्} %||1-41-20||

\twolineshloka
{मनोरथो महानेष भगीरथ महारथ}
{एवं भवतु भद्रं ते इक्ष्वाकुकुलवर्धन} %||1-41-21||

\twolineshloka
{इयं हैमवती गङ्गा ज्येष्ठा हिमवतः सुता}
{तां वै धारयितुं राजन्हरस्तत्र नियुज्यताम्} %||1-41-22||

\twolineshloka
{गङ्गायाः पतनं राजन्पृथिवी न सहिष्यते}
{तौ वै धारयितुं वीर नान्यं पश्यामि शूलिनः} %||1-41-23||

\twolineshloka
{तमेवमुक्त्वा राजानं गङ्गां चाभाष्य लोककृत्}
{जगाम त्रिदिवं देवः सह सर्वैर्मरुद्गणैः} %||1-42-24||


{॥इत्यार्षे श्रीमद्रामायणे वाल्मीकीये आदिकाव्ये बालकाण्डे एकचत्वारिंशः सर्गः॥}

\sect{द्विचत्वारिंशः सर्गः}

\twolineshloka
{देवदेवे गते तस्मिन्सोऽङ्गुष्ठाग्रनिपीडिताम्}
{कृत्वा वसुमतीं राम संवत्सरमुपासत} %||1-42-1||

\twolineshloka
{अथ संवत्सरे पूर्णे सर्वलोकनमस्कृतः}
{उमापतिः पशुपती राजानमिदमब्रवीत्} %||1-42-2||

\twolineshloka
{प्रीतस्तेऽहं नरश्रेष्ठ करिष्यामि तव प्रियम्}
{शिरसा धारयिष्यामि शैलराजसुतामहम्} %||1-42-3||

\threelineshloka
{ततो हैमवती ज्येष्ठा सर्वलोकनमस्कृता}
{तदा सातिमहद्रूपं कृत्वा वेगं च दुःसहम्}
{आकाशादपतद्राम शिवे शिवशिरस्युत} %||1-42-4||

\twolineshloka
{नैव सा निर्गमं लेखे जटामण्डलमोहिता}
{तत्रैवाबभ्रमद्देवी संवत्सरगणान्बहून्} %||1-42-5||

\twolineshloka
{अनेन तोषितश्चासीदत्यर्थं रघुनन्दन}
{विससर्ज ततो गङ्गां हरो बिन्दुसरः प्रति} %||1-42-6||

\twolineshloka
{गगनाच्छङ्करशिरस्ततो धरणिमागता}
{व्यसर्पत जलं तत्र तीव्रशब्दपुरस्कृतम्} %||1-42-7||

\twolineshloka
{ततो देवर्षिगन्धर्वा यक्षाः सिद्धगणास्तथा}
{व्यलोकयन्त ते तत्र गगनाद्गां गतां तदा} %||1-42-8||

\twolineshloka
{विमानैर्नगराकारैर्हयैर्गजवरैस्तथा}
{पारिप्लवगताश्चापि देवतास्तत्र विष्ठिताः} %||1-42-9||

\twolineshloka
{तदद्भुततमं लोके गङ्गा पतनमुत्तमम्}
{दिदृक्षवो देवगणाः समेयुरमितौजसः} %||1-42-10||

\twolineshloka
{सम्पतद्भिः सुरगणैस्तेषां चाभरणौजसा}
{शतादित्यमिवाभाति गगनं गततोयदम्} %||1-42-11||

\twolineshloka
{शिंशुमारोरगगणैर्मीनैरपि च चञ्चलैः}
{विद्युद्भिरिव विक्षिप्तैराकाशमभवत्तदा} %||1-42-12||

\twolineshloka
{पाण्डुरैः सलिलोत्पीडैः कीर्यमाणैः सहस्रधा}
{शारदाभ्रैरिवाक्रीत्णं गगनं हंससम्प्लवैः} %||1-42-13||

\twolineshloka
{क्वचिद्द्रुततरं याति कुटिलं क्वचिदायतम्}
{विनतं क्वचिदुद्धूतं क्वचिद्याति शनैः शनैः} %||1-42-14||

\twolineshloka
{सलिलेनैव सलिलं क्वचिदभ्याहतं पुनः}
{मुहुरूर्ध्वपथं गत्वा पपात वसुधां पुनः} %||1-42-15||

\twolineshloka
{तच्छङ्करशिरोभ्रष्टं भ्रष्टं भूमितले पुनः}
{व्यरोचत तदा तोयं निर्मलं गतकल्मषम्} %||1-42-16||

\twolineshloka
{तत्रर्षिगणगन्धर्वा वसुधातलवासिनः}
{भवाङ्गपतितं तोयं पवित्रमिति पस्पृशुः} %||1-42-17||

\twolineshloka
{शापात्प्रपतिता ये च गगनाद्वसुधातलम्}
{कृत्वा तत्राभिषेकं ते बभूवुर्गतकल्मषाः} %||1-42-18||

\twolineshloka
{धूपपापाः पुनस्तेन तोयेनाथ सुभास्वता}
{पुनराकाशमाविश्य स्वाँल्लोकान्प्रतिपेदिरे} %||1-42-19||

\twolineshloka
{मुमुदे मुदितो लोकस्तेन तोयेन भास्वता}
{कृताभिषेको गङ्गायां बभूव विगतक्लमः} %||1-42-20||

\twolineshloka
{भगीरथोऽपि राजर्षिर्दिव्यं स्यन्दनमास्थितः}
{प्रायादग्रे महातेजास्तं गङ्गा पृष्ठतोऽन्वगात्} %||1-42-21||

\twolineshloka
{देवाः सर्षिगणाः सर्वे दैत्यदानवराक्षसाः}
{गन्धर्वयक्षप्रवराः सकिंनरमहोरगाः} %||1-42-22||

\twolineshloka
{सर्वाश्चाप्सरसो राम भगीरथरथानुगाः}
{गङ्गामन्वगमन्प्रीताः सर्वे जलचराश्च ये} %||1-42-23||

\twolineshloka
{यतो भगीरथो राजा ततो गङ्गा यशस्विनी}
{जगाम सरितां श्रेष्ठा सर्वपापविनाशिनी} %||1-43-24||


{॥इत्यार्षे श्रीमद्रामायणे वाल्मीकीये आदिकाव्ये बालकाण्डे द्विचत्वारिंशः सर्गः॥}

\sect{त्रिचत्वारिंशः सर्गः}

\twolineshloka
{स गत्वा सागरं राजा गङ्गयानुगतस्तदा}
{प्रविवेश तलं भूमेर्यत्र ते भस्मसात्कृताः} %||1-43-1||

\twolineshloka
{भस्मन्यथाप्लुते राम गङ्गायाः सलिलेन वै}
{सर्व लोकप्रभुर्ब्रह्मा राजानमिदमब्रवीत्} %||1-43-2||

\twolineshloka
{तारिता नरशार्दूल दिवं याताश्च देववत्}
{षष्टिः पुत्रसहस्राणि सगरस्य महात्मनः} %||1-43-3||

\twolineshloka
{सागरस्य जलं लोके यावत्स्थास्यति पार्थिव}
{सगरस्यात्मजास्तावत्स्वर्गे स्थास्यन्ति देववत्} %||1-43-4||

\twolineshloka
{इयं च दुहिता ज्येष्ठा तव गङ्गा भविष्यति}
{त्वत्कृतेन च नाम्ना वै लोके स्थास्यति विश्रुता} %||1-43-5||

\twolineshloka
{गङ्गा त्रिपथगा नाम दिव्या भागीरथीति च}
{त्रिपथो भावयन्तीति ततस्त्रिपथगा स्मृता} %||1-43-6||

\twolineshloka
{पितामहानां सर्वेषां त्वमत्र मनुजाधिप}
{कुरुष्व सलिलं राजन्प्रतिज्ञामपवर्जय} %||1-43-7||

\twolineshloka
{पूर्वकेण हि ते राजंस्तेनातियशसा तदा}
{धर्मिणां प्रवरेणाथ नैष प्राप्तो मनोरथः} %||1-43-8||

\twolineshloka
{तथैवांशुमता तात लोकेऽप्रतिमतेजसा}
{गङ्गां प्रार्थयता नेतुं प्रतिज्ञा नापवर्जिता} %||1-43-9||

\twolineshloka
{राजर्षिणा गुणवता महर्षिसमतेजसा}
{मत्तुल्यतपसा चैव क्षत्रधर्मस्थितेन च} %||1-43-10||

\twolineshloka
{दिलीपेन महाभाग तव पित्रातितेजसा}
{पुनर्न शङ्किता नेतुं गङ्गां प्रार्थयतानघ} %||1-43-11||

\twolineshloka
{सा त्वया समतिक्रान्ता प्रतिज्ञा पुरुषर्षभ}
{प्राप्तोऽसि परमं लोके यशः परमसम्मतम्} %||1-43-12||

\twolineshloka
{यच्च गङ्गावतरणं त्वया कृतमरिन्दम}
{अनेन च भवान्प्राप्तो धर्मस्यायतनं महत्} %||1-43-13||

\twolineshloka
{प्लावयस्व त्वमात्मानं नरोत्तम सदोचिते}
{सलिले पुरुषव्याघ्र शुचिः पुण्यफलो भव} %||1-43-14||

\twolineshloka
{पितामहानां सर्वेषां कुरुष्व सलिलक्रियाम्}
{स्वस्ति तेऽस्तु गमिष्यामि स्वं लोकं गम्यतां नृप} %||1-43-15||

\twolineshloka
{इत्येवमुक्त्वा देवेशः सर्वलोकपितामहः}
{यथागतं तथागच्छद्देवलोकं महायशाः} %||1-43-16||

\threelineshloka
{भगीरथोऽपि राजर्षिः कृत्वा सलिलमुत्तमम्}
{यथाक्रमं यथान्यायं सागराणां महायशाः}
{कृतोदकः शुची राजा स्वपुरं प्रविवेश ह} %||1-43-17||

\threelineshloka
{समृद्धार्थो नरश्रेष्ठ स्वराज्यं प्रशशास ह}
{प्रमुमोद च लोकस्तं नृपमासाद्य राघव}
{नष्टशोकः समृद्धार्थो बभूव विगतज्वरः} %||1-43-18||

\twolineshloka
{एष ते राम गङ्गाया विस्तरोऽभिहितो मया}
{स्वस्ति प्राप्नुहि भद्रं ते सन्ध्याकालोऽतिवर्तते} %||1-43-19||

\twolineshloka
{धन्यं यशस्यमायुष्यं स्वर्ग्यं पुत्र्यमथापि च}
{इदमाख्यानमाख्यातं गङ्गावतरणं मया} %||1-44-20||


{॥इत्यार्षे श्रीमद्रामायणे वाल्मीकीये आदिकाव्ये बालकाण्डे त्रिचत्वारिंशः सर्गः॥}

\sect{चतुश्चत्वारिंशः सर्गः}

\twolineshloka
{विश्वामित्रवचः श्रुत्वा राघवः सहलक्ष्मणः}
{विस्मयं परमं गत्वा विश्वामित्रमथाब्रवीत्} %||1-44-1||

\twolineshloka
{अत्यद्भुतमिदं ब्रह्मन्कथितं परमं त्वया}
{गङ्गावतरणं पुण्यं सागरस्य च पूरणम्} %||1-44-2||

\twolineshloka
{तस्य सा शर्वरी सर्वा सह सौमित्रिणा तदा}
{जगाम चिन्तयानस्य विश्वामित्रकथां शुभाम्} %||1-44-3||

\twolineshloka
{ततः प्रभाते विमले विश्वामित्रं महामुनिम्}
{उवाच राघवो वाक्यं कृताह्निकमरिन्दमः} %||1-44-4||

\threelineshloka
{गता भगवती रात्रिः श्रोतव्यं परमं श्रुतम्}
{क्षणभूतेव सा रात्रिः संवृत्तेयं महातपः}
{इमां चिन्तयतः सर्वां निखिलेन कथां तव} %||1-44-5||

\threelineshloka
{तराम सरितां श्रेष्ठां पुण्यां त्रिपथगां नदीम्}
{नौरेषा हि सुखास्तीर्णा ऋषीणां पुण्यकर्मणाम्}
{भगवन्तमिह प्राप्तं ज्ञात्वा त्वरितमागता} %||1-44-6||

\twolineshloka
{तस्य तद्वचनं श्रुत्वा राघवस्य महात्मनः}
{सन्तारं कारयामास सर्षिसङ्घः सराघवः} %||1-44-7||

\twolineshloka
{उत्तरं तीरमासाद्य सम्पूज्यर्षिगणं तथ}
{गङ्गाकूले निविष्टास्ते विशालां ददृशुः पुरीम्} %||1-44-8||

\twolineshloka
{ततो मुनिवरस्तूर्णं जगाम सहराघवः}
{विशालां नगरीं रम्यां दिव्यां स्वर्गोपमां तदा} %||1-44-9||

\twolineshloka
{अथ रामो महाप्राज्ञो विश्वामित्रं महामुनिम्}
{पप्रच्छ प्राञ्जलिर्भूत्वा विशालामुत्तमां पुरीम्} %||1-44-10||

\twolineshloka
{कतरो राजवंशोऽयं विशालायां महामुने}
{श्रोतुमिच्छामि भद्रं ते परं कौतूहलं हि मे} %||1-44-11||

\twolineshloka
{तस्य तद्वचनं श्रुत्वा रामस्य मुनिपुङ्गवः}
{आख्यातुं तत्समारेभे विशालस्य पुरातनम्} %||1-44-12||

\twolineshloka
{श्रूयतां राम शक्रस्य कथां कथयतः शुभाम्}
{अस्मिन्देशे हि यद्वृत्तं शृणु तत्त्वेन राघव} %||1-44-13||

\twolineshloka
{पूर्वं कृतयुगे राम दितेः पुत्रा महाबलाः}
{अदितेश्च महाभागा वीर्यवन्तः सुधार्मिकाः} %||1-44-14||

\twolineshloka
{ततस्तेषां नरश्रेष्ठ बुद्धिरासीन्महात्मनाम्}
{अमरा निर्जराश्चैव कथं स्याम निरामयाः} %||1-44-15||

\twolineshloka
{तेषां चिन्तयतां राम बुद्धिरासीद्विपश्चिताम्}
{क्षीरोदमथनं कृत्वा रसं प्राप्स्याम तत्र वै} %||1-44-16||

\twolineshloka
{ततो निश्चित्य मथनं योक्त्रं कृत्वा च वासुकिम्}
{मन्थानं मन्दरं कृत्वा ममन्थुरमितौजसः} %||1-44-17||

\threelineshloka
{अथ धन्वन्तरिर्नाम अप्सराश्च सुवर्चसः}
{अप्सु निर्मथनादेव रसात्तस्माद्वरस्त्रियः}
{उत्पेतुर्मनुजश्रेष्ठ तस्मादप्सरसोऽभवन्} %||1-44-18||

\twolineshloka
{षष्टिः कोट्योऽभवंस्तासामप्सराणां सुवर्चसाम्}
{असङ्ख्येयास्तु काकुत्स्थ यास्तासां परिचारिकाः} %||1-44-19||

\twolineshloka
{न ताः स्म प्रतिगृह्णन्ति सर्वे ते देवदानवाः}
{अप्रतिग्रहणाच्चैव तेन साधारणाः स्मृताः} %||1-44-20||

\twolineshloka
{वरुणस्य ततः कन्या वारुणी रघुनन्दन}
{उत्पपात महाभागा मार्गमाणा परिग्रहम्} %||1-44-21||

\twolineshloka
{दितेः पुत्रा न तां राम जगृहुर्वरुणात्मजाम्}
{अदितेस्तु सुता वीर जगृहुस्तामनिन्दिताम्} %||1-44-22||

\twolineshloka
{असुरास्तेन दैतेयाः सुरास्तेनादितेः सुताः}
{हृष्टाः प्रमुदिताश्चासन्वारुणी ग्रहणात्सुराः} %||1-44-23||

\twolineshloka
{उच्चैःश्रवा हयश्रेष्ठो मणिरत्नं च कौस्तुभम्}
{उदतिष्ठन्नरश्रेष्ठ तथैवामृतमुत्तमम्} %||1-44-24||

\twolineshloka
{अथ तस्य कृते राम महानासीत्कुलक्षयः}
{अदितेस्तु ततः पुत्रा दितेः पुत्राण सूदयन्} %||1-44-25||

\twolineshloka
{अदितेरात्मजा वीरा दितेः पुत्रान्निजघ्निरे}
{तस्मिन्घोरे महायुद्धे दैतेयादित्ययोर्भृशम्} %||1-44-26||

\twolineshloka
{निहत्य दितिपुत्रांस्तु राज्यं प्राप्य पुरन्दरः}
{शशास मुदितो लोकान्सर्षिसङ्घान्सचारणान्} %||1-45-27||


{॥इत्यार्षे श्रीमद्रामायणे वाल्मीकीये आदिकाव्ये बालकाण्डे चतुश्चत्वारिंशः सर्गः॥}

\sect{पञ्चचत्वारिंशः सर्गः}

\twolineshloka
{हतेषु तेषु पुत्रेषु दितिः परमदुःखिता}
{मारीचं काश्यपं राम भर्तारमिदमब्रवीत्} %||1-45-1||

\twolineshloka
{हतपुत्रास्मि भगवंस्तव पुत्रैर्महाबलैः}
{शक्रहन्तारमिच्छामि पुत्रं दीर्घतपोऽर्जितम्} %||1-45-2||

\twolineshloka
{साहं तपश्चरिष्यामि गर्भं मे दातुमर्हसि}
{ईदृशं शक्रहन्तारं त्वमनुज्ञातुमर्हसि} %||1-45-3||

\twolineshloka
{तस्यास्तद्वचनं श्रुत्वा मारीचः काश्यपस्तदा}
{प्रत्युवाच महातेजा दितिं परमदुःखिताम्} %||1-45-4||

\twolineshloka
{एवं भवतु भद्रं ते शुचिर्भव तपोधने}
{जनयिष्यसि पुत्रं त्वं शक्र हन्तारमाहवे} %||1-45-5||

\twolineshloka
{पूर्णे वर्षसहस्रे तु शुचिर्यदि भविष्यसि}
{पुत्रं त्रैलोक्य हन्तारं मत्तस्त्वं जनयिष्यसि} %||1-45-6||

\twolineshloka
{एवमुक्त्वा महातेजाः पाणिना स ममार्ज ताम्}
{समालभ्य ततः स्वस्तीत्युक्त्वा स तपसे ययौ} %||1-45-7||

\twolineshloka
{गते तस्मिन्नरश्रेष्ठ दितिः परमहर्षिता}
{कुशप्लवनमासाद्य तपस्तेपे सुदारुणम्} %||1-45-8||

\twolineshloka
{तपस्तस्यां हि कुर्वत्यां परिचर्यां चकार ह}
{सहस्राक्षो नरश्रेष्ठ परया गुणसम्पदा} %||1-45-9||

\twolineshloka
{अग्निं कुशान्काष्ठमपः फलं मूलं तथैव च}
{न्यवेदयत्सहस्राक्षो यच्चान्यदपि काङ्क्षितम्} %||1-45-10||

\twolineshloka
{गात्रसंवाहनैश्चैव श्रमापनयनैस्तथा}
{शक्रः सर्वेषु कालेषु दितिं परिचचार ह} %||1-45-11||

\twolineshloka
{अथ वर्षसहस्रेतु दशोने रघु नन्दन}
{दितिः परमसम्प्रीता सहस्राक्षमथाब्रवीत्} %||1-45-12||

\twolineshloka
{तपश्चरन्त्या वर्षाणि दश वीर्यवतां वर}
{अवशिष्टानि भद्रं ते भ्रातरं द्रक्ष्यसे ततः} %||1-45-13||

\twolineshloka
{तमहं त्वत्कृते पुत्र समाधास्ये जयोत्सुकम्}
{त्रैलोक्यविजयं पुत्र सह भोक्ष्यसि विज्वरः} %||1-45-14||

\twolineshloka
{एवमुक्त्वा दितिः शक्रं प्राप्ते मध्यं दिवाकरे}
{निद्रयापहृता देवी पादौ कृत्वाथ शीर्षतः} %||1-45-15||

\twolineshloka
{दृष्ट्वा तामशुचिं शक्रः पादतः कृतमूर्धजाम्}
{शिरःस्थाने कृतौ पादौ जहास च मुमोद च} %||1-45-16||

\twolineshloka
{तस्याः शरीरविवरं विवेश च पुरन्दरः}
{गर्भं च सप्तधा राम बिभेद परमात्मवान्} %||1-45-17||

\twolineshloka
{बिध्यमानस्ततो गर्भो वज्रेण शतपर्वणा}
{रुरोद सुस्वरं राम ततो दितिरबुध्यत} %||1-45-18||

\twolineshloka
{मा रुदो मा रुदश्चेति गर्भं शक्रोऽभ्यभाषत}
{बिभेद च महातेजा रुदन्तमपि वासवः} %||1-45-19||

\twolineshloka
{न हन्तव्यो न हन्तव्य इत्येवं दितिरब्रवीत्}
{निष्पपात ततः शक्रो मातुर्वचनगौरवात्} %||1-45-20||

\twolineshloka
{प्राञ्जलिर्वज्रसहितो दितिं शक्रोऽभ्यभाषत}
{अशुचिर्देवि सुप्तासि पादयोः कृतमूर्धजा} %||1-45-21||

\twolineshloka
{तदन्तरमहं लब्ध्वा शक्रहन्तारमाहवे}
{अभिन्दं सप्तधा देवि तन्मे त्वं क्षन्तुमर्हसि} %||1-46-22||


{॥इत्यार्षे श्रीमद्रामायणे वाल्मीकीये आदिकाव्ये बालकाण्डे पञ्चचत्वारिंशः सर्गः॥}

\sect{षट्चत्वारिंशः सर्गः}

\twolineshloka
{सप्तधा तु कृते गर्भे दितिः परमदुःखिता}
{सहस्राक्षं दुराधर्षं वाक्यं सानुनयाब्रवीत्} %||1-46-1||

\twolineshloka
{ममापराधाद्गर्भोऽयं सप्तधा विफलीकृतः}
{नापराधोऽस्ति देवेश तवात्र बलसूदन} %||1-46-2||

\twolineshloka
{प्रियं तु कृतमिच्छामि मम गर्भविपर्यये}
{मरुतां सप्तं सप्तानां स्थानपाला भवन्त्विमे} %||1-46-3||

\twolineshloka
{वातस्कन्धा इमे सप्त चरन्तु दिवि पुत्रकाः}
{मारुता इति विख्याता दिव्यरूपा ममात्मजाः} %||1-46-4||

\twolineshloka
{ब्रह्मलोकं चरत्वेक इन्द्रलोकं तथापरः}
{दिवि वायुरिति ख्यातस्तृतीयोऽपि महायशाः} %||1-46-5||

\threelineshloka
{चत्वारस्तु सुरश्रेष्ठ दिशो वै तव शासनात्}
{सञ्चरिष्यन्ति भद्रं ते देवभूता ममात्मजाः}
{त्वत्कृतेनैव नाम्ना च मारुता इति विश्रुताः} %||1-46-6||

\twolineshloka
{तस्यास्तद्वचनं श्रुत्वा सहस्राक्षः पुरन्दरः}
{उवाच प्राञ्जलिर्वाक्यं दितिं बलनिषूदनः} %||1-46-7||

\twolineshloka
{सर्वमेतद्यथोक्तं ते भविष्यति न संशयः}
{विचरिष्यन्ति भद्रं ते देवभूतास्तवात्मजाः} %||1-46-8||

\twolineshloka
{एवं तौ निश्चयं कृत्वा मातापुत्रौ तपोवने}
{जग्मतुस्त्रिदिवं राम कृतार्थाविति नः श्रुतम्} %||1-46-9||

\twolineshloka
{एष देशः स काकुत्स्थ महेन्द्राध्युषितः पुरा}
{दितिं यत्र तपः सिद्धामेवं परिचचार सः} %||1-46-10||

\twolineshloka
{इक्ष्वाकोस्तु नरव्याघ्र पुत्रः परमधार्मिकः}
{अलम्बुषायामुत्पन्नो विशाल इति विश्रुतः} %||1-46-11||

{तेन चासीदिह स्थाने विशालेति पुरी कृता॥\devanumber{12}॥} %||1-46-12|| (Check)

\addtocounter{shlokacount}{1}

\twolineshloka
{विशालस्य सुतो राम हेमचन्द्रो महाबलः}
{सुचन्द्र इति विख्यातो हेमचन्द्रादनन्तरः} %||1-46-13||

\twolineshloka
{सुचन्द्रतनयो राम धूम्राश्व इति विश्रुतः}
{धूम्राश्वतनयश्चापि सृञ्जयः समपद्यत} %||1-46-14||

\twolineshloka
{सृञ्जयस्य सुतः श्रीमान्सहदेवः प्रतापवान्}
{कुशाश्वः सहदेवस्य पुत्रः परमधार्मिकः} %||1-46-15||

\twolineshloka
{कुशाश्वस्य महातेजाः सोमदत्तः प्रतापवान्}
{सोमदत्तस्य पुत्रस्तु काकुत्स्थ इति विश्रुतः} %||1-46-16||

\twolineshloka
{तस्य पुत्रो महातेजाः सम्प्रत्येष पुरीमिमाम्}
{आवसत्यमरप्रख्यः सुमतिर्नाम दुर्जयः} %||1-46-17||

\twolineshloka
{इक्ष्वाकोस्तु प्रसादेन सर्वे वैशालिका नृपाः}
{दीर्घायुषो महात्मानो वीर्यवन्तः सुधार्मिकाः} %||1-46-18||

\twolineshloka
{इहाद्य रजनीं राम सुखं वत्स्यामहे वयम्}
{श्वः प्रभाते नरश्रेष्ठ जनकं द्रष्टुमर्हसि} %||1-46-19||

\twolineshloka
{सुमतिस्तु महातेजा विश्वामित्रमुपागतम्}
{श्रुत्वा नरवरश्रेष्ठः प्रत्युद्गच्छन्महायशाः} %||1-46-20||

\twolineshloka
{पूजां च परमां कृत्वा सोपाध्यायः सबान्धवः}
{प्राञ्जलिः कुशलं पृष्ट्वा विश्वामित्रमथाब्रवीत्} %||1-46-21||

\twolineshloka
{धन्योऽस्म्यनुगृहीतोऽस्मि यस्य मे विषयं मुने}
{सम्प्राप्तो दर्शनं चैव नास्ति धन्यतरो मम} %||1-47-22||


{॥इत्यार्षे श्रीमद्रामायणे वाल्मीकीये आदिकाव्ये बालकाण्डे षट्चत्वारिंशः सर्गः॥}

\sect{सप्तचत्वारिंशः सर्गः}

\twolineshloka
{पृष्ट्वा तु कुशलं तत्र परस्परसमागमे}
{कथान्ते सुमतिर्वाक्यं व्याजहार महामुनिम्} %||1-47-1||

\twolineshloka
{इमौ कुमारौ भद्रं ते देवतुल्यपराक्रमौ}
{गजसिंहगती वीरौ शार्दूलवृषभोपमौ} %||1-47-2||

\twolineshloka
{पद्मपत्रविशालाक्षौ खड्गतूणीधनुर्धरौ}
{अश्विनाविव रूपेण समुपस्थितयौवनौ} %||1-47-3||

\twolineshloka
{यदृच्छयैव गां प्राप्तौ देवलोकादिवामरौ}
{कथं पद्भ्यामिह प्राप्तौ किमर्थं कस्य वा मुने} %||1-47-4||

\twolineshloka
{भूषयन्ताविमं देशं चन्द्रसूर्याविवाम्बरम्}
{परस्परस्य सदृशौ प्रमाणेङ्गितचेष्टितैः} %||1-47-5||

\twolineshloka
{किमर्थं च नरश्रेष्ठौ सम्प्राप्तौ दुर्गमे पथि}
{वरायुधधरौ वीरौ श्रोतुमिच्छामि तत्त्वतः} %||1-47-6||

\twolineshloka
{तस्य तद्वचनं श्रुत्वा यथावृत्तं न्यवेदयत्}
{सिद्धाश्रमनिवासं च राक्षसानां वधं तथा} %||1-47-7||

\threelineshloka
{विश्वामित्रवचः श्रुत्वा राजा परमहर्षितः}
{अतिथी परमौ प्राप्तौ पुत्रौ दशरथस्य तौ}
{पूजयामास विधिवत्सत्कारार्हौ महाबलौ} %||1-47-8||

\twolineshloka
{ततः परमसत्कारं सुमतेः प्राप्य राघवौ}
{उष्य तत्र निशामेकां जग्मतुर्मिथिलां ततः} %||1-47-9||

\twolineshloka
{तां दृष्ट्वा मुनयः सर्वे जनकस्य पुरीं शुभाम्}
{साधु साध्विति शंसन्तो मिथिलां समपूजयन्} %||1-47-10||

\twolineshloka
{मिथिलोपवने तत्र आश्रमं दृश्य राघवः}
{पुराणं निर्जनं रम्यं पप्रच्छ मुनिपुङ्गवम्} %||1-47-11||

\twolineshloka
{श्रीमदाश्रमसङ्काशं किं न्विदं मुनिवर्जितम्}
{श्रोतुमिच्छामि भगवन्कस्यायं पूर्व आश्रमः} %||1-47-12||

\twolineshloka
{तच्छ्रुता राघवेणोक्तं वाक्यं वाक्यविशारदः}
{प्रत्युवाच महातेजा विश्वमित्रो महामुनिः} %||1-47-13||

\twolineshloka
{हन्त ते कथयिष्यामि शृणु तत्त्वेन राघव}
{यस्यैतदाश्रमपदं शप्तं कोपान्महात्मना} %||1-47-14||

\twolineshloka
{गौतमस्य नरश्रेष्ठ पूर्वमासीन्महात्मनः}
{आश्रमो दिव्यसङ्काशः सुरैरपि सुपूजितः} %||1-47-15||

\twolineshloka
{स चेह तप आतिष्ठदहल्यासहितः पुरा}
{वर्षपूगान्यनेकानि राजपुत्र महायशः} %||1-47-16||

\twolineshloka
{तस्यान्तरं विदित्वा तु सहस्राक्षः शचीपतिः}
{मुनिवेषधरोऽहल्यामिदं वचनमब्रवीत्} %||1-47-17||

\twolineshloka
{ऋतुकालं प्रतीक्षन्ते नार्थिनः सुसमाहिते}
{सङ्गमं त्वहमिच्छामि त्वया सह सुमध्यमे} %||1-47-18||

\twolineshloka
{मुनिवेषं सहस्राक्षं विज्ञाय रघुनन्दन}
{मतिं चकार दुर्मेधा देवराजकुतूहलात्} %||1-47-19||

\threelineshloka
{अथाब्रवीत्सुरश्रेष्ठं कृतार्थेनान्तरात्मना}
{कृतार्थोऽसि सुरश्रेष्ठ गच्छ शीघ्रमितः प्रभो}
{आत्मानं मां च देवेश सर्वदा रक्ष मानदः} %||1-47-20||

\twolineshloka
{इन्द्रस्तु प्रहसन्वाक्यमहल्यामिदमब्रवीत्}
{सुश्रोणि परितुष्टोऽस्मि गमिष्यामि यथागतम्} %||1-47-21||

\twolineshloka
{एवं सङ्गम्य तु तया निश्चक्रामोटजात्ततः}
{स सम्भ्रमात्त्वरन्राम शङ्कितो गौतमं प्रति} %||1-47-22||

\fourlineshloka
{गौतमं स ददर्शाथ प्रविशन्तं महामुनिम्}
{देवदानवदुर्धर्षं तपोबलसमन्वितम्}
{तीर्थोदकपरिक्लिन्नं दीप्यमानमिवानलम्}
{गृहीतसमिधं तत्र सकुशं मुनिपुङ्गवम्} %||1-47-23||

{दृष्ट्वा सुरपतिस्त्रस्तो विषण्णवदनोऽभवत्॥\devanumber{24}॥} %||1-47-24|| (Check)

\addtocounter{shlokacount}{1}

\twolineshloka
{अथ दृष्ट्वा सहस्राक्षं मुनिवेषधरं मुनिः}
{दुर्वृत्तं वृत्तसम्पन्नो रोषाद्वचनमब्रवीत्} %||1-47-25||

\twolineshloka
{मम रूपं समास्थाय कृतवानसि दुर्मते}
{अकर्तव्यमिदं यस्माद्विफलस्त्वं भविष्यति} %||1-47-26||

\twolineshloka
{गौतमेनैवमुक्तस्य सरोषेण महात्मना}
{पेततुर्वृषणौ भूमौ सहस्राक्षस्य तत्क्षणात्} %||1-47-27||

\twolineshloka
{तथा शप्त्वा स वै शक्रं भार्यामपि च शप्तवान्}
{इह वर्षसहस्राणि बहूनि त्वं निवत्स्यसि} %||1-47-28||

\twolineshloka
{वायुभक्षा निराहारा तप्यन्ती भस्मशायिनी}
{अदृश्या सर्वभूतानामाश्रमेऽस्मिन्निवत्स्यसि} %||1-47-29||

\twolineshloka
{यदा चैतद्वनं घोरं रामो दशरथात्मजः}
{आगमिष्यति दुर्धर्षस्तदा पूता भविष्यसि} %||1-47-30||

\twolineshloka
{तस्यातिथ्येन दुर्वृत्ते लोभमोहविवर्जिता}
{मत्सकाशे मुदा युक्ता स्वं वपुर्धारयिष्यसि} %||1-47-31||

\threelineshloka
{एवमुक्त्वा महातेजा गौतमो दुष्टचारिणीम्}
{इममाश्रममुत्सृज्य सिद्धचारणसेविते}
{हिमवच्छिखरे रम्ये तपस्तेपे महातपाः} %||1-48-32||


{॥इत्यार्षे श्रीमद्रामायणे वाल्मीकीये आदिकाव्ये बालकाण्डे सप्तचत्वारिंशः सर्गः॥}

\sect{अष्टचत्वारिंशः सर्गः}

\twolineshloka
{अफलस्तु ततः शक्रो देवानग्निपुरोगमान्}
{अब्रवीत्त्रस्तवदनः सर्षिसङ्घान्सचारणान्} %||1-48-1||

\twolineshloka
{कुर्वता तपसो विघ्नं गौतमस्य महात्मनः}
{क्रोधमुत्पाद्य हि मया सुरकार्यमिदं कृतम्} %||1-48-2||

\twolineshloka
{अफलोऽस्मि कृतस्तेन क्रोधात्सा च निराकृता}
{शापमोक्षेण महता तपोऽस्यापहृतं मया} %||1-48-3||

\twolineshloka
{तन्मां सुरवराः सर्वे सर्षिसङ्घाः सचारणाः}
{सुरसाह्यकरं सर्वे सफलं कर्तुमर्हथ} %||1-48-4||

\twolineshloka
{शतक्रतोर्वचः श्रुत्वा देवाः साग्निपुरोगमाः}
{पितृदेवानुपेत्याहुः सह सर्वैर्मरुद्गणैः} %||1-48-5||

\twolineshloka
{अयं मेषः सवृषणः शक्रो ह्यवृषणः कृतः}
{मेषस्य वृषणौ गृह्य शक्रायाशु प्रयच्छत} %||1-48-6||

\twolineshloka
{अफलस्तु कृतो मेषः परां तुष्टिं प्रदास्यति}
{भवतां हर्षणार्थाय ये च दास्यन्ति मानवाः} %||1-48-7||

\twolineshloka
{अग्नेस्तु वचनं श्रुत्वा पितृदेवाः समागताः}
{उत्पाट्य मेषवृषणौ सहस्राक्षे न्यवेदयन्} %||1-48-8||

\twolineshloka
{तदा प्रभृति काकुत्स्थ पितृदेवाः समागताः}
{अफलान्भुञ्जते मेषान्फलैस्तेषामयोजयन्} %||1-48-9||

\twolineshloka
{इन्द्रस्तु मेषवृषणस्तदा प्रभृति राघव}
{गौतमस्य प्रभावेन तपसश्च महात्मनः} %||1-48-10||

\twolineshloka
{तदागच्छ महातेज आश्रमं पुण्यकर्मणः}
{तारयैनां महाभागामहल्यां देवरूपिणीम्} %||1-48-11||

\twolineshloka
{विश्वामित्रवचः श्रुत्वा राघवः सहलक्ष्मणः}
{विश्वामित्रं पुरस्कृत्य आश्रमं प्रविवेश ह} %||1-48-12||

\twolineshloka
{ददर्श च महाभागां तपसा द्योतितप्रभाम्}
{लोकैरपि समागम्य दुर्निरीक्ष्यां सुरासुरैः} %||1-48-13||

\twolineshloka
{प्रयत्नान्निर्मितां धात्रा दिव्यां मायामयीमिव}
{धूमेनाभिपरीताङ्गीं पूर्णचन्द्रप्रभामिव} %||1-48-14||

\twolineshloka
{सतुषारावृतां साभ्रां पूर्णचन्द्रप्रभामिव}
{मध्येऽम्भसो दुराधर्षां दीप्तां सूर्यप्रभामिव} %||1-48-15||

\twolineshloka
{स हि गौतमवाक्येन दुर्निरीक्ष्या बभूव ह}
{त्रयाणामपि लोकानां यावद्रामस्य दर्शनम्} %||1-48-16||

\twolineshloka
{राघवौ तु ततस्तस्याः पादौ जगृहतुस्तदा}
{स्मरन्ती गौतमवचः प्रतिजग्राह सा च तौ} %||1-48-17||

\twolineshloka
{पाद्यमर्घ्यं तथातिथ्यं चकार सुसमाहिता}
{प्रतिजग्राह काकुत्स्थो विधिदृष्टेन कर्मणा} %||1-48-18||

\twolineshloka
{पुष्पवृष्टिर्महत्यासीद्देवदुन्दुभिनिस्वनैः}
{गन्धर्वाप्सरसां चापि महानासीत्समागमः} %||1-48-19||

\twolineshloka
{साधु साध्विति देवास्तामहल्यां समपूजयन्}
{तपोबलविशुद्धाङ्गीं गौतमस्य वशानुगाम्} %||1-48-20||

\twolineshloka
{गौतमोऽपि महातेजा अहल्यासहितः सुखी}
{रामं सम्पूज्य विधिवत्तपस्तेपे महातपाः} %||1-48-21||

\twolineshloka
{रामोऽपि परमां पूजां गौतमस्य महामुनेः}
{सकाशाद्विधिवत्प्राप्य जगाम मिथिलां ततः} %||1-49-22||


{॥इत्यार्षे श्रीमद्रामायणे वाल्मीकीये आदिकाव्ये बालकाण्डे अष्टचत्वारिंशः सर्गः॥}

\sect{एकोनपञ्चाशत्तमः सर्गः}

\twolineshloka
{ततः प्रागुत्तरां गत्वा रामः सौमित्रिणा सह}
{विश्वामित्रं पुरस्कृत्य यज्ञवाटमुपागमत्} %||1-49-1||

\twolineshloka
{रामस्तु मुनिशार्दूलमुवाच सहलक्ष्मणः}
{साध्वी यज्ञसमृद्धिर्हि जनकस्य महात्मनः} %||1-49-2||

\twolineshloka
{बहूनीह सहस्राणि नानादेशनिवासिनाम्}
{ब्राह्मणानां महाभाग वेदाध्ययनशालिनाम्} %||1-49-3||

\twolineshloka
{ऋषिवाटाश्च दृश्यन्ते शकटीशतसङ्कुलाः}
{देशो विधीयतां ब्रह्मन्यत्र वत्स्यामहे वयम्} %||1-49-4||

\twolineshloka
{रामस्य वचनं श्रुत्वा विश्वामित्रो महामुनिः}
{निवेशमकरोद्देशे विविक्ते सलिलायुते} %||1-49-5||

\twolineshloka
{विश्वामित्रं मुनिश्रेष्ठं श्रुत्वा स नृपतिस्तदा}
{शतानन्दं पुरस्कृत्य पुरोहितमनिन्दितम्} %||1-49-6||

\twolineshloka
{ऋत्विजोऽपि महात्मानस्त्वर्घ्यमादाय सत्वरम्}
{विश्वामित्राय धर्मेण ददुर्मन्त्रपुरस्कृतम्} %||1-49-7||

\twolineshloka
{प्रतिगृह्य तु तां पूजां जनकस्य महात्मनः}
{पप्रच्छ कुशलं राज्ञो यज्ञस्य च निरामयम्} %||1-49-8||

\twolineshloka
{स तांश्चापि मुनीन्पृष्ट्वा सोपाध्याय पुरोधसः}
{यथान्यायं ततः सर्वैः समागच्छत्प्रहृष्टवान्} %||1-49-9||

\twolineshloka
{अथ राजा मुनिश्रेष्ठं कृताञ्जलिरभाषत}
{आसने भगवानास्तां सहैभिर्मुनिसत्तमैः} %||1-49-10||

\twolineshloka
{जनकस्य वचः श्रुत्वा निषसाद महामुनिः}
{पुरोधा ऋत्विजश्चैव राजा च सह मन्त्रिभिः} %||1-49-11||

\twolineshloka
{आसनेषु यथान्यायमुपविष्टान्समन्ततः}
{दृष्ट्वा स नृपतिस्तत्र विश्वामित्रमथाब्रवीत्} %||1-49-12||

\twolineshloka
{अद्य यज्ञसमृद्धिर्मे सफला दैवतैः कृता}
{अद्य यज्ञफलं प्राप्तं भगवद्दर्शनान्मया} %||1-49-13||

\twolineshloka
{धन्योऽस्म्यनुगृहीतोऽस्मि यस्य मे मुनिपुङ्गव}
{यज्ञोपसदनं ब्रह्मन्प्राप्तोऽसि मुनिभिः सह} %||1-49-14||

\twolineshloka
{द्वादशाहं तु ब्रह्मर्षे शेषमाहुर्मनीषिणः}
{ततो भागार्थिनो देवान्द्रष्टुमर्हसि कौशिक} %||1-49-15||

\twolineshloka
{इत्युक्त्वा मुनिशार्दूलं प्रहृष्टवदनस्तदा}
{पुनस्तं परिपप्रच्छ प्राञ्जलिः प्रयतो नृपः} %||1-49-16||

\twolineshloka
{इमौ कुमारौ भद्रं ते देवतुल्यपराक्रमौ}
{गजसिंहगती वीरौ शार्दूलवृषभोपमौ} %||1-49-17||

\twolineshloka
{पद्मपत्रविशालाक्षौ खड्गतूणीधनुर्धरौ}
{अश्विनाविव रूपेण समुपस्थितयौवनौ} %||1-49-18||

\twolineshloka
{यदृच्छयैव गां प्राप्तौ देवलोकादिवामरौ}
{कथं पद्भ्यामिह प्राप्तौ किमर्थं कस्य वा मुने} %||1-49-19||

\twolineshloka
{वरायुधधरौ वीरौ कस्य पुत्रौ महामुने}
{भूषयन्ताविमं देशं चन्द्रसूर्याविवाम्बरम्} %||1-49-20||

\twolineshloka
{परस्परस्य सदृशौ प्रमाणेङ्गितचेष्टितैः}
{काकपक्षधरौ वीरौ श्रोतुमिच्छामि तत्त्वतः} %||1-49-21||

\twolineshloka
{तस्य तद्वचनं श्रुत्वा जनकस्य महात्मनः}
{न्यवेदयन्महात्मानौ पुत्रौ दशरथस्य तौ} %||1-49-22||

\twolineshloka
{सिद्धाश्रमनिवासं च राक्षसानां वधं तथा}
{तच्चागमनमव्यग्रं विशालायाश्च दर्शनम्} %||1-49-23||

\twolineshloka
{अहल्यादर्शनं चैव गौतमेन समागमम्}
{महाधनुषि जिज्ञासां कर्तुमागमनं तथा} %||1-49-24||

\twolineshloka
{एतत्सर्वं महातेजा जनकाय महात्मने}
{निवेद्य विररामाथ विश्वामित्रो महामुनिः} %||1-50-25||


{॥इत्यार्षे श्रीमद्रामायणे वाल्मीकीये आदिकाव्ये बालकाण्डे एकोनपञ्चाशत्तमः सर्गः॥}

\sect{पञ्चाशत्तमः सर्गः}

\twolineshloka
{तस्य तद्वचनं श्रुत्वा विश्वामित्रस्य धीमतः}
{हृष्टरोमा महातेजाः शतानन्दो महातपाः} %||1-50-1||

\twolineshloka
{गौतमस्य सुतो ज्येष्ठस्तपसा द्योतितप्रभः}
{रामसन्दर्शनादेव परं विस्मयमागतः} %||1-50-2||

\twolineshloka
{स तौ निषण्णौ सम्प्रेक्ष्य सुखासीनौ नृपात्मजौ}
{शतानन्दो मुनिश्रेष्ठं विश्वामित्रमथाब्रवीत्} %||1-50-3||

\twolineshloka
{अपि ते मुनिशार्दूल मम माता यशस्विनी}
{दर्शिता राजपुत्राय तपो दीर्घमुपागता} %||1-50-4||

\twolineshloka
{अपि रामे महातेजो मम माता यशस्विनी}
{वन्यैरुपाहरत्पूजां पूजार्हे सर्वदेहिनाम्} %||1-50-5||

\twolineshloka
{अपि रामाय कथितं यथावृत्तं पुरातनम्}
{मम मातुर्महातेजो देवेन दुरनुष्ठितम्} %||1-50-6||

\twolineshloka
{अपि कौशिक भद्रं ते गुरुणा मम सङ्गता}
{माता मम मुनिश्रेष्ठ रामसन्दर्शनादितः} %||1-50-7||

\twolineshloka
{अपि मे गुरुणा रामः पूजितः कुशिकात्मज}
{इहागतो महातेजाः पूजां प्राप्य महात्मनः} %||1-50-8||

\twolineshloka
{अपि शान्तेन मनसा गुरुर्मे कुशिकात्मज}
{इहागतेन रामेण प्रयतेनाभिवादितः} %||1-50-9||

\twolineshloka
{तच्छ्रुत्वा वचनं तस्य विश्वामित्रो महामुनिः}
{प्रत्युवाच शतानन्दं वाक्यज्ञो वाक्यकोविदम्} %||1-50-10||

\twolineshloka
{नातिक्रान्तं मुनिश्रेष्ठ यत्कर्तव्यं कृतं मया}
{सङ्गता मुनिना पत्नी भार्गवेणेव रेणुका} %||1-50-11||

\twolineshloka
{तच्छ्रुत्वा वचनं तस्य विश्वामित्रस्य धीमतः}
{शतानन्दो महातेजा रामं वचनमब्रवीत्} %||1-50-12||

\twolineshloka
{स्वागतं ते नरश्रेष्ठ दिष्ट्या प्राप्तोऽसि राघव}
{विश्वामित्रं पुरस्कृत्य महर्षिमपराजितम्} %||1-50-13||

\twolineshloka
{अचिन्त्यकर्मा तपसा ब्रह्मर्षिरमितप्रभः}
{विश्वामित्रो महातेजा वेत्स्येनं परमां गतिम्} %||1-50-14||

\twolineshloka
{नास्ति धन्यतरो राम त्वत्तोऽन्यो भुवि कश्चन}
{गोप्ता कुशिकपुत्रस्ते येन तप्तं महत्तपः} %||1-50-15||

\twolineshloka
{श्रूयतां चाभिदास्यामि कौशिकस्य महात्मनः}
{यथाबलं यथावृत्तं तन्मे निगदतः शृणु} %||1-50-16||

\twolineshloka
{राजाभूदेष धर्मात्मा दीर्घ कालमरिन्दमः}
{धर्मज्ञः कृतविद्यश्च प्रजानां च हिते रतः} %||1-50-17||

\twolineshloka
{प्रजापतिसुतस्त्वासीत्कुशो नाम महीपतिः}
{कुशस्य पुत्रो बलवान्कुशनाभः सुधार्मिकः} %||1-50-18||

\twolineshloka
{कुशनाभसुतस्त्वासीद्गाधिरित्येव विश्रुतः}
{गाधेः पुत्रो महातेजा विश्वामित्रो महामुनिः} %||1-50-19||

\twolineshloka
{विश्वमित्रो महातेजाः पालयामास मेदिनीम्}
{बहुवर्षसहस्राणि राजा राज्यमकारयत्} %||1-50-20||

\twolineshloka
{कदाचित्तु महातेजा योजयित्वा वरूथिनीम्}
{अक्षौहिणीपरिवृतः परिचक्राम मेदिनीम्} %||1-50-21||

\twolineshloka
{नगराणि च राष्ट्राणि सरितश्च तथा गिरीन्}
{आश्रमान्क्रमशो राजा विचरन्नाजगामह} %||1-50-22||

\twolineshloka
{वसिष्ठस्याश्रमपदं नानापुष्पफलद्रुमम्}
{नानामृगगणाकीर्णं सिद्धचारणसेवितम्} %||1-50-23||

\twolineshloka
{देवदानवगन्धर्वैः किंनरैरुपशोभितम्}
{प्रशान्तहरिणाकीर्णं द्विजसङ्घनिषेवितम्} %||1-50-24||

\twolineshloka
{ब्रह्मर्षिगणसङ्कीर्णं देवर्षिगणसेवितम्}
{तपश्चरणसंसिद्धैरग्निकल्पैर्महात्मभिः} %||1-50-25||

\twolineshloka
{सततं सङ्कुलं श्रीमद्ब्रह्मकल्पैर्महात्मभिः}
{अब्भक्षैर्वायुभक्षैश्च शीर्णपर्णाशनैस्तथा} %||1-50-26||

\twolineshloka
{फलमूलाशनैर्दान्तैर्जितरोषैर्जितेन्द्रियैः}
{ऋषिभिर्वालखिल्यैश्च जपहोमपरायणैः} %||1-50-27||

\twolineshloka
{वसिष्ठस्याश्रमपदं ब्रह्मलोकमिवापरम्}
{ददर्श जयतां श्रेष्ठ विश्वामित्रो महाबलः} %||1-51-28||


{॥इत्यार्षे श्रीमद्रामायणे वाल्मीकीये आदिकाव्ये बालकाण्डे पञ्चाशत्तमः सर्गः॥}

\sect{एकपञ्चाशत्तमः सर्गः}

\twolineshloka
{स दृष्ट्वा परमप्रीतो विश्वामित्रो महाबलः}
{प्रणतो विनयाद्वीरो वसिष्ठं जपतां वरम्} %||1-51-1||

\twolineshloka
{स्वागतं तव चेत्युक्तो वसिष्ठेन महात्मना}
{आसनं चास्य भगवान्वसिष्ठो व्यादिदेश ह} %||1-51-2||

\twolineshloka
{उपविष्टाय च तदा विश्वामित्राय धीमते}
{यथान्यायं मुनिवरः फलमूलमुपाहरत्} %||1-51-3||

\twolineshloka
{प्रतिगृह्य च तां पूजां वसिष्ठाद्राजसत्तमः}
{तपोऽग्निहोत्रशिष्येषु कुशलं पर्यपृच्छत} %||1-51-4||

\twolineshloka
{विश्वामित्रो महातेजा वनस्पतिगणे तथा}
{सर्वत्र कुशलं चाह वसिष्ठो राजसत्तमम्} %||1-51-5||

\twolineshloka
{सुखोपविष्टं राजानं विश्वामित्रं महातपाः}
{पप्रच्छ जपतां श्रेष्ठो वसिष्ठो ब्रह्मणः सुतः} %||1-51-6||

\twolineshloka
{कच्चित्ते कुशलं राजन्कच्चिद्धर्मेण रञ्जयन्}
{प्रजाः पालयसे राजन्राजवृत्तेन धार्मिक} %||1-51-7||

\twolineshloka
{कच्चित्ते सुभृता भृत्याः कच्चित्तिष्ठन्ति शासने}
{कच्चित्ते विजिताः सर्वे रिपवो रिपुसूदन} %||1-51-8||

\twolineshloka
{कच्चिद्बले च कोशे च मित्रेषु च परन्तप}
{कुशलं ते नरव्याघ्र पुत्रपौत्रे तथानघ} %||1-51-9||

\twolineshloka
{सर्वत्र कुशलं राजा वसिष्ठं प्रत्युदाहरत्}
{विश्वामित्रो महातेजा वसिष्ठं विनयान्वितः} %||1-51-10||

\twolineshloka
{कृत्वोभौ सुचिरं कालं धर्मिष्ठौ ताः कथाः शुभाः}
{मुदा परमया युक्तौ प्रीयेतां तौ परस्परम्} %||1-51-11||

\twolineshloka
{ततो वसिष्ठो भगवान्कथान्ते रघुनन्दन}
{विश्वामित्रमिदं वाक्यमुवाच प्रहसन्निव} %||1-51-12||

\twolineshloka
{आतिथ्यं कर्तुमिच्छामि बलस्यास्य महाबल}
{तव चैवाप्रमेयस्य यथार्हं सम्प्रतीच्छ मे} %||1-51-13||

\twolineshloka
{सत्क्रियां तु भवानेतां प्रतीच्छतु मयोद्यताम्}
{राजंस्त्वमतिथिश्रेष्ठः पूजनीयः प्रयत्नतः} %||1-51-14||

\twolineshloka
{एवमुक्तो वसिष्ठेन विश्वामित्रो महामतिः}
{कृतमित्यब्रवीद्राजा पूजावाक्येन मे त्वया} %||1-51-15||

\twolineshloka
{फलमूलेन भगवन्विद्यते यत्तवाश्रमे}
{पाद्येनाचमनीयेन भगवद्दर्शनेन च} %||1-51-16||

\twolineshloka
{सर्वथा च महाप्राज्ञ पूजार्हेण सुपूजितः}
{गमिष्यामि नमस्तेऽस्तु मैत्रेणेक्षस्व चक्षुषा} %||1-51-17||

\twolineshloka
{एवं ब्रुवन्तं राजानं वसिष्ठः पुनरेव हि}
{न्यमन्त्रयत धर्मात्मा पुनः पुनरुदारधीः} %||1-51-18||

\twolineshloka
{बाढमित्येव गाधेयो वसिष्ठं प्रत्युवाच ह}
{यथा प्रियं भगवतस्तथास्तु मुनिसत्तम} %||1-51-19||

\twolineshloka
{एवमुक्तो महातेजा वसिष्ठो जपतां वरः}
{आजुहाव ततः प्रीतः कल्माषीं धूतकल्मषः} %||1-51-20||

\threelineshloka
{एह्येहि शबले क्षिप्रं शृणु चापि वचो मम}
{सबलस्यास्य राजर्षेः कर्तुं व्यवसितोऽस्म्यहम्}
{भोजनेन महार्हेण सत्कारं संविधत्स्व मे} %||1-51-21||

\twolineshloka
{यस्य यस्य यथाकामं षड्रसेष्वभिपूजितम्}
{तत्सर्वं कामधुग्दिव्ये अभिवर्षकृते मम} %||1-51-22||

\twolineshloka
{रसेनान्नेन पानेन लेह्यचोष्येण संयुतम्}
{अन्नानां निचयं सर्वं सृजस्व शबले त्वर} %||1-52-23||


{॥इत्यार्षे श्रीमद्रामायणे वाल्मीकीये आदिकाव्ये बालकाण्डे एकपञ्चाशत्तमः सर्गः॥}

\sect{द्विपञ्चाशत्तमः सर्गः}

\twolineshloka
{एवमुक्ता वसिष्ठेन शबला शत्रुसूदन}
{विदधे कामधुक्कामान्यस्य यस्य यथेप्सितम्} %||1-52-1||

\twolineshloka
{इक्षून्मधूंस्तथा लाजान्मैरेयांश्च वरासवान्}
{पानानि च महार्हाणि भक्ष्यांश्चोच्चावचांस्तथा} %||1-52-2||

\twolineshloka
{उष्णाढ्यस्यौदनस्यापि राशयः पर्वतोपमाः}
{मृष्टान्नानि च सूपाश्च दधिकुल्यास्तथैव च} %||1-52-3||

\twolineshloka
{नानास्वादुरसानां च षाडवानां तथैव च}
{भाजनानि सुपूर्णानि गौडानि च सहस्रशः} %||1-52-4||

\twolineshloka
{सर्वमासीत्सुसन्तुष्टं हृष्टपुष्टजनाकुलम्}
{विश्वामित्रबलं राम वसिष्ठेनाभितर्पितम्} %||1-52-5||

\twolineshloka
{विश्वामित्रोऽपि राजर्षिर्हृष्टपुष्टस्तदाभवत्}
{सान्तः पुरवरो राजा सब्राह्मणपुरोहितः} %||1-52-6||

\twolineshloka
{सामात्यो मन्त्रिसहितः सभृत्यः पूजितस्तदा}
{युक्तः परेण हर्षेण वसिष्ठमिदमब्रवीत्} %||1-52-7||

\twolineshloka
{पूजितोऽहं त्वया ब्रह्मन्पूजार्हेण सुसत्कृतः}
{श्रूयतामभिधास्यामि वाक्यं वाक्यविशारद} %||1-52-8||

\threelineshloka
{गवां शतसहस्रेण दीयतां शबला मम}
{रत्नं हि भगवन्नेतद्रत्नहारी च पार्थिवः}
{तस्मान्मे शबलां देहि ममैषा धर्मतो द्विज} %||1-52-9||

\twolineshloka
{एवमुक्तस्तु भगवान्वसिष्ठो मुनिसत्तमः}
{विश्वामित्रेण धर्मात्मा प्रत्युवाच महीपतिम्} %||1-52-10||

\twolineshloka
{नाहं शतसहस्रेण नापि कोटिशतैर्गवाम्}
{राजन्दास्यामि शबलां राशिभी रजतस्य वा} %||1-52-11||

\twolineshloka
{न परित्यागमर्हेयं मत्सकाशादरिन्दम}
{शाश्वती शबला मह्यं कीर्तिरात्मवतो यथा} %||1-52-12||

\twolineshloka
{अस्यां हव्यं च कव्यं च प्राणयात्रा तथैव च}
{आयत्तमग्निहोत्रं च बलिर्होमस्तथैव च} %||1-52-13||

\twolineshloka
{स्वाहाकारवषट्कारौ विद्याश्च विविधास्तथा}
{आयत्तमत्र राजर्षे सर्वमेतन्न संशयः} %||1-52-14||

\twolineshloka
{सर्व स्वमेतत्सत्येन मम तुष्टिकरी सदा}
{कारणैर्बहुभी राजन्न दास्ये शबलां तव} %||1-52-15||

\twolineshloka
{वसिष्ठेनैवमुक्तस्तु विश्वामित्रोऽब्रवीत्ततः}
{संरब्धतरमत्यर्थं वाक्यं वाक्यविशारदः} %||1-52-16||

\twolineshloka
{हैरण्यकक्ष्याग्रैवेयान्सुवर्णाङ्कुशभूषितान्}
{ददामि कुञ्जराणां ते सहस्राणि चतुर्दश} %||1-52-17||

\twolineshloka
{हैरण्यानां रथानां च श्वेताश्वानां चतुर्युजाम्}
{ददामि ते शतान्यष्टौ किङ्किणीकविभूषितान्} %||1-52-18||

\twolineshloka
{हयानां देशजातानां कुलजानां महौजसाम्}
{सहस्रमेकं दश च ददामि तव सुव्रत} %||1-52-19||

\twolineshloka
{नानावर्णविभक्तानां वयःस्थानां तथैव च}
{ददाम्येकां गवां कोटिं शबला दीयतां मम} %||1-52-20||

\twolineshloka
{एवमुक्तस्तु भगवान्विश्वामित्रेण धीमता}
{न दास्यामीति शबलां प्राह राजन्कथञ्चन} %||1-52-21||

\twolineshloka
{एतदेव हि मे रत्नमेतदेव हि मे धनम्}
{एतदेव हि सर्वस्वमेतदेव हि जीवितम्} %||1-52-22||

\twolineshloka
{दर्शश्च पूर्णमासश्च यज्ञाश्चैवाप्तदक्षिणाः}
{एतदेव हि मे राजन्विविधाश्च क्रियास्तथा} %||1-52-23||

\twolineshloka
{अदोमूलाः क्रियाः सर्वा मम राजन्न संशयः}
{बहूनां किं प्रलापेन न दास्ये कामदोहिनीम्} %||1-53-24||


{॥इत्यार्षे श्रीमद्रामायणे वाल्मीकीये आदिकाव्ये बालकाण्डे द्विपञ्चाशत्तमः सर्गः॥}

\sect{त्रिपञ्चाशत्तमः सर्गः}

\twolineshloka
{कामधेनुं वसिष्ठोऽपि यदा न त्यजते मुनिः}
{तदास्य शबलां राम विश्वामित्रोऽन्वकर्षत} %||1-53-1||

\twolineshloka
{नीयमाना तु शबला राम राज्ञा महात्मना}
{दुःखिता चिन्तयामास रुदन्ती शोककर्शिता} %||1-53-2||

\twolineshloka
{परित्यक्ता वसिष्ठेन किमहं सुमहात्मना}
{याहं राजभृतैर्दीना ह्रियेयं भृशदुःखिता} %||1-53-3||

\twolineshloka
{किं मयापकृतं तस्य महर्षेर्भावितात्मनः}
{यन्मामनागसं भक्तामिष्टां त्यजति धार्मिकः} %||1-53-4||

\twolineshloka
{इति सा चिन्तयित्वा तु निःश्वस्य च पुनः पुनः}
{जगाम वेगेन तदा वसिष्ठं परमौजसम्} %||1-53-5||

\twolineshloka
{निर्धूय तांस्तदा भृत्याञ्शतशः शत्रुसूदन}
{जगामानिलवेगेन पादमूलं महात्मनः} %||1-53-6||

\twolineshloka
{शबला सा रुदन्ती च क्रोशन्ती चेदमब्रवीत्}
{वसिष्ठस्याग्रतः स्थित्वा मेघदुन्दुभिराविणी} %||1-53-7||

\twolineshloka
{भगवन्किं परित्यक्ता त्वयाहं ब्रह्मणः सुत}
{यस्माद्राजभृता मां हि नयन्ते त्वत्सकाशतः} %||1-53-8||

\twolineshloka
{एवमुक्तस्तु ब्रह्मर्षिरिदं वचनमब्रवीत्}
{शोकसन्तप्तहृदयां स्वसारमिव दुःखिताम्} %||1-53-9||

\twolineshloka
{न त्वां त्यजामि शबले नापि मेऽपकृतं त्वया}
{एष त्वां नयते राजा बलान्मत्तो महाबलः} %||1-53-10||

\twolineshloka
{न हि तुल्यं बलं मह्यं राजा त्वद्य विशेषतः}
{बली राजा क्षत्रियश्च पृथिव्याः पतिरेव च} %||1-53-11||

\twolineshloka
{इयमक्षौहिणीपूर्णा सवाजिरथसङ्कुला}
{हस्तिध्वजसमाकीर्णा तेनासौ बलवत्तरः} %||1-53-12||

\twolineshloka
{एवमुक्ता वसिष्ठेन प्रत्युवाच विनीतवत्}
{वचनं वचनज्ञा सा ब्रह्मर्षिममितप्रभम्} %||1-53-13||

\twolineshloka
{न बलं क्षत्रियस्याहुर्ब्राह्मणो बलवत्तरः}
{ब्रह्मन्ब्रह्मबलं दिव्यं क्षत्रात्तु बलवत्तरम्} %||1-53-14||

\twolineshloka
{अप्रमेयबलं तुभ्यं न त्वया बलवत्तरः}
{विश्वामित्रो महावीर्यस्तेजस्तव दुरासदम्} %||1-53-15||

\twolineshloka
{नियुङ्क्ष्व मां महातेजस्त्वद्ब्रह्मबलसम्भृताम्}
{तस्य दर्पं बलं यत्तन्नाशयामि दुरात्मनः} %||1-53-16||

\twolineshloka
{इत्युक्तस्तु तया राम वसिष्ठः सुमहायशाः}
{सृजस्वेति तदोवाच बलं परबलारुजम्} %||1-53-17||

\twolineshloka
{तस्या हुम्भारवोत्सृष्टाः पह्लवाः शतशो नृप}
{नाशयन्ति बलं सर्वं विश्वामित्रस्य पश्यतः} %||1-53-18||

\twolineshloka
{स राजा परमक्रुद्धः क्रोधविस्फारितेक्षणः}
{पह्लवान्नाशयामास शस्त्रैरुच्चावचैरपि} %||1-53-19||

\twolineshloka
{विश्वामित्रार्दितान्दृष्ट्वा पह्लवाञ्शतशस्तदा}
{भूय एवासृजद्घोराञ्शकान्यवनमिश्रितान्} %||1-53-20||

\twolineshloka
{तैरासीत्संवृता भूमिः शकैर्यवनमिश्रितैः}
{प्रभावद्भिर्महावीर्यैर्हेमकिञ्जल्कसंनिभैः} %||1-53-21||

\twolineshloka
{दीर्घासिपट्टिशधरैर्हेमवर्णाम्बरावृतैः}
{निर्दग्धं तद्बलं सर्वं प्रदीप्तैरिव पावकैः} %||1-53-22||

{ततोऽस्त्राणि महातेजा विश्वामित्रो मुमोच ह॥\devanumber{23}॥} %||1-54-23|| (Check)

\addtocounter{shlokacount}{1}


{॥इत्यार्षे श्रीमद्रामायणे वाल्मीकीये आदिकाव्ये बालकाण्डे त्रिपञ्चाशत्तमः सर्गः॥}

\sect{चतुःपञ्चाशत्तमः सर्गः}

\twolineshloka
{ततस्तानाकुलान्दृष्ट्वा विश्वामित्रास्त्रमोहितान्}
{वसिष्ठश्चोदयामास कामधुक्सृज योगतः} %||1-54-1||

\twolineshloka
{तस्या हुम्भारवाज्जाताः काम्बोजा रविसंनिभाः}
{ऊधसस्त्वथ संजाताः पह्लवाः शस्त्रपाणयः} %||1-54-2||

\twolineshloka
{योनिदेशाच्च यवनः शकृद्देशाच्छकास्तथा}
{रोमकूपेषु मेच्छाश्च हरीताः सकिरातकाः} %||1-54-3||

\twolineshloka
{तैस्तन्निषूदितं सैन्यं विश्वमित्रस्य तत्क्षणात्}
{सपदातिगजं साश्वं सरथं रघुनन्दन} %||1-54-4||

\twolineshloka
{दृष्ट्वा निषूदितं सैन्यं वसिष्ठेन महात्मना}
{विश्वामित्रसुतानां तु शतं नानाविधायुधम्} %||1-54-5||

\twolineshloka
{अभ्यधावत्सुसङ्क्रुद्धं वसिष्ठं जपतां वरम्}
{हुङ्कारेणैव तान्सर्वान्निर्ददाह महानृषिः} %||1-54-6||

\twolineshloka
{ते साश्वरथपादाता वसिष्ठेन महात्मना}
{भस्मीकृता मुहूर्तेन विश्वामित्रसुतास्तदा} %||1-54-7||

\twolineshloka
{दृष्ट्वा विनाशितान्पुत्रान्बलं च सुमहायशाः}
{सव्रीडश्चिन्तयाविष्टो विश्वामित्रोऽभवत्तदा} %||1-54-8||

\twolineshloka
{सन्दुर इव निर्वेगो भग्नदंष्ट्र इवोरगः}
{उपरक्त इवादित्यः सद्यो निष्प्रभतां गतः} %||1-54-9||

\twolineshloka
{हतपुत्रबलो दीनो लूनपक्ष इव द्विजः}
{हतदर्पो हतोत्साहो निर्वेदं समपद्यत} %||1-54-10||

\twolineshloka
{स पुत्रमेकं राज्याय पालयेति नियुज्य च}
{पृथिवीं क्षत्रधर्मेण वनमेवान्वपद्यत} %||1-54-11||

\twolineshloka
{स गत्वा हिमवत्पार्श्वं किंनरोरगसेवितम्}
{महादेवप्रसादार्थं तपस्तेपे महातपाः} %||1-54-12||

\twolineshloka
{केनचित्त्वथ कालेन देवेशो वृषभध्वजः}
{दर्शयामास वरदो विश्वामित्रं महामुनिम्} %||1-54-13||

\twolineshloka
{किमर्थं तप्यसे राजन्ब्रूहि यत्ते विवक्षितम्}
{वरदोऽस्मि वरो यस्ते काङ्क्षितः सोऽभिधीयताम्} %||1-54-14||

\twolineshloka
{एवमुक्तस्तु देवेन विश्वामित्रो महातपाः}
{प्रणिपत्य महादेवमिदं वचनमब्रवीत्} %||1-54-15||

\twolineshloka
{यदि तुष्टो महादेव धनुर्वेदो ममानघ}
{साङ्गोपाङ्गोपनिषदः सरहस्यः प्रदीयताम्} %||1-54-16||

\twolineshloka
{यानि देवेषु चास्त्राणि दानवेषु महर्षिषु}
{गन्धर्वयक्षरक्षःसु प्रतिभान्तु ममानघ} %||1-54-17||

\twolineshloka
{तव प्रसादाद्भवतु देवदेव ममेप्सितम्}
{एवमस्त्विति देवेशो वाक्यमुक्त्वा दिवं गतः} %||1-54-18||

\twolineshloka
{प्राप्य चास्त्राणि राजर्षिर्विश्वामित्रो महाबलः}
{दर्पेण महता युक्तो दर्पपूर्णोऽभवत्तदा} %||1-54-19||

\twolineshloka
{विवर्धमानो वीर्येण समुद्र इव पर्वणि}
{हतमेव तदा मेने वसिष्ठमृषिसत्तमम्} %||1-54-20||

\twolineshloka
{ततो गत्वाश्रमपदं मुमोचास्त्राणि पार्थिवः}
{यैस्तत्तपोवनं सर्वं निर्दग्धं चास्त्रतेजसा} %||1-54-21||

\twolineshloka
{उदीर्यमाणमस्त्रं तद्विश्वामित्रस्य धीमतः}
{दृष्ट्वा विप्रद्रुता भीता मुनयः शतशो दिशः} %||1-54-22||

\twolineshloka
{वसिष्ठस्य च ये शिष्यास्तथैव मृगपक्षिणः}
{विद्रवन्ति भयाद्भीता नानादिग्भ्यः सहस्रशः} %||1-54-23||

\twolineshloka
{वसिष्ठस्याश्रमपदं शून्यमासीन्महात्मनः}
{मुहूर्तमिव निःशब्दमासीदीरिणसंनिभम्} %||1-54-24||

\twolineshloka
{वदतो वै वसिष्ठस्य मा भैष्टेति मुहुर्मुहुः}
{नाशयाम्यद्य गाधेयं नीहारमिव भास्करः} %||1-54-25||

\twolineshloka
{एवमुक्त्वा महातेजा वसिष्ठो जपतां वरः}
{विश्वामित्रं तदा वाक्यं सरोषमिदमब्रवीत्} %||1-54-26||

\twolineshloka
{आश्रमं चिरसंवृद्धं यद्विनाशितवानसि}
{दुराचारोऽसि यन्मूढ तस्मात्त्वं न भविष्यसि} %||1-54-27||

\twolineshloka
{इत्युक्त्वा परमक्रुद्धो दण्डमुद्यम्य सत्वरः}
{विधूम इव कालाग्निर्यमदण्डमिवापरम्} %||1-55-28||


{॥इत्यार्षे श्रीमद्रामायणे वाल्मीकीये आदिकाव्ये बालकाण्डे चतुःपञ्चाशत्तमः सर्गः॥}

\sect{पञ्चपञ्चाशत्तमः सर्गः}

\twolineshloka
{एवमुक्तो वसिष्ठेन विश्वामित्रो महाबलः}
{आग्नेयमस्त्रमुत्क्षिप्य तिष्ठ तिष्ठेति चाब्रवीत्} %||1-55-1||

{वसिष्ठो भगवान्क्रोधादिदं वचनमब्रवीत्॥\devanumber{2}॥} %||1-55-2|| (Check)

\addtocounter{shlokacount}{1}

\twolineshloka
{क्षत्रबन्धो स्थितोऽस्म्येष यद्बलं तद्विदर्शय}
{नाशयाम्येष ते दर्पं शस्त्रस्य तव गाधिज} %||1-55-3||

\twolineshloka
{क्व च ते क्षत्रियबलं क्व च ब्रह्मबलं महत्}
{पश्य ब्रह्मबलं दिव्यं मम क्षत्रियपांसन} %||1-55-4||

\twolineshloka
{तस्यास्त्रं गाधिपुत्रस्य घोरमाग्नेयमुत्तमम्}
{ब्रह्मदण्डेन तच्छान्तमग्नेर्वेग इवाम्भसा} %||1-55-5||

\twolineshloka
{वारुणं चैव रौद्रं च ऐन्द्रं पाशुपतं तथा}
{ऐषीकं चापि चिक्षेप रुषितो गाधिनन्दनः} %||1-55-6||

\twolineshloka
{मानवं मोहनं चैव गान्धर्वं स्वापनं तथा}
{जृम्भणं मोहनं चैव सन्तापनविलापने} %||1-55-7||

\twolineshloka
{शोषणं दारणं चैव वज्रमस्त्रं सुदुर्जयम्}
{ब्रह्मपाशं कालपाशं वारुणं पाशमेव च} %||1-55-8||

\twolineshloka
{पिनाकास्त्रं च दयितं शुष्कार्द्रे अशनी तथा}
{दण्डास्त्रमथ पैशाचं क्रौञ्चमस्त्रं तथैव च} %||1-55-9||

\twolineshloka
{धर्मचक्रं कालचक्रं विष्णुचक्रं तथैव च}
{वायव्यं मथनं चैव अस्त्रं हयशिरस्तथा} %||1-55-10||

\twolineshloka
{शक्तिद्वयं च चिक्षेप कङ्कालं मुसलं तथा}
{वैद्याधरं महास्त्रं च कालास्त्रमथ दारुणम्} %||1-55-11||

\twolineshloka
{त्रिशूलमस्त्रं घोरं च कापालमथ कङ्कणम्}
{एतान्यस्त्राणि चिक्षेप सर्वाणि रघुनन्दन} %||1-55-12||

\twolineshloka
{वसिष्ठे जपतां श्रेष्ठे तदद्भुतमिवाभवत्}
{तानि सर्वाणि दण्डेन ग्रसते ब्रह्मणः सुतः} %||1-55-13||

\twolineshloka
{तेषु शान्तेषु ब्रह्मास्त्रं क्षिप्तवान्गाधिनन्दनः}
{तदस्त्रमुद्यतं दृष्ट्वा देवाः साग्निपुरोगमाः} %||1-55-14||

\twolineshloka
{देवर्षयश्च सम्भ्रान्ता गन्धर्वाः समहोरगाः}
{त्रैलोक्यमासीत्सन्त्रस्तं ब्रह्मास्त्रे समुदीरिते} %||1-55-15||

\twolineshloka
{तदप्यस्त्रं महाघोरं ब्राह्मं ब्राह्मेण तेजसा}
{वसिष्ठो ग्रसते सर्वं ब्रह्मदण्डेन राघव} %||1-55-16||

\twolineshloka
{ब्रह्मास्त्रं ग्रसमानस्य वसिष्ठस्य महात्मनः}
{त्रैलोक्यमोहनं रौद्रं रूपमासीत्सुदारुणम्} %||1-55-17||

\twolineshloka
{रोमकूपेषु सर्वेषु वसिष्ठस्य महात्मनः}
{मरीच्य इव निष्पेतुरग्नेर्धूमाकुलार्चिषः} %||1-55-18||

\twolineshloka
{प्राज्वलद्ब्रह्मदण्डश्च वसिष्ठस्य करोद्यतः}
{विधूम इव कालाग्निर्यमदण्ड इवापरः} %||1-55-19||

\twolineshloka
{ततोऽस्तुवन्मुनिगणा वसिष्ठं जपतां वरम्}
{अमोघं ते बलं ब्रह्मंस्तेजो धारय तेजसा} %||1-55-20||

\twolineshloka
{निगृहीतस्त्वया ब्रह्मन्विश्वामित्रो महातपाः}
{प्रसीद जपतां श्रेष्ठ लोकाः सन्तु गतव्यथाः} %||1-55-21||

\twolineshloka
{एवमुक्तो महातेजाः शमं चक्रे महातपाः}
{विश्वामित्रोऽपि निकृतो विनिःश्वस्येदमब्रवीत्} %||1-55-22||

\twolineshloka
{धिग्बलं क्षत्रियबलं ब्रह्मतेजोबलं बलम्}
{एकेन ब्रह्मदण्डेन सर्वास्त्राणि हतानि मे} %||1-55-23||

\twolineshloka
{तदेतत्समवेक्ष्याहं प्रसन्नेन्द्रियमानसः}
{तपो महत्समास्थास्ये यद्वै ब्रह्मत्वकारकम्} %||1-56-24||


{॥इत्यार्षे श्रीमद्रामायणे वाल्मीकीये आदिकाव्ये बालकाण्डे पञ्चपञ्चाशत्तमः सर्गः॥}

\sect{षट्पञ्चाशत्तमः सर्गः}

\twolineshloka
{ततः सन्तप्तहृदयः स्मरन्निग्रहमात्मनः}
{विनिःश्वस्य विनिःश्वस्य कृतवैरो महात्मना} %||1-56-1||

\threelineshloka
{स दक्षिणां दिशं गत्वा महिष्या सह राघव}
{तताप परमं घोरं विश्वामित्रो महातपाः}
{फलमूलाशनो दान्तश्चचार परमं तपः} %||1-56-2||

\twolineshloka
{अथास्य जज्ञिरे पुत्राः सत्यधर्मपरायणाः}
{हविष्पन्दो मधुष्पन्दो दृढनेत्रो महारथः} %||1-56-3||

\twolineshloka
{पूर्णे वर्षसहस्रे तु ब्रह्मा लोकपितामहः}
{अब्रवीन्मधुरं वाक्यं विश्वामित्रं तपोधनम्} %||1-56-4||

\twolineshloka
{जिता राजर्षिलोकास्ते तपसा कुशिकात्मज}
{अनेन तपसा त्वां हि राजर्षिरिति विद्महे} %||1-56-5||

\twolineshloka
{एवमुक्त्वा महातेजा जगाम सह दैवतैः}
{त्रिविष्टपं ब्रह्मलोकं लोकानां परमेश्वरः} %||1-56-6||

\twolineshloka
{विश्वामित्रोऽपि तच्छ्रुत्वा ह्रिया किञ्चिदवाङ्मुखः}
{दुःखेन महताविष्टः समन्युरिदमब्रवीत्} %||1-56-7||

\twolineshloka
{तपश्च सुमहत्तप्तं राजर्षिरिति मां विदुः}
{देवाः सर्षिगणाः सर्वे नास्ति मन्ये तपःफलम्} %||1-56-8||

\twolineshloka
{एवं निश्चित्य मनसा भूय एव महातपाः}
{तपश्चचार काकुत्स्थ परमं परमात्मवान्} %||1-56-9||

\twolineshloka
{एतस्मिन्नेव काले तु सत्यवादी जितेन्द्रियः}
{त्रिशङ्कुरिति विख्यात इक्ष्वाकु कुलनन्दनः} %||1-56-10||

\twolineshloka
{तस्य बुद्धिः समुत्पन्ना यजेयमिति राघव}
{गच्छेयं स्वशरीरेण देवानां परमां गतिम्} %||1-56-11||

\twolineshloka
{स वसिष्ठं समाहूय कथयामास चिन्तितम्}
{अशक्यमिति चाप्युक्तो वसिष्ठेन महात्मना} %||1-56-12||

\twolineshloka
{प्रत्याख्यातो वसिष्ठेन स ययौ दक्षिणां दिशम्}
{वसिष्ठा दीर्घ तपसस्तपो यत्र हि तेपिरे} %||1-56-13||

\twolineshloka
{त्रिशङ्कुः सुमहातेजाः शतं परमभास्वरम्}
{वसिष्ठपुत्रान्ददृशे तप्यमानान्यशस्विनः} %||1-56-14||

\threelineshloka
{सोऽभिगम्य महात्मानः सर्वानेव गुरोः सुतान्}
{अभिवाद्यानुपूर्व्येण ह्रिया किञ्चिदवाङ्मुखः}
{अब्रवीत्सुमहातेजाः सर्वानेव कृताञ्जलिः} %||1-56-15||

\twolineshloka
{शरणं वः प्रपद्येऽहं शरण्याञ्शरणागतः}
{प्रत्याख्यातोऽस्मि भद्रं वो वसिष्ठेन महात्मना} %||1-56-16||

\twolineshloka
{यष्टुकामो महायज्ञं तदनुज्ञातुमर्थथ}
{गुरुपुत्रानहं सर्वान्नमस्कृत्य प्रसादये} %||1-56-17||

\threelineshloka
{शिरसा प्रणतो याचे ब्राह्मणांस्तपसि स्थितान्}
{ते मां भवन्तः सिद्ध्यर्थं याजयन्तु समाहिताः}
{सशरीरो यथाहं हि देवलोकमवाप्नुयाम्} %||1-56-18||

\twolineshloka
{प्रत्याख्यातो वसिष्ठेन गतिमन्यां तपोधनाः}
{गुरुपुत्रानृते सर्वान्नाहं पश्यामि काञ्चन} %||1-56-19||

\twolineshloka
{इक्ष्वाकूणां हि सर्वेषां पुरोधाः परमा गतिः}
{तस्मादनन्तरं सर्वे भवन्तो दैवतं मम} %||1-57-20||


{॥इत्यार्षे श्रीमद्रामायणे वाल्मीकीये आदिकाव्ये बालकाण्डे षट्पञ्चाशत्तमः सर्गः॥}

\sect{सप्तपञ्चाशत्तमः सर्गः}

\twolineshloka
{ततस्त्रिशङ्कोर्वचनं श्रुत्वा क्रोधसमन्वितम्}
{ऋषिपुत्रशतं राम राजानमिदमब्रवीत्} %||1-57-1||

\twolineshloka
{प्रत्याख्यातोऽसि दुर्बुद्धे गुरुणा सत्यवादिना}
{तं कथं समतिक्रम्य शाखान्तरमुपेयिवान्} %||1-57-2||

\twolineshloka
{इक्ष्वाकूणां हि सर्वेषां पुरोधाः परमा गतिः}
{न चातिक्रमितुं शक्यं वचनं सत्यवादिनः} %||1-57-3||

\twolineshloka
{अशक्यमिति चोवाच वसिष्ठो भगवानृषिः}
{तं वयं वै समाहर्तुं क्रतुं शक्ताः कथं तव} %||1-57-4||

\twolineshloka
{बालिशस्त्वं नरश्रेष्ठ गम्यतां स्वपुरं पुनः}
{याजने भगवाञ्शक्तस्त्रैलोक्यस्यापि पार्थिव} %||1-57-5||

\twolineshloka
{तेषां तद्वचनं श्रुत्वा क्रोधपर्याकुलाक्षरम्}
{स राजा पुनरेवैतानिदं वचनमब्रवीत्} %||1-57-6||

\twolineshloka
{प्रत्याख्यातोऽस्मि गुरुणा गुरुपुत्रैस्तथैव च}
{अन्यां गतिं गमिष्यामि स्वस्ति वोऽस्तु तपोधनाः} %||1-57-7||

\threelineshloka
{ऋषिपुत्रास्तु तच्छ्रुत्वा वाक्यं घोराभिसंहितम्}
{शेपुः परमसङ्क्रुद्धाश्चण्डालत्वं गमिष्यसि}
{एवमुक्त्वा महात्मानो विविशुस्ते स्वमाश्रमम्} %||1-57-8||

\threelineshloka
{अथ रात्र्यां व्यतीतायां राजा चण्डालतां गतः}
{नीलवस्त्रधरो नीलः परुषो ध्वस्तमूर्धजः}
{चित्यमाल्यानुलेपश्च आयसाभरणोऽभवत्} %||1-57-9||

\twolineshloka
{तं दृष्ट्वा मन्त्रिणः सर्वे त्यक्त्वा चण्डालरूपिणम्}
{प्राद्रवन्सहिता राम पौरा येऽस्यानुगामिनः} %||1-57-10||

\twolineshloka
{एको हि राजा काकुत्स्थ जगाम परमात्मवान्}
{दह्यमानो दिवारात्रं विश्वामित्रं तपोधनम्} %||1-57-11||

\twolineshloka
{विश्वामित्रस्तु तं दृष्ट्वा राजानं विफलीकृतम्}
{चण्डालरूपिणं राम मुनिः कारुण्यमागतः} %||1-57-12||

\twolineshloka
{कारुण्यात्स महातेजा वाक्यं परम धार्मिकः}
{इदं जगाद भद्रं ते राजानं घोरदर्शनम्} %||1-57-13||

\twolineshloka
{किमागमनकार्यं ते राजपुत्र महाबल}
{अयोध्याधिपते वीर शापाच्चण्डालतां गतः} %||1-57-14||

\twolineshloka
{अथ तद्वाक्यमाकर्ण्य राजा चण्डालतां गतः}
{अब्रवीत्प्राञ्जलिर्वाक्यं वाक्यज्ञो वाक्यकोविदम्} %||1-57-15||

\twolineshloka
{प्रत्याख्यातोऽस्मि गुरुणा गुरुपुत्रैस्तथैव च}
{अनवाप्यैव तं कामं मया प्राप्तो विपर्ययः} %||1-57-16||

\twolineshloka
{सशरीरो दिवं यायामिति मे सौम्यदर्शनम्}
{मया चेष्टं क्रतुशतं तच्च नावाप्यते फलम्} %||1-57-17||

\twolineshloka
{अनृतं नोक्त पूर्वं मे न च वक्ष्ये कदाचन}
{कृच्छ्रेष्वपि गतः सौम्य क्षत्रधर्मेण ते शपे} %||1-57-18||

\twolineshloka
{यज्ञैर्बहुविधैरिष्टं प्रजा धर्मेण पालिताः}
{गुरवश्च महात्मानः शीलवृत्तेन तोषिताः} %||1-57-19||

\twolineshloka
{धर्मे प्रयतमानस्य यज्ञं चाहर्तुमिच्छतः}
{परितोषं न गच्छन्ति गुरवो मुनिपुङ्गव} %||1-57-20||

\twolineshloka
{दैवमेव परं मन्ये पौरुषं तु निरर्थकम्}
{दैवेनाक्रम्यते सर्वं दैवं हि परमा गतिः} %||1-57-21||

\twolineshloka
{तस्य मे परमार्तस्य प्रसादमभिकाङ्क्षतः}
{कर्तुमर्हसि भद्रं ते दैवोपहतकर्मणः} %||1-57-22||

\twolineshloka
{नान्यां गतिं गमिष्यामि नान्यः शरणमस्ति मे}
{दैवं पुरुषकारेण निवर्तयितुमर्हसि} %||1-58-23||


{॥इत्यार्षे श्रीमद्रामायणे वाल्मीकीये आदिकाव्ये बालकाण्डे सप्तपञ्चाशत्तमः सर्गः॥}

\sect{अष्टपञ्चाशत्तमः सर्गः}

\twolineshloka
{उक्तवाक्यं तु राजानं कृपया कुशिकात्मजः}
{अब्रवीन्मधुरं वाक्यं साक्षाच्चण्डालरूपिणम्} %||1-58-1||

\twolineshloka
{इक्ष्वाको स्वागतं वत्स जानामि त्वां सुधार्मिकम्}
{शरणं ते भविष्यामि मा भैषीर्नृपपुङ्गव} %||1-58-2||

\twolineshloka
{अहमामन्त्रये सर्वान्महर्षीन्पुण्यकर्मणः}
{यज्ञसाह्यकरान्राजंस्ततो यक्ष्यसि निर्वृतः} %||1-58-3||

\twolineshloka
{गुरुशापकृतं रूपं यदिदं त्वयि वर्तते}
{अनेन सह रूपेण सशरीरो गमिष्यसि} %||1-58-4||

\twolineshloka
{हस्तप्राप्तमहं मन्ये स्वर्गं तव नरेश्वर}
{यस्त्वं कौशिकमागम्य शरण्यं शरणं गतः} %||1-58-5||

\twolineshloka
{एवमुक्त्वा महातेजाः पुत्रान्परमधार्मिकान्}
{व्यादिदेश महाप्राज्ञान्यज्ञसम्भारकारणात्} %||1-58-6||

{सर्वाञ्शिष्यान्समाहूय वाक्यमेतदुवाच ह॥\devanumber{7}॥} %||1-58-7|| (Check)

\addtocounter{shlokacount}{1}

\twolineshloka
{सर्वानृषिवरान्वत्सा आनयध्वं ममाज्ञया}
{सशिष्यान्सुहृदश्चैव सर्त्विजः सुबहुश्रुतान्} %||1-58-8||

\twolineshloka
{यदन्यो वचनं ब्रूयान्मद्वाक्यबलचोदितः}
{तत्सर्वमखिलेनोक्तं ममाख्येयमनादृतम्} %||1-58-9||

\twolineshloka
{तस्य तद्वचनं श्रुत्वा दिशो जग्मुस्तदाज्ञया}
{आजग्मुरथ देशेभ्यः सर्वेभ्यो ब्रह्मवादिनः} %||1-58-10||

\twolineshloka
{ते च शिष्याः समागम्य मुनिं ज्वलिततेजसम्}
{ऊचुश्च वचनं सर्वे सर्वेषां ब्रह्मवादिनाम्} %||1-58-11||

\twolineshloka
{श्रुत्वा ते वचनं सर्वे समायान्ति द्विजातयः}
{सर्वदेशेषु चागच्छन्वर्जयित्वा महोदयम्} %||1-58-12||

\twolineshloka
{वासिष्ठं तच्छतं सर्वं क्रोधपर्याकुलाक्षरम्}
{यदाह वचनं सर्वं शृणु त्वं मुनिपुङ्गव} %||1-58-13||

\twolineshloka
{क्षत्रियो याजको यस्य चण्डालस्य विशेषतः}
{कथं सदसि भोक्तारो हविस्तस्य सुरर्षयः} %||1-58-14||

\twolineshloka
{ब्राह्मणा वा महात्मानो भुक्त्वा चण्डालभोजनम्}
{कथं स्वर्गं गमिष्यन्ति विश्वामित्रेण पालिताः} %||1-58-15||

\twolineshloka
{एतद्वचनं नैष्ठुर्यमूचुः संरक्तलोचनाः}
{वासिष्ठा मुनिशार्दूल सर्वे ते समहोदयाः} %||1-58-16||

\twolineshloka
{तेषां तद्वचनं श्रुत्वा सर्वेषां मुनिपुङ्गवः}
{क्रोधसंरक्तनयनः सरोषमिदमब्रवीत्} %||1-58-17||

\twolineshloka
{यद्दूषयन्त्यदुष्टं मां तप उग्रं समास्थितम्}
{भस्मीभूता दुरात्मानो भविष्यन्ति न संशयः} %||1-58-18||

\twolineshloka
{अद्य ते कालपाशेन नीता वैवस्वतक्षयम्}
{सप्तजातिशतान्येव मृतपाः सन्तु सर्वशः} %||1-58-19||

\twolineshloka
{श्वमांसनियताहारा मुष्टिका नाम निर्घृणाः}
{विकृताश्च विरूपाश्च लोकाननुचरन्त्विमान्} %||1-58-20||

\twolineshloka
{महोदयश्च दुर्बुद्धिर्मामदूष्यं ह्यदूषयत्}
{दूषिटः सर्वलोकेषु निषादत्वं गमिष्यति} %||1-58-21||

\twolineshloka
{प्राणातिपातनिरतो निरनुक्रोशतां गतः}
{दीर्घकालं मम क्रोधाद्दुर्गतिं वर्तयिष्यति} %||1-58-22||

\twolineshloka
{एतावदुक्त्वा वचनं विश्वामित्रो महातपाः}
{विरराम महातेजा ऋषिमध्ये महामुनिः} %||1-59-23||


{॥इत्यार्षे श्रीमद्रामायणे वाल्मीकीये आदिकाव्ये बालकाण्डे अष्टपञ्चाशत्तमः सर्गः॥}

\sect{एकोनषष्टितमः सर्गः}

\twolineshloka
{तपोबलहतान्कृत्वा वासिष्ठान्समहोदयान्}
{ऋषिमध्ये महातेजा विश्वामित्रोऽभ्यभाषत} %||1-59-1||

\threelineshloka
{अयमिक्ष्वाकुदायादस्त्रिशङ्कुरिति विश्रुतः}
{धर्मिष्ठश्च वदान्यश्च मां चैव शरणं गतः}
{स्वेनानेन शरीरेण देवलोकजिगीषया} %||1-59-2||

\twolineshloka
{यथायं स्वशरीरेण देवलोकं गमिष्यति}
{तथा प्रवर्त्यतां यज्ञो भवद्भिश्च मया सह} %||1-59-3||

\twolineshloka
{विश्वामित्रवचः श्रुत्वा सर्व एव महर्षयः}
{ऊचुः समेत्य सहिता धर्मज्ञा धर्मसंहितम्} %||1-59-4||

\twolineshloka
{अयं कुशिकदायादो मुनिः परमकोपनः}
{यदाह वचनं सम्यगेतत्कार्यं न संशयः} %||1-59-5||

\threelineshloka
{अग्निकल्पो हि भगवाञ्शापं दास्यति रोषितः}
{तस्मात्प्रवर्त्यतां यज्ञः सशरीरो यथा दिवम्}
{गच्छेदिक्ष्वाकुदायादो विश्वामित्रस्य तेजसा} %||1-59-6||

{ततः प्रवर्त्यतां यज्ञः सर्वे समधितिष्ठते॥\devanumber{7}॥} %||1-59-7|| (Check)

\addtocounter{shlokacount}{1}

\twolineshloka
{एवमुक्त्वा महर्षयः संजह्रुस्ताः क्रियास्तदा}
{याजकाश्च महातेजा विश्वामित्रोऽभवत्क्रतौ} %||1-59-8||

\twolineshloka
{ऋत्विजश्चानुपूर्व्येण मन्त्रवन्मन्त्रकोविदाः}
{चक्रुः सर्वाणि कर्माणि यथाकल्पं यथाविधि} %||1-59-9||

\twolineshloka
{ततः कालेन महता विश्वामित्रो महातपाः}
{चकारावाहनं तत्र भागार्थं सर्वदेवताः} %||1-59-10||

\twolineshloka
{नाह्यागमंस्तदाहूता भागार्थं सर्वदेवताः}
{ततः क्रोधसमाविष्टो विश्वमित्रो महामुनिः} %||1-59-11||

\twolineshloka
{स्रुवमुद्यम्य सक्रोधस्त्रिशङ्कुमिदमब्रवीत्}
{पश्य मे तपसो वीर्यं स्वार्जितस्य नरेश्वर} %||1-59-12||

\twolineshloka
{एष त्वां स्वशरीरेण नयामि स्वर्गमोजसा}
{दुष्प्रापं स्वशरीरेण दिवं गच्छ नराधिप} %||1-59-13||

\twolineshloka
{स्वार्जितं किञ्चिदप्यस्ति मया हि तपसः फलम्}
{राजंस्त्वं तेजसा तस्य सशरीरो दिवं व्रज} %||1-59-14||

\twolineshloka
{उक्तवाक्ये मुनौ तस्मिन्सशरीरो नरेश्वरः}
{दिवं जगाम काकुत्स्थ मुनीनां पश्यतां तदा} %||1-59-15||

\twolineshloka
{देवलोकगतं दृष्ट्वा त्रिशङ्कुं पाकशासनः}
{सह सर्वैः सुरगणैरिदं वचनमब्रवीत्} %||1-59-16||

\twolineshloka
{त्रिशङ्को गच्छ भूयस्त्वं नासि स्वर्गकृतालयः}
{गुरुशापहतो मूढ पत भूमिमवाक्शिराः} %||1-59-17||

\twolineshloka
{एवमुक्तो महेन्द्रेण त्रिशङ्कुरपतत्पुनः}
{विक्रोशमानस्त्राहीति विश्वामित्रं तपोधनम्} %||1-59-18||

\twolineshloka
{तच्छ्रुत्वा वचनं तस्य क्रोशमानस्य कौशिकः}
{रोषमाहारयत्तीव्रं तिष्ठ तिष्ठेति चाब्रवीत्} %||1-59-19||

\twolineshloka
{ऋषिमध्ये स तेजस्वी प्रजापतिरिवापरः}
{सृजन्दक्षिणमार्गस्थान्सप्तर्षीनपरान्पुनः} %||1-59-20||

\twolineshloka
{नक्षत्रमालामपरामसृजत्क्रोधमूर्छितः}
{दक्षिणां दिशमास्थाय मुनिमध्ये महायशाः} %||1-59-21||

\threelineshloka
{सृष्ट्वा नक्षत्रवंशं च क्रोधेन कलुषीकृतः}
{अन्यमिन्द्रं करिष्यामि लोको वा स्यादनिन्द्रकः}
{दैवतान्यपि स क्रोधात्स्रष्टुं समुपचक्रमे} %||1-59-22||

\twolineshloka
{ततः परमसम्भ्रान्ताः सर्षिसङ्घाः सुरर्षभाः}
{विश्वामित्रं महात्मानमूचुः सानुनयं वचः} %||1-59-23||

\twolineshloka
{अयं राजा महाभाग गुरुशापपरिक्षतः}
{सशरीरो दिवं यातुं नार्हत्येव तपोधन} %||1-59-24||

\twolineshloka
{तेषां तद्वचनं श्रुत्वा देवानां मुनिपुङ्गवः}
{अब्रवीत्सुमहद्वाक्यं कौशिकः सर्वदेवताः} %||1-59-25||

\twolineshloka
{सशरीरस्य भद्रं वस्त्रिशङ्कोरस्य भूपतेः}
{आरोहणं प्रतिज्ञाय नानृतं कर्तुमुत्सहे} %||1-59-26||

\twolineshloka
{सर्गोऽस्तु सशरीरस्य त्रिशङ्कोरस्य शाश्वतः}
{नक्षत्राणि च सर्वाणि मामकानि ध्रुवाण्यथ} %||1-59-27||

\twolineshloka
{यावल्लोका धरिष्यन्ति तिष्ठन्त्वेतानि सर्वशः}
{मत्कृतानि सुराः सर्वे तदनुज्ञातुमर्हथ} %||1-59-28||

{एवमुक्ताः सुराः सर्वे प्रत्यूचुर्मुनिपुङ्गवम्॥\devanumber{29}॥} %||1-59-29|| (Check)

\addtocounter{shlokacount}{1}

\twolineshloka
{एवं भवतु भद्रं ते तिष्ठन्त्वेतानि सर्वशः}
{गगने तान्यनेकानि वैश्वानरपथाद्बहिः} %||1-59-30||

\twolineshloka
{नक्षत्राणि मुनिश्रेष्ठ तेषु ज्योतिःषु जाज्वलन्}
{अवाक्शिरास्त्रिशङ्कुश्च तिष्ठत्वमरसंनिभः} %||1-59-31||

\twolineshloka
{विश्वामित्रस्तु धर्मात्मा सर्वदेवैरभिष्टुतः}
{ऋषिभिश्च महातेजा बाढमित्याह देवताः} %||1-59-32||

\twolineshloka
{ततो देवा महात्मानो मुनयश्च तपोधनाः}
{जग्मुर्यथागतं सर्वे यज्ञस्यान्ते नरोत्तम} %||1-60-33||


{॥इत्यार्षे श्रीमद्रामायणे वाल्मीकीये आदिकाव्ये बालकाण्डे एकोनषष्टितमः सर्गः॥}

\sect{षष्टितमः सर्गः}

\twolineshloka
{विश्वामित्रो महात्माथ प्रस्थितान्प्रेक्ष्य तानृषीन्}
{अब्रवीन्नरशार्दूल सर्वांस्तान्वनवासिनः} %||1-60-1||

\twolineshloka
{महाविघ्नः प्रवृत्तोऽयं दक्षिणामास्थितो दिशम्}
{दिशमन्यां प्रपत्स्यामस्तत्र तप्स्यामहे तपः} %||1-60-2||

\twolineshloka
{पश्चिमायां विशालायां पुष्करेषु महात्मनः}
{सुखं तपश्चरिष्यामः परं तद्धि तपोवनम्} %||1-60-3||

\twolineshloka
{एवमुक्त्वा महातेजाः पुष्करेषु महामुनिः}
{तप उग्रं दुराधर्षं तेपे मूलफलाशनः} %||1-60-4||

\twolineshloka
{एतस्मिन्नेव काले तु अयोध्याधिपतिर्नृपः}
{अम्बरीष इति ख्यातो यष्टुं समुपचक्रमे} %||1-60-5||

\twolineshloka
{तस्य वै यजमानस्य पशुमिन्द्रो जहार ह}
{प्रनष्टे तु पशौ विप्रो राजानमिदमब्रवीत्} %||1-60-6||

\twolineshloka
{पशुरद्य हृतो राजन्प्रनष्टस्तव दुर्नयात्}
{अरक्षितारं राजानं घ्नन्ति दोषा नरेश्वर} %||1-60-7||

\twolineshloka
{प्रायश्चित्तं महद्ध्येतन्नरं वा पुरुषर्षभ}
{आनयस्व पशुं शीघ्रं यावत्कर्म प्रवर्तते} %||1-60-8||

\twolineshloka
{उपाध्याय वचः श्रुत्वा स राजा पुरुषर्षभ}
{अन्वियेष महाबुद्धिः पशुं गोभिः सहस्रशः} %||1-60-9||

\twolineshloka
{देशाञ्जनपदांस्तांस्तान्नगराणि वनानि च}
{आश्रमाणि च पुण्यानि मार्गमाणो महीपतिः} %||1-60-10||

\twolineshloka
{स पुत्रसहितं तात सभार्यं रघुनन्दन}
{भृगुतुन्दे समासीनमृचीकं सन्ददर्श ह} %||1-60-11||

\threelineshloka
{तमुवाच महातेजाः प्रणम्याभिप्रसाद्य च}
{ब्रह्मर्षिं तपसा दीप्तं राजर्षिरमितप्रभः}
{पृष्ट्वा सर्वत्र कुशलमृचीकं तमिदं वचः} %||1-60-12||

\twolineshloka
{गवां शतसहस्रेण विक्रिणीषे सुतं यदि}
{पशोरर्थे महाभाग कृतकृत्योऽस्मि भार्गव} %||1-60-13||

\twolineshloka
{सर्वे परिसृता देशा यज्ञियं न लभे पशुम्}
{दातुमर्हसि मूल्येन सुतमेकमितो मम} %||1-60-14||

\twolineshloka
{एवमुक्तो महातेजा ऋचीकस्त्वब्रवीद्वचः}
{नाहं ज्येष्ठं नरश्रेष्ठं विक्रीणीयां कथञ्चन} %||1-60-15||

\twolineshloka
{ऋचीकस्य वचः श्रुत्वा तेषां माता महात्मनाम्}
{उवाच नरशार्दूलमम्बरीषं तपस्विनी} %||1-60-16||

{ममापि दयितं विद्धि कनिष्ठं शुनकं नृप॥\devanumber{17}॥} %||1-60-17|| (Check)

\addtocounter{shlokacount}{1}

\twolineshloka
{प्रायेण हि नरश्रेष्ठ ज्येष्ठाः पितृषु वल्लभाः}
{मातॄणां च कनीयांसस्तस्माद्रक्षे कनीयसम्} %||1-60-18||

\twolineshloka
{उक्तवाक्ये मुनौ तस्मिन्मुनिपत्न्यां तथैव च}
{शुनःशेपः स्वयं राम मध्यमो वाक्यमब्रवीत्} %||1-60-19||

\twolineshloka
{पिता ज्येष्ठमविक्रेयं माता चाह कनीयसम्}
{विक्रीतं मध्यमं मन्ये राजन्पुत्रं नयस्व माम्} %||1-60-20||

\twolineshloka
{गवां शतसहस्रेण शुनःशेपं नरेश्वरः}
{गृहीत्वा परमप्रीतो जगाम रघुनन्दन} %||1-60-21||

\twolineshloka
{अम्बरीषस्तु राजर्षी रथमारोप्य सत्वरः}
{शुनःशेपं महातेजा जगामाशु महायशाः} %||1-61-22||


{॥इत्यार्षे श्रीमद्रामायणे वाल्मीकीये आदिकाव्ये बालकाण्डे षष्टितमः सर्गः॥}

\sect{एकषष्टितमः सर्गः}

\twolineshloka
{शुनःशेपं नरश्रेष्ठ गृहीत्वा तु महायशाः}
{व्यश्राम्यत्पुष्करे राजा मध्याह्ने रघुनन्दन} %||1-61-1||

\twolineshloka
{तस्य विश्रममाणस्य शुनःशेपो महायशाः}
{पुष्करं श्रेष्ठमागम्य विश्वामित्रं ददर्श ह} %||1-61-2||

\twolineshloka
{विषण्णवदनो दीनस्तृष्णया च श्रमेण च}
{पपाताङ्के मुने राम वाक्यं चेदमुवाच ह} %||1-61-3||

\twolineshloka
{न मेऽस्ति माता न पिता ज्ञातयो बान्धवाः कुतः}
{त्रातुमर्हसि मां सौम्य धर्मेण मुनिपुङ्गव} %||1-61-4||

\twolineshloka
{त्राता त्वं हि मुनिश्रेष्ठ सर्वेषां त्वं हि भावनः}
{राजा च कृतकार्यः स्यादहं दीर्घायुरव्ययः} %||1-61-5||

\threelineshloka
{स्वर्गलोकमुपाश्नीयां तपस्तप्त्वा ह्यनुत्तमम्}
{स मे नाथो ह्यनाथस्य भव भव्येन चेतसा}
{पितेव पुत्रं धर्मात्मंस्त्रातुमर्हसि किल्बिषात्} %||1-61-6||

\twolineshloka
{तस्य तद्वचनं श्रुत्वा विश्वामित्रो महातपाः}
{सान्त्वयित्वा बहुविधं पुत्रानिदमुवाच ह} %||1-61-7||

\twolineshloka
{यत्कृते पितरः पुत्राञ्जनयन्ति शुभार्थिनः}
{परलोकहितार्थाय तस्य कालोऽयमागतः} %||1-61-8||

\twolineshloka
{अयं मुनिसुतो बालो मत्तः शरणमिच्छति}
{अस्य जीवितमात्रेण प्रियं कुरुत पुत्रकाः} %||1-61-9||

\twolineshloka
{सर्वे सुकृतकर्माणः सर्वे धर्मपरायणाः}
{पशुभूता नरेन्द्रस्य तृप्तिमग्नेः प्रयच्छत} %||1-61-10||

\twolineshloka
{नाथवांश्च शुनःशेपो यज्ञश्चाविघ्नतो भवेत्}
{देवतास्तर्पिताश्च स्युर्मम चापि कृतं वचः} %||1-61-11||

\twolineshloka
{मुनेस्तु वचनं श्रुत्वा मधुष्यन्दादयः सुताः}
{साभिमानं नरश्रेष्ठ सलीलमिदमब्रुवन्} %||1-61-12||

\twolineshloka
{कथमात्मसुतान्हित्वा त्रायसेऽन्यसुतं विभो}
{अकार्यमिव पश्यामः श्वमांसमिव भोजने} %||1-61-13||

\twolineshloka
{तेषां तद्वचनं श्रुत्वा पुत्राणां मुनिपुङ्गवः}
{क्रोधसंरक्तनयनो व्याहर्तुमुपचक्रमे} %||1-61-14||

\twolineshloka
{निःसाध्वसमिदं प्रोक्तं धर्मादपि विगर्हितम्}
{अतिक्रम्य तु मद्वाक्यं दारुणं रोमहर्षणम्} %||1-61-15||

\twolineshloka
{श्वमांसभोजिनः सर्वे वासिष्ठा इव जातिषु}
{पूर्णं वर्षसहस्रं तु पृथिव्यामनुवत्स्यथ} %||1-61-16||

\twolineshloka
{कृत्वा शापसमायुक्तान्पुत्रान्मुनिवरस्तदा}
{शुनःशेपमुवाचार्तं कृत्वा रक्षां निरामयाम्} %||1-61-17||

\twolineshloka
{पवित्रपाशैरासक्तो रक्तमाल्यानुलेपनः}
{वैष्णवं यूपमासाद्य वाग्भिरग्निमुदाहर} %||1-61-18||

\twolineshloka
{इमे तु गाथे द्वे दिव्ये गायेथा मुनिपुत्रक}
{अम्बरीषस्य यज्ञेऽस्मिंस्ततः सिद्धिमवाप्स्यसि} %||1-61-19||

\twolineshloka
{शुनःशेपो गृहीत्वा ते द्वे गाथे सुसमाहितः}
{त्वरया राजसिंहं तमम्बरीषमुवाच ह} %||1-61-20||

\twolineshloka
{राजसिंह महासत्त्व शीघ्रं गच्छावहे सदः}
{निवर्तयस्व राजेन्द्र दीक्षां च समुपाहर} %||1-61-21||

\twolineshloka
{तद्वाक्यमृषिपुत्रस्य श्रुत्वा हर्षं समुत्सुकः}
{जगाम नृपतिः शीघ्रं यज्ञवाटमतन्द्रितः} %||1-61-22||

\twolineshloka
{सदस्यानुमते राजा पवित्रकृतलक्षणम्}
{पशुं रक्ताम्बरं कृत्वा यूपे तं समबन्धयत्} %||1-61-23||

\twolineshloka
{स बद्धो वाग्भिरग्र्याभिरभितुष्टाव वै सुरौ}
{इन्द्रमिन्द्रानुजं चैव यथावन्मुनिपुत्रकः} %||1-61-24||

\twolineshloka
{ततः प्रीतः सहस्राक्षो रहस्यस्तुतितर्पितः}
{दीर्घमायुस्तदा प्रादाच्छुनःशेपाय राघव} %||1-61-25||

\twolineshloka
{स च राजा नरश्रेष्ठ यज्ञस्य च समाप्तवान्}
{फलं बहुगुणं राम सहस्राक्षप्रसादजम्} %||1-61-26||

\twolineshloka
{विश्वामित्रोऽपि धर्मात्मा भूयस्तेपे महातपाः}
{पुष्करेषु नरश्रेष्ठ दशवर्षशतानि च} %||1-62-27||


{॥इत्यार्षे श्रीमद्रामायणे वाल्मीकीये आदिकाव्ये बालकाण्डे एकषष्टितमः सर्गः॥}

\sect{द्विषष्टितमः सर्गः}

\twolineshloka
{पूर्णे वर्षसहस्रे तु व्रतस्नातं महामुनिम्}
{अभ्यागच्छन्सुराः सर्वे तपःफलचिकीर्षवः} %||1-62-1||

\twolineshloka
{अब्रवीत्सुमहातेजा ब्रह्मा सुरुचिरं वचः}
{ऋषिस्त्वमसि भद्रं ते स्वार्जितैः कर्मभिः शुभैः} %||1-62-2||

\twolineshloka
{तमेवमुक्त्वा देवेशस्त्रिदिवं पुनरभ्यगात्}
{विश्वामित्रो महातेजा भूयस्तेपे महत्तपः} %||1-62-3||

\twolineshloka
{ततः कालेन महता मेनका परमाप्सराः}
{पुष्करेषु नरश्रेष्ठ स्नातुं समुपचक्रमे} %||1-62-4||

\twolineshloka
{तां ददर्श महातेजा मेनकां कुशिकात्मजः}
{रूपेणाप्रतिमां तत्र विद्युतं जलदे यथा} %||1-62-5||

\threelineshloka
{दृष्ट्वा कन्दर्पवशगो मुनिस्तामिदमब्रवीत्}
{अप्सरः स्वागतं तेऽस्तु वस चेह ममाश्रमे}
{अनुगृह्णीष्व भद्रं ते मदनेन सुमोहितम्} %||1-62-6||

\twolineshloka
{इत्युक्ता सा वरारोहा तत्रावासमथाकरोत्}
{तपसो हि महाविघ्नो विश्वामित्रमुपागतः} %||1-62-7||

\twolineshloka
{तस्यां वसन्त्यां वर्षाणि पञ्च पञ्च च राघव}
{विश्वामित्राश्रमे सौम्य सुखेन व्यतिचक्रमुः} %||1-62-8||

\twolineshloka
{अथ काले गते तस्मिन्विश्वामित्रो महामुनिः}
{सव्रीड इव संवृत्तश्चिन्ताशोकपरायणः} %||1-62-9||

\twolineshloka
{बुद्धिर्मुनेः समुत्पन्ना सामर्षा रघुनन्दन}
{सर्वं सुराणां कर्मैतत्तपोऽपहरणं महत्} %||1-62-10||

\twolineshloka
{अहोरात्रापदेशेन गताः संवत्सरा दश}
{काममोहाभिभूतस्य विघ्नोऽयं प्रत्युपस्थितः} %||1-62-11||

{विनिःश्वसन्मुनिवरः पश्चात्तापेन दुःखितः॥\devanumber{12}॥} %||1-62-12|| (Check)

\addtocounter{shlokacount}{1}

\threelineshloka
{भीतामप्सरसं दृष्ट्वा वेपन्तीं प्राञ्जलिं स्थिताम्}
{मेनकां मधुरैर्वाक्यैर्विसृज्य कुशिकात्मजः}
{उत्तरं पर्वतं राम विश्वामित्रो जगाम ह} %||1-62-13||

\twolineshloka
{स कृत्वा नैष्ठिकीं बुद्धिं जेतुकामो महायशाः}
{कौशिकीतीरमासाद्य तपस्तेपे सुदारुणम्} %||1-62-14||

\twolineshloka
{तस्य वर्षसहस्रं तु घोरं तप उपासतः}
{उत्तरे पर्वते राम देवतानामभूद्भयम्} %||1-62-15||

\twolineshloka
{अमन्त्रयन्समागम्य सर्वे सर्षिगणाः सुराः}
{महर्षिशब्दं लभतां साध्वयं कुशिकात्मजः} %||1-62-16||

\twolineshloka
{देवतानां वचः श्रुत्वा सर्वलोकपितामहः}
{अब्रवीन्मधुरं वाक्यं विश्वामित्रं तपोधनम्} %||1-62-17||

\twolineshloka
{महर्षे स्वागतं वत्स तपसोग्रेण तोषितः}
{महत्त्वमृषिमुख्यत्वं ददामि तव कौशिक} %||1-62-18||

\twolineshloka
{ब्रह्मणः स वचः श्रुत्वा विश्वामित्रस्तपोधनः}
{प्राञ्जलिः प्रणतो भूत्वा प्रत्युवाच पितामहम्} %||1-62-19||

\twolineshloka
{ब्रह्मर्षि शब्दमतुलं स्वार्जितैः कर्मभिः शुभैः}
{यदि मे भगवानाह ततोऽहं विजितेन्द्रियः} %||1-62-20||

\twolineshloka
{तमुवाच ततो ब्रह्मा न तावत्त्वं जितेन्द्रियः}
{यतस्व मुनिशार्दूल इत्युक्त्वा त्रिदिवं गतः} %||1-62-21||

\twolineshloka
{विप्रस्थितेषु देवेषु विश्वामित्रो महामुनिः}
{ऊर्ध्वबाहुर्निरालम्बो वायुभक्षस्तपश्चरन्} %||1-62-22||

\twolineshloka
{धर्मे पञ्चतपा भूत्वा वर्षास्वाकाशसंश्रयः}
{शिशिरे सलिलस्थायी रात्र्यहानि तपोधनः} %||1-62-23||

\twolineshloka
{एवं वर्षसहस्रं हि तपो घोरमुपागमत्}
{तस्मिन्सन्तप्यमाने तु विश्वामित्रे महामुनौ} %||1-62-24||

\twolineshloka
{सम्भ्रमः सुमहानासीत्सुराणां वासवस्य च}
{रम्भामप्सरसं शक्रः सह सर्वैर्मरुद्गणैः} %||1-62-25||

{उवाचात्महितं वाक्यमहितं कौशिकस्य च॥\devanumber{26}॥} %||1-63-26|| (Check)

\addtocounter{shlokacount}{1}


{॥इत्यार्षे श्रीमद्रामायणे वाल्मीकीये आदिकाव्ये बालकाण्डे द्विषष्टितमः सर्गः॥}

\sect{त्रिषष्टितमः सर्गः}

\twolineshloka
{सुरकार्यमिदं रम्भे कर्तव्यं सुमहत्त्वया}
{लोभनं कौशिकस्येह काममोहसमन्वितम्} %||1-63-1||

\twolineshloka
{तथोक्ता साप्सरा राम सहस्राक्षेण धीमता}
{व्रीडिता प्राञ्जलिर्भूत्वा प्रत्युवाच सुरेश्वरम्} %||1-63-2||

\threelineshloka
{अयं सुरपते घोरो विश्वामित्रो महामुनिः}
{क्रोधमुत्स्रक्ष्यते घोरं मयि देव न संशयः}
{ततो हि मे भयं देव प्रसादं कर्तुमर्हसि} %||1-63-3||

\twolineshloka
{तामुवाच सहस्राक्षो वेपमानां कृताञ्जलिम्}
{मा भैषि रम्भे भद्रं ते कुरुष्व मम शासनम्} %||1-63-4||

\twolineshloka
{कोकिलो हृदयग्राही माधवे रुचिरद्रुमे}
{अहं कन्दर्पसहितः स्थास्यामि तव पार्श्वतः} %||1-63-5||

\twolineshloka
{त्वं हि रूपं बहुगुणं कृत्वा परमभास्वरम्}
{तमृषिं कौशिकं रम्भे भेदयस्व तपस्विनम्} %||1-63-6||

\twolineshloka
{सा श्रुत्वा वचनं तस्य कृत्वा रूपमनुत्तमम्}
{लोभयामास ललिता विश्वामित्रं शुचिस्मिता} %||1-63-7||

\twolineshloka
{कोकिलस्य तु शुश्राव वल्गु व्याहरतः स्वनम्}
{सम्प्रहृष्टेन मनसा तत एनामुदैक्षत} %||1-63-8||

\twolineshloka
{अथ तस्य च शब्देन गीतेनाप्रतिमेन च}
{दर्शनेन च रम्भाया मुनिः सन्देहमागतः} %||1-63-9||

\twolineshloka
{सहस्राक्षस्य तत्कर्म विज्ञाय मुनिपुङ्गवः}
{रम्भां क्रोधसमाविष्टः शशाप कुशिकात्मजः} %||1-63-10||

\twolineshloka
{यन्मां लोभयसे रम्भे कामक्रोधजयैषिणम्}
{दशवर्षसहस्राणि शैली स्थास्यसि दुर्भगे} %||1-63-11||

\twolineshloka
{ब्राह्मणः सुमहातेजास्तपोबलसमन्वितः}
{उद्धरिष्यति रम्भे त्वां मत्क्रोधकलुषीकृताम्} %||1-63-12||

\twolineshloka
{एवमुक्त्वा महातेजा विश्वामित्रो महामुनिः}
{अशक्नुवन्धारयितुं कोपं सन्तापमागतः} %||1-63-13||

\twolineshloka
{तस्य शापेन महता रम्भा शैली तदाभवत्}
{वचः श्रुत्वा च कन्दर्पो महर्षेः स च निर्गतः} %||1-63-14||

\twolineshloka
{कोपेन स महातेजास्तपोऽपहरणे कृते}
{इन्द्रियैरजितै राम न लेभे शान्तिमात्मनः} %||1-64-15||


{॥इत्यार्षे श्रीमद्रामायणे वाल्मीकीये आदिकाव्ये बालकाण्डे त्रिषष्टितमः सर्गः॥}

\sect{चतुःषष्टितमः सर्गः}

\twolineshloka
{अथ हैमवतीं राम दिशं त्यक्त्वा महामुनिः}
{पूर्वां दिशमनुप्राप्य तपस्तेपे सुदारुणम्} %||1-64-1||

\twolineshloka
{मौनं वर्षसहस्रस्य कृत्वा व्रतमनुत्तमम्}
{चकाराप्रतिमं राम तपः परमदुष्करम्} %||1-64-2||

\twolineshloka
{पूर्णे वर्षसहस्रे तु काष्ठभूतं महामुनिम्}
{विघ्नैर्बहुभिराधूतं क्रोधो नान्तरमाविशत्} %||1-64-3||

\threelineshloka
{ततो देवाः सगन्धर्वाः पन्नगासुरराक्षसाः}
{मोहितास्तेजसा तस्य तपसा मन्दरश्मयः}
{कश्मलोपहताः सर्वे पितामहमथाब्रुवन्} %||1-64-4||

\twolineshloka
{बहुभिः कारणैर्देव विश्वामित्रो महामुनिः}
{लोभितः क्रोधितश्चैव तपसा चाभिवर्धते} %||1-64-5||

\fourlineshloka
{न ह्यस्य वृजिनं किञ्चिद्दृश्यते सूक्ष्ममप्यथ}
{न दीयते यदि त्वस्य मनसा यदभीप्सितम्}
{विनाशयति त्रैलोक्यं तपसा सचराचरम्}
{व्याकुलाश्च दिशः सर्वा न च किञ्चित्प्रकाशते} %||1-64-6||

\twolineshloka
{सागराः क्षुभिताः सर्वे विशीर्यन्ते च पर्वताः}
{प्रकम्पते च पृथिवी वायुर्वाति भृशाकुलः} %||1-64-7||

\twolineshloka
{बुद्धिं न कुरुते यावन्नाशे देव महामुनिः}
{तावत्प्रसाद्यो भगवानग्निरूपो महाद्युतिः} %||1-64-8||

\twolineshloka
{कालाग्निना यथा पूर्वं त्रैलोक्यं दह्यतेऽखिलम्}
{देवराज्ये चिकीर्षेत दीयतामस्य यन्मतम्} %||1-64-9||

\twolineshloka
{ततः सुरगणाः सर्वे पितामहपुरोगमाः}
{विश्वामित्रं महात्मानं वाक्यं मधुरमब्रुवन्} %||1-64-10||

\twolineshloka
{ब्रह्मर्षे स्वागतं तेऽस्तु तपसा स्म सुतोषिताः}
{ब्राह्मण्यं तपसोग्रेण प्राप्तवानसि कौशिक} %||1-64-11||

\twolineshloka
{दीर्घमायुश्च ते ब्रह्मन्ददामि समरुद्गणः}
{स्वस्ति प्राप्नुहि भद्रं ते गच्छ सौम्य यथासुखम्} %||1-64-12||

\twolineshloka
{पितामहवचः श्रुत्वा सर्वेषां च दिवौकसाम्}
{कृत्वा प्रणामं मुदितो व्याजहार महामुनिः} %||1-64-13||

\twolineshloka
{ब्राह्मण्यं यदि मे प्राप्तं दीर्घमायुस्तथैव च}
{ओङ्कारोऽथ वषट्कारो वेदाश्च वरयन्तु माम्} %||1-64-14||

\threelineshloka
{क्षत्रवेदविदां श्रेष्ठो ब्रह्मवेदविदामपि}
{ब्रह्मपुत्रो वसिष्ठो मामेवं वदतु देवताः}
{यद्ययं परमः कामः कृतो यान्तु सुरर्षभाः} %||1-64-15||

\twolineshloka
{ततः प्रसादितो देवैर्वसिष्ठो जपतां वरः}
{सख्यं चकार ब्रह्मर्षिरेवमस्त्विति चाब्रवीत्} %||1-64-16||

\twolineshloka
{ब्रह्मर्षित्वं न सन्देहः सर्वं सम्पत्स्यते तव}
{इत्युक्त्वा देवताश्चापि सर्वा जग्मुर्यथागतम्} %||1-64-17||

\twolineshloka
{विश्वामित्रोऽपि धर्मात्मा लब्ध्वा ब्राह्मण्यमुत्तमम्}
{पूजयामास ब्रह्मर्षिं वसिष्ठं जपतां वरम्} %||1-64-18||

\twolineshloka
{कृतकामो महीं सर्वां चचार तपसि स्थितः}
{एवं त्वनेन ब्राह्मण्यं प्राप्तं राम महात्मना} %||1-64-19||

\twolineshloka
{एष राम मुनिश्रेष्ठ एष विग्रहवांस्तपः}
{एष धर्मः परो नित्यं वीर्यस्यैष परायणम्} %||1-64-20||

\twolineshloka
{शतानन्दवचः श्रुत्वा रामलक्ष्मणसंनिधौ}
{जनकः प्राञ्जलिर्वाक्यमुवाच कुशिकात्मजम्} %||1-64-21||

\twolineshloka
{धन्योऽस्म्यनुगृहीतोऽस्मि यस्य मे मुनिपुङ्गव}
{यज्ञं काकुत्स्थ सहितः प्राप्तवानसि धार्मिक} %||1-64-22||

\twolineshloka
{पावितोऽहं त्वया ब्रह्मन्दर्शनेन महामुने}
{गुणा बहुविधाः प्राप्तास्तव सन्दर्शनान्मया} %||1-64-23||

\twolineshloka
{विस्तरेण च ते ब्रह्मन्कीर्त्यमानं महत्तपः}
{श्रुतं मया महातेजो रामेण च महात्मना} %||1-64-24||

{सदस्यैः प्राप्य च सदः श्रुतास्ते बहवो गुणाः॥\devanumber{25}॥} %||1-64-25|| (Check)

\addtocounter{shlokacount}{1}

\twolineshloka
{अप्रमेयं तपस्तुभ्यमप्रमेयं च ते बलम्}
{अप्रमेया गुणाश्चैव नित्यं ते कुशिकात्मज} %||1-64-26||

\twolineshloka
{तृप्तिराश्चर्यभूतानां कथानां नास्ति मे विभो}
{कर्मकालो मुनिश्रेष्ठ लम्बते रविमण्डलम्} %||1-64-27||

\twolineshloka
{श्वः प्रभाते महातेजो द्रष्टुमर्हसि मां पुनः}
{स्वागतं तपसां श्रेष्ठ मामनुज्ञातुमर्हसि} %||1-64-28||

\twolineshloka
{एवमुक्त्वा मुनिश्रेष्ठं वैदेहो मिथिलाधिपः}
{प्रदक्षिणं चकाराशु सोपाध्यायः सबान्धवः} %||1-64-29||

\twolineshloka
{विश्वामित्रोऽपि धर्मात्मा सहरामः सलक्ष्मणः}
{स्वं वाटमभिचक्राम पूज्यमानो महर्षिभिः} %||1-65-30||


{॥इत्यार्षे श्रीमद्रामायणे वाल्मीकीये आदिकाव्ये बालकाण्डे चतुःषष्टितमः सर्गः॥}

\sect{पञ्चषष्टितमः सर्गः}

\twolineshloka
{ततः प्रभाते विमले कृतकर्मा नराधिपः}
{विश्वामित्रं महात्मानमाजुहाव सराघवम्} %||1-65-1||

\twolineshloka
{तमर्चयित्वा धर्मात्मा शास्त्रदृष्ट्तेन कर्मणा}
{राघवौ च महात्मानौ तदा वाक्यमुवाच ह} %||1-65-2||

\twolineshloka
{भगवन्स्वागतं तेऽस्तु किं करोमि तवानघ}
{भवानाज्ञापयतु मामाज्ञाप्यो भवता ह्यहम्} %||1-65-3||

\twolineshloka
{एवमुक्तः स धर्मात्मा जनकेन महात्मना}
{प्रत्युवाच मुनिर्वीरं वाक्यं वाक्यविशारदः} %||1-65-4||

\twolineshloka
{पुत्रौ दशरथस्येमौ क्षत्रियौ लोकविश्रुतौ}
{द्रष्टुकामौ धनुः श्रेष्ठं यदेतत्त्वयि तिष्ठति} %||1-65-5||

\twolineshloka
{एतद्दर्शय भद्रं ते कृतकामौ नृपात्मजौ}
{दर्शनादस्य धनुषो यथेष्टं प्रतियास्यतः} %||1-65-6||

\twolineshloka
{एवमुक्तस्तु जनकः प्रत्युवाच महामुनिम्}
{श्रूयतामस्य धनुषो यदर्थमिह तिष्ठति} %||1-65-7||

\twolineshloka
{देवरात इति ख्यातो निमेः षष्ठो महीपतिः}
{न्यासोऽयं तस्य भगवन्हस्ते दत्तो महात्मना} %||1-65-8||

\twolineshloka
{दक्षयज्ञवधे पूर्वं धनुरायम्य वीर्यवान्}
{रुद्रस्तु त्रिदशान्रोषात्सलीलमिदमब्रवीत्} %||1-65-9||

\twolineshloka
{यस्माद्भागार्थिनो भागान्नाकल्पयत मे सुराः}
{वराङ्गानि महार्हाणि धनुषा शातयामि वः} %||1-65-10||

\twolineshloka
{ततो विमनसः सर्वे देवा वै मुनिपुङ्गव}
{प्रसादयन्ति देवेशं तेषां प्रीतोऽभवद्भवः} %||1-65-11||

{प्रीतियुक्तः स सर्वेषां ददौ तेषां महात्मनाम्॥\devanumber{12}॥} %||1-65-12|| (Check)

\addtocounter{shlokacount}{1}

\twolineshloka
{तदेतद्देवदेवस्य धनूरत्नं महात्मनः}
{न्यासभूतं तदा न्यस्तमस्माकं पूर्वके विभो} %||1-65-13||

\twolineshloka
{अथ मे कृषतः क्षेत्रं लाङ्गलादुत्थिता मम}
{क्षेत्रं शोधयता लब्ध्वा नाम्ना सीतेति विश्रुता} %||1-65-14||

\twolineshloka
{भूतलादुत्थिता सा तु व्यवर्धत ममात्मजा}
{वीर्यशुल्केति मे कन्या स्थापितेयमयोनिजा} %||1-65-15||

\twolineshloka
{भूतलादुत्थितां तां तु वर्धमानां ममात्मजाम्}
{वरयामासुरागम्य राजानो मुनिपुङ्गव} %||1-65-16||

\twolineshloka
{तेषां वरयतां कन्यां सर्वेषां पृथिवीक्षिताम्}
{वीर्यशुल्केति भगवन्न ददामि सुतामहम्} %||1-65-17||

\twolineshloka
{ततः सर्वे नृपतयः समेत्य मुनिपुङ्गव}
{मिथिलामभ्युपागम्य वीर्यं जिज्ञासवस्तदा} %||1-65-18||

\twolineshloka
{तेषां जिज्ञासमानानां वीर्यं धनुरुपाहृतम्}
{न शेकुर्ग्रहणे तस्य धनुषस्तोलनेऽपि वा} %||1-65-19||

\twolineshloka
{तेषां वीर्यवतां वीर्यमल्पं ज्ञात्वा महामुने}
{प्रत्याख्याता नृपतयस्तन्निबोध तपोधन} %||1-65-20||

\twolineshloka
{ततः परमकोपेन राजानो मुनिपुङ्गव}
{अरुन्धन्मिथिलां सर्वे वीर्यसन्देहमागताः} %||1-65-21||

\twolineshloka
{आत्मानमवधूतं ते विज्ञाय मुनिपुङ्गव}
{रोषेण महताविष्टाः पीडयन्मिथिलां पुरीम्} %||1-65-22||

\twolineshloka
{ततः संवत्सरे पूर्णे क्षयं यातानि सर्वशः}
{साधनानि मुनिरेष्ठ ततोऽहं भृशदुःखितः} %||1-65-23||

\twolineshloka
{ततो देवगणान्सर्वांस्तपसाहं प्रसादयम्}
{ददुश्च परमप्रीताश्चतुरङ्गबलं सुराः} %||1-65-24||

\twolineshloka
{ततो भग्ना नृपतयो हन्यमाना दिशो ययुः}
{अवीर्या वीर्यसन्दिग्धा सामात्याः पापकारिणः} %||1-65-25||

\twolineshloka
{तदेतन्मुनिशार्दूल धनुः परमभास्वरम्}
{रामलक्ष्मणयोश्चापि दर्शयिष्यामि सुव्रत} %||1-65-26||

\twolineshloka
{यद्यस्य धनुषो रामः कुर्यादारोपणं मुने}
{सुतामयोनिजां सीतां दद्यां दाशरथेरहम्} %||1-66-27||


{॥इत्यार्षे श्रीमद्रामायणे वाल्मीकीये आदिकाव्ये बालकाण्डे पञ्चषष्टितमः सर्गः॥}

\sect{षट्षष्टितमः सर्गः}

\twolineshloka
{जनकस्य वचः श्रुत्वा विश्वामित्रो महामुनिः}
{धनुर्दर्शय रामाय इति होवाच पार्थिवम्} %||1-66-1||

\twolineshloka
{ततः स राजा जनकः सचिवान्व्यादिदेश ह}
{धनुरानीयतां दिव्यं गन्धमाल्यविभूषितम्} %||1-66-2||

\twolineshloka
{जनकेन समादिष्ठाः सचिवाः प्राविशन्पुरीम्}
{तद्धनुः पुरतः कृत्वा निर्जग्मुः पार्थिवाज्ञया} %||1-66-3||

\twolineshloka
{नृपां शतानि पञ्चाशद्व्यायतानां महात्मनाम्}
{मञ्जूषामष्टचक्रां तां समूहुस्ते कथञ्चन} %||1-66-4||

\twolineshloka
{तामादाय तु मञ्जूषामायतीं यत्र तद्धनुः}
{सुरोपमं ते जनकमूचुर्नृपतिमन्त्रिणः} %||1-66-5||

\twolineshloka
{इदं धनुर्वरं राजन्पूजितं सर्वराजभिः}
{मिथिलाधिप राजेन्द्र दर्शनीयं यदीच्छसि} %||1-66-6||

\twolineshloka
{तेषां नृपो वचः श्रुत्वा कृताञ्जलिरभाषत}
{विश्वामित्रं महात्मानं तौ चोभौ रामलक्ष्मणौ} %||1-66-7||

\twolineshloka
{इदं धनुर्वरं ब्रह्मञ्जनकैरभिपूजितम्}
{राजभिश्च महावीर्यैरशक्यं पूरितुं तदा} %||1-66-8||

\twolineshloka
{नैतत्सुरगणाः सर्वे नासुरा न च राक्षसाः}
{गन्धर्वयक्षप्रवराः सकिंनरमहोरगाः} %||1-66-9||

\twolineshloka
{क्व गतिर्मानुषाणां च धनुषोऽस्य प्रपूरणे}
{आरोपणे समायोगे वेपने तोलनेऽपि वा} %||1-66-10||

\twolineshloka
{तदेतद्धनुषां श्रेष्ठमानीतं मुनिपुङ्गव}
{दर्शयैतन्महाभाग अनयो राजपुत्रयोः} %||1-66-11||

\twolineshloka
{विश्वामित्रस्तु धर्मात्मा श्रुत्वा जनकभाषितम्}
{वत्स राम धनुः पश्य इति राघवमब्रवीत्} %||1-66-12||

\twolineshloka
{महर्षेर्वचनाद्रामो यत्र तिष्ठति तद्धनुः}
{मञ्जूषां तामपावृत्य दृष्ट्वा धनुरथाब्रवीत्} %||1-66-13||

\twolineshloka
{इदं धनुर्वरं ब्रह्मन्संस्पृशामीह पाणिना}
{यत्नवांश्च भविष्यामि तोलने पूरणेऽपि वा} %||1-66-14||

\twolineshloka
{बाढमित्येव तं राजा मुनिश्च समभाषत}
{लीलया स धनुर्मध्ये जग्राह वचनान्मुनेः} %||1-66-15||

\twolineshloka
{पश्यतां नृषहस्राणां बहूनां रघुनन्दनः}
{आरोपयत्स धर्मात्मा सलीलमिव तद्धनुः} %||1-66-16||

\twolineshloka
{आरोपयित्वा मौर्वीं च पूरयामास वीर्यवान्}
{तद्बभञ्ज धनुर्मध्ये नरश्रेष्ठो महायशाः} %||1-66-17||

\twolineshloka
{तस्य शब्दो महानासीन्निर्घातसमनिःस्वनः}
{भूमिकम्पश्च सुमहान्पर्वतस्येव दीर्यतः} %||1-66-18||

\twolineshloka
{निपेतुश्च नराः सर्वे तेन शब्देन मोहिताः}
{वर्जयित्वा मुनिवरं राजानं तौ च राघवौ} %||1-66-19||

\twolineshloka
{प्रत्याश्वस्ते जने तस्मिन्राजा विगतसाध्वसः}
{उवाच प्राञ्जलिर्वाक्यं वाक्यज्ञो मुनिपुङ्गवम्} %||1-66-20||

\twolineshloka
{भगवन्दृष्टवीर्यो मे रामो दशरथात्मजः}
{अत्यद्भुतमचिन्त्यं च अतर्कितमिदं मया} %||1-66-21||

\twolineshloka
{जनकानां कुले कीर्तिमाहरिष्यति मे सुता}
{सीता भर्तारमासाद्य रामं दशरथात्मजम्} %||1-66-22||

\twolineshloka
{मम सत्या प्रतिज्ञा च वीर्यशुल्केति कौशिक}
{सीता प्राणैर्बहुमता देया रामाय मे सुता} %||1-66-23||

\twolineshloka
{भवतोऽनुमते ब्रह्मञ्शीघ्रं गच्छन्तु मन्त्रिणः}
{मम कौशिक भद्रं ते अयोध्यां त्वरिता रथैः} %||1-66-24||

\twolineshloka
{राजानं प्रश्रितैर्वाक्यैरानयन्तु पुरं मम}
{प्रदानं वीर्यशुल्कायाः कथयन्तु च सर्वशः} %||1-66-25||

\twolineshloka
{मुनिगुप्तौ च काकुत्स्थौ कथयन्तु नृपाय वै}
{प्रीयमाणं तु राजानमानयन्तु सुशीघ्रगाः} %||1-66-26||

\twolineshloka
{कौशिकश्च तथेत्याह राजा चाभाष्य मन्त्रिणः}
{अयोध्यां प्रेषयामास धर्मात्मा कृतशासनात्} %||1-67-27||


{॥इत्यार्षे श्रीमद्रामायणे वाल्मीकीये आदिकाव्ये बालकाण्डे षट्षष्टितमः सर्गः॥}

\sect{सप्तषष्टितमः सर्गः}

\twolineshloka
{जनकेन समादिष्टा दूतास्ते क्लान्तवाहनाः}
{त्रिरात्रमुषित्वा मार्गे तेऽयोध्यां प्राविशन्पुरीम्} %||1-67-1||

\twolineshloka
{ते राजवचनाद्दूता राजवेश्मप्रवेशिताः}
{ददृशुर्देवसङ्काशं वृद्धं दशरथं नृपम्} %||1-67-2||

\twolineshloka
{बद्धाञ्जलिपुटाः सर्वे दूता विगतसाध्वसाः}
{राजानं प्रयता वाक्यमब्रुवन्मधुराक्षरम्} %||1-67-3||

\twolineshloka
{मैथिलो जनको राजा साग्निहोत्रपुरस्कृतः}
{कुशलं चाव्ययं चैव सोपाध्यायपुरोहितम्} %||1-67-4||

\twolineshloka
{मुहुर्मुहुर्मधुरया स्नेहसंयुक्तया गिरा}
{जनकस्त्वां महाराज पृच्छते सपुरःसरम्} %||1-67-5||

\twolineshloka
{पृष्ट्वा कुशलमव्यग्रं वैदेहो मिथिलाधिपः}
{कौशिकानुमते वाक्यं भवन्तमिदमब्रवीत्} %||1-67-6||

\twolineshloka
{पूर्वं प्रतिज्ञा विदिता वीर्यशुल्का ममात्मजा}
{राजानश्च कृतामर्षा निर्वीर्या विमुखीकृताः} %||1-67-7||

\twolineshloka
{सेयं मम सुता राजन्विश्वामित्र पुरःसरैः}
{यदृच्छयागतैर्वीरैर्निर्जिता तव पुत्रकैः} %||1-67-8||

\twolineshloka
{तच्च राजन्धनुर्दिव्यं मध्ये भग्नं महात्मना}
{रामेण हि महाराज महत्यां जनसंसदि} %||1-67-9||

\twolineshloka
{अस्मै देया मया सीता वीर्यशुल्का महात्मने}
{प्रतिज्ञां तर्तुमिच्छामि तदनुज्ञातुमर्हसि} %||1-67-10||

\twolineshloka
{सोपाध्यायो महाराज पुरोहितपुरस्कृतः}
{शीघ्रमागच्छ भद्रं ते द्रष्टुमर्हसि राघवौ} %||1-67-11||

\twolineshloka
{प्रीतिं च मम राजेन्द्र निर्वर्तयितुमर्हसि}
{पुत्रयोरुभयोरेव प्रीतिं त्वमपि लप्स्यसे} %||1-67-12||

\twolineshloka
{एवं विदेहाधिपतिर्मधुरं वाक्यमब्रवीत्}
{विश्वामित्राभ्यनुज्ञातः शतानन्दमते स्थितः} %||1-67-13||

\twolineshloka
{दूतवाक्यं तु तच्छ्रुत्वा राजा परमहर्षितः}
{वसिष्ठं वामदेवं च मन्त्रिणोऽन्यांश्च सोऽब्रवीत्} %||1-67-14||

\twolineshloka
{गुप्तः कुशिकपुत्रेण कौसल्यानन्दवर्धनः}
{लक्ष्मणेन सह भ्रात्रा विदेहेषु वसत्यसौ} %||1-67-15||

\twolineshloka
{दृष्टवीर्यस्तु काकुत्स्थो जनकेन महात्मना}
{सम्प्रदानं सुतायास्तु राघवे कर्तुमिच्छति} %||1-67-16||

\twolineshloka
{यदि वो रोचते वृत्तं जनकस्य महात्मनः}
{पुरीं गच्छामहे शीघ्रं मा भूत्कालस्य पर्ययः} %||1-67-17||

\twolineshloka
{मन्त्रिणो बाढमित्याहुः सह सर्वैर्महर्षिभिः}
{सुप्रीतश्चाब्रवीद्राजा श्वो यात्रेति स मन्त्रिणः} %||1-67-18||

\twolineshloka
{मन्त्रिणस्तु नरेन्द्रस्य रात्रिं परमसत्कृताः}
{ऊषुः प्रमुदिताः सर्वे गुणैः सर्वैः समन्विताः} %||1-68-19||


{॥इत्यार्षे श्रीमद्रामायणे वाल्मीकीये आदिकाव्ये बालकाण्डे सप्तषष्टितमः सर्गः॥}

\sect{अष्टषष्टितमः सर्गः}

\twolineshloka
{ततो रात्र्यां व्यतीतायां सोपाध्यायः सबान्धवः}
{राजा दशरथो हृष्टः सुमन्त्रमिदमब्रवीत्} %||1-68-1||

\twolineshloka
{अद्य सर्वे धनाध्यक्षा धनमादाय पुष्कलम्}
{व्रजन्त्वग्रे सुविहिता नानारत्नसमन्विताः} %||1-68-2||

\twolineshloka
{चतुरङ्गबलं चापि शीघ्रं निर्यातु सर्वशः}
{ममाज्ञासमकालं च यानयुग्यमनुत्तमम्} %||1-68-3||

\twolineshloka
{वसिष्ठो वामदेवश्च जाबालिरथ काश्यपः}
{मार्कण्डेयश्च दीर्घायुरृषिः कात्यायनस्तथा} %||1-68-4||

\twolineshloka
{एते द्विजाः प्रयान्त्वग्रे स्यन्दनं योजयस्व मे}
{यथा कालात्ययो न स्याद्दूता हि त्वरयन्ति माम्} %||1-68-5||

\twolineshloka
{वचनाच्च नरेन्द्रस्य सा सेना चतुरङ्गिणी}
{राजानमृषिभिः सार्धं व्रजन्तं पृष्ठतोऽन्वगात्} %||1-68-6||

\twolineshloka
{गत्वा चतुरहं मार्गं विदेहानभ्युपेयिवान्}
{राजा तु जनकः श्रीमाञ्श्रुत्वा पूजामकल्पयत्} %||1-68-7||

\threelineshloka
{ततो राजानमासाद्य वृद्धं दशरथं नृपम्}
{जनको मुदितो राजा हर्षं च परमं ययौ}
{उवाच न नरश्रेष्ठो नरश्रेष्ठं मुदान्वितम्} %||1-68-8||

\twolineshloka
{स्वागतं ते महाराज दिष्ट्या प्राप्तोऽसि राघव}
{पुत्रयोरुभयोः प्रीतिं लप्स्यसे वीर्यनिर्जिताम्} %||1-68-9||

\twolineshloka
{दिष्ट्या प्राप्तो महातेजा वसिष्ठो भगवानृषिः}
{सह सर्वैर्द्विजश्रेष्ठैर्देवैरिव शतक्रतुः} %||1-68-10||

\twolineshloka
{दिष्ट्या मे निर्जिता विघ्ना दिष्ट्या मे पूजितं कुलम्}
{राघवैः सह सम्बन्धाद्वीर्यश्रेष्ठैर्महात्मभिः} %||1-68-11||

\twolineshloka
{श्वः प्रभाते नरेन्द्रेन्द्र निर्वर्तयितुमर्हसि}
{यज्ञस्यान्ते नरश्रेष्ठ विवाहमृषिसम्मतम्} %||1-68-12||

\twolineshloka
{तस्य तद्वचनं श्रुत्वा ऋषिमध्ये नराधिपः}
{वाक्यं वाक्यविदां श्रेष्ठः प्रत्युवाच महीपतिम्} %||1-68-13||

\twolineshloka
{प्रतिग्रहो दातृवशः श्रुतमेतन्मया पुरा}
{यथा वक्ष्यसि धर्मज्ञ तत्करिष्यामहे वयम्} %||1-68-14||

\twolineshloka
{तद्धर्मिष्ठं यशस्यं च वचनं सत्यवादिनः}
{श्रुत्वा विदेहाधिपतिः परं विस्मयमागतः} %||1-68-15||

\twolineshloka
{ततः सर्वे मुनिगणाः परस्परसमागमे}
{हर्षेण महता युक्तास्तां निशामवसन्सुखम्} %||1-68-16||

\twolineshloka
{राजा च राघवौ पुत्रौ निशाम्य परिहर्षितः}
{उवास परमप्रीतो जनकेन सुपूजितः} %||1-68-17||

\twolineshloka
{जनकोऽपि महातेजाः क्रिया धर्मेण तत्त्ववित्}
{यज्ञस्य च सुताभ्यां च कृत्वा रात्रिमुवास ह} %||1-69-18||


{॥इत्यार्षे श्रीमद्रामायणे वाल्मीकीये आदिकाव्ये बालकाण्डे अष्टषष्टितमः सर्गः॥}

\sect{एकोनसप्ततितमः सर्गः}

\twolineshloka
{ततः प्रभाते जनकः कृतकर्मा महर्षिभिः}
{उवाच वाक्यं वाक्यज्ञः शतानन्दं पुरोहितम्} %||1-69-1||

\twolineshloka
{भ्राता मम महातेजा यवीयानतिधार्मिकः}
{कुशध्वज इति ख्यातः पुरीमध्यवसच्छुभाम्} %||1-69-2||

\twolineshloka
{वार्याफलकपर्यन्तां पिबन्निक्षुमतीं नदीम्}
{साङ्काश्यां पुण्यसङ्काशां विमानमिव पुष्पकम्} %||1-69-3||

\twolineshloka
{तमहं द्रष्टुमिच्छामि यज्ञगोप्ता स मे मतः}
{प्रीतिं सोऽपि महातेजा इम्मां भोक्ता मया सह} %||1-69-4||

\twolineshloka
{शासनात्तु नरेन्द्रस्य प्रययुः शीघ्रवाजिभिः}
{समानेतुं नरव्याघ्रं विष्णुमिन्द्राज्ञया यथा} %||1-69-5||

{आज्ञया तु नरेन्द्रस्य आजगाम कुशध्वजः॥\devanumber{6}॥} %||1-69-6|| (Check)

\addtocounter{shlokacount}{1}

\twolineshloka
{स ददर्श महात्मानं जनकं धर्मवत्सलम्}
{सोऽभिवाद्य शतानन्दं राजानं चापि धार्मिकम्} %||1-69-7||

\twolineshloka
{राजार्हं परमं दिव्यमासनं चाध्यरोहत}
{उपविष्टावुभौ तौ तु भ्रातरावमितौजसौ} %||1-69-8||

\threelineshloka
{प्रेषयामासतुर्वीरौ मन्त्रिश्रेष्ठं सुदामनम्}
{गच्छ मन्त्रिपते शीघ्रमैक्ष्वाकममितप्रभम्}
{आत्मजैः सह दुर्धर्षमानयस्व समन्त्रिणम्} %||1-69-9||

\twolineshloka
{औपकार्यां स गत्वा तु रघूणां कुलवर्धनम्}
{ददर्श शिरसा चैनमभिवाद्येदमब्रवीत्} %||1-69-10||

\twolineshloka
{अयोध्याधिपते वीर वैदेहो मिथिलाधिपः}
{स त्वां द्रष्टुं व्यवसितः सोपाध्यायपुरोहितम्} %||1-69-11||

\twolineshloka
{मन्त्रिश्रेष्ठवचः श्रुत्वा राजा सर्षिगणस्तदा}
{सबन्धुरगमत्तत्र जनको यत्र वर्तते} %||1-69-12||

\twolineshloka
{स राजा मन्त्रिसहितः सोपाध्यायः सबान्धवः}
{वाक्यं वाक्यविदां श्रेष्ठो वैदेहमिदमब्रवीत्} %||1-69-13||

\twolineshloka
{विदितं ते महाराज इक्ष्वाकुकुलदैवतम्}
{वक्ता सर्वेषु कृत्येषु वसिष्ठो भगवानृषिः} %||1-69-14||

\twolineshloka
{विश्वामित्राभ्यनुज्ञातः सह सर्वैर्महर्षिभिः}
{एष वक्ष्यति धर्मात्मा वसिष्ठो मे यथाक्रमम्} %||1-69-15||

\twolineshloka
{तूष्णीम्भूते दशरथे वसिष्ठो भगवानृषिः}
{उवाच वाक्यं वाक्यज्ञो वैदेहं सपुरोहितम्} %||1-69-16||

\twolineshloka
{अव्यक्तप्रभवो ब्रह्मा शाश्वतो नित्य अव्ययः}
{तस्मान्मरीचिः संजज्ञे मरीचेः कश्यपः सुतः} %||1-69-17||

\twolineshloka
{विवस्वान्कश्यपाज्जज्ञे मनुर्वैवैस्वतः स्मृतः}
{मनुः प्रजापतिः पूर्वमिक्ष्वाकुस्तु मनोः सुतः} %||1-69-18||

\twolineshloka
{तमिक्ष्वाकुमयोध्यायां राजानं विद्धि पूर्वकम्}
{इक्ष्वाकोस्तु सुतः श्रीमान्विकुक्षिरुदपद्यत} %||1-69-19||

\twolineshloka
{विकुक्षेस्तु महातेजा बाणः पुत्रः प्रतापवान्}
{बाणस्य तु महातेजा अनरण्यः प्रतापवान्} %||1-69-20||

\twolineshloka
{अनरण्यात्पृथुर्जज्ञे त्रिशङ्कुस्तु पृथोः सुतः}
{त्रिशङ्कोरभवत्पुत्रो धुन्धुमारो महायशाः} %||1-69-21||

\twolineshloka
{धुन्धुमारान्महातेजा युवनाश्वो महारथः}
{युवनाश्वसुतः श्रीमान्मान्धाता पृथिवीपतिः} %||1-69-22||

\twolineshloka
{मान्धातुस्तु सुतः श्रीमान्सुसन्धिरुदपद्यत}
{सुसन्धेरपि पुत्रौ द्वौ ध्रुवसन्धिः प्रसेनजित्} %||1-69-23||

\twolineshloka
{यशस्वी ध्रुवसन्धेस्तु भरतो नाम नामतः}
{भरतात्तु महातेजा असितो नाम जायत} %||1-69-24||

\twolineshloka
{सह तेन गरेणैव जातः स सगरोऽभवत्}
{सगरस्यासमञ्जस्तु असमञ्जादथांशुमान्} %||1-69-25||

\twolineshloka
{दिलीपोंऽशुमतः पुत्रो दिलीपस्य भगीरथः}
{भगीरथात्ककुत्स्थश्च ककुत्स्थस्य रघुस्तथा} %||1-69-26||

\twolineshloka
{रघोस्तु पुत्रस्तेजस्वी प्रवृद्धः पुरुषादकः}
{कल्माषपादो ह्यभवत्तस्माज्जातस्तु शङ्खणः} %||1-69-27||

\twolineshloka
{सुदर्शनः शङ्खणस्य अग्निवर्णः सुदर्शनात्}
{शीघ्रगस्त्वग्निवर्णस्य शीघ्रगस्य मरुः सुतः} %||1-69-28||

\twolineshloka
{मरोः प्रशुश्रुकस्त्वासीदम्बरीषः प्रशुश्रुकात्}
{अम्बरीषस्य पुत्रोऽभून्नहुषः पृथिवीपतिः} %||1-69-29||

\threelineshloka
{नहुषस्य ययातिस्तु नाभागस्तु ययातिजः}
{नाभागस्य भभूवाज अजाद्दशरथोऽभवत्}
{तस्माद्दशरथाज्जातौ भ्रातरौ रामलक्ष्मणौ} %||1-69-30||

\twolineshloka
{आदिवंशविशुद्धानां राज्ञां परमधर्मिणाम्}
{इक्ष्वाकुकुलजातानां वीराणां सत्यवादिनाम्} %||1-69-31||

\twolineshloka
{रामलक्ष्मणयोरर्थे त्वत्सुते वरये नृप}
{सदृशाभ्यां नरश्रेष्ठ सदृशे दातुमर्हसि} %||1-70-32||


{॥इत्यार्षे श्रीमद्रामायणे वाल्मीकीये आदिकाव्ये बालकाण्डे एकोनसप्ततितमः सर्गः॥}

\sect{सप्ततितमः सर्गः}

\twolineshloka
{एवं ब्रुवाणं जनकः प्रत्युवाच कृताञ्जलिः}
{श्रोतुमर्हसि भद्रं ते कुलं नः कीर्तितं परम्} %||1-70-1||

\twolineshloka
{प्रदाने हि मुनिश्रेष्ठ कुलं निरवशेषतः}
{वक्तव्यं कुलजातेन तन्निबोध महामुने} %||1-70-2||

\twolineshloka
{राजाभूत्त्रिषु लोकेषु विश्रुतः स्वेन कर्मणा}
{निमिः परमधर्मात्मा सर्वसत्त्ववतां वरः} %||1-70-3||

\twolineshloka
{तस्य पुत्रो मिथिर्नाम जनको मिथि पुत्रकः}
{प्रथमो जनको नाम जनकादप्युदावसुः} %||1-70-4||

\twolineshloka
{उदावसोस्तु धर्मात्मा जातो वै नन्दिवर्धनः}
{नन्दिवर्धन पुत्रस्तु सुकेतुर्नाम नामतः} %||1-70-5||

\twolineshloka
{सुकेतोरपि धर्मात्मा देवरातो महाबलः}
{देवरातस्य राजर्षेर्बृहद्रथ इति श्रुतः} %||1-70-6||

\twolineshloka
{बृहद्रथस्य शूरोऽभून्महावीरः प्रतापवान्}
{महावीरस्य धृतिमान्सुधृतिः सत्यविक्रमः} %||1-70-7||

\twolineshloka
{सुधृतेरपि धर्मात्मा धृष्टकेतुः सुधार्मिकः}
{धृष्टकेतोस्तु राजर्षेर्हर्यश्व इति विश्रुतः} %||1-70-8||

\twolineshloka
{हर्यश्वस्य मरुः पुत्रो मरोः पुत्रः प्रतीन्धकः}
{प्रतीन्धकस्य धर्मात्मा राजा कीर्तिरथः सुतः} %||1-70-9||

\twolineshloka
{पुत्रः कीर्तिरथस्यापि देवमीढ इति स्मृतः}
{देवमीढस्य विबुधो विबुधस्य महीध्रकः} %||1-70-10||

\twolineshloka
{महीध्रकसुतो राजा कीर्तिरातो महाबलः}
{कीर्तिरातस्य राजर्षेर्महारोमा व्यजायत} %||1-70-11||

\twolineshloka
{महारोंणस्तु धर्मात्मा स्वर्णरोमा व्यजायत}
{स्वर्णरोंणस्तु राजर्षेर्ह्रस्वरोमा व्यजायत} %||1-70-12||

\twolineshloka
{तस्य पुत्रद्वयं जज्ञे धर्मज्ञस्य महात्मनः}
{ज्येष्ठोऽहमनुजो भ्राता मम वीरः कुशध्वजः} %||1-70-13||

\twolineshloka
{मां तु ज्येष्ठं पिता राज्ये सोऽभिषिच्य नराधिपः}
{कुशध्वजं समावेश्य भारं मयि वनं गतः} %||1-70-14||

\twolineshloka
{वृद्धे पितरि स्वर्याते धर्मेण धुरमावहम्}
{भ्रातरं देवसङ्काशं स्नेहात्पश्यन्कुशध्वजम्} %||1-70-15||

\twolineshloka
{कस्यचित्त्वथ कालस्य साङ्काश्यादगमत्पुरात्}
{सुधन्वा वीर्यवान्राजा मिथिलामवरोधकः} %||1-70-16||

\twolineshloka
{स च मे प्रेषयामास शैवं धनुरनुत्तमम्}
{सीता कन्या च पद्माक्षी मह्यं वै दीयतामिति} %||1-70-17||

\twolineshloka
{तस्याप्रदानाद्ब्रह्मर्षे युद्धमासीन्मया सह}
{स हतोऽभिमुखो राजा सुधन्वा तु मया रणे} %||1-70-18||

\twolineshloka
{निहत्य तं मुनिश्रेष्ठ सुधन्वानं नराधिपम्}
{साङ्काश्ये भ्रातरं शूरमभ्यषिञ्चं कुशध्वजम्} %||1-70-19||

\twolineshloka
{कनीयानेष मे भ्राता अहं ज्येष्ठो महामुने}
{ददामि परमप्रीतो वध्वौ ते मुनिपुङ्गव} %||1-70-20||

\twolineshloka
{सीतां रामाय भद्रं ते ऊर्मिलां लक्ष्मणाय च}
{वीर्यशुल्कां मम सुतां सीतां सुरसुतोपमाम्} %||1-70-21||

\twolineshloka
{द्वितीयामूर्मिलां चैव त्रिर्वदामि न संशयः}
{ददामि परमप्रीतो वध्वौ ते रघुनन्दन} %||1-70-22||

\twolineshloka
{रामलक्ष्मणयो राजन्गोदानं कारयस्व ह}
{पितृकार्यं च भद्रं ते ततो वैवाहिकं कुरु} %||1-70-23||

\threelineshloka
{मघा ह्यद्य महाबाहो तृतीये दिवसे प्रभो}
{फल्गुन्यामुत्तरे राजंस्तस्मिन्वैवाहिकं कुरु}
{रामलक्ष्मणयोरर्थे दानं कार्यं सुखोदयम्} %||1-71-24||


{॥इत्यार्षे श्रीमद्रामायणे वाल्मीकीये आदिकाव्ये बालकाण्डे सप्ततितमः सर्गः॥}

\sect{एकसप्ततितमः सर्गः}

\twolineshloka
{तमुक्तवन्तं वैदेहं विश्वामित्रो महामुनिः}
{उवाच वचनं वीरं वसिष्ठसहितो नृपम्} %||1-71-1||

\twolineshloka
{अचिन्त्यान्यप्रमेयानि कुलानि नरपुङ्गव}
{इक्ष्वाकूणां विदेहानां नैषां तुल्योऽस्ति कश्चन} %||1-71-2||

\twolineshloka
{सदृशो धर्मसम्बन्धः सदृशो रूपसम्पदा}
{रामलक्ष्मणयो राजन्सीता चोर्मिलया सह} %||1-71-3||

{वक्तव्यं न नरश्रेष्ठ श्रूयतां वचनं मम॥\devanumber{4}॥} %||1-71-4|| (Check)

\addtocounter{shlokacount}{1}

\threelineshloka
{भ्राता यवीयान्धर्मज्ञ एष राजा कुशध्वजः}
{अस्य धर्मात्मनो राजन्रूपेणाप्रतिमं भुवि}
{सुता द्वयं नरश्रेष्ठ पत्न्यर्थं वरयामहे} %||1-71-5||

\twolineshloka
{भरतस्य कुमारस्य शत्रुघ्नस्य च धीमतः}
{वरयेम सुते राजंस्तयोरर्थे महात्मनोः} %||1-71-6||

\twolineshloka
{पुत्रा दशरथस्येमे रूपयौवनशालिनः}
{लोकपालोपमाः सर्वे देवतुल्यपराक्रमाः} %||1-71-7||

\twolineshloka
{उभयोरपि राजेन्द्र सम्बन्धेनानुबध्यताम्}
{इक्ष्वाकुकुलमव्यग्रं भवतः पुण्यकर्मणः} %||1-71-8||

\twolineshloka
{विश्वामित्रवचः श्रुत्वा वसिष्ठस्य मते तदा}
{जनकः प्राञ्जलिर्वाक्यमुवाच मुनिपुङ्गवौ} %||1-71-9||

\threelineshloka
{सदृशं कुलसम्बन्धं यदाज्ञापयथः स्वयम्}
{एवं भवतु भद्रं वः कुशध्वजसुते इमे}
{पत्न्यौ भजेतां सहितौ शत्रुघ्नभरतावुभौ} %||1-71-10||

\twolineshloka
{एकाह्ना राजपुत्रीणां चतसॄणां महामुने}
{पाणीन्गृह्णन्तु चत्वारो राजपुत्रा महाबलाः} %||1-71-11||

\twolineshloka
{उत्तरे दिवसे ब्रह्मन्फल्गुनीभ्यां मनीषिणः}
{वैवाहिकं प्रशंसन्ति भगो यत्र प्रजापतिः} %||1-71-12||

\twolineshloka
{एवमुक्त्वा वचः सौम्यं प्रत्युत्थाय कृताञ्जलिः}
{उभौ मुनिवरौ राजा जनको वाक्यमब्रवीत्} %||1-71-13||

\twolineshloka
{परो धर्मः कृतो मह्यं शिष्योऽस्मि भवतोः सदा}
{इमान्यासनमुख्यानि आसेतां मुनिपुङ्गवौ} %||1-71-14||

\twolineshloka
{यथा दशरथस्येयं तथायोध्या पुरी मम}
{प्रभुत्वे नासित्सन्देहो यथार्हं कर्तुमर्हथः} %||1-71-15||

\twolineshloka
{तथा ब्रुवति वैदेहे जनके रघुनन्दनः}
{राजा दशरथो हृष्टः प्रत्युवाच महीपतिम्} %||1-71-16||

\twolineshloka
{युवामसङ्ख्येय गुणौ भ्रातरौ मिथिलेश्वरौ}
{ऋषयो राजसङ्घाश्च भवद्भ्यामभिपूजिताः} %||1-71-17||

\twolineshloka
{स्वस्ति प्राप्नुहि भद्रं ते गमिष्यामि स्वमालयम्}
{श्राद्धकर्माणि सर्वाणि विधास्य इति चाब्रवीत्} %||1-71-18||

\twolineshloka
{तमापृष्ट्वा नरपतिं राजा दशरथस्तदा}
{मुनीन्द्रौ तौ पुरस्कृत्य जगामाशु महायशाः} %||1-71-19||

\twolineshloka
{स गत्वा निलयं राजा श्राद्धं कृत्वा विधानतः}
{प्रभाते काल्यमुत्थाय चक्रे गोदानमुत्तमम्} %||1-71-20||

\twolineshloka
{गवां शतसहस्राणि ब्राह्मणेभ्यो नराधिपः}
{एकैकशो ददौ राजा पुत्रानुद्धिश्य धर्मतः} %||1-71-21||

\twolineshloka
{सुवर्णशृङ्गाः सम्पन्नाः सवत्साः कांस्यदोहनाः}
{गवां शतसहस्राणि चत्वारि पुरुषर्षभः} %||1-71-22||

\twolineshloka
{वित्तमन्यच्च सुबहु द्विजेभ्यो रघुनन्दनः}
{ददौ गोदानमुद्दिश्य पुत्राणां पुत्रवत्सलः} %||1-71-23||

\twolineshloka
{स सुतैः कृतगोदानैर्वृतश्च नृपतिस्तदा}
{लोकपालैरिवाभाति वृतः सौम्यः प्रजापतिः} %||1-72-24||


{॥इत्यार्षे श्रीमद्रामायणे वाल्मीकीये आदिकाव्ये बालकाण्डे एकसप्ततितमः सर्गः॥}

\sect{द्विसप्ततितमः सर्गः}

\twolineshloka
{यस्मिंस्तु दिवसे राजा चक्रे गोदानमुत्तमम्}
{तस्मिंस्तु दिवसे शूरो युधाजित्समुपेयिवान्} %||1-72-1||

\twolineshloka
{पुत्रः केकयराजस्य साक्षाद्भरतमातुलः}
{दृष्ट्वा पृष्ट्वा च कुशलं राजानमिदमब्रवीत्} %||1-72-2||

\twolineshloka
{केकयाधिपती राजा स्नेहात्कुशलमब्रवीत्}
{येषां कुशलकामोऽसि तेषां सम्प्रत्यनामयम्} %||1-72-3||

\twolineshloka
{स्वस्रीयं मम राजेन्द्र द्रष्टुकामो महीपते}
{तदर्थमुपयातोऽहमयोध्यां रघुनन्दन} %||1-72-4||

\twolineshloka
{श्रुत्वा त्वहमयोध्यायां विवाहार्थं तवात्मजान्}
{मिथिलामुपयातास्तु त्वया सह महीपते} %||1-72-5||

\twolineshloka
{त्वरयाभुपयातोऽहं द्रष्टुकामः स्वसुः सुतम्}
{अथ राजा दशरथः प्रियातिथिमुपस्थिम} %||1-72-6||

\twolineshloka
{दृष्ट्वा परमसत्कारैः पूजार्हं समपूजयत्}
{ततस्तामुषितो रात्रिं सह पुत्रैर्महात्मभिः} %||1-72-7||

\threelineshloka
{ऋषींस्तदा पुरस्कृत्य यज्ञवाटमुपागमत्}
{युक्ते मुहूर्ते विजये सर्वाभरणभूषितैः}
{भ्रातृभिः सहितो रामः कृतकौतुकमङ्गलः} %||1-72-8||

{वसिष्ठं पुरतः कृत्वा महर्षीनपरानपि॥\devanumber{9}॥} %||1-72-9|| (Check)

\addtocounter{shlokacount}{1}

\twolineshloka
{राजा रशरथो राजन्कृतकौतुकमङ्गलैः}
{पुत्रैर्नरवरश्रेष्ठ दातारमभिकाङ्क्षते} %||1-72-10||

\twolineshloka
{दातृप्रतिग्रहीतृभ्यां सर्वार्थाः प्रभवन्ति हि}
{स्वधर्मं प्रतिपद्यस्व कृत्वा वैवाह्यमुत्तमम्} %||1-72-11||

\twolineshloka
{इत्युक्तः परमोदारो वसिष्ठेन महात्मना}
{प्रत्युवाच महातेजा वाक्यं परमधर्मवित्} %||1-72-12||

\twolineshloka
{कः स्थितः प्रतिहारो मे कस्याज्ञा सम्प्रतीक्ष्यते}
{स्वगृहे को विचारोऽस्ति यथा राज्यमिदं तव} %||1-72-13||

\twolineshloka
{कृतकौतुकसर्वस्वा वेदिमूलमुपागताः}
{मम कन्या मुनिश्रेष्ठ दीप्ता वह्नेरिवार्चिषः} %||1-72-14||

\twolineshloka
{सज्जोऽहं त्वत्प्रतीक्षोऽस्मि वेद्यामस्यां प्रतिष्हितः}
{अविघ्नं कुरुतां राजा किमर्थं हि विलम्ब्यते} %||1-72-15||

\twolineshloka
{तद्वाक्यं जनकेनोक्तं श्रुत्वा दशरथस्तदा}
{प्रवेशयामास सुतान्सर्वानृषिगणानपि} %||1-72-16||

\threelineshloka
{अब्रवीज्जनको राजा कौसल्यानन्दवर्धनम्}
{इयं सीता मम सुता सहधर्मचरी तव}
{प्रतीच्छ चैनां भद्रं ते पाणिं गृह्णीष्व पाणिना} %||1-72-17||

\twolineshloka
{लक्ष्मणागच्छ भद्रं ते ऊर्मिलामुद्यतां मया}
{प्रतीच्छ पाणिं गृह्णीष्व मा भूत्कालस्य पर्ययः} %||1-72-18||

\twolineshloka
{तमेवमुक्त्वा जनको भरतं चाभ्यभाषत}
{गृहाण पाणिं माण्डव्याः पाणिना रघुनन्दन} %||1-72-19||

\twolineshloka
{शत्रुघ्नं चापि धर्मात्मा अब्रवीज्जनकेश्वरः}
{श्रुतकीर्त्या महाबाहो पाणिं गृह्णीष्व पाणिना} %||1-72-20||

\twolineshloka
{सर्वे भवन्तः संयाश्च सर्वे सुचरितव्रताः}
{पत्नीभिः सन्तु काकुत्स्था मा भूत्कालस्य पर्ययः} %||1-72-21||

\twolineshloka
{जनकस्य वचः श्रुत्वा पाणीन्पाणिभिरस्पृशन्}
{चत्वारस्ते चतसृणां वसिष्ठस्य मते स्थिताः} %||1-72-22||

\threelineshloka
{अग्निं प्रदक्षिणं कृत्वा वेदिं राजानमेव च}
{ऋषींश्चैव महात्मानः सह भार्या रघूत्तमाः}
{यथोक्तेन तथा चक्रुर्विवाहं विधिपूर्वकम्} %||1-72-23||

\twolineshloka
{पुष्पवृष्टिर्महत्यासीदन्तरिक्षात्सुभास्वरा}
{दिव्यदुन्दुभिनिर्घोषैर्गीतवादित्रनिस्वनैः} %||1-72-24||

\twolineshloka
{ननृतुश्चाप्सरःसङ्घा गन्धर्वाश्च जगुः कलम्}
{विवाहे रघुमुख्यानां तदद्भुतमिवाभवत्} %||1-72-25||

\twolineshloka
{ईदृशे वर्तमाने तु तूर्योद्घुष्टनिनादिते}
{त्रिरग्निं ते परिक्रम्य ऊहुर्भार्या महौजसः} %||1-72-26||

\twolineshloka
{अथोपकार्यां जग्मुस्ते सदारा रघुनन्दनः}
{राजाप्यनुययौ पश्यन्सर्षिसङ्घः सबान्धवः} %||1-73-27||


{॥इत्यार्षे श्रीमद्रामायणे वाल्मीकीये आदिकाव्ये बालकाण्डे द्विसप्ततितमः सर्गः॥}

\sect{त्रिसप्ततितमः सर्गः}

\twolineshloka
{अथ रात्र्यां व्यतीतायां विश्वामित्रो महामुनिः}
{आपृच्छ्य तौ च राजानौ जगामोत्तरपर्वतम्} %||1-73-1||

\twolineshloka
{विश्वामित्रो गते राजा वैदेहं मिथिलाधिपम्}
{आपृच्छ्याथ जगामाशु राजा दशरथः पुरीम्} %||1-73-2||

\twolineshloka
{अथ राजा विदेहानां ददौ कन्याधनं बहु}
{गवां शतसहस्राणि बहूनि मिथिलेश्वरः} %||1-73-3||

\twolineshloka
{कम्बलानां च मुख्यानां क्षौमकोट्यम्बराणि च}
{हस्त्यश्वरथपादातं दिव्यरूपं स्वलङ्कृतम्} %||1-73-4||

\twolineshloka
{ददौ कन्या पिता तासां दासीदासमनुत्तमम्}
{हिरण्यस्य सुवर्णस्य मुक्तानां विद्रुमस्य च} %||1-73-5||

\twolineshloka
{ददौ परमसंहृष्टः कन्याधनमनुत्तमम्}
{दत्त्वा बहुधनं राजा समनुज्ञाप्य पार्थिवम्} %||1-73-6||

\twolineshloka
{प्रविवेश स्वनिलयं मिथिलां मिथिलेश्वरः}
{राजाप्ययोध्याधिपतिः सह पुत्रैर्महात्मभिः} %||1-73-7||

\twolineshloka
{ऋषीन्सर्वान्पुरस्कृत्य जगाम सबलानुगः}
{गच्छन्तं तु नरव्याघ्रं सर्षिसङ्घं सराघवम्} %||1-73-8||

\twolineshloka
{घोराः स्म पक्षिणो वाचो व्याहरन्ति ततस्ततः}
{भौमाश्चैव मृगाः सर्वे गच्छन्ति स्म प्रदक्षिणम्} %||1-73-9||

\threelineshloka
{तान्दृष्ट्वा राजशार्दूलो वसिष्ठं पर्यपृच्छत}
{असौम्याः पक्षिणो घोरा मृगाश्चापि प्रदक्षिणाः}
{किमिदं हृदयोत्कम्पि मनो मम विषीदति} %||1-73-10||

\twolineshloka
{राज्ञो दशरथस्यैतच्छ्रुत्वा वाक्यं महानृषिः}
{उवाच मधुरां वाणीं श्रूयतामस्य यत्फलम्} %||1-73-11||

\twolineshloka
{उपस्थितं भयं घोरं दिव्यं पक्षिमुखाच्च्युतम्}
{मृगाः प्रशमयन्त्येते सन्तापस्त्यज्यतामयम्} %||1-73-12||

\twolineshloka
{तेषां संवदतां तत्र वायुः प्रादुर्बभूव ह}
{कम्पयन्मेदिनीं सर्वां पातयंश्च द्रुमाञ्शुभान्} %||1-73-13||

\twolineshloka
{तमसा संवृतः सूर्यः सर्वा न प्रबभुर्दिशः}
{भस्मना चावृतं सर्वं सम्मूढमिव तद्बलम्} %||1-73-14||

\twolineshloka
{वसिष्ठ ऋषयश्चान्ये राजा च ससुतस्तदा}
{ससंज्ञा इव तत्रासन्सर्वमन्यद्विचेतनम्} %||1-73-15||

\twolineshloka
{तस्मिंस्तमसि घोरे तु भस्मच्छन्नेव सा चमूः}
{ददर्श भीमसङ्काशं जटामण्डलधारिणम्} %||1-73-16||

\twolineshloka
{कैलासमिव दुर्धर्षं कालाग्निमिव दुःसहम्}
{ज्वलन्तमिव तेजोभिर्दुर्निरीक्ष्यं पृथग्जनैः} %||1-73-17||

\twolineshloka
{स्कन्धे चासज्य परशुं धनुर्विद्युद्गणोपमम्}
{प्रगृह्य शरमुख्यं च त्रिपुरघ्नं यथा हरम्} %||1-73-18||

\threelineshloka
{तं दृष्ट्वा भीमसङ्काशं ज्वलन्तमिव पावकम्}
{वसिष्ठप्रमुखा विप्रा जपहोमपरायणाः}
{सङ्गता मुनयः सर्वे संजजल्पुरथो मिथः} %||1-73-19||

\threelineshloka
{कच्चित्पितृवधामर्षी क्षत्रं नोत्सादयिष्यति}
{पूर्वं क्षत्रवधं कृत्वा गतमन्युर्गतज्वरः}
{क्षत्रस्योत्सादनं भूयो न खल्वस्य चिकीर्षितम्} %||1-73-20||

\twolineshloka
{एवमुक्त्वार्घ्यमादाय भार्गवं भीमदर्शनम्}
{ऋषयो राम रामेति मधुरां वाचमब्रुवन्} %||1-73-21||

\twolineshloka
{प्रतिगृह्य तु तां पूजामृषिदत्तां प्रतापवान्}
{रामं दाशरथिं रामो जामदग्न्योऽभ्यभाषत} %||1-74-22||


{॥इत्यार्षे श्रीमद्रामायणे वाल्मीकीये आदिकाव्ये बालकाण्डे त्रिसप्ततितमः सर्गः॥}

\sect{चतुःसप्ततितमः सर्गः}

\twolineshloka
{राम दाशरथे वीर वीर्यं ते श्रूयतेऽधुतम्}
{धनुषो भेदनं चैव निखिलेन मया श्रुतम्} %||1-74-1||

\twolineshloka
{तदद्भुतमचिन्त्यं च भेदनं धनुषस्त्वया}
{तच्छ्रुत्वाहमनुप्राप्तो धनुर्गृह्यापरं शुभम्} %||1-74-2||

\twolineshloka
{तदिदं घोरसङ्काशं जामदग्न्यं महद्धनुः}
{पूरयस्व शरेणैव स्वबलं दर्शयस्व च} %||1-74-3||

\twolineshloka
{तदहं ते बलं दृष्ट्वा धनुषोऽस्य प्रपूरणे}
{द्वन्द्वयुद्धं प्रदास्यामि वीर्यश्लाघ्यमिदं तव} %||1-74-4||

\twolineshloka
{तस्य तद्वचनं श्रुत्वा राजा दशरतःस्तदा}
{विषण्णवदनो दीनः प्राञ्जलिर्वाक्यमब्रवीत्} %||1-74-5||

\twolineshloka
{क्षत्ररोषात्प्रशान्तस्त्वं ब्राह्मणस्य महायशाः}
{बालानां मम पुत्राणामभयं दातुमर्हसि} %||1-74-6||

\twolineshloka
{भार्गवाणां कुले जातः स्वाध्यायव्रतशालिनाम्}
{सहस्राक्षे प्रतिज्ञाय शस्त्रं निक्षिप्तवानसि} %||1-74-7||

\twolineshloka
{स त्वं धर्मपरो भूत्वा काश्यपाय वसुन्धराम्}
{दत्त्वा वनमुपागम्य महेन्द्रकृतकेतनः} %||1-74-8||

\twolineshloka
{मम सर्वविनाशाय सम्प्राप्तस्त्वं महामुने}
{न चैकस्मिन्हते रामे सर्वे जीवामहे वयम्} %||1-74-9||

\twolineshloka
{ब्रुवत्येवं दशरथे जामदग्न्यः प्रतापवान्}
{अनादृत्यैव तद्वाक्यं राममेवाभ्यभाषत} %||1-74-10||

\twolineshloka
{इमे द्वे धनुषी श्रेष्ठे दिव्ये लोकाभिविश्रुते}
{दृढे बलवती मुख्ये सुकृते विश्वकर्मणा} %||1-74-11||

\twolineshloka
{अतिसृष्टं सुरैरेकं त्र्यम्बकाय युयुत्सवे}
{त्रिपुरघ्नं नरश्रेष्ठ भग्नं काकुत्स्ह यत्त्वया} %||1-74-12||

\twolineshloka
{इदं द्वितीयं दुर्धर्षं विष्णोर्दत्तं सुरोत्तमैः}
{समानसारं काकुत्स्थ रौद्रेण धनुषा त्विदम्} %||1-74-13||

\twolineshloka
{तदा तु देवताः सर्वाः पृच्छन्ति स्म पितामहम्}
{शितिकण्ठस्य विष्णोश्च बलाबलनिरीक्षया} %||1-74-14||

\twolineshloka
{अभिप्रायं तु विज्ञाय देवतानां पितामहः}
{विरोधं जनयामास तयोः सत्यवतां वरः} %||1-74-15||

\twolineshloka
{विरोधे च महद्युद्धमभवद्रोमहर्षणम्}
{शितिकण्ठस्य विष्णोश्च परस्परजयैषिणोः} %||1-74-16||

\twolineshloka
{तदा तज्जृम्भितं शैवं धनुर्भीमपराक्रमम्}
{हुङ्कारेण महादेवः स्तम्भितोऽथ त्रिलोचनः} %||1-74-17||

\twolineshloka
{देवैस्तदा समागम्य सर्षिसङ्घैः सचारणैः}
{याचितौ प्रशमं तत्र जग्मतुस्तौ सुरोत्तमौ} %||1-74-18||

\twolineshloka
{जृम्भितं तद्धनुर्दृष्ट्वा शैवं विष्णुपराक्रमैः}
{अधिकं मेनिरे विष्णुं देवाः सर्षिगणास्तदा} %||1-74-19||

\twolineshloka
{धनू रुद्रस्तु सङ्क्रुद्धो विदेहेषु महायशाः}
{देवरातस्य राजर्षेर्ददौ हस्ते ससायकम्} %||1-74-20||

\twolineshloka
{इदं च विष्णवं राम धनुः परपुरंजयम्}
{ऋचीके भार्गवे प्रादाद्विष्णुः स न्यासमुत्तमम्} %||1-74-21||

\twolineshloka
{ऋचीकस्तु महातेजाः पुत्रस्याप्रतिकर्मणः}
{पितुर्मम ददौ दिव्यं जमदग्नेर्महात्मनः} %||1-74-22||

\twolineshloka
{न्यस्तशस्त्रे पितरि मे तपोबलसमन्विते}
{अर्जुनो विदधे मृत्युं प्राकृतां बुद्धिमास्थितः} %||1-74-23||

\twolineshloka
{वधमप्रतिरूपं तु पितुः श्रुत्वा सुदारुणम्}
{क्षत्रमुत्सादयं रोषाज्जातं जातमनेकशः} %||1-74-24||

\twolineshloka
{पृथिवीं चाखिलां प्राप्य काश्यपाय महात्मने}
{यज्ञस्यान्ते तदा राम दक्षिणां पुण्यकर्मणे} %||1-74-25||

\twolineshloka
{दत्त्वा महेन्द्रनिलयस्तपोबलसमन्वितः}
{श्रुतवान्धनुषो भेदं ततोऽहं द्रुतमागतः} %||1-74-26||

\twolineshloka
{तदिदं वैष्णवं राम पितृपैतामहं महत्}
{क्षत्रधर्मं पुरस्कृत्य गृह्णीष्व धनुरुत्तमम्} %||1-74-27||

\twolineshloka
{योजयस्व धनुः श्रेष्ठे शरं परपुरंजयम्}
{यदि शक्नोषि काकुत्स्थ द्वन्द्वं दास्यामि ते ततः} %||1-75-28||


{॥इत्यार्षे श्रीमद्रामायणे वाल्मीकीये आदिकाव्ये बालकाण्डे चतुःसप्ततितमः सर्गः॥}

\sect{पञ्चसप्ततितमः सर्गः}

\twolineshloka
{श्रुत्वा तज्जामदग्न्यस्य वाक्यं दाशरथिस्तदा}
{गौरवाद्यन्त्रितकथः पितू राममथाब्रवीत्} %||1-75-1||

\twolineshloka
{श्रुतवानस्मि यत्कर्म कृतवानसि भार्गव}
{अनुरुन्ध्यामहे ब्रह्मन्पितुरानृण्यमास्थितः} %||1-75-2||

\twolineshloka
{वीर्यहीनमिवाशक्तं क्षत्रधर्मेण भार्गव}
{अवजानामि मे तेजः पश्य मेऽद्य पराक्रमम्} %||1-75-3||

\twolineshloka
{इत्युक्त्वा राघवः क्रुद्धो भार्गवस्य वरायुधम्}
{शरं च प्रतिसङ्गृह्य हस्ताल्लघुपराक्रमः} %||1-75-4||

\twolineshloka
{आरोप्य स धनू रामः शरं सज्यं चकार ह}
{जामदग्न्यं ततो रामं रामः क्रुद्धोऽब्रवीद्वचः} %||1-75-5||

\twolineshloka
{ब्राह्मणोऽसीति पूज्यो मे विश्वामित्रकृतेन च}
{तस्माच्छक्तो न ते राम मोक्तुं प्राणहरं शरम्} %||1-75-6||

\twolineshloka
{इमां वा त्वद्गतिं राम तपोबलसमार्जितान्}
{लोकानप्रतिमान्वापि हनिष्यामि यदिच्छसि} %||1-75-7||

\twolineshloka
{न ह्ययं वैष्णवो दिव्यः शरः परपुरंजयः}
{मोघः पतति वीर्येण बलदर्पविनाशनः} %||1-75-8||

\twolineshloka
{वरायुधधरं राम द्रष्टुं सर्षिगणाः सुराः}
{पितामहं पुरस्कृत्य समेतास्तत्र सङ्घशः} %||1-75-9||

\twolineshloka
{गन्धर्वाप्सरसश्चैव सिद्धचारणकिंनराः}
{यक्षराक्षसनागाश्च तद्द्रष्टुं महदद्भुतम्} %||1-75-10||

\twolineshloka
{जडीकृते तदा लोके रामे वरधनुर्धरे}
{निर्वीर्यो जामदग्न्योऽसौ रमो राममुदैक्षत} %||1-75-11||

\twolineshloka
{तेजोभिर्हतवीर्यत्वाज्जामदग्न्यो जडीकृतः}
{रामं कमल पत्राक्षं मन्दं मन्दमुवाच ह} %||1-75-12||

\twolineshloka
{काश्यपाय मया दत्ता यदा पूर्वं वसुन्धरा}
{विषये मे न वस्तव्यमिति मां काश्यपोऽब्रवीत्} %||1-75-13||

\twolineshloka
{सोऽहं गुरुवचः कुर्वन्पृथिव्यां न वसे निशाम्}
{इति प्रतिज्ञा काकुत्स्थ कृता वै काश्यपस्य ह} %||1-75-14||

\twolineshloka
{तदिमां त्वं गतिं वीर हन्तुं नार्हसि राघव}
{मनोजवं गमिष्यामि महेन्द्रं पर्वतोत्तमम्} %||1-75-15||

\twolineshloka
{लोकास्त्वप्रतिमा राम निर्जितास्तपसा मया}
{जहि ताञ्शरमुख्येन मा भूत्कालस्य पर्ययः} %||1-75-16||

\twolineshloka
{अक्षय्यं मधुहन्तारं जानामि त्वां सुरेश्वरम्}
{धनुषोऽस्य परामर्शात्स्वस्ति तेऽस्तु परन्तप} %||1-75-17||

\twolineshloka
{एते सुरगणाः सर्वे निरीक्षन्ते समागताः}
{त्वामप्रतिमकर्माणमप्रतिद्वन्द्वमाहवे} %||1-75-18||

\twolineshloka
{न चेयं मम काकुत्स्थ व्रीडा भवितुमर्हति}
{त्वया त्रैलोक्यनाथेन यदहं विमुखीकृतः} %||1-75-19||

\twolineshloka
{शरमप्रतिमं राम मोक्तुमर्हसि सुव्रत}
{शरमोक्षे गमिष्यामि महेन्द्रं पर्वतोत्तमम्} %||1-75-20||

\twolineshloka
{तथा ब्रुवति रामे तु जामदग्न्ये प्रतापवान्}
{रामो दाशरथिः श्रीमांश्चिक्षेप शरमुत्तमम्} %||1-75-21||

\twolineshloka
{ततो वितिमिराः सर्वा दिशा चोपदिशस्तथा}
{सुराः सर्षिगणा रामं प्रशशंसुरुदायुधम्} %||1-75-22||

\twolineshloka
{रामं दाशरथिं रामो जामदग्न्यः प्रशस्य च}
{ततः प्रदक्षिणीकृत्य जगामात्मगतिं प्रभुः} %||1-76-23||


{॥इत्यार्षे श्रीमद्रामायणे वाल्मीकीये आदिकाव्ये बालकाण्डे पञ्चसप्ततितमः सर्गः॥}

\sect{षट्सप्ततितमः सर्गः}

\twolineshloka
{गते रामे प्रशान्तात्मा रामो दाशरथिर्धनुः}
{वरुणायाप्रमेयाय ददौ हस्ते ससायकम्} %||1-76-1||

\twolineshloka
{अभिवाद्य ततो रामो वसिष्ठ प्रमुखानृषीन्}
{पितरं विह्वलं दृष्ट्वा प्रोवाच रघुनन्दनः} %||1-76-2||

\twolineshloka
{जामदग्न्यो गतो रामः प्रयातु चतुरङ्गिणी}
{अयोध्याभिमुखी सेना त्वया नाथेन पालिता} %||1-76-3||

\twolineshloka
{रामस्य वचनं श्रुत्वा राजा दशरथः सुतम्}
{बाहुभ्यां सम्परिष्वज्य मूर्ध्नि चाघ्राय राघवम्} %||1-76-4||

\twolineshloka
{गतो राम इति श्रुत्वा हृष्टः प्रमुदितो नृपः}
{चोदयामास तां सेनां जगामाशु ततः पुरीम्} %||1-76-5||

\twolineshloka
{पताकाध्वजिनीं रम्यां तूर्योद्घुष्टनिनादिताम्}
{सिक्तराजपथां रम्यां प्रकीर्णकुसुमोत्कराम्} %||1-76-6||

\twolineshloka
{राजप्रवेशसुमुखैः पौरैर्मङ्गलवादिभिः}
{सम्पूर्णां प्राविशद्राजा जनौघैः समलङ्कृताम्} %||1-76-7||

\twolineshloka
{कौसल्या च सुमित्रा च कैकेयी च सुमध्यमा}
{वधूप्रतिग्रहे युक्ता याश्चान्या राजयोषितः} %||1-76-8||

\twolineshloka
{ततः सीतां महाभागामूर्मिलां च यशस्विनीम्}
{कुशध्वजसुते चोभे जगृहुर्नृपपत्नयः} %||1-76-9||

\twolineshloka
{मङ्गलालापनैश्चैव शोभिताः क्षौमवाससः}
{देवतायतनान्याशु सर्वास्ताः प्रत्यपूजयन्} %||1-76-10||

\twolineshloka
{अभिवाद्याभिवाद्यांश्च सर्वा राजसुतास्तदा}
{रेमिरे मुदिताः सर्वा भर्तृभिः सहिता रहः} %||1-76-11||

\twolineshloka
{कृतदाराः कृतास्त्राश्च सधनाः ससुहृज्जनाः}
{शुश्रूषमाणाः पितरं वर्तयन्ति नरर्षभाः} %||1-76-12||

\twolineshloka
{तेषामतियशा लोके रामः सत्यपराक्रमः}
{स्वयम्भूरिव भूतानां बभूव गुणवत्तरः} %||1-76-13||

\twolineshloka
{रामस्तु सीतया सार्धं विजहार बहूनृतून्}
{मनस्वी तद्गतस्तस्या नित्यं हृदि समर्पितः} %||1-76-14||

\twolineshloka
{प्रिया तु सीता रामस्य दाराः पितृकृता इति}
{गुणाद्रूपगुणाच्चापि प्रीतिर्भूयो व्यवर्धत} %||1-76-15||

\twolineshloka
{तस्याश्च भर्ता द्विगुणं हृदये परिवर्तते}
{अन्तर्जातमपि व्यक्तमाख्याति हृदयं हृदा} %||1-76-16||

\twolineshloka
{तस्य भूयो विशेषेण मैथिली जनकात्मजा}
{देवताभिः समा रूपे सीता श्रीरिव रूपिणी} %||1-76-17||

\fourlineindentedshloka
{तया स राजर्षिसुतोऽभिरामया}
{ समेयिवानुत्तमराजकन्यया}
{अतीव रामः शुशुभेऽतिकामया}
{ विभुः श्रिया विष्णुरिवामरेश्वरः} %||1-77-18||


{॥इत्यार्षे श्रीमद्रामायणे वाल्मीकीये आदिकाव्ये बालकाण्डे षट्सप्ततितमः सर्गः॥}

